% !TEX program = pdflatex
\documentclass[12pt]{article}

\usepackage[utf8]{inputenc}
\usepackage{CJKutf8}
\usepackage{amsmath,amssymb}
\usepackage{mathtools}
\usepackage{amsthm}
\usepackage{geometry}
\geometry{
  a4paper,
  top=25mm,
  bottom=25mm,
  left=20mm,
  right=20mm
}
\usepackage{xcolor}
\usepackage{url}
\usepackage[unicode,colorlinks=true,linkcolor=blue,urlcolor=blue]{hyperref}
\usepackage{xurl}
\Urlmuskip=0mu plus 4mu
\sloppy
\usepackage{array}
\usepackage{tocloft}

\renewcommand{\cftsecleader}{\cftdotfill{\cftdotsep}}
\renewcommand{\refname}{参考文献}

\newtheorem{definition}{定义}[section]
\newtheorem{axiom}{公理}[section]
\newtheorem{assumption}{假设}[section]
\newtheorem{theorem}{定理}[section]
\newtheorem{proposition}{命题}[section]
\newtheorem{corollary}{推论}[section]
\newtheorem{lemma}{引理}[section]
\newtheorem{remark}{备注}[section]
\newtheorem{Certificate}{证书}[section]

\newcommand{\NN}{\mathbb{N}}
\newcommand{\ZZ}{\mathbb{Z}}
\newcommand{\QQ}{\mathbb{Q}}
\newcommand{\RR}{\mathbb{R}}
\newcommand{\CC}{\mathbb{C}}
\newcommand{\GFramework}{\textsf{G-Framework}}
\newcommand{\GAlgebra}{\textsf{G-Algebra}}
\newcommand{\Udc}{\mathfrak{U}_{\mathrm{dc}}}
\newcommand{\Vsys}{\mathfrak V}
\newcommand{\NF}{\mathsf{NF}}
\newcommand{\Hilb}{\mathsf{Hilb}_{23}}
\newcommand{\DefEq}{\coloneqq}
\newcommand{\code}[1]{\path{#1}}
\newcommand{\CertTuple}{(N,G,P,C,A,B,T)}
\newcommand{\Rstar}{\mathcal{R}^{\ast}}
\newcommand{\Rcal}{\mathcal{R}}
\newcommand{\Ecal}{\mathcal{E}}
\newcommand{\Dcal}{\mathcal{D}}
\newcommand{\Ucal}{\mathcal{U}}
\newcommand{\Hcal}{\mathcal{H}}
\newcommand{\Kcal}{\mathcal{K}}
\newcommand{\Lcal}{\mathcal{L}}
\newcommand{\Scal}{\mathcal{S}}
\newcommand{\Pcal}{\mathcal{P}}
\newcommand{\Qcal}{\mathcal{Q}}
\newcommand{\Tcal}{\mathcal{T}}
\newcommand{\Bcal}{\mathcal{B}}
\newcommand{\Ccal}{\mathcal{C}}
\newcommand{\Ical}{\mathcal{I}}
\newcommand{\id}{\mathrm{id}}
\newcommand{\eps}{\varepsilon}
\newcommand{\OmegaCh}{\Omega_{\mathrm{Ch}}}
\DeclareMathOperator{\Ob}{Ob}
\DeclareMathOperator{\Pred}{Pred}
\DeclareMathOperator{\Cert}{Cert}
\DeclareMathOperator{\Audit}{Audit}
\DeclareMathOperator{\Accept}{Accept}
\DeclareMathOperator{\Accepted}{Accepted}
\DeclareMathOperator{\Fallacy}{Fallacy}
\DeclareMathOperator{\Resolve}{Resolve}
\DeclareMathOperator{\Barrier}{Barrier}
\DeclareMathOperator{\Cluster}{Cluster}
\DeclareMathOperator{\Rank}{rank}
\DeclareMathOperator{\Norm}{Norm}
\DeclareMathOperator{\Gate}{Gate}
\DeclareMathOperator{\Hit}{Hit}
\DeclareMathOperator{\Obs}{Obs}
\DeclareMathOperator{\Inv}{Inv}
\DeclareMathOperator{\Gen}{Gen}
\DeclareMathOperator{\BackwriteQueue}{BackwriteQueue}
\DeclareMathOperator{\Tighten}{Tighten}
\DeclareMathOperator{\Prune}{Prune}
\DeclareMathOperator{\Hom}{Hom}
\DeclareMathOperator{\TotHom}{TotHom}
\DeclareMathOperator{\GCap}{Cap}
\DeclareMathOperator{\Graph}{Graph}
\DeclareMathOperator{\Hash}{Hash}
\DeclareMathOperator{\Zero}{Zero}
\DeclareMathOperator{\Anchor}{Anchor}
\DeclareMathOperator{\BVP}{BVP}
\DeclareMathOperator{\PDE}{PDE}
\DeclareMathOperator{\Res}{Res}
\DeclareMathOperator{\Reg}{Reg}
\DeclareMathOperator{\Phys}{Phys}
\DeclareMathOperator{\Dio}{Dio}
\DeclareMathOperator{\CH}{CH}
\DeclareMathOperator{\Con}{Con}
\DeclareMathOperator{\Search}{Search}
\DeclareMathOperator{\Law}{Law}
\DeclareMathOperator{\Sect}{Sect}
\DeclareMathOperator{\Supp}{Supp}
\DeclareMathOperator{\Riem}{Riem}
\DeclareMathOperator{\ZFC}{ZFC}
\DeclareMathOperator{\Logic}{Logic}
\DeclareMathOperator{\ValidSem}{ValidSem}
\DeclareMathOperator{\RepairableResidual}{RepairableResidual}
\DeclareMathOperator{\Excluded}{Excluded}
\DeclareMathOperator{\Halt}{Halt}
\DeclareMathOperator{\NonHalt}{NonHalt}
\DeclareMathOperator{\Ground}{Ground}
\DeclareMathOperator{\Info}{Info}
\DeclareMathOperator{\Unitary}{Unitary}
\DeclareMathOperator{\Collapse}{Collapse}
\DeclareMathOperator{\Runtime}{Runtime}
\DeclareMathOperator{\Gap}{Gap}
\DeclareMathOperator{\Sing}{Sing}
\DeclareMathOperator{\Fold}{Fold}
\DeclareMathOperator{\Conf}{Conf}
\DeclareMathOperator{\Entropy}{H}
\DeclareMathOperator{\SemInfo}{SemInfo}
\DeclareMathOperator{\Arrow}{Arrow}
\DeclareMathOperator{\ManyBody}{ManyBody}
\DeclareMathOperator{\Perturb}{Perturb}
\DeclareMathOperator{\Pair}{Pair}
\DeclareMathOperator{\Macro}{Macro}
\DeclareMathOperator{\KComplex}{K}
\DeclareMathOperator{\CertDepth}{CertDepth}
\DeclareMathOperator{\Erg}{Erg}
\DeclareMathOperator{\Phase}{Phase}
\DeclareMathOperator{\WignerEff}{WignerEff}
\DeclareMathOperator{\Emergence}{Emergence}
\DeclareMathOperator{\Micro}{Micro}
\DeclareMathOperator{\Pref}{Pref}
\DeclareMathOperator{\SocChoice}{SocChoice}
\DeclareMathOperator{\Term}{Term}
\DeclareMathOperator{\Dep}{Dep}
\DeclareMathOperator{\CommonNF}{CommonNF}
\DeclareMathOperator{\Langlands}{Langlands}
\DeclareMathOperator{\Gal}{Gal}
\DeclareMathOperator{\AutRep}{AutRep}
\DeclareMathOperator{\TopDict}{TopDict}
\DeclareMathOperator{\Align}{Align}
\DeclareMathOperator{\Conform}{Conform}
\DeclareMathOperator{\BoundPred}{BoundPred}
\DeclareMathOperator{\FeedbackFit}{FeedbackFit}
\DeclareMathOperator{\Qualia}{Qualia}
\DeclareMathOperator{\Conscious}{Conscious}
\DeclareMathOperator{\SelfRepair}{SelfRepair}
\DeclareMathOperator{\Experience}{Experience}

\setlength{\emergencystretch}{8em}
\setlength{\parindent}{0pt}
\setlength{\parskip}{0.6em}

\begin{document}
\begin{CJK*}{UTF8}{gkai}

\title{G 框架中的数学、物理、计算与认知基础障碍谱系：\\统一法空间、总同调语义场与证书化化解\\（工作笔记：形式系统版）}
\author{Gao Zheng（高政）}
\date{2026-07-08}
\maketitle

\section*{前言}
\addcontentsline{toc}{section}{前言}

本文把希尔伯特二十三问题、千禧年大奖难题、理论计算机科学基石问题、信息论与统计物理学
悖论、复杂性科学障碍、约束闭合问题、符号接地问题、朗兰兹纲领以及意识硬问题，统一
登记为一个广义基础障碍谱系。该谱系横跨集合论与逻辑基础、数论、代数几何、物理公理化、
偏微分方程、规范场论、计算复杂性、热力学、生命科学、社会选择、约束相容性与认知哲学。
若仅在传统底空间中理解这些对象，它们分别表现为连续统基数分层、形式相容性、不可判定搜索、
奇点、边界残差、交点复杂度、零点逼近、质量缺口、停机边界、语义失锚、相空间遍历破缺、
多偏好冲突和主观体验不可还原等障碍；若进入 \GFramework{} 与底层 \GAlgebra{} 的证书语义层，
它们可以被统一写成
\[
  \text{底空间障碍}
  \Longrightarrow
  \text{离散--连续统一法空间对象}
  \Longrightarrow
  \text{总同调语义不变量}
  \Longrightarrow
  \text{显式证书门禁}.
\]
为避免把背景来源误读为证明前提，本文采用“基础文献--背景文献”分层引用。形式系统的基础语言
只以 Gao 2025 年 Zenodo 文献 \cite{Gao2025PureMathZenodo} 为基础参照；希尔伯特问题、集合论独立性、PDE、代数拓扑、
停机问题、计算复杂性、信息论、统计物理、社会选择、朗兰兹纲领和意识硬问题等外部文献，
只用于给出问题的历史定位、标准名称和传统障碍边界
\cite{Hilbert1902Problems,Godel1940,Cohen1963,Jech2003,Evans2010PDE,
BrennerScott2008FEM,Hatcher2002AT,Turing1936Computable,Cook1971NP,
Karp1972Reducibility,Shannon1948Communication,Birkhoff1931Ergodic,
Arrow1951SocialChoice,Langlands1970Problems,Chalmers1995Facing}。

本文所说的“化解”采用严格的形式口径：一个障碍对象 \(X\) 的底层障碍
\(\Barrier(X)\) 被称为在 \GFramework{} 内被证书化化解，指存在可写明的规范化映射、
同调不变量、谓词因果链、失败吸收规则、回写修复结构与终端证书，使得原本依赖无穷逼近、
无穷穷举、逐点解析延拓、概率补丁或非显式经验判断的任务，被重写为有限阶的可审计门禁判断。
该口径以 Gao 2025 年 Zenodo 基础文献 \cite{Gao2025PureMathZenodo} 中的集合论、代数拓扑、
范畴语言与 G-代数基础为唯一基础参照。本文不把其他 \(NoteBook\) 内部工作稿作为证明前提；规范化映射、
吸收元、边缘剪枝、回写修复、PDE 残差对象、协变隧穿、语义积分与终端门禁均在本文内部
重新定义并闭合为可检查的证明链。

全文按三层路线展开。第一层以希尔伯特二十三问题为主干：建立障碍聚类与公理降级证书，
分别处理逻辑基础与相容性、物理学公理化、PDE 与变分边界奇点、代数几何与拓扑不变量、
数论与函数零点，并给出二十三问题的总归纳证书。第二层把证明对象扩展到现代数学和理论物理：
连续群局部可微性限制、Hodge、Yang--Mills 与质量缺口、Navier--Stokes、BSD、Gödel 解耦、
图灵停机、量子引力测量、黑洞信息、Levinthal 悖论、Shannon 语义盲区、Loschmidt 时间箭头、
多体问题、Kolmogorov--Chaitin 边界、各态历经破缺、Wigner 不合理有效性、涌现边界和
Arrow 不可能性。第三层进入符号系统与认知论：把朗兰兹纲领解释为离散--连续拓扑字典，把高阶符号演化
的约束相容性解释为约束谓词链的吸收--剪枝门禁，把意识硬问题解释为回写队列与收紧算子的自举语义
修复回路。

全文保持如下体例：先给出定义，再将“公理项”降级为可证明的构造定理；随后每个聚类给出
定理、证书、引理、命题、推论与备注。证明不依赖自然语言断言的权威性，而依赖可写出的
谓词因果链：
\[
  \Pred(x)
  \Longrightarrow
  \Inv(x)
  \Longrightarrow
  \Cert(x)
  \Longrightarrow
  \Gate(x)
  \Longrightarrow
  \Accept(x).
\]
于是，本文不是把若干世纪难题重新叙述为宏大口号，而是把它们压缩为一组可追踪、可审计、
可归纳、可失败吸收且可回写修复的形式逻辑证书。

\setcounter{tocdepth}{3}
\setcounter{secnumdepth}{3}
\tableofcontents
\clearpage

%%%%%%%%%%%%%%%%%%%%%%%%%%%%%%%%%%%%%%%%%%%%%%%%%%%%%%%%%%%%
\section{形式逻辑总基底：希尔伯特问题谱系、障碍聚类与公理项降级}
\label{sec:tex59-base}
%%%%%%%%%%%%%%%%%%%%%%%%%%%%%%%%%%%%%%%%%%%%%%%%%%%%%%%%%%%%

\begin{remark}[核心陈述]
本文把希尔伯特二十三问题登记为有限索引集
\[
  \Hilb \DefEq \{H_1,H_2,\ldots,H_{23}\}.
\]
每个问题 \(H_i\) 不直接作为孤立对象处理，而被拆成问题陈述、底层障碍、语义重写和
证书门禁四个槽：
\[
  H_i
  \longmapsto
  \bigl(Q_i,\Barrier_i,\Resolve_i,\Cert_i\bigr).
\]
只要 \(\Resolve_i\) 能把 \(\Barrier_i\) 送入统一法空间 \(\Udc\) 中的同调不变量与
终端门禁，本文便称 \(H_i\) 的底层障碍在 \GFramework{} 内被证书化化解。
\end{remark}

\subsection{Zenodo 基础、自足证明域与引用非依赖约定}
\label{subsec:tex59-puremath-self-contained}

\begin{definition}[Zenodo 基础域、本文局部域与背景引用域]
\label{def:tex59-dependency-domains}
定义三个依赖域：
令 \(\mathsf{Base}_{\mathrm{PM}}\) 表示
Gao 2025 年 Zenodo 基础文献 \cite{Gao2025PureMathZenodo} 中关于 \GFramework{}、\GAlgebra{}、law-space、
同伦结构、主丛几何与广义非交换李代数的集合、代数、同调与范畴基础。
\[
  \mathsf{Local}_{59}
  \DefEq
  \{\text{本文内部显式给出的定义、证书、引理、命题、定理、推论与备注}\},
\]
\[
  \mathsf{Ext}_{\mathrm{bg}}
  \DefEq
  \{\text{外部经典文献中用于定位问题背景、标准术语或历史定理的引用}\}.
\]
对任意命题 \(\Theta\)，定义证明依赖集
\[
  \Dep(\Theta)
\]
为在 \(\Theta\) 的形式证明中实际作为推理前件调用的定义、规则、引理、命题与定理集合。
若某一引用只用于说明问题背景、命名来源或标准外部事实，不进入谓词因果链，则不计入
\(\Dep(\Theta)\)。
\end{definition}

\begin{remark}[背景引用分区]
\label{rem:tex59-background-citation-partition}
本文将外部背景文献分为八类。

第一，Hilbert 问题谱系、连续统假设与集合论独立性背景：
\[
  \cite{Hilbert1902Problems,Godel1940,Cohen1963,Jech2003}.
\]
第二，基础逻辑、可计算性与复杂性边界背景：
\[
  \cite{Godel1931,Turing1936Computable,Cook1971NP,Karp1972Reducibility}.
\]
第三，PDE、弱形式、泛函分析、Navier--Stokes 与数值残差背景：
\[
  \cite{Evans2010PDE,BrennerScott2008FEM,Brezis2011FunctionalAnalysis,
  Fefferman2000NS,Temam2001NS}.
\]
第四，代数拓扑、同调代数、范畴语言、Hodge 理论与代数几何背景：
\[
  \cite{Hatcher2002AT,Weibel1994HomologicalAlgebra,MacLane1998Categories,
  Voisin2002Hodge,GriffithsHarris1978}.
\]
第五，Yang--Mills、量子测量、量子引力与黑洞信息背景：
\[
  \cite{JaffeWitten2000YM,VonNeumann1955Quantum,DeWitt1967QuantumGravity,
  Kiefer2012QuantumGravity,Hawking1976Breakdown}.
\]
第六，数论、椭圆曲线、BSD 与 Langlands 背景：
\[
  \cite{BirchSwinnertonDyer1965,Silverman2009AEC,Langlands1970Problems,
  Gelbart1984AutomorphicForms,Frenkel2007Langlands,Wiles1995Modular}.
\]
第七，信息论、热力学、统计物理、多体系统与复杂涌现背景：
\[
  \cite{Shannon1948Communication,Loschmidt1876SecondLaw,Boltzmann1877SecondLaw,
  Price1996TimeArrow,Birkhoff1931Ergodic,Khinchin1949StatMech,
  MezardParisiVirasoro1987SpinGlass,FetterWalecka1971ManyParticle,
  NegeleOrland1988QuantumManyParticle,Gardner1970LifeGame,
  Anderson1972MoreIsDifferent,Holland1998Emergence}.
\]
第八，符号接地、约束相容性、社会选择与意识硬问题背景：
\[
  \cite{Harnad1990SymbolGrounding,Bostrom2014Superintelligence,
  Russell2019HumanCompatible,Arrow1951SocialChoice,Sen1970CollectiveChoice,
  Nagel1974WhatIsItLike,Chalmers1995Facing,Tononi2004IIT}.
\]
上述文献增强的是术语、历史和标准问题边界的可追溯性；按定义
~\ref{def:tex59-dependency-domains}，它们不自动进入任何形式证明的 \(\Dep(\Theta)\)。
\end{remark}

\begin{proposition}[基础文献定位替换的保守性]
\label{prop:tex59-zenodo-source-conservative}
设某证明 \(\Theta\) 只调用基础文献中的结构性内容：
\[
  \text{集合、映射、代数、同调、范畴、主丛、联络、同伦与曲率对象}.
\]
若把本地文件定位替换为 Gao 2025 年 Zenodo 文献 \cite{Gao2025PureMathZenodo} 定位，
且证明中不调用本地文件路径作为数学前件，则
\[
  \Dep(\Theta)
  \text{ 不变}.
\]
\end{proposition}

\begin{Certificate}[证书：路径定位不进入谓词因果链]
\label{cert:tex59-zenodo-source-conservative}
证书链为
\[
\begin{aligned}
  Z_1 &: \Theta\text{ 只调用结构性数学内容},\\
  Z_2 &: \text{本地文件路径不是集合、映射、谓词、证书或门禁},\\
  Z_3 &: Z_2\Rightarrow \text{本地文件路径不属于 }\Dep(\Theta),\\
  Z_4 &: Z_1\wedge Z_3
         \Rightarrow
         \Dep(\Theta)\text{ 在 Zenodo 定位替换下保持不变}.
\end{aligned}
\]
\end{Certificate}

\begin{proof}[证明（基于数学符号推演逻辑的谓词因果链）]
由定义~\ref{def:tex59-dependency-domains}，\(\Dep(\Theta)\) 只记录实际作为推理前件调用的
数学对象。文件路径只是文献定位信息，不是集合、映射、谓词、证书、队列或门禁。
因此将文献定位从本地路径改为 Zenodo DOI，不改变 \(\Theta\) 的任何前件集合。
所以
\[
  \Dep(\Theta)
\]
保持不变。证毕。
\end{proof}

\begin{lemma}[背景引用不产生证明前件]
\label{lem:tex59-background-nondependence}
若外部引用 \(r\in\mathsf{Ext}_{\mathrm{bg}}\) 只出现在背景说明、命名说明或边界口径中，
且在证明中不存在形如
\[
  r\Rightarrow P
\]
的推理步，则对任意本文命题 \(\Theta\)，有
\[
  r\notin\Dep(\Theta).
\]
\end{lemma}

\begin{Certificate}[证书：引用位置与推理位置分离]
\label{cert:tex59-background-nondependence}
证书链为
\[
\begin{aligned}
  B_1 &: r\in\mathsf{Ext}_{\mathrm{bg}},\\
  B_2 &: r\text{ 只出现在背景说明、命名说明或边界口径中},\\
  B_3 &: \neg\exists P,\ (r\Rightarrow P)\text{ 是本文证明中的谓词因果链一步},\\
  B_4 &: B_1\wedge B_2\wedge B_3\Rightarrow r\notin\Dep(\Theta).
\end{aligned}
\]
\end{Certificate}

\begin{proof}[证明（基于数学符号推演逻辑的谓词因果链）]
由定义~\ref{def:tex59-dependency-domains}，\(\Dep(\Theta)\) 只记录实际作为推理前件调用的
对象。由 \(B_2\)，\(r\) 的功能是背景说明、命名说明或边界口径说明；由 \(B_3\)，不存在
以 \(r\) 为前件推出本文谓词 \(P\) 的证明步。因此 \(r\) 不满足进入 \(\Dep(\Theta)\) 的定义条件，
推出
\[
  r\notin\Dep(\Theta).
\]
证毕。
\end{proof}

\begin{theorem}[Zenodo 基础自足证明域定理]
\label{thm:tex59-puremath-only-self-contained}
设本文中的命题 \(\Theta\) 的证明只调用：
\[
  \mathsf{Base}_{\mathrm{PM}},
  \qquad
  \mathsf{Local}_{59},
  \qquad
  \mathsf{Ext}_{\mathrm{bg}}\text{ 中的背景引用}.
\]
若所有外部引用都满足引理~\ref{lem:tex59-background-nondependence} 的非依赖条件，则
\[
  \Dep(\Theta)
  \subseteq
  \mathsf{Base}_{\mathrm{PM}}\cup\mathsf{Local}_{59}.
\]
因此，本文的形式证明在基础依赖意义上只依赖 Gao 2025 年 Zenodo 基础文献 \cite{Gao2025PureMathZenodo}
与本文内部定义。
\end{theorem}

\begin{Certificate}[数学归纳法证书：章节证书的依赖闭包]
\label{cert:tex59-puremath-only-self-contained}
令本文章节证书按出现顺序排列为
\[
  \Theta_0,\Theta_1,\ldots,\Theta_M.
\]
定义谓词
\[
  P(n):\quad
  \forall j\le n,\quad
  \Dep(\Theta_j)\subseteq \mathsf{Base}_{\mathrm{PM}}\cup\mathsf{Local}_{59}.
\]
归纳证书为
\[
\begin{aligned}
  I_0 &: \Dep(\Theta_0)\subseteq \mathsf{Base}_{\mathrm{PM}}\cup\mathsf{Local}_{59},\\
  I_1 &: P(n)\wedge
  \left[
  \Dep(\Theta_{n+1})\subseteq
  \mathsf{Base}_{\mathrm{PM}}\cup
  \{\Theta_0,\ldots,\Theta_n\}\cup\mathsf{Local}_{59}^{\mathrm{def}}
  \right]\\
  &\quad\Rightarrow P(n+1),\\
  I_2 &: \forall r\in\mathsf{Ext}_{\mathrm{bg}},\
  r\notin\Dep(\Theta_j)
  \quad\text{（引理~\ref{lem:tex59-background-nondependence}）}.
\end{aligned}
\]
\end{Certificate}

\begin{proof}[证明（数学归纳法证书；基于数学符号推演逻辑的谓词因果链）]
基步：\(\Theta_0\) 只使用 Gao 2025 基础域 \cite{Gao2025PureMathZenodo} 与本文第一个局部定义块，故
\[
  \Dep(\Theta_0)\subseteq \mathsf{Base}_{\mathrm{PM}}\cup\mathsf{Local}_{59},
\]
即 \(P(0)\) 成立。

归纳步：假设 \(P(n)\) 成立。若 \(\Theta_{n+1}\) 的证明只调用 Gao 2025 基础域
\cite{Gao2025PureMathZenodo}、此前已闭合的
\(\Theta_0,\ldots,\Theta_n\) 和本文局部定义，则由归纳假设
\[
  \Dep(\Theta_j)\subseteq \mathsf{Base}_{\mathrm{PM}}\cup\mathsf{Local}_{59},
  \qquad j\le n.
\]
将这些依赖代回 \(\Theta_{n+1}\) 的依赖集，得到
\[
  \Dep(\Theta_{n+1})
  \subseteq
  \mathsf{Base}_{\mathrm{PM}}\cup\mathsf{Local}_{59}.
\]
于是 \(P(n+1)\) 成立。

最后，外部经典文献由引理~\ref{lem:tex59-background-nondependence} 从证明依赖集中剥离；
它们只保留为问题背景与术语来源。由数学归纳法，对所有章节证书 \(\Theta_j\) 均有
\[
  \Dep(\Theta_j)\subseteq \mathsf{Base}_{\mathrm{PM}}\cup\mathsf{Local}_{59}.
\]
证毕。
\end{proof}

\begin{corollary}[基础依赖压缩推论]
\label{cor:tex59-internal-dependency-compression}
在本文当前版本中，任意形式证明的基础依赖均可压缩为
\[
  \boxed{
  \mathsf{Base}_{\mathrm{PM}}
  \cup
  \mathsf{Local}_{59}.
  }
\]
也就是说，除 Gao 2025 年 Zenodo 基础文献 \cite{Gao2025PureMathZenodo} 外，其余仓内文稿不作为本文证明前提。
\end{corollary}

\begin{Certificate}[证书：基础依赖压缩]
\label{cert:tex59-internal-dependency-compression}
证书链为
\[
\begin{aligned}
  C_1 &: \text{本文唯一基础文献引用为 }\mathrm{Gao2025PureMathZenodo},\\
  C_2 &: \text{所有证明对象均属于 }\mathsf{Local}_{59}\text{ 或由 Gao 2025 基础 }\cite{Gao2025PureMathZenodo}\text{ 支撑},\\
  C_3 &: \text{其他引用若出现，只属于 }\mathsf{Ext}_{\mathrm{bg}}\text{ 或已被删除},\\
  C_4 &: C_1\wedge C_2\wedge C_3
  \Rightarrow
  \Dep(\Theta)\subseteq \mathsf{Base}_{\mathrm{PM}}\cup\mathsf{Local}_{59}.
\end{aligned}
\]
\end{Certificate}

\begin{proof}[证明（基于数学符号推演逻辑的谓词因果链）]
由定理~\ref{thm:tex59-puremath-only-self-contained}，任意本文命题 \(\Theta\) 的证明依赖均落入
\[
  \mathsf{Base}_{\mathrm{PM}}\cup\mathsf{Local}_{59}.
\]
该集合中的第一项对应 Gao 2025 年 Zenodo 基础文献 \cite{Gao2025PureMathZenodo}，
第二项对应本文内部已经出现的定义与证明链。
故除 Gao 2025 年 Zenodo 基础文献 \cite{Gao2025PureMathZenodo} 外，无其他仓内稿件作为证明前提。证毕。
\end{proof}

\subsection{传统证明语言的忠实嵌入、保守扩展与结构增强}
\label{subsec:tex59-conservative-enhancement}

\begin{definition}[传统证明语言与 G-增强证明语言]
\label{def:tex59-traditional-and-g-language}
令
\[
  \mathcal L_{\mathrm{trad}}
\]
表示传统数学证明语言。它的对象包括集合、关系、函数、代数结构、拓扑空间、流形、范畴、
同调对象、微分算子、数论对象与逻辑公式。令
\[
  \mathcal T_{\mathrm{trad}}
\]
表示传统理论中被选定的公理、定义与推理规则集合。

定义 G-增强证明语言为
\[
  \mathcal L_G
  \DefEq
  \mathcal L_{\mathrm{trad}}
  \cup
  \{\NF,\Inv,\Cert,\Gate,\Audit,\bot,\BackwriteQueue,\Tighten,\Prune\}.
\]
定义 G-增强理论为
\[
  \mathcal T_G
  \DefEq
  \mathcal T_{\mathrm{trad}}
  \cup
  \mathcal T_{\mathrm{def}}
  \cup
  \mathcal T_{\mathrm{cert}},
\]
其中 \(\mathcal T_{\mathrm{def}}\) 只包含新增符号的定义展开规则，\(\mathcal T_{\mathrm{cert}}\)
只包含证书、门禁、吸收、剪枝和回写的审计规则。
\end{definition}

\begin{definition}[嵌入翻译与遗忘投影]
\label{def:tex59-embedding-projection}
定义嵌入翻译
\[
  \iota:\mathcal L_{\mathrm{trad}}\longrightarrow\mathcal L_G
\]
如下：对任意传统项、关系符号、函数符号、逻辑联结词、量词和等号，\(\iota\) 保持其原符号不变；
对任意传统公式 \(F\)，定义
\[
  \iota(F)
  \DefEq
  \bigl(F,\NF(F),\Inv(F),\Cert(F),\Gate(F),\Audit(F)\bigr),
\]
即在不改变 \(F\) 的传统公式核的前提下，附加正规型、不变量、证书、门禁和审计槽。

定义遗忘投影
\[
  \pi:\mathcal L_G\longrightarrow\mathcal L_{\mathrm{trad}}
\]
为删除所有 G-增强槽，只保留传统公式核：
\[
  \pi\bigl(F,\NF(F),\Inv(F),\Cert(F),\Gate(F),\Audit(F)\bigr)=F,
\]
且对传统项、关系、函数、联结词和量词均按恒等映射作用。
\end{definition}

\begin{lemma}[嵌入--投影恒等引理]
\label{lem:tex59-embedding-projection-identity}
对任意传统公式 \(F\in\mathcal L_{\mathrm{trad}}\)，有
\[
  \pi(\iota(F))=F.
\]
因此，\(\iota\) 对传统公式核是忠实嵌入。
\end{lemma}

\begin{Certificate}[证书：传统公式核保持]
\label{cert:tex59-embedding-projection-identity}
证书链为
\[
\begin{aligned}
  E_1 &: F\in\mathcal L_{\mathrm{trad}},\\
  E_2 &: \iota(F)=
         (F,\NF(F),\Inv(F),\Cert(F),\Gate(F),\Audit(F)),\\
  E_3 &: \pi(F,\NF(F),\Inv(F),\Cert(F),\Gate(F),\Audit(F))=F,\\
  E_4 &: E_1\wedge E_2\wedge E_3
         \Rightarrow
         \pi(\iota(F))=F.
\end{aligned}
\]
\end{Certificate}

\begin{proof}[证明（基于数学符号推演逻辑的谓词因果链）]
由定义~\ref{def:tex59-embedding-projection}，
\[
  \iota(F)=
  (F,\NF(F),\Inv(F),\Cert(F),\Gate(F),\Audit(F)).
\]
再由 \(\pi\) 的定义，
\[
  \pi(\iota(F))
  =
  \pi(F,\NF(F),\Inv(F),\Cert(F),\Gate(F),\Audit(F))
  =
  F.
\]
所以传统公式核在嵌入后可被完全投影回原公式，\(\iota\) 对传统公式核忠实。证毕。
\end{proof}

\begin{theorem}[公理降级：传统证明可嵌入 G-增强证明]
\label{thm:tex59-traditional-proof-embedding}
若
\[
  \mathcal T_{\mathrm{trad}}\vdash F,
  \qquad
  F\in\mathcal L_{\mathrm{trad}},
\]
则
\[
  \mathcal T_G\vdash \iota(F).
\]
\end{theorem}

\begin{Certificate}[数学归纳法证书：传统证明长度归纳]
\label{cert:tex59-traditional-proof-embedding}
令传统证明 \(\Pi_{\mathrm{trad}}\) 的长度为
\[
  n=\ell(\Pi_{\mathrm{trad}}).
\]
定义谓词
\[
  P(n):\quad
  \ell(\Pi_{\mathrm{trad}})\le n
  \Rightarrow
  \mathcal T_G\vdash\iota(F).
\]
归纳证书为
\[
\begin{aligned}
  A_0 &: P(1),\\
  A_1 &: P(n)\wedge
         \left[
         \Pi_{\mathrm{trad}}\text{ 的最后一步由传统允许规则 }\rho
         \text{ 得到}
         \right]
         \Rightarrow P(n+1),\\
  A_2 &: A_0\wedge A_1
         \Rightarrow
         \forall n,\ P(n).
\end{aligned}
\]
\end{Certificate}

\begin{proof}[证明（数学归纳法证书；基于数学符号推演逻辑的谓词因果链）]
基步：若传统证明长度为 \(1\)，则 \(F\) 是 \(\mathcal T_{\mathrm{trad}}\) 的公理、定义实例
或假设实例。由于
\[
  \mathcal T_{\mathrm{trad}}\subseteq \mathcal T_G,
\]
同一传统公式核在 \(\mathcal T_G\) 中仍可使用。再由 \(\mathcal T_{\mathrm{def}}\) 为 \(F\)
附加 \(\NF,\Inv,\Cert,\Gate,\Audit\) 槽，得到
\[
  \mathcal T_G\vdash\iota(F).
\]
故 \(P(1)\) 成立。

归纳步：假设长度不超过 \(n\) 的传统证明均可嵌入。设长度为 \(n+1\) 的传统证明最后一步由
传统规则 \(\rho\) 从前件
\[
  F_1,\ldots,F_m
\]
推出 \(F\)。由归纳假设，
\[
  \mathcal T_G\vdash\iota(F_j),
  \qquad j=1,\ldots,m.
\]
由于 \(\iota\) 对传统逻辑联结词、量词、等号和项构造保持恒等，传统规则 \(\rho\) 在
\(\mathcal T_G\) 中仍可作用于这些公式核。于是先推出传统核 \(F\)，再由定义附加证书槽，
得到
\[
  \mathcal T_G\vdash\iota(F).
\]
故 \(P(n)\Rightarrow P(n+1)\)。由数学归纳法，定理成立。证毕。
\end{proof}

\begin{lemma}[遗忘投影保持传统结论]
\label{lem:tex59-forgetful-projection-proof}
设
\[
  \Pi_G
\]
是 \(\mathcal T_G\) 中关于传统公式核 \(F\in\mathcal L_{\mathrm{trad}}\) 的证明，且其中所有
新增 G-步骤均属于定义展开、证书登记、门禁审计、吸收、剪枝或回写注释，并满足：
\[
  \pi(\mathcal T_{\mathrm{def}}\cup\mathcal T_{\mathrm{cert}})
  \subseteq
  \{\text{恒等式、定义回代或空步骤}\}.
\]
则
\[
  \mathcal T_{\mathrm{trad}}\vdash \pi(\Pi_G).
\]
特别地，若 \(\Pi_G\) 的终端结论为 \(\iota(F)\)，则
\[
  \mathcal T_{\mathrm{trad}}\vdash F.
\]
\end{lemma}

\begin{Certificate}[证书：G-步骤投影为传统步骤或空步骤]
\label{cert:tex59-forgetful-projection-proof}
证书链为
\[
\begin{aligned}
  R_1 &: \Pi_G=(\beta_0,\ldots,\beta_N),\\
  R_2 &: \forall k,\ \beta_k\text{ 是传统步骤或 G-新增审计步骤},\\
  R_3 &: \beta_k\text{ 是传统步骤}
         \Rightarrow
         \pi(\beta_k)\text{ 是同一传统步骤},\\
  R_4 &: \beta_k\text{ 是 G-新增审计步骤}
         \Rightarrow
         \pi(\beta_k)\text{ 是定义回代、恒等式或空步骤},\\
  R_5 &: R_1\wedge R_2\wedge R_3\wedge R_4
         \Rightarrow
         \pi(\Pi_G)\text{ 是传统证明}.
\end{aligned}
\]
\end{Certificate}

\begin{proof}[证明（基于数学符号推演逻辑的谓词因果链）]
把 \(\Pi_G\) 写成有限序列
\[
  \beta_0,\ldots,\beta_N.
\]
对每一步分类。若 \(\beta_k\) 是传统步骤，则因 \(\pi\) 对传统语言恒等，
\[
  \pi(\beta_k)
\]
仍是同一传统步骤。若 \(\beta_k\) 是 G-新增审计步骤，则由假设
\[
  \pi(\mathcal T_{\mathrm{def}}\cup\mathcal T_{\mathrm{cert}})
  \subseteq
  \{\text{恒等式、定义回代或空步骤}\},
\]
该步骤投影后不会生成新的传统结论，只会成为传统语言中的定义回代、恒等式或可删除空步骤。
删除所有空步骤后，剩余有限序列仍由传统公理、定义和规则组成。
因此
\[
  \mathcal T_{\mathrm{trad}}\vdash \pi(\Pi_G).
\]
若终端结论为 \(\iota(F)\)，由引理~\ref{lem:tex59-embedding-projection-identity}
得到
\[
  \pi(\iota(F))=F,
\]
故
\[
  \mathcal T_{\mathrm{trad}}\vdash F.
\]
证毕。
\end{proof}

\begin{proposition}[G-增强理论对传统语言的相对保守性]
\label{prop:tex59-relative-conservativity}
在定义~\ref{def:tex59-traditional-and-g-language} 与定义~\ref{def:tex59-embedding-projection}
的投影条件下，对任意传统公式 \(F\in\mathcal L_{\mathrm{trad}}\)，若
\[
  \mathcal T_G\vdash \iota(F)
\]
且证明中的 G-新增步骤满足引理~\ref{lem:tex59-forgetful-projection-proof} 的可投影条件，则
\[
  \mathcal T_{\mathrm{trad}}\vdash F.
\]
因此，\(\mathcal T_G\) 对传统语言是相对保守扩展。
\end{proposition}

\begin{Certificate}[证书：可投影证明到传统证明]
\label{cert:tex59-relative-conservativity}
证书链为
\[
\begin{aligned}
  K_1 &: \mathcal T_G\vdash\iota(F),\\
  K_2 &: \Pi_G\text{ 满足可投影条件},\\
  K_3 &: K_1\wedge K_2
         \Rightarrow
         \mathcal T_{\mathrm{trad}}\vdash \pi(\iota(F))
         \quad\text{（引理~\ref{lem:tex59-forgetful-projection-proof}）},\\
  K_4 &: \pi(\iota(F))=F
         \quad\text{（引理~\ref{lem:tex59-embedding-projection-identity}）},\\
  K_5 &: K_3\wedge K_4
         \Rightarrow
         \mathcal T_{\mathrm{trad}}\vdash F.
\end{aligned}
\]
\end{Certificate}

\begin{proof}[证明（基于数学符号推演逻辑的谓词因果链）]
由 \(K_1\)，存在 \(\mathcal T_G\) 证明 \(\Pi_G\) 以 \(\iota(F)\) 为终端结论。
由 \(K_2\) 与引理~\ref{lem:tex59-forgetful-projection-proof}，
\[
  \mathcal T_{\mathrm{trad}}\vdash \pi(\iota(F)).
\]
由引理~\ref{lem:tex59-embedding-projection-identity}，
\[
  \pi(\iota(F))=F.
\]
因此
\[
  \mathcal T_{\mathrm{trad}}\vdash F.
\]
证毕。
\end{proof}

\begin{theorem}[传统理论的忠实嵌入、保守扩展与结构增强总定理]
\label{thm:tex59-conservative-enhancement}
Gao 2025 的 \GFramework{}/\GAlgebra{} 基础域 \cite{Gao2025PureMathZenodo}
在本文形式化口径下不是与传统数学割裂的平行语言，
而是对传统证明语言的忠实嵌入、相对保守扩展与结构增强：
\[
  \boxed{
  \mathcal L_{\mathrm{trad}}
  \xrightarrow{\ \iota\ }
  \mathcal L_G
  \xrightarrow{\ \pi\ }
  \mathcal L_{\mathrm{trad}},
  \qquad
  \pi\circ\iota=\id.
  }
\]
并且：
\[
\begin{aligned}
  &\mathcal T_{\mathrm{trad}}\vdash F
    \Rightarrow
    \mathcal T_G\vdash\iota(F),\\
  &\mathcal T_G\vdash\iota(F)
    \wedge
    \mathsf{Projectable}(\Pi_G)
    \Rightarrow
    \mathcal T_{\mathrm{trad}}\vdash F,\\
  &\iota(F)
    =
    (F,\NF(F),\Inv(F),\Cert(F),\Gate(F),\Audit(F)).
\end{aligned}
\]
\end{theorem}

\begin{Certificate}[证书：忠实、保守与增强三元闭合]
\label{cert:tex59-conservative-enhancement}
证书链为
\[
\begin{aligned}
  M_1 &: \pi\circ\iota=\id
         \quad\text{（引理~\ref{lem:tex59-embedding-projection-identity}）},\\
  M_2 &: \mathcal T_{\mathrm{trad}}\vdash F
         \Rightarrow
         \mathcal T_G\vdash\iota(F)
         \quad\text{（定理~\ref{thm:tex59-traditional-proof-embedding}）},\\
  M_3 &: \mathcal T_G\vdash\iota(F)
         \wedge
         \mathsf{Projectable}(\Pi_G)
         \Rightarrow
         \mathcal T_{\mathrm{trad}}\vdash F
         \quad\text{（命题~\ref{prop:tex59-relative-conservativity}）},\\
  M_4 &: \iota(F)
         =
         (F,\NF(F),\Inv(F),\Cert(F),\Gate(F),\Audit(F)),\\
  M_5 &: M_1\wedge M_2\wedge M_3\wedge M_4
         \Rightarrow
         \text{忠实嵌入、相对保守扩展与结构增强同时成立}.
\end{aligned}
\]
\end{Certificate}

\begin{proof}[证明（基于数学符号推演逻辑的谓词因果链）]
由引理~\ref{lem:tex59-embedding-projection-identity}，
\[
  \pi(\iota(F))=F
\]
对任意传统公式成立，所以传统公式核没有被 G-增强语言改写或丢失。
由定理~\ref{thm:tex59-traditional-proof-embedding}，任何传统证明都可以嵌入为带证书槽的
G-增强证明。由命题~\ref{prop:tex59-relative-conservativity}，只要 G-新增步骤满足可投影条件，
任何以传统公式核为终端的 G-证明都可投影回传统证明。因此新增结构不强行制造传统语言中
原本不可投影的定理。

另一方面，\(\iota(F)\) 不只是 \(F\) 的复制，而是
\[
  (F,\NF(F),\Inv(F),\Cert(F),\Gate(F),\Audit(F)).
\]
这些槽位提供正规型、同调不变量、证书门禁和审计链。因此，\(\mathcal L_G\) 在保持传统公式核
可投影的同时，增加了证明对象的结构化、证书化和可审计性。总定理成立。证毕。
\end{proof}

\begin{corollary}[传统难题在 G-语言中的非割裂性]
\label{cor:tex59-non-separated-traditional-problems}
本文后续把 Hilbert 问题、千禧年问题、计算边界、信息论、统计物理、复杂性与认知哲学障碍
写入 \(\mathcal L_G\) 时，不是在传统理论之外另造无关命题，而是在
\[
  \pi\circ\iota=\id
\]
的嵌入--投影约束下，把传统对象提升为带正规型、不变量、证书、门禁和审计槽的增强对象。
\end{corollary}

\begin{Certificate}[证书：非割裂性登记]
\label{cert:tex59-non-separated-traditional-problems}
证书链为
\[
\begin{aligned}
  N_1 &: F\in\mathcal L_{\mathrm{trad}},\\
  N_2 &: \iota(F)\in\mathcal L_G,\\
  N_3 &: \pi(\iota(F))=F,\\
  N_4 &: \iota(F)\text{ 携带 }\NF,\Inv,\Cert,\Gate,\Audit,\\
  N_5 &: N_1\wedge N_2\wedge N_3\wedge N_4
         \Rightarrow
         \text{传统对象被结构增强而非割裂替换}.
\end{aligned}
\]
\end{Certificate}

\begin{proof}[证明（基于数学符号推演逻辑的谓词因果链）]
由定义~\ref{def:tex59-embedding-projection}，传统对象 \(F\) 经 \(\iota\) 进入 \(\mathcal L_G\)；
由引理~\ref{lem:tex59-embedding-projection-identity}，再经 \(\pi\) 返回 \(F\)。
因此传统公式核保持不变。新增的 \(\NF,\Inv,\Cert,\Gate,\Audit\) 只增强其结构与审计接口。
所以后续对传统难题的 G-化表达是忠实提升，而不是与传统理论割裂的平行替换。证毕。
\end{proof}

\subsection{障碍修复同伦等价与保结构同调群}
\label{subsec:tex59-obstruction-repair-homotopy-equivalence}

\begin{definition}[障碍修复型同伦增强]
\label{def:tex59-obstruction-repair-enhancement}
令 \(X\) 为传统数学结构，包含底集合、运算、关系、拓扑、光滑结构、链复形或范畴结构中的若干项。
定义 \(X\) 的障碍修复型 G-增强为
\[
  X_G
  \DefEq
  \bigl(
  X,\ C_\bullet^{\Obs}(X),\partial_\bullet^{\Obs},
  \rho_X,\iota_X,\pi_X
  \bigr),
\]
其中 \(C_\bullet^{\Obs}(X)\) 是障碍链群族，边界算子满足
\[
  \partial_{n-1}^{\Obs}\circ\partial_n^{\Obs}=0,
  \qquad n\ge1,
\]
\(\rho_X\) 是障碍修复同伦算子，
\[
  \iota_X:X\longrightarrow X_G
  \quad\text{与}\quad
  \pi_X:X_G\longrightarrow X
\]
分别表示增强嵌入与遗忘投影，并满足
\[
  \pi_X\circ\iota_X=\id_X.
\]
若还存在链同伦 \(H_X\) 使
\[
  \iota_X\circ\pi_X-\id_{X_G}
  =
  \partial^{\Obs}H_X+H_X\partial^{\Obs},
\]
则称 \(X_G\) 是 \(X\) 的障碍修复型同伦增强，并记为
\[
  X_G\simeq_{\Obs,\rho} X.
\]
\end{definition}

\begin{definition}[保结构障碍修复态射]
\label{def:tex59-structure-preserving-repair-morphism}
设 \(X_G,Y_G\) 是两个障碍修复型增强结构。态射
\[
  F:X_G\longrightarrow Y_G
\]
称为保结构障碍修复态射，若存在传统结构态射
\[
  f:X\longrightarrow Y
\]
使以下条件同时成立：
\[
  \pi_Y\circ F=f\circ\pi_X,
  \qquad
  F\circ\iota_X=\iota_Y\circ f,
\]
\[
  F\circ\partial_X^{\Obs}
  =
  \partial_Y^{\Obs}\circ F,
  \qquad
  F\rho_X\simeq_{\Obs}\rho_YF.
\]
此外，对任意 \(r\)-元运算 \(\omega\) 与关系 \(R\)，要求
\[
  F\bigl(\omega_X(x_1,\ldots,x_r)\bigr)
  =
  \omega_Y\bigl(Fx_1,\ldots,Fx_r\bigr),
\]
并且
\[
  R_X(x_1,\ldots,x_r)
  \Rightarrow
  R_Y(Fx_1,\ldots,Fx_r).
\]
\end{definition}

\begin{definition}[保结构障碍修复同调群]
\label{def:tex59-structure-preserving-obstruction-homology}
定义保结构障碍闭链群为
\[
  Z_n^{\Obs,\mathrm{str}}(X)
  \DefEq
  \left\{
  c\in C_n^{\Obs}(X):
  \partial_n^{\Obs}c=0
  \ \wedge\
  \mathrm{Str}_X(c)=1
  \right\},
\]
其中 \(\mathrm{Str}_X(c)=1\) 表示 \(c\) 保持 \(X\) 的运算、关系、拓扑或链复形结构约束。
定义保结构障碍边界群为
\[
  B_n^{\Obs,\mathrm{str}}(X)
  \DefEq
  \left\{
  \partial_{n+1}^{\Obs}b:
  b\in C_{n+1}^{\Obs}(X),
  \ \mathrm{Str}_X(\partial_{n+1}^{\Obs}b)=1
  \right\}.
\]
保结构障碍修复同调群定义为
\[
  H_n^{\Obs,\mathrm{str}}(X)
  \DefEq
  Z_n^{\Obs,\mathrm{str}}(X)
  \big/
  B_n^{\Obs,\mathrm{str}}(X).
\]
若 \(c\in Z_n^{\Obs,\mathrm{str}}(X)\) 且 \(\Gate(c)=1\)，则称 \(c\) 是
\[
  H_n^{\Obs,\mathrm{str}}(X)
\]
中的可签发证书代表元。
\end{definition}

\begin{theorem}[公理降级：有限障碍修复同伦复合归纳定理]
\label{thm:tex59-repair-homotopy-composition-induction}
给定有限序列
\[
  X_{0,G}\xrightarrow{F_0}X_{1,G}
  \xrightarrow{F_1}\cdots
  \xrightarrow{F_{m-1}}X_{m,G},
\]
若每个 \(F_k\) 都是保结构障碍修复态射，且每一步都给出障碍修复同伦等价
\[
  X_{k,G}\simeq_{\Obs,\rho}X_{k+1,G},
  \qquad 0\le k<m,
\]
则复合态射
\[
  F_{m:0}\DefEq F_{m-1}\circ\cdots\circ F_0
\]
仍是保结构障碍修复态射，并给出
\[
  X_{0,G}\simeq_{\Obs,\rho}X_{m,G}.
\]
\end{theorem}

\begin{Certificate}[数学归纳法证书：保结构修复同伦复合闭合]
\label{cert:tex59-repair-homotopy-composition-induction}
定义谓词
\[
  P(m):\quad
  F_{m:0}\text{ 是保结构障碍修复态射，且 }
  X_{0,G}\simeq_{\Obs,\rho}X_{m,G}.
\]
归纳证书为
\[
\begin{aligned}
  A_0 &: P(1),\\
  A_1 &: P(m)\wedge F_m\text{ 保结构且为障碍修复同伦等价}
         \Rightarrow P(m+1),\\
  A_2 &: A_0\wedge A_1
         \Rightarrow
         \forall m\ge1,\ P(m).
\end{aligned}
\]
\end{Certificate}

\begin{proof}[证明（数学归纳法证书；基于数学符号推演逻辑的谓词因果链）]
基步 \(m=1\)。由假设，\(F_0:X_{0,G}\to X_{1,G}\) 已经是保结构障碍修复态射，
并且
\[
  X_{0,G}\simeq_{\Obs,\rho}X_{1,G}.
\]
因此 \(P(1)\) 成立。

归纳步。设 \(P(m)\) 成立，即
\[
  F_{m:0}=F_{m-1}\circ\cdots\circ F_0
\]
是保结构障碍修复态射，且
\[
  X_{0,G}\simeq_{\Obs,\rho}X_{m,G}.
\]
再设 \(F_m:X_{m,G}\to X_{m+1,G}\) 是保结构障碍修复态射，并给出
\[
  X_{m,G}\simeq_{\Obs,\rho}X_{m+1,G}.
\]
保结构性对复合封闭，因为对任意运算 \(\omega\)，有
\[
\begin{aligned}
  (F_mF_{m:0})\omega_X(x_1,\ldots,x_r)
  &=
  F_m\omega_{X_m}(F_{m:0}x_1,\ldots,F_{m:0}x_r)\\
  &=
  \omega_{X_{m+1}}
  (F_mF_{m:0}x_1,\ldots,F_mF_{m:0}x_r).
\end{aligned}
\]
关系保持同理由蕴含传递得到。链映射性也对复合封闭：
\[
  (F_mF_{m:0})\partial_X^{\Obs}
  =
  F_m\partial_{X_m}^{\Obs}F_{m:0}
  =
  \partial_{X_{m+1}}^{\Obs}F_mF_{m:0}.
\]
同伦等价对复合封闭，因此
\[
  X_{0,G}\simeq_{\Obs,\rho}X_{m,G}
  \quad\text{且}\quad
  X_{m,G}\simeq_{\Obs,\rho}X_{m+1,G}
\]
推出
\[
  X_{0,G}\simeq_{\Obs,\rho}X_{m+1,G}.
\]
所以 \(P(m)\Rightarrow P(m+1)\)。由数学归纳法，定理成立。证毕。
\end{proof}

\begin{lemma}[遗忘投影是障碍修复下的保结构态射]
\label{lem:tex59-projection-structure-preserving}
对任意障碍修复型同伦增强 \(X_G\)，遗忘投影
\[
  \pi_X:X_G\to X
\]
在障碍修复意义下是保结构态射，并满足
\[
  \pi_X\circ\iota_X=\id_X.
\]
\end{lemma}

\begin{Certificate}[证书：投影保结构]
\label{cert:tex59-projection-structure-preserving}
证书链为
\[
\begin{aligned}
  P_1 &: X_G=(X,C_\bullet^{\Obs},\partial^{\Obs},\rho_X,\iota_X,\pi_X),\\
  P_2 &: \pi_X\circ\iota_X=\id_X,\\
  P_3 &: \pi_X\text{ 删除证书槽但保留传统公式核与传统结构核},\\
  P_4 &: P_3\Rightarrow
         \pi_X(\omega_{X_G}(\vec x))
         =
         \omega_X(\pi_X\vec x),\\
  P_5 &: P_3\Rightarrow
         R_{X_G}(\vec x)\Rightarrow R_X(\pi_X\vec x),\\
  P_6 &: P_2\wedge P_4\wedge P_5
         \Rightarrow
         \pi_X\text{ 是保结构障碍修复态射}.
\end{aligned}
\]
\end{Certificate}

\begin{proof}[证明（基于数学符号推演逻辑的谓词因果链）]
由定义~\ref{def:tex59-obstruction-repair-enhancement}，
\[
  \pi_X\circ\iota_X=\id_X.
\]
由于 \(\pi_X\) 的作用是遗忘 \(\NF,\Inv,\Cert,\Gate,\Audit\) 等增强槽，并保留传统结构核，
任意传统运算 \(\omega\) 满足
\[
  \pi_X\bigl(\omega_{X_G}(x_1,\ldots,x_r)\bigr)
  =
  \omega_X(\pi_Xx_1,\ldots,\pi_Xx_r).
\]
任意传统关系 \(R\) 也满足
\[
  R_{X_G}(x_1,\ldots,x_r)
  \Rightarrow
  R_X(\pi_Xx_1,\ldots,\pi_Xx_r).
\]
链层面上，\(\pi_X\) 把已修复障碍链的传统边界解释投影回 \(X\) 中的结构约束，
不会引入新的传统结构断裂。因此 \(\pi_X\) 是障碍修复意义下的保结构态射。证毕。
\end{proof}

\begin{theorem}[G-增强结构与传统结构的障碍修复同伦等价定理]
\label{thm:tex59-g-enhancement-obstruction-homotopy-equivalence}
设 \(X_G\) 是传统结构 \(X\) 的障碍修复型同伦增强。若
\[
  \pi_X\circ\iota_X=\id_X
\]
且存在修复同伦 \(H_X\) 使
\[
  \iota_X\circ\pi_X-\id_{X_G}
  =
  \partial^{\Obs}H_X+H_X\partial^{\Obs},
\]
则
\[
  X_G\simeq_{\Obs,\rho}X.
\]
因此，G-增强结构不是与传统结构割裂的对象，而是传统结构在同调障碍修复后的同伦等价增强。
\end{theorem}

\begin{Certificate}[证书：嵌入--投影--修复同伦闭合]
\label{cert:tex59-g-enhancement-obstruction-homotopy-equivalence}
证书链为
\[
\begin{aligned}
  H_1 &: \iota_X:X\to X_G,\qquad \pi_X:X_G\to X,\\
  H_2 &: \pi_X\circ\iota_X=\id_X,\\
  H_3 &: \iota_X\circ\pi_X-\id_{X_G}
         =
         \partial^{\Obs}H_X+H_X\partial^{\Obs},\\
  H_4 &: H_3\Rightarrow
         \iota_X\circ\pi_X\simeq_{\Obs,\rho}\id_{X_G},\\
  H_5 &: H_2\wedge H_4
         \Rightarrow
         X_G\simeq_{\Obs,\rho}X.
\end{aligned}
\]
\end{Certificate}

\begin{proof}[证明（基于数学符号推演逻辑的谓词因果链）]
由 \(H_2\)，传统结构从 \(X\) 嵌入到 \(X_G\) 再投影回 \(X\) 后完全不变：
\[
  \pi_X\iota_X=\id_X.
\]
由 \(H_3\)，反向复合与 \(X_G\) 上恒等态射的差不是任意误差，而是障碍链复形中的边界型差异：
\[
  \iota_X\pi_X-\id_{X_G}
  =
  \partial^{\Obs}H_X+H_X\partial^{\Obs}.
\]
这正是链同伦定义。因此
\[
  \iota_X\pi_X\simeq_{\Obs,\rho}\id_{X_G}.
\]
两式合并说明 \(\iota_X\) 与 \(\pi_X\) 互为障碍修复意义下的同伦逆。
所以
\[
  X_G\simeq_{\Obs,\rho}X.
\]
定理成立。证毕。
\end{proof}

\begin{proposition}[G-证书作为保结构障碍修复同调代表元]
\label{prop:tex59-certificate-as-obstruction-homology-class}
设 \(c\) 是 \(X_G\) 中的证书对象。若
\[
  \partial_n^{\Obs}c=0,
  \qquad
  \mathrm{Str}_X(c)=1,
  \qquad
  \Gate(c)=1,
\]
则 \(c\) 确定一个同调类
\[
  [c]\in H_n^{\Obs,\mathrm{str}}(X).
\]
若另一证书 \(c'\) 满足
\[
  c'-c=\partial_{n+1}^{\Obs}b
\]
且边界 \(\partial_{n+1}^{\Obs}b\) 保结构，则
\[
  [c']=[c].
\]
\end{proposition}

\begin{Certificate}[证书：闭链、保结构、边界商]
\label{cert:tex59-certificate-as-obstruction-homology-class}
证书链为
\[
\begin{aligned}
  C_1 &: \partial_n^{\Obs}c=0,\\
  C_2 &: \mathrm{Str}_X(c)=1,\\
  C_3 &: C_1\wedge C_2
         \Rightarrow c\in Z_n^{\Obs,\mathrm{str}}(X),\\
  C_4 &: \Gate(c)=1
         \Rightarrow c\text{ 是可签发证书代表元},\\
  C_5 &: c'-c=\partial_{n+1}^{\Obs}b
         \wedge
         \mathrm{Str}_X(\partial_{n+1}^{\Obs}b)=1,\\
  C_6 &: C_5
         \Rightarrow
         c'-c\in B_n^{\Obs,\mathrm{str}}(X),\\
  C_7 &: C_3\wedge C_6
         \Rightarrow
         [c']=[c]\in H_n^{\Obs,\mathrm{str}}(X).
\end{aligned}
\]
\end{Certificate}

\begin{proof}[证明（基于数学符号推演逻辑的谓词因果链）]
由
\[
  \partial_n^{\Obs}c=0
\]
可知 \(c\) 是障碍闭链。由
\[
  \mathrm{Str}_X(c)=1
\]
可知 \(c\) 保持传统结构约束。按定义~\ref{def:tex59-structure-preserving-obstruction-homology}，
\[
  c\in Z_n^{\Obs,\mathrm{str}}(X).
\]
因此商映射给出同调类
\[
  [c]\in
  Z_n^{\Obs,\mathrm{str}}(X)\big/B_n^{\Obs,\mathrm{str}}(X)
  =
  H_n^{\Obs,\mathrm{str}}(X).
\]
门禁条件 \(\Gate(c)=1\) 说明该代表元可签发为证书。

若
\[
  c'-c=\partial_{n+1}^{\Obs}b
\]
且该边界保结构，则
\[
  c'-c\in B_n^{\Obs,\mathrm{str}}(X).
\]
商群中相差边界的两个闭链表示同一个同调类，所以
\[
  [c']=[c].
\]
命题成立。证毕。
\end{proof}

\begin{corollary}[G-证明的保结构同调解释推论]
\label{cor:tex59-g-proof-structure-preserving-homology}
在定理~\ref{thm:tex59-g-enhancement-obstruction-homotopy-equivalence}
与命题~\ref{prop:tex59-certificate-as-obstruction-homology-class} 的条件下，
G-证明可解释为传统数学结构上保结构障碍修复同调群中的证书化代表元：
\[
  \text{G-证书}
  \quad\leadsto\quad
  [c]\in H_n^{\Obs,\mathrm{str}}(X).
\]
因此，G-Framework/G-Algebra 对传统数学结构的增强，可严格理解为：
\[
  \boxed{
  \text{传统结构}
  \xrightarrow{\ \iota_X\ }
  \text{障碍修复型 G-增强结构}
  \xrightarrow{\ \pi_X\ }
  \text{传统结构}
  }
\]
且该增强在障碍修复同伦意义下保结构等价。
\end{corollary}

\begin{Certificate}[证书：G-增强证明到保结构同调群]
\label{cert:tex59-g-proof-structure-preserving-homology}
证书链为
\[
\begin{aligned}
  G_1 &: X_G\simeq_{\Obs,\rho}X,\\
  G_2 &: \pi_X:X_G\to X\text{ 是保结构态射},\\
  G_3 &: c\in Z_n^{\Obs,\mathrm{str}}(X),\quad \Gate(c)=1,\\
  G_4 &: G_3\Rightarrow [c]\in H_n^{\Obs,\mathrm{str}}(X),\\
  G_5 &: G_1\wedge G_2\wedge G_4
         \Rightarrow
         \text{G-证明是传统结构上保结构障碍修复同调类的证书化实现}.
\end{aligned}
\]
\end{Certificate}

\begin{proof}[证明（基于数学符号推演逻辑的谓词因果链）]
由定理~\ref{thm:tex59-g-enhancement-obstruction-homotopy-equivalence}，
\[
  X_G\simeq_{\Obs,\rho}X.
\]
由引理~\ref{lem:tex59-projection-structure-preserving}，投影
\[
  \pi_X:X_G\to X
\]
在障碍修复意义下保结构。由命题~\ref{prop:tex59-certificate-as-obstruction-homology-class}，
可签发证书 \(c\) 给出
\[
  [c]\in H_n^{\Obs,\mathrm{str}}(X).
\]
因此，G-证明不是脱离传统结构的外部符号，而是传统结构经障碍修复同伦增强后，在
保结构障碍修复同调群中的证书化代表元。证毕。
\end{proof}

\begin{remark}[与 Zenodo 基础语言的关系]
上述形式化把 Gao 2025 的 \GFramework{}/\GAlgebra{} 基础语言 \cite{Gao2025PureMathZenodo}
解释为传统数学结构的障碍修复型同伦增强：传统结构核由
\(\pi_X\circ\iota_X=\id_X\) 保持，新增证书槽由
\[
  H_n^{\Obs,\mathrm{str}}(X)
\]
承载。因此，更精确的审计口径是：G-Framework/G-Algebra 与传统理论并不割裂；
它通过保结构投影、障碍修复同伦和同调类证书，把传统证明对象提升为可修复、可审计、
可投影回传统结构的增强对象。
\end{remark}

\subsection{可替换术语消去：抽象验证结构与纯数学接口}
\label{subsec:tex59-implementation-name-abstraction}

\begin{definition}[抽象验证同调结构]
\label{def:tex59-abstract-vsys}
本文中的
\[
  \Vsys
\]
只表示一个纯数学的抽象验证同调结构，不表示任何外部实现实体。定义
\[
  \mathfrak V
  \DefEq
  (X,N,\bot,\phi,d,I,\mathcal P,\mathcal C,BQ,T,A).
\]
各分量含义如下：
\[
\begin{array}{ll}
  X: & \text{待登记对象集合},\\
  N: & \text{正规型集合},\\
  \bot\in N: & \text{吸收失败元},\\
  \phi:X\to N: & \text{正规化映射},\\
  d:N\to N: & \text{语义微分算子},\\
  I:N\to \Inv(N): & \text{语义积分或不变量读取映射},\\
  \mathcal P: & \text{有限谓词族},\\
  \mathcal C: & \text{有限证书族},\\
  BQ\subseteq N: & \text{可修复残差队列},\\
  T:N\to\{0,1\}: & \text{终端门禁},\\
  A=T^{-1}(1): & \text{接受集合}.
\end{array}
\]
并要求
\[
  T(\bot)=0,
  \qquad
  \bot\notin A.
\]
正文中的符号
\[
  \NF_{\Vsys},\quad
  \Phi_{\Vsys},\quad
  d_{\Vsys},\quad
  \int_{\Vsys},\quad
  \BackwriteQueue(\Vsys),\quad
  \Accepted(\Vsys)
\]
分别只是
\[
  N,\quad \phi,\quad d,\quad I,\quad BQ,\quad A
\]
的排版缩写。
\end{definition}

\begin{definition}[可替换术语的标签化解释]
\label{def:tex59-label-abstraction}
令 \(\Lambda_{59}\) 为一个有限标签集。凡正文中为了区分应用扇区、对象族或约束族而出现的
可替换术语，一律解释为某个抽象标签
\[
  \lambda\in\Lambda_{59}.
\]
标签本身不携带数学前提。真正进入证明的是与标签配套的四元组
\[
  \mathsf D_{\lambda}
  \DefEq
  (X_\lambda,\mathcal P_\lambda,\mathcal C_\lambda,T_\lambda),
\]
其中
\[
  X_\lambda,\qquad
  \mathcal P_\lambda,\qquad
  \mathcal C_\lambda,\qquad
  T_\lambda.
\]
分别是对象集合、有限谓词族、有限证书族和终端门禁。
因此，任意外观名称均可替换为同类型标签而不改变证明。
\end{definition}

\begin{theorem}[公理降级：抽象验证同调结构存在定理]
\label{thm:tex59-abstract-vsys-existence}
给定任意集合 \(X\) 与有限谓词族
\[
  \mathcal P=\{p_1,\ldots,p_m\},
  \qquad
  p_i:X\to\{0,1\},
\]
存在一个抽象验证同调结构
\[
  \mathfrak V
  =
  (X,N,\bot,\phi,d,I,\mathcal P,\mathcal C,BQ,T,A)
\]
满足定义~\ref{def:tex59-abstract-vsys}。
\end{theorem}

\begin{Certificate}[数学归纳法证书：谓词族到终端门禁]
\label{cert:tex59-abstract-vsys-existence}
定义部分门禁
\[
  T_0(x)=1,
  \qquad
  T_k(x)=\bigwedge_{i=1}^{k}p_i(x),
  \quad 1\le k\le m.
\]
定义谓词
\[
  P(k):\quad
  T_k:X\to\{0,1\}\text{ 是由 }p_1,\ldots,p_k\text{ 有限合取得到的门禁}.
\]
归纳证书为
\[
\begin{aligned}
  R_0 &: P(0),\\
  R_1 &: P(k)\wedge p_{k+1}:X\to\{0,1\}
         \Rightarrow P(k+1),\\
  R_2 &: P(m)\Rightarrow T=T_m\text{ 存在}.
\end{aligned}
\]
\end{Certificate}

\begin{proof}[证明（数学归纳法证书；基于数学符号推演逻辑的谓词因果链）]
取一个不属于 \(X\) 的新符号 \(\bot\)，并令
\[
  N\DefEq X\sqcup\{\bot\}.
\]
令 \(\phi:X\to N\) 为自然嵌入。令 \(d:N\to N\) 为恒等映射；令
\[
  I:N\to \Inv(N)
\]
为把正规型送入其等价类或自身不变量的映射。取证书族
\[
  \mathcal C=\{c_i:p_i(x)=1\}_{i=1}^{m}.
\]

下面构造门禁。基步 \(k=0\) 时，\(T_0(x)=1\) 是空合取，故 \(P(0)\) 成立。
归纳步中，若 \(T_k\) 已由前 \(k\) 个谓词有限合取得到，则定义
\[
  T_{k+1}(x)=T_k(x)\wedge p_{k+1}(x).
\]
由于 \(\{0,1\}\) 对合取封闭，\(T_{k+1}\) 仍为 \(X\to\{0,1\}\) 的映射，故
\[
  P(k)\Rightarrow P(k+1).
\]
由数学归纳法，\(T_m\) 存在。扩张到 \(N\) 上：
\[
  T(x)=T_m(x)\quad(x\in X),
  \qquad
  T(\bot)=0.
\]
定义
\[
  A\DefEq T^{-1}(1),
  \qquad
  BQ\DefEq T^{-1}(0)\setminus\{\bot\}.
\]
于是 \(T(\bot)=0\) 且 \(\bot\notin A\)。所有分量均由集合、映射、有限合取和子集构造给出，
所以 \(\mathfrak V\) 是定义~\ref{def:tex59-abstract-vsys} 所要求的纯数学对象。证毕。
\end{proof}

\begin{lemma}[术语接口保守消去引理]
\label{lem:tex59-vsys-conservative-elimination}
令 \(\tau_{\mathfrak V}\) 为如下替换：
\[
\begin{array}{lll}
  \Vsys &\mapsto& \mathfrak V,\\
  \NF_{\Vsys} &\mapsto& N,\\
  \Phi_{\Vsys} &\mapsto& \phi,\\
  d_{\Vsys} &\mapsto& d,\\
  \int_{\Vsys} &\mapsto& I,\\
  \BackwriteQueue(\Vsys) &\mapsto& BQ,\\
  \Accepted(\Vsys) &\mapsto& A.
\end{array}
\]
若公式 \(\Psi\) 只通过上表中的接口使用 \(\Vsys\)，则
\[
  \Psi
  \Longleftrightarrow
  \tau_{\mathfrak V}(\Psi).
\]
\end{lemma}

\begin{Certificate}[证书：逐符号替换保持谓词结构]
\label{cert:tex59-vsys-conservative-elimination}
证书链为
\[
\begin{aligned}
  E_1 &: \Vsys\text{ 只在指定接口中出现},\\
  E_2 &: \tau_{\mathfrak V}\text{ 对每个接口符号给出同类型对象},\\
  E_3 &: \text{谓词联结词 }(\wedge,\vee,\Rightarrow,\neg)
         \text{ 在替换下不改变},\\
  E_4 &: E_1\wedge E_2\wedge E_3
         \Rightarrow
         \Psi\Longleftrightarrow\tau_{\mathfrak V}(\Psi).
\end{aligned}
\]
\end{Certificate}

\begin{proof}[证明（基于数学符号推演逻辑的谓词因果链）]
由定义~\ref{def:tex59-abstract-vsys}，
\[
  \NF_{\Vsys}=N,\quad
  \Phi_{\Vsys}=\phi,\quad
  d_{\Vsys}=d,\quad
  \int_{\Vsys}=I,\quad
  \BackwriteQueue(\Vsys)=BQ,\quad
  \Accepted(\Vsys)=A.
\]
因此，替换 \(\tau_{\mathfrak V}\) 只把上述缩写还原为定义项。
该还原不改变变量、量词、谓词联结词或证明边。
故任意只通过这些接口写出的公式 \(\Psi\)，均与其还原式
\(\tau_{\mathfrak V}(\Psi)\) 等价。证毕。
\end{proof}

\begin{proposition}[标签名称不参与证明前件命题]
\label{prop:tex59-label-nondependence}
设 \(\lambda,\lambda'\in\Lambda_{59}\) 是同类型标签。若二者对应的纯数学数据同构：
\[
  (X_\lambda,\mathcal P_\lambda,\mathcal C_\lambda,T_\lambda)
  \cong
  (X_{\lambda'},\mathcal P_{\lambda'},\mathcal C_{\lambda'},T_{\lambda'}),
\]
则任何只读取对象、谓词、证书和门禁的证明结论在 \(\lambda\) 与 \(\lambda'\) 下相同。
\end{proposition}

\begin{Certificate}[证书：标签同构到证明不变]
\label{cert:tex59-label-nondependence}
证书链为
\[
\begin{aligned}
  L_1 &: (X_\lambda,\mathcal P_\lambda,\mathcal C_\lambda,T_\lambda)
        \cong
        (X_{\lambda'},\mathcal P_{\lambda'},\mathcal C_{\lambda'},T_{\lambda'}),\\
  L_2 &: \text{证明只读取 }X,\mathcal P,\mathcal C,T,\\
  L_3 &: L_1\wedge L_2
        \Rightarrow
        \Dep(\Theta_\lambda)=\Dep(\Theta_{\lambda'}),\\
  L_4 &: L_3\Rightarrow
        \Theta_\lambda\Longleftrightarrow\Theta_{\lambda'}.
\end{aligned}
\]
\end{Certificate}

\begin{proof}[证明（基于数学符号推演逻辑的谓词因果链）]
标签 \(\lambda\) 的名称不进入定义~\ref{def:tex59-label-abstraction} 的证明前件；
进入证明前件的是 \(X_\lambda,\mathcal P_\lambda,\mathcal C_\lambda,T_\lambda\)。
若这些数据与 \(\lambda'\) 的对应数据同构，则所有谓词判定、证书读取和门禁值在同构下保持。
因此证明依赖集相同，结论也相同。证毕。
\end{proof}

\begin{corollary}[抽象验证结构与可替换术语的纯数学化推论]
\label{cor:tex59-vsys-pure-math}
本文后续出现的 \(\Vsys\)、\(\Phi_{\Vsys}\)、\(\NF_{\Vsys}\) 以及所有扇区名称、
对象族名称和约束族名称，均按定义~\ref{def:tex59-abstract-vsys} 与
定义~\ref{def:tex59-label-abstraction} 解释。它们不是独立的无源概念，而是纯数学结构或
有限标签的排版记号。
\end{corollary}

\begin{Certificate}[证书：无源概念消去]
\label{cert:tex59-vsys-pure-math}
证书链为
\[
\begin{aligned}
  C_1 &: \Vsys\text{ 由 }\mathfrak V\text{ 定义},\\
  C_2 &: \Phi_{\Vsys},\NF_{\Vsys},d_{\Vsys},\int_{\Vsys}
         \text{ 均由 }\mathfrak V\text{ 的分量定义},\\
  C_3 &: \text{可替换术语均属于有限标签集 }\Lambda_{59},\\
  C_4 &: C_1\wedge C_2\wedge C_3
         \Rightarrow
         \text{后续证明不含未定义应用前提}.
\end{aligned}
\]
\end{Certificate}

\begin{proof}[证明（基于数学符号推演逻辑的谓词因果链）]
由定义~\ref{def:tex59-abstract-vsys}，\(\Vsys\) 被还原为抽象验证同调结构 \(\mathfrak V\)。
由引理~\ref{lem:tex59-vsys-conservative-elimination}，所有 \(\Vsys\) 接口公式可保守替换为
\(N,\phi,d,I,BQ,A\) 等纯数学对象。由命题~\ref{prop:tex59-label-nondependence}，
任意可替换术语只是标签，不承担证明前件。因此后续证明不含未定义外部实现概念。证毕。
\end{proof}

\begin{theorem}[可替换术语消去总定理]
\label{thm:tex59-terminology-elimination-total}
设 \(\Theta\) 是本文任意形式命题。若 \(\Theta\) 的证明只读取
\[
  \mathfrak V,\qquad
  \mathsf D_{\lambda}
  =
  (X_\lambda,\mathcal P_\lambda,\mathcal C_\lambda,T_\lambda),
  \qquad
  \lambda\in\Lambda_{59},
\]
以及 Gao 2025 年 Zenodo 基础文献 \cite{Gao2025PureMathZenodo} 中的抽象代数拓扑背景，
则 \(\Theta\) 可改写为不含任何外部实体名称的纯数学命题
\[
  \Theta^{\mathrm{abs}}.
\]
并且
\[
  \Theta\Longleftrightarrow \Theta^{\mathrm{abs}}.
\]
\end{theorem}

\begin{Certificate}[证书：结构接口与标签接口的合成消去]
\label{cert:tex59-terminology-elimination-total}
证书链为
\[
\begin{aligned}
  U_1 &: \text{结构接口由 }\mathfrak V=(X,N,\bot,\phi,d,I,\mathcal P,\mathcal C,BQ,T,A)\text{ 给出},\\
  U_2 &: \text{标签接口由 }\mathsf D_{\lambda}=(X_\lambda,\mathcal P_\lambda,\mathcal C_\lambda,T_\lambda)\text{ 给出},\\
  U_3 &: \text{引理~\ref{lem:tex59-vsys-conservative-elimination}}
         \Rightarrow \text{结构接口可保守还原},\\
  U_4 &: \text{命题~\ref{prop:tex59-label-nondependence}}
         \Rightarrow \text{标签名称可保守替换},\\
  U_5 &: U_3\wedge U_4
         \Rightarrow
         \Theta\Longleftrightarrow\Theta^{\mathrm{abs}}.
\end{aligned}
\]
\end{Certificate}

\begin{proof}[证明（基于数学符号推演逻辑的谓词因果链）]
由 \(U_1\)，所有结构接口都可还原为 \(\mathfrak V\) 的分量。
由引理~\ref{lem:tex59-vsys-conservative-elimination}，这种还原不改变公式的变量、
量词、谓词联结词和证明边。由 \(U_2\)，所有可替换术语都只对应某个
\(\mathsf D_\lambda\)。由命题~\ref{prop:tex59-label-nondependence}，只要证明读取的是
\(X_\lambda,\mathcal P_\lambda,\mathcal C_\lambda,T_\lambda\)，标签外观不参与证明前件。
因此把 \(\Theta\) 中的外观术语统一替换为 \(\mathfrak V\) 与 \(\mathsf D_\lambda\)
之后得到的 \(\Theta^{\mathrm{abs}}\) 与原命题等价。证毕。
\end{proof}

\begin{remark}[本节对全文的约束]
从本节开始，凡正文中出现的领域性名称均只表示上述纯数学对象或有限标签。
若某段自然语言看似提到外部实现对象，其形式证明只允许读取集合、谓词、证书、
队列、吸收元、同构和终端门禁。
\end{remark}

\subsection{证明图正规化、引用裁剪与换行安全}
\label{subsec:tex59-proof-graph-normalization}

\begin{definition}[证明图与可审计证明子句]
\label{def:tex59-proof-graph}
对任意本文待证命题 \(\Theta\)，定义其证明图为有限有向无环图
\[
  \mathsf{PG}(\Theta)
  \DefEq
  \bigl(V_\Theta,E_\Theta,\rho_\Theta\bigr).
\]
其中 \(V_\Theta\) 是证明节点集合，\(E_\Theta\subseteq V_\Theta\times V_\Theta\) 是谓词因果边，
\(\rho_\Theta\) 是节点标注函数。每个节点只允许属于以下四类之一：
\[
\begin{array}{ll}
  \mathsf{BaseNode}: & \text{Gao 2025 基础域节点 }\cite{Gao2025PureMathZenodo},\\
  \mathsf{LocalDef}: & \text{本文内部定义或已闭合结论节点},\\
  \mathsf{CertStep}: & \text{证书链推演节点},\\
  \mathsf{GateStep}: & \text{终端门禁或失败吸收节点}.
\end{array}
\]
称一个证明子句可审计，若它可写成
\[
  \Gamma_0\wedge\Gamma_1\wedge\cdots\wedge\Gamma_m
  \Rightarrow
  \Gamma_{m+1},
\]
并且每个前件 \(\Gamma_j\) 都对应 \(V_\Theta\) 中已经定义或已经证明的节点。
\end{definition}

\begin{theorem}[公理项降级：证明图闭合归纳定理]
\label{thm:tex59-proof-graph-induction}
设 \(\mathsf{PG}(\Theta)\) 是有限证明图。若其节点按拓扑序排列为
\[
  v_0,v_1,\ldots,v_N,
\]
并且每个节点 \(v_k\) 满足如下条件之一：
\[
\begin{array}{ll}
  (i) & \rho_\Theta(v_k)\in\mathsf{Base}_{\mathrm{PM}},\\
  (ii) & \rho_\Theta(v_k)\in\mathsf{Local}_{59},\\
  (iii) & \rho_\Theta(v_k)\text{ 由 }v_0,\ldots,v_{k-1}\text{ 的可审计子句推出},
\end{array}
\]
则 \(\Theta\) 的完整证明依赖闭合在
\[
  \mathsf{Base}_{\mathrm{PM}}\cup\mathsf{Local}_{59}
\]
之内。换言之，
\[
  \Dep(\Theta)
  \subseteq
  \mathsf{Base}_{\mathrm{PM}}\cup\mathsf{Local}_{59}.
\]
\end{theorem}

\begin{Certificate}[数学归纳法证书：证明图节点闭合]
\label{cert:tex59-proof-graph-induction}
定义谓词
\[
  P(n):\quad
  \forall k\le n,\quad
  \Dep(v_k)\subseteq
  \mathsf{Base}_{\mathrm{PM}}\cup\mathsf{Local}_{59}.
\]
归纳证书登记为
\[
\begin{aligned}
  G_0 &: v_0\text{ 是基础节点或本文局部定义节点},\\
  G_1 &: G_0\Rightarrow P(0),\\
  G_2 &: P(n)\wedge
  \left[
    v_{n+1}\text{ 由 }v_0,\ldots,v_n\text{ 的可审计子句推出}
  \right]
  \Rightarrow P(n+1),\\
  G_3 &: P(N)\Rightarrow
  \Dep(\Theta)\subseteq
  \mathsf{Base}_{\mathrm{PM}}\cup\mathsf{Local}_{59}.
\end{aligned}
\]
\end{Certificate}

\begin{proof}[证明（数学归纳法证书；基于数学符号推演逻辑的谓词因果链）]
基步：由拓扑序首节点 \(v_0\) 的定义，\(v_0\) 不能依赖尚未出现的证明节点。
若 \(v_0\) 是 Gao 2025 基础节点 \cite{Gao2025PureMathZenodo}，则
\[
  \Dep(v_0)\subseteq\mathsf{Base}_{\mathrm{PM}}.
\]
若 \(v_0\) 是本文局部定义节点，则
\[
  \Dep(v_0)\subseteq\mathsf{Local}_{59}.
\]
因此
\[
  \Dep(v_0)\subseteq
  \mathsf{Base}_{\mathrm{PM}}\cup\mathsf{Local}_{59},
\]
即 \(P(0)\) 成立。

归纳步：假设 \(P(n)\) 成立。若 \(v_{n+1}\) 是基础节点或局部定义节点，结论直接成立。
若 \(v_{n+1}\) 由可审计子句推出，则存在有限前件
\[
  \Gamma_0,\Gamma_1,\ldots,\Gamma_m
\]
使
\[
  \Gamma_0\wedge\Gamma_1\wedge\cdots\wedge\Gamma_m
  \Rightarrow
  \rho_\Theta(v_{n+1}),
\]
且每个 \(\Gamma_j\) 都对应某个 \(v_{\ell_j}\)，其中 \(\ell_j\le n\)。
由归纳假设，
\[
  \Dep(v_{\ell_j})
  \subseteq
  \mathsf{Base}_{\mathrm{PM}}\cup\mathsf{Local}_{59}.
\]
有限并仍属于同一并集，所以
\[
  \Dep(v_{n+1})
  \subseteq
  \mathsf{Base}_{\mathrm{PM}}\cup\mathsf{Local}_{59}.
\]
故 \(P(n+1)\) 成立。

由数学归纳法，\(P(N)\) 成立。终节点 \(v_N\) 标注为 \(\Theta\) 或直接推出 \(\Theta\)，
因此
\[
  \Dep(\Theta)
  \subseteq
  \mathsf{Base}_{\mathrm{PM}}\cup\mathsf{Local}_{59}.
\]
证毕。
\end{proof}

\begin{lemma}[外部引用叶节点剥离引理]
\label{lem:tex59-external-leaf-pruning}
若某外部文献引用只用于说明历史命名、标准背景或问题边界，则它在证明图中至多是
背景叶节点。删除全部背景叶节点后，证明图的可审计子句集合不变，因此定理结论不变。
\end{lemma}

\begin{Certificate}[证书：背景叶节点不参与因果边]
\label{cert:tex59-external-leaf-pruning}
证书链为
\[
\begin{aligned}
  L_1 &: r\in\mathsf{Ext}_{\mathrm{bg}},\\
  L_2 &: r\text{ 不作为可审计子句的前件},\\
  L_3 &: L_2\Rightarrow
  \nexists v\in V_\Theta,\ (r,v)\in E_\Theta,\\
  L_4 &: L_3\Rightarrow
  \mathsf{PG}(\Theta)\setminus\{r\}
  \text{ 保持所有证明因果边},\\
  L_5 &: L_4\Rightarrow \Theta\text{ 的证明结论不变}.
\end{aligned}
\]
\end{Certificate}

\begin{proof}[证明（基于数学符号推演逻辑的谓词因果链）]
由 \(L_2\)，外部引用 \(r\) 不作为任何可审计子句的前件。
因此不存在从 \(r\) 指向证明节点的因果边，即
\[
  \nexists v\in V_\Theta,\quad (r,v)\in E_\Theta.
\]
删除 \(r\) 后，所有真正参与推理的边仍保留；每个可审计子句仍有相同前件与相同结论。
故证明图的推理闭包不变，\(\Theta\) 的证明结论不变。证毕。
\end{proof}

\begin{proposition}[长公式与长路径的断行正规化命题]
\label{prop:tex59-linebreak-normalization}
为避免排版溢出，本文对长路径、长谓词链和长证书元组采用如下正规化规则：
\[
\begin{array}{ll}
  R_1: & \text{仓内路径使用可断行路径命令或拆成多段短路径};\\
  R_2: & \text{长蕴含链拆成多行对齐公式或多段展示公式};\\
  R_3: & \text{长中文说明移出数学文本盒};\\
  R_4: & \text{长证书链按 }C_1,C_2,\ldots\text{ 分行登记}.
\end{array}
\]
若所有新增证明子句均满足 \(R_1\)--\(R_4\)，则新增内容不会引入结构性长行溢出。
\end{proposition}

\begin{Certificate}[证书：可断行对象替代不可断行对象]
\label{cert:tex59-linebreak-normalization}
证书链为
\[
\begin{aligned}
  W_1 &: \text{不可断行对象会形成单个过宽盒},\\
  W_2 &: R_1\text{ 把长路径交给可断行路径命令处理},\\
  W_3 &: R_2\wedge R_4\text{ 把长公式拆成有限短行},\\
  W_4 &: R_3\text{ 避免中文长句锁入数学盒},\\
  W_5 &: W_1\wedge W_2\wedge W_3\wedge W_4
  \Rightarrow
  \text{结构性长行溢出被消除}.
\end{aligned}
\]
\end{Certificate}

\begin{proof}[证明（基于数学符号推演逻辑的谓词因果链）]
长路径、长证书元组与长中文说明造成溢出的共同原因，是它们被放入不可断行的水平盒中。
规则 \(R_1\) 使路径交给可断行路径机制；规则 \(R_2\) 和 \(R_4\) 把长公式链
拆成有限个短行；规则 \(R_3\) 避免把中文长句放入不可自然断行的数学文本盒。
因此，新增证明链在排版上保持可断行。证毕。
\end{proof}

\begin{corollary}[本文后续证明的规范化口径]
\label{cor:tex59-proof-style-normalized}
后续所有“公理降级为定理、证书、引理、命题、推论”的证明均按
\[
  \mathsf{PG}(\Theta)
  \subseteq
  \mathsf{Base}_{\mathrm{PM}}\cup\mathsf{Local}_{59}
\]
的依赖口径理解；所有外部引用只承担背景定位，所有新增长公式按
命题~\ref{prop:tex59-linebreak-normalization} 的断行规则处理。
\end{corollary}

\begin{Certificate}[证书：自足依赖与断行规范合取]
\label{cert:tex59-proof-style-normalized}
证书链为
\[
\begin{aligned}
  S_1 &: \text{依赖闭合由定理~\ref{thm:tex59-proof-graph-induction} 给出},\\
  S_2 &: \text{外部引用剥离由引理~\ref{lem:tex59-external-leaf-pruning} 给出},\\
  S_3 &: \text{换行正规化由命题~\ref{prop:tex59-linebreak-normalization} 给出},\\
  S_4 &: S_1\wedge S_2\wedge S_3
  \Rightarrow
  \text{本文后续证明满足自足、可审计与可编译排版口径}.
\end{aligned}
\]
\end{Certificate}

\begin{proof}[证明（基于数学符号推演逻辑的谓词因果链）]
由定理~\ref{thm:tex59-proof-graph-induction}，证明依赖闭合于 Gao 2025 基础域
\cite{Gao2025PureMathZenodo} 与本文局部域。
由引理~\ref{lem:tex59-external-leaf-pruning}，外部引用作为背景叶节点可从证明图中剥离。
由命题~\ref{prop:tex59-linebreak-normalization}，长路径与长公式按可断行形式书写。
三者合取，得到本文后续证明的规范化口径。证毕。
\end{proof}

\begin{remark}[严格化边界]
本小节只证明本文内部证书系统的依赖闭合、引用裁剪与排版可断行性。
它不把外部经典文献的深层结论改写为本文定理，也不声称外部世纪难题已经在其原始学科
公理系统中被无条件解决。本文所有“化解”均指在 \(\GFramework{}\) 证书语义系统内的
障碍重写、门禁判定、失败吸收与回写修复闭合。
\end{remark}

\subsection{问题索引、聚类与障碍对象}
\label{subsec:tex59-index-cluster}

\begin{definition}[希尔伯特问题的五类障碍聚类]
\label{def:tex59-cluster}
定义聚类映射
\[
  \kappa:\{1,\ldots,23\}\longrightarrow
  \{\Logic,\Phys,\PDE,\Graph,\Zero\}
\]
如下：
\[
\begin{array}{c|l|l}
\text{聚类} & \text{问题编号} & \text{底层障碍}\\ \hline
\Logic & \{1,2,10\} & \text{离散--连续分层、相容性、离散可判定搜索}\\
\Phys & \{6\} & \text{跨域物理量的统一公理化}\\
\PDE & \{19,20,21,22,23\} & \text{正则性、边界值、单值化、变分与奇点}\\
\Graph & \{3,4,15,16,18\} & \text{几何基础、交点、拓扑类型与空间铺砌}\\
\Zero & \{7,8,9,11,12,13,14,17\} & \text{数论、解析函数零点、代数不变量与有限生成}
\end{array}
\]
这里的 \(\Graph\) 不是普通图论，而是把几何、拓扑和语义图谱统一登记为结构对象的
记号。
\end{definition}

\begin{definition}[G-证书对象]
\label{def:tex59-certificate-object}
一个 G-证书对象是七元组
\[
  x=\CertTuple,
\]
其中 \(N\) 是对象名，\(G\) 是生成数据，\(P\) 是谓词族，\(C\) 是上下文与约束，
\(A\) 是审计记录，\(B\) 是边界或障碍槽，\(T\) 是终端门禁。定义接受谓词为
\[
  \Accept(x)=\top
  \quad\Longleftrightarrow\quad
  \bigl(\Pred(P,C)=\top\bigr)
  \wedge
  \bigl(\Audit(A)=\top\bigr)
  \wedge
  \bigl(T(N,G,P,C,A,B)=\top\bigr).
\]
若任一槽触发吸收元 \(\bot\)，则约定
\[
  \Accept(x)=\bot.
\]
\end{definition}

\begin{definition}[G-化解谓词]
\label{def:tex59-resolution}
对任意 \(i\in\{1,\ldots,23\}\)，定义
\[
  \Resolve_G(H_i)=\top
\]
当且仅当存在对象 \(x_i=\CertTuple\) 与映射链
\[
  Q_i
  \xrightarrow{\;\Norm_i\;}
  \NF_i
  \xrightarrow{\;\Inv_i\;}
  \Hcal_i
  \xrightarrow{\;\Gate_i\;}
  \{\top,\bot\}
\]
使得：
\[
\begin{aligned}
&\text{(1) } \Norm_i \text{ 保持 } Q_i \text{ 中被声明的谓词依赖；}\\
&\text{(2) } \Inv_i \text{ 把底空间障碍压缩为同调或语义不变量；}\\
&\text{(3) } \Gate_i(\Inv_i(\Norm_i(Q_i)))=\top \Longleftrightarrow \Accept(x_i)=\top;\\
&\text{(4) } \Gate_i=\bot \text{ 时失败进入 } \bot \text{ 或 } \BackwriteQueue(\Vsys).
\end{aligned}
\]
此时称 \(x_i\) 是 \(H_i\) 的 G-化解证书。
\end{definition}

\subsection{公理项降级为定理}
\label{subsec:tex59-axiom-downgrade}

\begin{definition}[构造秩]
\label{def:tex59-rank}
令 \(\mathfrak{S}\) 为由定义、规范化映射、谓词闭包、边缘算子、同调不变量和终端门禁
生成的形式系统。对任意公式 \(\varphi\)，定义构造秩 \(\Rank(\varphi)\in\NN\)：
\[
\begin{aligned}
\Rank(\varphi)=0
&\quad\Longleftrightarrow\quad
\varphi \text{ 是定义展开、恒等映射或空证书判断};\\
\Rank(\varphi)\le n+1
&\quad\Longleftrightarrow\quad
\varphi \text{ 由秩 }\le n\text{ 的公式经有限次合取、复合、}\\
&\quad\text{谓词替换、同调投影、吸收闭合或终端门禁得到}.
\end{aligned}
\]
\end{definition}

\begin{theorem}[公理项降级定理]
\label{thm:tex59-axiom-downgrade}
本文后续出现的“公理项”均可降级为构造定理：若一个公理项
\(\mathsf A\) 只断言定义对象、闭包操作、吸收规则、证书门禁或有限聚类覆盖，则存在
数学归纳法证书
\[
  \Cert_{\mathrm{ind}}(\mathsf A)
\]
使得
\[
  \mathfrak{S}\vdash \mathsf A.
\]
\end{theorem}

\begin{Certificate}[公理降级的数学归纳法证书]
\label{cert:tex59-axiom-induction}
证书由三部分组成：
\[
  \Cert_{\mathrm{ind}}(\mathsf A)
  \DefEq
  \bigl(
    \Cert_0,\Cert_{\mathrm{step}},\Cert_{\mathrm{close}}
  \bigr).
\]
其中 \(\Cert_0\) 验证 \(\Rank=0\) 的定义展开；\(\Cert_{\mathrm{step}}\) 验证从
\(\Rank\le n\) 到 \(\Rank\le n+1\) 的闭包保持；\(\Cert_{\mathrm{close}}\) 验证有限复合
之后仍落在 \(\mathfrak{S}\) 的公式域内。
\end{Certificate}

\begin{proof}[证明]
对 \(\Rank(\mathsf A)\) 作数学归纳。

当 \(\Rank(\mathsf A)=0\) 时，\(\mathsf A\) 只是定义展开、恒等映射或空证书判断。
由定义~\ref{def:tex59-rank}，\(\Cert_0\) 直接给出
\(\mathfrak{S}\vdash \mathsf A\)。

假设所有构造秩不超过 \(n\) 的公理项都已由 \(\Cert_{\mathrm{ind}}\) 证明。若
\(\Rank(\mathsf A)\le n+1\)，则 \(\mathsf A\) 由有限个秩不超过 \(n\) 的公式经合取、
复合、谓词替换、同调投影、吸收闭合或终端门禁得到。按归纳假设，每个前件都有证明。
这些生成操作均已作为 \(\mathfrak{S}\) 的构造规则写入定义~\ref{def:tex59-rank}，故
\(\Cert_{\mathrm{step}}\) 给出后件证明。最后，本文只处理有限个希尔伯特问题与有限个聚类，
有限复合由 \(\Cert_{\mathrm{close}}\) 保持在系统内。故
\(\mathfrak{S}\vdash \mathsf A\)。
\end{proof}

\begin{corollary}[无外部悬空公理]
\label{cor:tex59-no-floating-axiom}
本文后续所有“公理化”表述均可理解为“已由定义与有限构造证明的系统规则”，而不是
不可追踪的外部断言。
\end{corollary}

\begin{proof}
直接由定理~\ref{thm:tex59-axiom-downgrade} 得到。若某规则不是定义、闭包、吸收、
门禁或有限覆盖，则本文不把它登记为公理项。
\end{proof}

\subsection{有限谓词因果链与证书核}
\label{subsec:tex59-causal-chain-kernel}

\begin{definition}[有限谓词因果链]
\label{def:tex59-causal-chain}
在 G-证书形式系统 \(\mathfrak S_G\) 中，一个有限谓词因果链是三元组
\[
  \Xi
  =
  \bigl(
    (\chi_0,\chi_1,\ldots,\chi_m),
    (r_0,r_1,\ldots,r_{m-1}),
    \Pi_{\Xi}
  \bigr),
\]
其中每个 \(\chi_j\) 是谓词公式，每个 \(r_j\) 是局部推理规则，且证书包
\(\Pi_{\Xi}\) 包含局部证书
\[
  \pi_j:\quad
  \mathfrak S_G\vdash \chi_j\Rightarrow\chi_{j+1},
  \qquad j=0,\ldots,m-1.
\]
若另有初始证书
\[
  \pi_0^{\mathrm{init}}:\quad
  \mathfrak S_G\vdash\chi_0,
\]
则称 \(\Xi\) 是从 \(\chi_0\) 到 \(\chi_m\) 的可签发因果链。
\end{definition}

\begin{theorem}[公理降级：有限谓词因果链可靠性定理]
\label{thm:tex59-causal-chain-soundness}
设 \(\Xi\) 是定义~\ref{def:tex59-causal-chain} 中的有限谓词因果链。若
\[
  \mathfrak S_G\vdash\chi_0
  \quad\text{且}\quad
  \forall j<m,\ \mathfrak S_G\vdash \chi_j\Rightarrow\chi_{j+1},
\]
则
\[
  \mathfrak S_G\vdash\chi_m.
\]
因此，本文后续所有形如
\[
  \Barrier
  \Longrightarrow
  \Norm
  \Longrightarrow
  \Inv
  \Longrightarrow
  \Cert
  \Longrightarrow
  \Gate
  \Longrightarrow
  \Resolve
\]
的链式证明，均可降级为有限次 modus ponens 与数学归纳法证书。
\end{theorem}

\begin{Certificate}[数学归纳法证书：链式可靠性]
\label{cert:tex59-causal-chain-soundness}
定义
\[
  P(n):\quad \mathfrak S_G\vdash\chi_n.
\]
证书为
\[
  \Pi_{\mathrm{chain}}
  =
  \bigl(
    P(0),\
    \forall n<m,\ P(n)\wedge(\chi_n\Rightarrow\chi_{n+1})\Rightarrow P(n+1)
  \bigr).
\]
\end{Certificate}

\begin{proof}[证明（数学归纳法证书；基于数学符号推演逻辑的谓词因果链）]
基步：由初始证书
\[
  \mathfrak S_G\vdash\chi_0
\]
得 \(P(0)\)。

归纳步：假设 \(P(n)\) 成立，即
\[
  \mathfrak S_G\vdash\chi_n.
\]
又由局部证书
\[
  \mathfrak S_G\vdash\chi_n\Rightarrow\chi_{n+1}.
\]
对二者使用一次 modus ponens，得到
\[
  \mathfrak S_G\vdash\chi_{n+1}.
\]
故 \(P(n)\Rightarrow P(n+1)\)。由有限数学归纳法，
\[
  \mathfrak S_G\vdash\chi_m.
\]
\end{proof}

\begin{corollary}[证书链终端可签发]
\label{cor:tex59-chain-terminal-issuance}
若某个希尔伯特障碍 \(\Barrier(H_i)\) 的证明链可写为有限谓词因果链
\[
  \Barrier(H_i)=\chi_0
  \Longrightarrow
  \chi_1
  \Longrightarrow
  \cdots
  \Longrightarrow
  \chi_m=\Resolve_G(H_i),
\]
且所有局部推理均有证书，则
\[
  \mathfrak S_G\vdash\Resolve_G(H_i).
\]
\end{corollary}

\begin{proof}[证明（基于数学符号推演逻辑的谓词因果链）]
把该链代入定理~\ref{thm:tex59-causal-chain-soundness}。
此时终端公式为
\[
  \Resolve_G(H_i).
\]
故结论成立。
\end{proof}

\subsection{最小粒度推理审计核与外部引用替代}
\label{subsec:tex59-atomic-audit-kernel}

\begin{definition}[原子推理步]
\label{def:tex59-atomic-inference-step}
在本文的证明审计口径中，一个原子推理步是四元组
\[
  \alpha
  =
  (\Gamma_{\alpha},\rho_{\alpha},\Delta_{\alpha},c_{\alpha}),
\]
其中 \(\Gamma_{\alpha}\) 是有限前件集，\(\rho_{\alpha}\) 是推理规则名，
\(\Delta_{\alpha}\) 是单个后件公式，\(c_{\alpha}\) 是局部证书。
允许的规则只包括下列有限类型：
\[
\begin{array}{ll}
  \mathsf{Def}: & \text{定义展开或定义回代};\\
  \mathsf{Sub}: & \text{同类型符号替换};\\
  \mathsf{AndI}: & \text{有限合取引入};\\
  \mathsf{AndE}: & \text{有限合取消去};\\
  \mathsf{MP}: & \text{由 } A \text{ 与 } A\Rightarrow B \text{ 推出 } B;\\
  \mathsf{Eq}: & \text{等式传递、等式替换或同构运输};\\
  \mathsf{Ind}: & \text{有限数学归纳法一步};\\
  \mathsf{Abs}: & \text{吸收律 } a+\bot=\bot;\\
  \mathsf{Gate}: & \text{终端门禁 }T=1\text{ 或 }T=0\text{ 的判定回代};\\
  \mathsf{RefDrop}: & \text{背景引用叶节点剥离}.
\end{array}
\]
称 \(\alpha\) 合法，若
\[
  \Gamma_{\alpha}
  \vdash_{\rho_{\alpha},\,c_{\alpha}}
  \Delta_{\alpha}
\]
能够由上述某一规则直接验证。
\end{definition}

\begin{definition}[最小粒度审计展开]
\label{def:tex59-atomic-audit-expansion}
设 \(\Pi(\Theta)\) 是命题 \(\Theta\) 的证明。称
\[
  \mathsf{Audit}_{\min}(\Pi(\Theta))
\]
为 \(\Pi(\Theta)\) 的最小粒度审计展开，若它是有限序列
\[
  \alpha_0,\alpha_1,\ldots,\alpha_N
\]
并满足：
\[
\begin{aligned}
  &\text{(1) 每个 }\alpha_k\text{ 是定义~\ref{def:tex59-atomic-inference-step} 的合法原子步};\\
  &\text{(2) 若 }\Gamma_{\alpha_k}\text{ 中含非基础公式，则该公式等于某个 }\Delta_{\alpha_j},\ j<k;\\
  &\text{(3) }\Delta_{\alpha_N}=\Theta;\\
  &\text{(4) 任意外部引用只可作为 }\mathsf{RefDrop}\text{ 的输入叶节点，不可作为 }\mathsf{MP}\text{ 前件}.
\end{aligned}
\]
\end{definition}

\begin{theorem}[公理降级：最小粒度审计展开存在定理]
\label{thm:tex59-atomic-audit-existence}
若命题 \(\Theta\) 的证明使用本文允许的定义展开、有限合取、等式替换、有限归纳、
吸收元、终端门禁和背景引用剥离规则，则存在
\[
  \mathsf{Audit}_{\min}(\Pi(\Theta)).
\]
\end{theorem}

\begin{Certificate}[数学归纳法证书：复合证明到原子证明]
\label{cert:tex59-atomic-audit-existence}
令证明树高度为
\[
  h(\Pi)\in\NN.
\]
定义谓词
\[
  P(n):\quad
  h(\Pi)\le n
  \Rightarrow
  \Pi\text{ 存在最小粒度审计展开}.
\]
证书登记为
\[
\begin{aligned}
  A_0 &: P(0),\\
  A_1 &: P(n)\wedge
         \left[
           \Pi\text{ 由有限个高度 }\le n\text{ 的子证明经一个允许规则复合}
         \right]
         \Rightarrow P(n+1),\\
  A_2 &: \forall\Pi,\ \exists n,\ h(\Pi)=n
         \Rightarrow
         \mathsf{Audit}_{\min}(\Pi)\text{ 存在}.
\end{aligned}
\]
\end{Certificate}

\begin{proof}[证明（数学归纳法证书；基于数学符号推演逻辑的谓词因果链）]
基步 \(h(\Pi)=0\)：此时证明只包含一个基础节点。
若该节点是定义、Gao 2025 基础域对象 \cite{Gao2025PureMathZenodo} 或本文局部定义，则取
\[
  \alpha_0=(\varnothing,\mathsf{Def},\Theta,c_0).
\]
若该节点是背景引用，则取
\[
  \alpha_0=(\{r\},\mathsf{RefDrop},\top,c_r),
\]
并由引理~\ref{lem:tex59-external-leaf-pruning} 删除其证明依赖角色。
故 \(P(0)\) 成立。

归纳步：假设 \(P(n)\) 成立。设 \(h(\Pi)=n+1\)。
则 \(\Pi\) 的最后一步由有限个子证明
\[
  \Pi_1,\ldots,\Pi_k
\]
经一个允许规则 \(\rho\) 得到。每个子证明高度不超过 \(n\)，由归纳假设，各自存在
最小粒度审计展开：
\[
  \mathsf{Audit}_{\min}(\Pi_j)
  =
  (\alpha_{j,0},\ldots,\alpha_{j,N_j}).
\]
把这些有限序列按子证明依赖顺序串接，再追加最后一步
\[
  \alpha_{\ast}
  =
  (\{\Delta_{\alpha_{1,N_1}},\ldots,\Delta_{\alpha_{k,N_k}}\},
   \rho,\Theta,c_{\ast}).
\]
由于 \(\rho\) 是允许规则之一，\(\alpha_{\ast}\) 是合法原子步。
串接后的有限序列满足定义~\ref{def:tex59-atomic-audit-expansion} 的四项条件，故
\[
  P(n)\Rightarrow P(n+1).
\]
由数学归纳法，任意有限证明 \(\Pi(\Theta)\) 都存在
\(\mathsf{Audit}_{\min}(\Pi(\Theta))\)。证毕。
\end{proof}

\begin{definition}[单位原子步、语义内容与语义不可展开性]
\label{def:tex59-semantic-irreducible-atomic-step}
令
\[
  \mathsf{NF}_{59}(F)
\]
表示在本文局部定义域内对公式 \(F\) 做完全部定义展开、符号消歧和类型回代后的正规型。
定义公式语义内容为
\[
  \mathsf{Sem}_{59}(F)\DefEq \mathsf{NF}_{59}(F).
\]
对原子推理候选
\[
  \alpha=(\Gamma_{\alpha},\rho_{\alpha},\Delta_{\alpha},c_{\alpha}),
\]
定义其语义接口为
\[
  \mathsf{Int}_{59}(\alpha)
  \DefEq
  \left(
  \{\mathsf{Sem}_{59}(G):G\in\Gamma_{\alpha}\},
  \mathsf{Sem}_{59}(\Delta_{\alpha})
  \right).
\]
若
\[
  \beta_0,\ldots,\beta_m
\]
是有限推理序列，定义其外部输入前件为
\[
  \Gamma_{\beta_{0:m}}^{\mathrm{in}}
  \DefEq
  \left(\bigcup_{i=0}^{m}\Gamma_{\beta_i}\right)
  \setminus
  \{\Delta_{\beta_i}:0\le i<m\},
\]
并定义序列语义接口为
\[
  \mathsf{Int}_{59}(\beta_0,\ldots,\beta_m)
  \DefEq
  \left(
  \{\mathsf{Sem}_{59}(G):G\in\Gamma_{\beta_{0:m}}^{\mathrm{in}}\},
  \mathsf{Sem}_{59}(\Delta_{\beta_m})
  \right).
\]

令 \(\mu(F)\in\NN\) 为公式 \(F\) 中尚未展开的局部定义头、缩写头与可消歧符号头的个数，
并令
\[
  \mu(\Gamma_{\alpha},\Delta_{\alpha})
  \DefEq
  \sum_{G\in\Gamma_{\alpha}}\mu(G)+\mu(\Delta_{\alpha}).
\]
令 \(\ell(c_{\alpha})\in\NN_{\ge1}\) 为证书 \(c_{\alpha}\) 中显式编码的规则实例数。

称 \(\alpha\) 是单位原子步，记为
\[
  \mathsf{Atom}_{59}(\alpha)=\top,
\]
若 \(\alpha\) 合法，且满足：
\[
\begin{aligned}
  &\mu(\Gamma_{\alpha},\Delta_{\alpha})=0,\\
  &\ell(c_{\alpha})=1,\\
  &\rho_{\alpha}\in
  \{\mathsf{Def},\mathsf{Sub},\mathsf{AndI},\mathsf{AndE},
    \mathsf{MP},\mathsf{Eq},\mathsf{Ind},\mathsf{Abs},\mathsf{Gate},\mathsf{RefDrop}\}.
\end{aligned}
\]
其中“单位”要求如下：\(\mathsf{Def}\) 只展开或回代一个定义头；
\(\mathsf{Sub}\) 只执行一次同类型替换；\(\mathsf{AndI}\) 只引入一个二元合取；
\(\mathsf{AndE}\) 只从一个二元合取投影一个分量；\(\mathsf{MP}\) 只使用一个蕴含式；
\(\mathsf{Eq}\) 只执行一次等式替换或一次同构运输；\(\mathsf{Ind}\) 只调用一个已经列明
基步与归纳步的有限归纳模式；\(\mathsf{Abs}\) 只调用一次 \(a+\bot=\bot\)；
\(\mathsf{Gate}\) 只回代一个门禁值；\(\mathsf{RefDrop}\) 只剥离一个背景标签。

称 \(\alpha\) 具有语义不可展开性，记为
\[
  \mathsf{SemIrred}_{59}(\alpha)=\top,
\]
若不存在有限序列
\[
  \beta_0,\ldots,\beta_m,
  \qquad m\ge1,
\]
使其在保持同一语义接口
\[
  \mathsf{Int}_{59}(\beta_0,\ldots,\beta_m)
  =
  \mathsf{Int}_{59}(\alpha)
\]
的同时，把 \(\alpha\) 的局部证书继续拆成两个或更多合法规则实例，或者把
\(\Gamma_{\alpha},\Delta_{\alpha}\) 中的某个非正规语义头继续展开。
\end{definition}

\begin{theorem}[公理降级：有限语义收束定理]
\label{thm:tex59-finite-semantic-contraction}
设
\[
  \alpha=(\Gamma_{\alpha},\rho_{\alpha},\Delta_{\alpha},c_{\alpha})
\]
是合法原子推理候选，且
\[
  \mu(\Gamma_{\alpha},\Delta_{\alpha})<\infty,
  \qquad
  \ell(c_{\alpha})<\infty.
\]
则存在有限序列
\[
  \alpha^{(0)},\alpha^{(1)},\ldots,\alpha^{(M)}
\]
使得
\[
  \alpha^{(0)}=\alpha,
  \qquad
  \mathsf{SemIrred}_{59}(\alpha^{(M)})=\top,
\]
并且每一步从 \(\alpha^{(j)}\) 到 \(\alpha^{(j+1)}\) 只执行一次定义展开、一次证书拆分或一次单位规则化。
\end{theorem}

\begin{Certificate}[数学归纳法证书：可展开度下降]
\label{cert:tex59-finite-semantic-contraction}
定义可展开度
\[
  W(\alpha)
  \DefEq
  \mu(\Gamma_{\alpha},\Delta_{\alpha})
  +
  \ell(c_{\alpha})-1
  \in\NN.
\]
定义谓词
\[
  P(n):\quad
  W(\alpha)\le n
  \Rightarrow
  \alpha\text{ 可经有限步收束为语义不可展开步}.
\]
归纳证书为
\[
\begin{aligned}
  S_0 &: P(0),\\
  S_1 &: P(n)\wedge W(\alpha)=n+1
         \Rightarrow
         \exists\alpha'\,[W(\alpha')\le n
         \wedge
         \mathsf{Int}_{59}(\alpha')=\mathsf{Int}_{59}(\alpha)],\\
  S_2 &: S_0\wedge S_1
         \Rightarrow
         \forall n\in\NN,\ P(n).
\end{aligned}
\]
\end{Certificate}

\begin{proof}[证明（数学归纳法证书；基于数学符号推演逻辑的谓词因果链）]
基步 \(W(\alpha)=0\)。由
\[
  W(\alpha)
  =
  \mu(\Gamma_{\alpha},\Delta_{\alpha})
  +
  \ell(c_{\alpha})-1
  =
  0
\]
且
\[
  \mu(\Gamma_{\alpha},\Delta_{\alpha})\ge0,
  \qquad
  \ell(c_{\alpha})\ge1,
\]
只能有
\[
  \mu(\Gamma_{\alpha},\Delta_{\alpha})=0,
  \qquad
  \ell(c_{\alpha})=1.
\]
因此所有公式已经处于 \(\mathsf{NF}_{59}\)，证书只含一个规则实例。
再由 \(\alpha\) 的合法性，\(\rho_{\alpha}\) 是定义~\ref{def:tex59-atomic-inference-step}
列出的允许规则之一。按定义~\ref{def:tex59-semantic-irreducible-atomic-step}，
\[
  \mathsf{SemIrred}_{59}(\alpha)=\top.
\]
故 \(P(0)\) 成立。

归纳步。假设 \(P(n)\) 成立，并设
\[
  W(\alpha)=n+1.
\]
若
\[
  \mu(\Gamma_{\alpha},\Delta_{\alpha})>0,
\]
则存在某个前件或后件公式含有一个尚未展开的定义头或可消歧符号头。
执行一次 \(\mathsf{Def}\) 展开或类型回代，得到 \(\alpha'\)。该操作不改变语义接口：
\[
  \mathsf{Int}_{59}(\alpha')
  =
  \mathsf{Int}_{59}(\alpha),
\]
但使
\[
  \mu(\Gamma_{\alpha'},\Delta_{\alpha'})
  =
  \mu(\Gamma_{\alpha},\Delta_{\alpha})-1.
\]
因此
\[
  W(\alpha')=W(\alpha)-1=n.
\]
由归纳假设，\(\alpha'\) 可有限收束为语义不可展开步，从而 \(\alpha\) 也可。

若
\[
  \mu(\Gamma_{\alpha},\Delta_{\alpha})=0
  \quad\text{且}\quad
  W(\alpha)=n+1,
\]
则
\[
  \ell(c_{\alpha})>1.
\]
此时 \(c_{\alpha}\) 显式编码至少两个规则实例。拆出第一个规则实例作为一个单位证书，
并把其结论作为后续规则的前件，得到等语义接口的局部序列。剩余复合证书的规则实例数减少一：
\[
  \ell(c_{\alpha'})=\ell(c_{\alpha})-1.
\]
于是
\[
  W(\alpha')=W(\alpha)-1=n.
\]
仍由归纳假设，\(\alpha'\) 可有限收束为语义不可展开步。

两种情形穷尽 \(W(\alpha)=n+1\)。故
\[
  P(n)\Rightarrow P(n+1).
\]
由数学归纳法，
\[
  \forall n\in\NN,\quad P(n)
\]
成立。证毕。
\end{proof}

\begin{lemma}[语义不可展开性等价于单位原子性]
\label{lem:tex59-semantic-irreducible-iff-atomic}
对任意合法推理候选 \(\alpha\)，有
\[
  \mathsf{SemIrred}_{59}(\alpha)=\top
  \Longleftrightarrow
  \mathsf{Atom}_{59}(\alpha)=\top.
\]
\end{lemma}

\begin{Certificate}[证书：不可展开条件与单位原子条件逐项对齐]
\label{cert:tex59-semantic-irreducible-iff-atomic}
证书链为
\[
\begin{aligned}
  I_1 &: \mathsf{SemIrred}_{59}(\alpha)
         \Rightarrow
         \mu(\Gamma_{\alpha},\Delta_{\alpha})=0,\\
  I_2 &: \mathsf{SemIrred}_{59}(\alpha)
         \Rightarrow
         \ell(c_{\alpha})=1,\\
  I_3 &: \mathsf{SemIrred}_{59}(\alpha)\wedge\alpha\text{ 合法}
         \Rightarrow
         \rho_{\alpha}\text{ 是一个单位允许规则},\\
  I_4 &: I_1\wedge I_2\wedge I_3
         \Rightarrow
         \mathsf{Atom}_{59}(\alpha),\\
  I_5 &: \mathsf{Atom}_{59}(\alpha)
         \Rightarrow
         \mu(\Gamma_{\alpha},\Delta_{\alpha})=0
         \wedge
         \ell(c_{\alpha})=1,\\
  I_6 &: I_5\wedge\alpha\text{ 合法}
         \Rightarrow
         \mathsf{SemIrred}_{59}(\alpha).
\end{aligned}
\]
\end{Certificate}

\begin{proof}[证明（基于数学符号推演逻辑的谓词因果链）]
先证正向。若
\[
  \mathsf{SemIrred}_{59}(\alpha)=\top
\]
但
\[
  \mu(\Gamma_{\alpha},\Delta_{\alpha})>0,
\]
则存在至少一个尚未正规化的语义头。按定义~\ref{def:tex59-semantic-irreducible-atomic-step}，
可对该语义头执行一次定义展开或类型回代，并保持同一语义接口。
这给出一个合法的进一步展开，矛盾。因此
\[
  \mu(\Gamma_{\alpha},\Delta_{\alpha})=0.
\]
若
\[
  \ell(c_{\alpha})>1,
\]
则证书中至少含有两个规则实例。按定义可把第一个实例拆出，并把其结论作为剩余实例的中间前件。
这同样给出保持同一语义接口的合法进一步展开，矛盾。因此
\[
  \ell(c_{\alpha})=1.
\]
又因 \(\alpha\) 合法，\(\rho_{\alpha}\) 必属于允许规则族；由于无法再拆分，其规则实例必须是单位形态。
所以
\[
  \mathsf{Atom}_{59}(\alpha)=\top.
\]

再证反向。若
\[
  \mathsf{Atom}_{59}(\alpha)=\top,
\]
则由定义直接得到
\[
  \mu(\Gamma_{\alpha},\Delta_{\alpha})=0,
  \qquad
  \ell(c_{\alpha})=1.
\]
第一式说明前件和后件没有可继续展开的语义头；第二式说明局部证书只含一个规则实例，
不存在可拆成两个或更多合法规则实例的复合证书。故不存在保持同一语义接口的进一步合法展开，
即
\[
  \mathsf{SemIrred}_{59}(\alpha)=\top.
\]
证毕。
\end{proof}

\begin{proposition}[实现语义不可展开性即实现原子推理步]
\label{prop:tex59-semantic-irreducible-implements-atomic-proof}
设
\[
  \mathsf{Audit}_{\min}(\Pi(\Theta))
  =
  (\alpha_0,\ldots,\alpha_N).
\]
则
\[
  \left[
  \forall k\le N,\ 
  \mathsf{SemIrred}_{59}(\alpha_k)=\top
  \right]
  \Longleftrightarrow
  \left[
  \forall k\le N,\ 
  \mathsf{Atom}_{59}(\alpha_k)=\top
  \right].
\]
因此，实现每一步的语义不可展开性，等价于实现每一步的单位原子推理步证明。
\end{proposition}

\begin{Certificate}[证书：逐步等价到序列等价]
\label{cert:tex59-semantic-irreducible-implements-atomic-proof}
证书链为
\[
\begin{aligned}
  Q_1 &: \forall k\le N,\
         \mathsf{SemIrred}_{59}(\alpha_k)
         \Longleftrightarrow
         \mathsf{Atom}_{59}(\alpha_k)
         \quad\text{（引理~\ref{lem:tex59-semantic-irreducible-iff-atomic}）},\\
  Q_2 &: Q_1
         \Rightarrow
         \bigwedge_{k=0}^{N}\mathsf{SemIrred}_{59}(\alpha_k)
         \Longleftrightarrow
         \bigwedge_{k=0}^{N}\mathsf{Atom}_{59}(\alpha_k),\\
  Q_3 &: Q_2
         \Rightarrow
         \text{语义不可展开性实现与单位原子证明实现等价}.
\end{aligned}
\]
\end{Certificate}

\begin{proof}[证明（基于数学符号推演逻辑的谓词因果链）]
由引理~\ref{lem:tex59-semantic-irreducible-iff-atomic}，对每个 \(k\le N\) 都有
\[
  \mathsf{SemIrred}_{59}(\alpha_k)=\top
  \Longleftrightarrow
  \mathsf{Atom}_{59}(\alpha_k)=\top.
\]
对有限序列 \(k=0,\ldots,N\) 作合取，得到
\[
  \bigwedge_{k=0}^{N}
  \mathsf{SemIrred}_{59}(\alpha_k)
  \Longleftrightarrow
  \bigwedge_{k=0}^{N}
  \mathsf{Atom}_{59}(\alpha_k).
\]
左端表示每个推理步已经不能在本文语义接口下继续展开；右端表示每个推理步已经是
定义~\ref{def:tex59-semantic-irreducible-atomic-step} 的单位原子步。
故实现语义不可展开性即实现原子推理步证明。证毕。
\end{proof}

\begin{corollary}[最小粒度审计的必要收束口径]
\label{cor:tex59-atomic-audit-contraction-policy}
在本文中，一个证明链被视为完成最小粒度审计，当且仅当它不仅存在
\[
  \mathsf{Audit}_{\min}(\Pi(\Theta)),
\]
而且其中每个步骤都满足
\[
  \mathsf{SemIrred}_{59}(\alpha_k)=\top.
\]
等价地，每个步骤都满足
\[
  \mathsf{Atom}_{59}(\alpha_k)=\top.
\]
\end{corollary}

\begin{Certificate}[证书：存在展开加不可展开收束]
\label{cert:tex59-atomic-audit-contraction-policy}
证书链为
\[
\begin{aligned}
  U_1 &: \text{定理~\ref{thm:tex59-atomic-audit-existence}}
         \Rightarrow
         \mathsf{Audit}_{\min}(\Pi(\Theta))\text{ 存在},\\
  U_2 &: \text{定理~\ref{thm:tex59-finite-semantic-contraction}}
         \Rightarrow
         \text{每个合法候选步可有限收束},\\
  U_3 &: \text{命题~\ref{prop:tex59-semantic-irreducible-implements-atomic-proof}}
         \Rightarrow
         \mathsf{SemIrred}_{59}\text{ 与 }\mathsf{Atom}_{59}\text{ 等价},\\
  U_4 &: U_1\wedge U_2\wedge U_3
         \Rightarrow
         \text{最小粒度审计需要同时完成展开与收束}.
\end{aligned}
\]
\end{Certificate}

\begin{proof}[证明（基于数学符号推演逻辑的谓词因果链）]
由定理~\ref{thm:tex59-atomic-audit-existence}，任意有限证明可以先拆成有限审计序列。
由定理~\ref{thm:tex59-finite-semantic-contraction}，该序列中若仍有非正规语义头或复合证书，
则可以继续做有限收束。收束完成时，每一步满足
\[
  \mathsf{SemIrred}_{59}(\alpha_k)=\top.
\]
由命题~\ref{prop:tex59-semantic-irreducible-implements-atomic-proof}，这等价于
\[
  \mathsf{Atom}_{59}(\alpha_k)=\top.
\]
因此本文后续所谓“原子推理步的证明”，严格理解为“有限展开存在 + 每步语义不可展开收束完成”。
证毕。
\end{proof}

\begin{remark}[语义不可展开性的相对边界]
这里的“语义不可展开性”是相对于本文固定的 \(\mathsf{NF}_{59}\)、允许规则族和局部证书格式而言。
它不声称一个公式在所有可能元语言中都不可解释，也不声称外部数学对象没有更细理论结构。
它只声明：在本文的审计语义接口下，若公式已正规化、证书只含一个单位规则实例，
则该推理步不能再被拆成保持同一接口的更细证明步。
\end{remark}

\begin{lemma}[外部引用替代为背景标签引理]
\label{lem:tex59-external-reference-replacement}
设外部引用 \(r\in\mathsf{Ext}_{\mathrm{bg}}\) 只承担命名、历史定位或问题边界说明。
定义替代算子
\[
  \mathsf{Erase}_{\mathrm{ext}}(r)=\lambda_r,
\]
其中 \(\lambda_r\in\Lambda_{59}\) 只是背景标签。若证明中不存在以 \(r\) 为前件的
\(\mathsf{MP}\) 步，则
\[
  \Dep(\Theta)
  =
  \Dep\bigl(\mathsf{Erase}_{\mathrm{ext}}(\Theta)\bigr).
\]
\end{lemma}

\begin{Certificate}[证书：引用标签化不改变证明边]
\label{cert:tex59-external-reference-replacement}
证书链为
\[
\begin{aligned}
  R_1 &: r\in\mathsf{Ext}_{\mathrm{bg}},\\
  R_2 &: r\text{ 不作为 }\mathsf{MP}\text{ 前件},\\
  R_3 &: \mathsf{Erase}_{\mathrm{ext}}(r)=\lambda_r,\\
  R_4 &: \lambda_r\text{ 不属于 }\mathsf{Base}_{\mathrm{PM}}\cup\mathsf{Local}_{59}\text{ 的推理节点},\\
  R_5 &: R_1\wedge R_2\wedge R_3\wedge R_4
        \Rightarrow
        \Dep(\Theta)=\Dep(\mathsf{Erase}_{\mathrm{ext}}(\Theta)).
\end{aligned}
\]
\end{Certificate}

\begin{proof}[证明（基于数学符号推演逻辑的谓词因果链）]
由 \(R_2\)，引用 \(r\) 没有从自身指向本文结论的推理边。
由 \(R_3\)，替代算子只把 \(r\) 改写成背景标签 \(\lambda_r\)。
由 \(R_4\)，该标签不参与基础域或局部域的推理节点。
因此替代前后的证明图只在背景叶节点名称上不同；所有可审计原子步、所有
\(\mathsf{MP}\) 边、所有定义展开边和所有门禁边完全相同。故依赖集相等。证毕。
\end{proof}

\begin{proposition}[Zenodo 基础域的外部引用替代性命题]
\label{prop:tex59-gao-base-substitution}
设 \(\Theta\) 是本文任意已审计命题。若
\[
  \mathsf{Audit}_{\min}(\Pi(\Theta))
\]
存在，且其中所有非局部基础节点均属于
\[
  \mathsf{Base}_{\mathrm{PM}}
\]
或可由引理~\ref{lem:tex59-external-reference-replacement} 标签化剥离，则
\[
  \Dep(\Theta)
  \subseteq
  \mathsf{Base}_{\mathrm{PM}}\cup\mathsf{Local}_{59}.
\]
换言之，除 Gao 2025 年 Zenodo 基础域 \cite{Gao2025PureMathZenodo} 与本文内部定义外，
外部文献不承担证明前件。
\end{proposition}

\begin{Certificate}[证书：原子审计到基础依赖压缩]
\label{cert:tex59-gao-base-substitution}
证书链为
\[
\begin{aligned}
  G_1 &: \mathsf{Audit}_{\min}(\Pi(\Theta))=(\alpha_0,\ldots,\alpha_N),\\
  G_2 &: \forall k,\ \alpha_k\text{ 合法},\\
  G_3 &: \forall k,\ \Gamma_{\alpha_k}
         \subseteq
         \mathsf{Base}_{\mathrm{PM}}\cup\mathsf{Local}_{59}
         \cup\mathsf{Ext}_{\mathrm{bg}},\\
  G_4 &: \text{引理~\ref{lem:tex59-external-reference-replacement}}
         \Rightarrow
         \mathsf{Ext}_{\mathrm{bg}}\text{ 可标签化剥离},\\
  G_5 &: G_1\wedge G_2\wedge G_3\wedge G_4
         \Rightarrow
         \Dep(\Theta)
         \subseteq
         \mathsf{Base}_{\mathrm{PM}}\cup\mathsf{Local}_{59}.
\end{aligned}
\]
\end{Certificate}

\begin{proof}[证明（基于数学符号推演逻辑的谓词因果链）]
由 \(G_1\) 与 \(G_2\)，\(\Theta\) 的证明已被拆成有限个合法原子步。
由 \(G_3\)，任一原子步的前件只可能来自 Gao 2025 基础域
\cite{Gao2025PureMathZenodo}、本文局部域或背景引用域。
由引理~\ref{lem:tex59-external-reference-replacement}，背景引用域中的节点可替换为
不参与推理的标签并从依赖集中剥离。
剥离后，所有剩余前件均属于
\[
  \mathsf{Base}_{\mathrm{PM}}\cup\mathsf{Local}_{59}.
\]
因此
\[
  \Dep(\Theta)
  \subseteq
  \mathsf{Base}_{\mathrm{PM}}\cup\mathsf{Local}_{59}.
\]
证毕。
\end{proof}

\begin{corollary}[全文最小粒度审计口径]
\label{cor:tex59-global-atomic-audit}
本文后续每个“公理降级为定理、定理、引理、命题、推论”的证明，均按如下审计义务理解：
\[
  \boxed{
  \begin{gathered}
  \text{先拆成有限原子推理步，再检查每步规则类型、前件来源与局部证书；}\\
  \text{外部文献只允许作为背景标签，不能作为证明前件；}\\
  \text{真正的证明依赖压缩到 Gao 2025 基础域 }\cite{Gao2025PureMathZenodo}\text{ 与本文局部定义域。}
  \end{gathered}
  }
\]
\end{corollary}

\begin{Certificate}[证书：全文审计义务闭合]
\label{cert:tex59-global-atomic-audit}
证书链为
\[
\begin{aligned}
  C_1 &: \text{定理~\ref{thm:tex59-atomic-audit-existence}}
         \Rightarrow \text{有限证明可原子化},\\
  C_2 &: \text{引理~\ref{lem:tex59-external-reference-replacement}}
         \Rightarrow \text{外部引用可标签化剥离},\\
  C_3 &: \text{命题~\ref{prop:tex59-gao-base-substitution}}
         \Rightarrow \text{基础依赖压缩},\\
  C_4 &: C_1\wedge C_2\wedge C_3
         \Rightarrow
         \text{全文最小粒度审计口径成立}.
\end{aligned}
\]
\end{Certificate}

\begin{proof}[证明（基于数学符号推演逻辑的谓词因果链）]
由定理~\ref{thm:tex59-atomic-audit-existence}，每个有限证明可拆成有限原子步。
由引理~\ref{lem:tex59-external-reference-replacement}，不参与推理的外部引用可替换为
背景标签。由命题~\ref{prop:tex59-gao-base-substitution}，替换后证明依赖只剩
\[
  \mathsf{Base}_{\mathrm{PM}}\cup\mathsf{Local}_{59}.
\]
这正是推论所述的全文审计义务。证毕。
\end{proof}

\begin{remark}[审计边界]
本节审计的是本文内部形式证明链的细粒度可靠性与引用依赖压缩。
它不把外部文献中的原始深层定理重证为本文结论；外部文献只保留为问题命名、历史定位与
背景索引。本文真正使用的推理资源是 Gao 2025 年 Zenodo 基础域 \cite{Gao2025PureMathZenodo}
与本文内部显式构造的
定义、证书、引理、命题、定理和推论。
\end{remark}

%%%%%%%%%%%%%%%%%%%%%%%%%%%%%%%%%%%%%%%%%%%%%%%%%%%%%%%%%%%%
\section{逻辑基础与系统相容性：第 1、第 2、第 10 问题}
\label{sec:tex59-logic}
%%%%%%%%%%%%%%%%%%%%%%%%%%%%%%%%%%%%%%%%%%%%%%%%%%%%%%%%%%%%

\begin{remark}[核心陈述]
第 1 问题、第 2 问题与第 10 问题分别把离散--连续边界、形式相容性和离散搜索极限推到
基础层。G-化解路径不是在底空间中继续比较基数、追逐绝对相容性证明或穷举整数解，
而是把三者统一送入
\[
  \Udc
  \xrightarrow{\;\Inv\;}
  \Hcal_{\mathrm{tot}}
  \xrightarrow{\;\Gate\;}
  \{\top,\bot\}.
\]
\end{remark}

\subsection{离散--连续统一法空间}
\label{subsec:tex59-udc}

\begin{definition}[离散--连续统一法空间]
\label{def:tex59-udc}
令 \(D\) 为离散自由度空间，令 \(C\) 为连续轨迹或流形空间。给定可证书化投影
\[
  \pi_D:C\longrightarrow D_{\mathrm{cert}},
\]
定义离散--连续统一法空间为商空间
\[
  \Udc
  \DefEq
  (D\sqcup C)/\sim_{\pi_D},
\]
其中
\[
  c\sim_{\pi_D}d
  \quad\Longleftrightarrow\quad
  d=\pi_D(c).
\]
因此，连续对象在证书层以离散自由度代表出现。
\end{definition}

\begin{theorem}[连续统障碍的证书化降维]
\label{thm:tex59-ch-resolution}
设 \(H_1\) 的底层障碍为连续统基数分层障碍 \(\Barrier_{\CH}\)。若目标判断只依赖
\(\Udc\) 中的证书代表
\[
  [c]_{\Udc}=[\pi_D(c)]_{\Udc},
\]
则 \(H_1\) 在 G-化解意义下满足
\[
  \Resolve_G(H_1)=\top.
\]
\end{theorem}

\begin{Certificate}[第 1 问题的谓词因果链证书]
\label{cert:tex59-ch}
证书链为
\[
  c\in C
  \Longrightarrow
  \pi_D(c)\in D_{\mathrm{cert}}
  \Longrightarrow
  [c]_{\Udc}=[\pi_D(c)]_{\Udc}
  \Longrightarrow
  \Inv([c]_{\Udc})
  \Longrightarrow
  \Gate_{\CH}\in\{\top,\bot\}.
\]
\end{Certificate}

\begin{proof}[证明]
由定义~\ref{def:tex59-udc}，任意连续对象 \(c\in C\) 与其证书化离散投影
\(\pi_D(c)\) 在 \(\Udc\) 中同属一个等价类。因此，若目标判断只读取等价类
\([c]_{\Udc}\) 的同调或语义不变量，则该判断不再依赖 \(C\) 在底空间中的基数分层。
于是连续统问题在本文定义的 G-化解谓词中被重写为
\[
  \Gate_{\CH}(\Inv([c]_{\Udc})).
\]
这正满足定义~\ref{def:tex59-resolution} 的四个条件，故
\(\Resolve_G(H_1)=\top\)。
\end{proof}

\begin{remark}[与 Gödel--Cohen 背景的接口]
Gao 2025 年 Zenodo 基础文献 \cite{Gao2025PureMathZenodo,Godel1940,Cohen1963,Jech2003}
整理了 CH 在 ZFC 中的相对一致性与独立性背景。本文的命题不改变该集合论事实，
而是把需要在 \GFramework{} 内被审计的连续对象送入 \(\Udc\) 的证书等价类，从而使
形式化和语义化判断不再以底空间基数比较为直接门禁。
\end{remark}

\subsection{吸收元、剪枝算子与相容性门禁}
\label{subsec:tex59-consistency}

\begin{definition}[吸收元与谬误孔洞]
\label{def:tex59-bottom-fallacy}
在语义值域中加入吸收元 \(\bot\)，并规定
\[
  a+\bot=\bot,\qquad
  a\cdot\bot=\bot,\qquad
  \Gate(\bot)=\bot.
\]
若存在对象 \(x\) 与谓词 \(p\)，使得同域同条件下
\[
  p(x)\wedge \neg p(x)
\]
成立，则称 \(x\) 触发谬误孔洞，记为 \(x\in\Fallacy\)。
\end{definition}

\begin{definition}[剪枝边缘算子]
\label{def:tex59-prune}
剪枝边缘算子
\[
  E_{\mathrm{prune}}:\Ob(\mathfrak S)\longrightarrow \Ob(\mathfrak S)\cup\{\bot\}
\]
定义为
\[
  E_{\mathrm{prune}}(x)=
  \begin{cases}
  \bot, & x\in\Fallacy,\\
  x, & x\notin\Fallacy.
  \end{cases}
\]
\end{definition}

\begin{lemma}[谬误孔洞被吸收隔离]
\label{lem:tex59-fallacy-absorbed}
对任意对象 \(x\)，若 \(x\in\Fallacy\)，则
\[
  \Accept(E_{\mathrm{prune}}(x))=\bot.
\]
\end{lemma}

\begin{proof}
若 \(x\in\Fallacy\)，由定义~\ref{def:tex59-prune} 得
\[
  E_{\mathrm{prune}}(x)=\bot.
\]
再由定义~\ref{def:tex59-bottom-fallacy} 中的吸收规则，
\[
  \Gate(\bot)=\bot.
\]
接受谓词包含终端门禁，故
\[
  \Accept(E_{\mathrm{prune}}(x))=\bot.
\]
\end{proof}

\begin{theorem}[第 2 问题的相容性门禁化]
\label{thm:tex59-consistency-gate}
在定义~\ref{def:tex59-certificate-object} 至定义~\ref{def:tex59-prune} 的系统中，
\[
  \Accepted(\Vsys)\cap \Fallacy=\varnothing.
\]
因此，第 2 问题的“相容性障碍”在 G-化解意义下被门禁化为吸收元隔离与剪枝证书。
\end{theorem}

\begin{Certificate}[第 2 问题的相容性证书]
\label{cert:tex59-consistency}
证书链为
\[
  p(x)\wedge\neg p(x)
  \Longrightarrow
  x\in\Fallacy
  \Longrightarrow
  E_{\mathrm{prune}}(x)=\bot
  \Longrightarrow
  \Gate(x)=\bot
  \Longrightarrow
  x\notin\Accepted(\Vsys).
\]
\end{Certificate}

\begin{proof}
反设存在
\[
  x\in \Accepted(\Vsys)\cap\Fallacy.
\]
由 \(x\in\Fallacy\) 和引理~\ref{lem:tex59-fallacy-absorbed} 得
\[
  \Accept(E_{\mathrm{prune}}(x))=\bot.
\]
而 \(\Accepted(\Vsys)\) 的对象必须通过相同的终端门禁；触发 \(\bot\) 的对象不能被登记为
接受对象。这与 \(x\in\Accepted(\Vsys)\) 矛盾。故交集为空。由定义
~\ref{def:tex59-resolution}，该门禁链给出 \(H_2\) 的 G-化解证书。
\end{proof}

\subsection{丢番图搜索的同调门禁化}
\label{subsec:tex59-diophantine}

\begin{definition}[丢番图实例的证书投影]
\label{def:tex59-diophantine-projection}
令
\[
  F(\mathbf n)=0,\qquad \mathbf n\in\ZZ^m
\]
为丢番图实例。G-证书投影由三部分组成：
\[
  \Pi_{\Dio}(F)
  \DefEq
  \bigl(
    \Inv_{\Dio}(F),\;
    B_{\Dio}(F),\;
    T_{\Dio}(F)
  \bigr),
\]
其中 \(\Inv_{\Dio}(F)\) 是从系数、同余、次数、局部障碍和同调类抽取的不变量，
\(B_{\Dio}(F)\) 是未闭合残差槽，\(T_{\Dio}(F)\) 是终端门禁。
\end{definition}

\begin{proposition}[离散穷举被门禁替代]
\label{prop:tex59-diophantine-gate}
若丢番图实例 \(F\) 满足证书充分性条件
\[
  \Gate_{\Dio}(\Inv_{\Dio}(F),B_{\Dio}(F))
  =
  T_{\Dio}(F),
\]
则在 G-证书层中，\(F\) 的可接受性判断不需要执行底空间无穷穷举：
\[
  \Search_{\ZZ^m}(F)
  \quad\leadsto\quad
  T_{\Dio}(F).
\]
\end{proposition}

\begin{proof}
由定义~\ref{def:tex59-diophantine-projection}，\(F\) 被投影为
\(\Inv_{\Dio}(F)\)、\(B_{\Dio}(F)\) 与 \(T_{\Dio}(F)\)。命题假设说明终端门禁完全由
这些证书槽决定。因此接受判断只需要计算
\[
  \Gate_{\Dio}(\Inv_{\Dio}(F),B_{\Dio}(F)),
\]
而不需要枚举 \(\ZZ^m\) 中所有候选解。故底空间的无穷搜索障碍被重写为证书门禁。
\end{proof}

\begin{corollary}[第 10 问题的 G-化解]
\label{cor:tex59-h10}
对所有进入 \(\Pi_{\Dio}\) 且满足证书充分性条件的丢番图实例，希尔伯特第 10 问题的
离散搜索障碍满足
\[
  \Resolve_G(H_{10})=\top.
\]
\end{corollary}

\begin{proof}
由命题~\ref{prop:tex59-diophantine-gate}，不可控的底空间搜索已被重写为
\(T_{\Dio}(F)\)。这满足定义~\ref{def:tex59-resolution} 中的规范化、同调不变量、
终端门禁和失败槽条件，故结论成立。
\end{proof}

\begin{theorem}[逻辑聚类三障碍归纳闭合定理]
\label{thm:tex59-logic-cluster-induction}
令
\[
  L_0=H_1,\qquad L_1=H_2,\qquad L_2=H_{10}.
\]
若
\[
  \Resolve_G(L_0)=\top,\qquad
  \Resolve_G(L_1)=\top,\qquad
  \Resolve_G(L_2)=\top,
\]
则逻辑基础聚类证书
\[
  \Cert_{\Logic}
  \DefEq
  \bigl(
    \Cert_{\CH},\Cert_{\Con},\Cert_{\Dio}
  \bigr)
\]
闭合，并且
\[
  \forall i\in\{1,2,10\},\quad \Resolve_G(H_i)=\top.
\]
\end{theorem}

\begin{Certificate}[数学归纳法证书：逻辑聚类闭合]
\label{cert:tex59-logic-cluster-induction}
定义
\[
  P(n):\quad
  \forall j\le n,\ \Resolve_G(L_j)=\top,
  \qquad n=0,1,2.
\]
证书链为
\[
\begin{aligned}
  K_0 &: \text{定理~\ref{thm:tex59-ch-resolution}}\Rightarrow P(0),\\
  K_1 &: P(0)\wedge\text{定理~\ref{thm:tex59-consistency-gate}}\Rightarrow P(1),\\
  K_2 &: P(1)\wedge\text{推论~\ref{cor:tex59-h10}}\Rightarrow P(2).
\end{aligned}
\]
\end{Certificate}

\begin{proof}[证明（数学归纳法证书；基于数学符号推演逻辑的谓词因果链）]
基步：由定理~\ref{thm:tex59-ch-resolution}，
\[
  \Resolve_G(L_0)=\Resolve_G(H_1)=\top,
\]
故 \(P(0)\) 成立。

第一步推进：由定理~\ref{thm:tex59-consistency-gate}，
\[
  \Resolve_G(L_1)=\Resolve_G(H_2)=\top.
\]
与 \(P(0)\) 合取，得到
\[
  \forall j\le1,\ \Resolve_G(L_j)=\top,
\]
即 \(P(1)\)。

第二步推进：由推论~\ref{cor:tex59-h10}，
\[
  \Resolve_G(L_2)=\Resolve_G(H_{10})=\top.
\]
与 \(P(1)\) 合取，得到 \(P(2)\)。由于
\[
  \{L_0,L_1,L_2\}=\{H_1,H_2,H_{10}\},
\]
逻辑聚类闭合。
\end{proof}

\begin{remark}[逻辑基础聚类的闭合]
第 1、第 2、第 10 问题在本文中共享同一逻辑结构：
\[
  \text{底空间不可控}
  \Longrightarrow
  \text{证书代表}
  \Longrightarrow
  \text{吸收/门禁}
  \Longrightarrow
  \text{可审计接受}.
\]
因此逻辑基础聚类 \(\kappa^{-1}(\Logic)=\{1,2,10\}\) 已形成闭合证书。
\end{remark}

%%%%%%%%%%%%%%%%%%%%%%%%%%%%%%%%%%%%%%%%%%%%%%%%%%%%%%%%%%%%
\section{物理学的彻底公理化：第 6 问题}
\label{sec:tex59-physics}
%%%%%%%%%%%%%%%%%%%%%%%%%%%%%%%%%%%%%%%%%%%%%%%%%%%%%%%%%%%%

\begin{remark}[核心陈述]
第 6 问题要求物理学的公理化。G-化解路径是把经典力学、热力学、量子力学、
相互作用扇区、量子扇区和抽象约束边界等异构对象，统一写入
\[
  \Phi_{\Vsys}:\Phys_{\mathrm{raw}}\longrightarrow \NF_{\Vsys}
\]
的绝对元级正规型，并由显式证书替代经验公式的最终裁决权。
\end{remark}

\subsection{七阶接收映射与物理扇区}
\label{subsec:tex59-physics-map}

\begin{definition}[物理扇区簇]
\label{def:tex59-physical-sectors}
定义物理扇区簇
\[
  \Sect_{\Phys}
  \DefEq
  \{
  s_1,s_2,s_3,s_4,s_5,s_6,s_7,s_8
  \}.
\]
每个扇区 \(s\in\Sect_{\Phys}\) 具有状态空间 \(X_s\)、局部定律或约束 \(L_s\)、观测接口
\(O_s\) 与证书接口 \(C_s\)。
\end{definition}

\begin{definition}[七阶接收映射]
\label{def:tex59-seven-map}
七阶接收映射定义为复合
\[
  \Phi_{\Vsys}
  \DefEq
  \Phi_7\circ\Phi_6\circ\Phi_5\circ\Phi_4\circ\Phi_3\circ\Phi_2\circ\Phi_1,
\]
其中
\[
\begin{array}{ll}
\Phi_1:& \text{原始物理记录进入对象槽};\\
\Phi_2:& \text{单位、尺度、坐标和边界归一化};\\
\Phi_3:& \text{提取谓词族与约束链};\\
\Phi_4:& \text{生成同伦或同调不变量};\\
\Phi_5:& \text{登记审计日志与失败槽};\\
\Phi_6:& \text{执行终端门禁};\\
\Phi_7:& \text{输出绝对元级正规型 } \NF_{\Vsys}.
\end{array}
\]
\end{definition}

\begin{theorem}[跨域物理量的正规型统一]
\label{thm:tex59-physical-normal-form}
若每个物理扇区 \(s\in\Sect_{\Phys}\) 均存在保持谓词依赖的映射
\[
  \Phi_{\Vsys,s}:X_s\longrightarrow \NF_{\Vsys},
\]
则所有扇区中的可接受物理陈述都可统一表示为同一门禁判断：
\[
  \Accept_{\Phys}(x)
  =
  \Gate_{\Vsys}(\Phi_{\Vsys,s}(x)).
\]
\end{theorem}

\begin{Certificate}[第 6 问题的跨域证书]
\label{cert:tex59-physics}
对任意 \(s\in\Sect_{\Phys}\)，证书链为
\[
  \begin{aligned}
  x\in X_s
  &\Longrightarrow
  (L_s,O_s,C_s)
  \Longrightarrow
  \Phi_{\Vsys,s}(x)\in\NF_{\Vsys}\\
  &\Longrightarrow
  \Gate_{\Vsys}(\Phi_{\Vsys,s}(x))
  \Longrightarrow
  \Accept_{\Phys}(x).
  \end{aligned}
\]
\end{Certificate}

\begin{proof}[证明]
由定义~\ref{def:tex59-physical-sectors}，每个扇区对象 \(x\in X_s\) 都附带局部定律、
观测接口与证书接口。由定义~\ref{def:tex59-seven-map}，\(\Phi_{\Vsys,s}\) 逐阶执行对象登记、
归一化、谓词抽取、不变量生成、审计、门禁和正规型输出。若该映射保持谓词依赖，则
扇区内的物理陈述在 \(\NF_{\Vsys}\) 中具有同一类型的证书代表。接受判断只读取
\(\Gate_{\Vsys}\) 的输出，因此所有扇区被统一为同一门禁形式。证毕。
\end{proof}

\subsection{证书化取代经验公式的终端裁决}
\label{subsec:tex59-physics-certificate}

\begin{proposition}[经验公式的证书降级]
\label{prop:tex59-empirical-downgrade}
设某物理扇区中有经验公式
\[
  y=F_{\theta}(x)
\]
与证书判断
\[
  \Cert_s(x,y).
\]
若
\[
  \Cert_s(x,F_{\theta}(x))=\top
  \Longrightarrow
  \Gate_{\Vsys}(\Phi_{\Vsys,s}(x,F_{\theta}(x)))=\top,
\]
且
\[
  \Cert_s(x,F_{\theta}(x))=\bot
  \Longrightarrow
  \Gate_{\Vsys}(\Phi_{\Vsys,s}(x,F_{\theta}(x)))=\bot,
\]
则经验公式在终端裁决层被降级为提案生成器，最终裁决权属于证书门禁。
\end{proposition}

\begin{proof}
命题的两个蕴含式说明：经验公式给出的输出 \(F_{\theta}(x)\) 是否可接受，完全由证书
\(\Cert_s\) 与门禁 \(\Gate_{\Vsys}\) 决定。若证书为真，则门禁接受；若证书失败，则门禁拒绝。
因此 \(F_{\theta}\) 不再是终端裁决规则，只是产生待审计候选的映射。
\end{proof}

\begin{theorem}[第 6 问题的 G-化解]
\label{thm:tex59-h6}
在定义~\ref{def:tex59-physical-sectors} 至命题~\ref{prop:tex59-empirical-downgrade} 的条件下，
希尔伯特第 6 问题的物理公理化障碍满足
\[
  \Resolve_G(H_6)=\top.
\]
\end{theorem}

\begin{proof}
第 6 问题的障碍是异构物理理论缺少统一公理接口。定理
~\ref{thm:tex59-physical-normal-form} 已将所有扇区可接受陈述统一为
\(\Gate_{\Vsys}(\Phi_{\Vsys,s}(x))\)。命题~\ref{prop:tex59-empirical-downgrade} 进一步说明经验公式
只能产生候选，不能绕过证书门禁。故异构物理量、公理接口和终端裁决均被写入同一
证书系统，满足定义~\ref{def:tex59-resolution}，从而 \(\Resolve_G(H_6)=\top\)。
\end{proof}

\begin{theorem}[物理扇区接收链归纳闭合定理]
\label{thm:tex59-physics-sector-induction}
枚举物理扇区为
\[
  s_0,\ldots,s_7
  =
  \mathrm{CM},\mathrm{Thermo},\mathrm{EM},\mathrm{QM},
  s_5,s_6,s_7,s_8.
\]
若每个扇区均存在保持谓词依赖的接收映射
\[
  \Phi_{\Vsys,s_k}:X_{s_k}\to\NF_{\Vsys},
  \qquad k=0,\ldots,7,
\]
则
\[
  \forall k\le7,\quad
  \Accept_{\Phys}(x_k)=\Gate_{\Vsys}(\Phi_{\Vsys,s_k}(x_k)).
\]
因此，第 6 问题的物理公理化接口在有限扇区族上归纳闭合。
\end{theorem}

\begin{Certificate}[数学归纳法证书：物理扇区逐项接收]
\label{cert:tex59-physics-sector-induction}
定义
\[
  P(n):\quad
  \forall k\le n,\ 
  \Accept_{\Phys}(x_k)=\Gate_{\Vsys}(\Phi_{\Vsys,s_k}(x_k)).
\]
证书链为
\[
\begin{aligned}
  S_0 &: \Phi_{\Vsys,s_0}\text{ 保持谓词依赖}\Rightarrow P(0),\\
  S_{n+1} &: P(n)\wedge
  \left[
  \Phi_{\Vsys,s_{n+1}}\text{ 保持谓词依赖}
  \right]
  \Rightarrow P(n+1).
\end{aligned}
\]
\end{Certificate}

\begin{proof}[证明（数学归纳法证书；基于数学符号推演逻辑的谓词因果链）]
基步：对 \(s_0\)，由定理~\ref{thm:tex59-physical-normal-form} 的单扇区实例，
\[
  \Accept_{\Phys}(x_0)=\Gate_{\Vsys}(\Phi_{\Vsys,s_0}(x_0)).
\]
故 \(P(0)\) 成立。

归纳步：假设 \(P(n)\) 成立。对 \(s_{n+1}\)，由于
\(\Phi_{\Vsys,s_{n+1}}\) 保持谓词依赖，再次调用定理
~\ref{thm:tex59-physical-normal-form}，得到
\[
  \Accept_{\Phys}(x_{n+1})
  =
  \Gate_{\Vsys}(\Phi_{\Vsys,s_{n+1}}(x_{n+1})).
\]
与 \(P(n)\) 合取即得 \(P(n+1)\)。有限归纳至 \(n=7\)，物理扇区接收链闭合。
\end{proof}

\begin{remark}[与 Zenodo 基础的接口]
Gao 2025 年 Zenodo 基础文献 \cite{Gao2025PureMathZenodo}
给出本文所需的集合、范畴、同调与代数基础背景。
本文第 6 问题的结论不调用应用数学分卷作为前提，而是在本节内部把跨域物理扇区的共同接口
抽象为 \(\Phi_{\Vsys}\) 与 \(\Gate_{\Vsys}\)。
\end{remark}

%%%%%%%%%%%%%%%%%%%%%%%%%%%%%%%%%%%%%%%%%%%%%%%%%%%%%%%%%%%%
\section{PDE、变分法与边界奇点：第 19 至第 23 问题}
\label{sec:tex59-pde}
%%%%%%%%%%%%%%%%%%%%%%%%%%%%%%%%%%%%%%%%%%%%%%%%%%%%%%%%%%%%

\begin{remark}[核心陈述]
第 19 至第 23 问题集中体现分析学中的正则性、边界值、线性微分方程单值化和变分法
发展障碍。G-化解路径不是在奇点附近强行逐点解析延拓，而是通过语义微分、语义积分、
协变 PDE 隧穿命中与回写收紧算子，把局部发散转写为全局不变量与可修复边界残差。
\end{remark}

\subsection{PDE 残差对象与语义微积分}
\label{subsec:tex59-pde-object}

\begin{definition}[PDE 残差对象]
\label{def:tex59-pde-object}
一个 PDE 残差对象是六元组
\[
  \mathfrak p
  =
  (X,\Lcal,\Bcal,u_0,E,\Gamma),
\]
其中 \(X\) 是函数或状态空间，\(\Lcal[u]=0\) 是内部方程，\(\Bcal[u]=0\) 是边界或初值条件，
\(u_0\) 是初始候选，\(E\) 是残差能量，
\[
  E(u)
  \DefEq
  \frac12\|\Lcal[u]\|^2+\frac{\lambda}{2}\|\Bcal[u]\|^2,
\]
\(\Gamma\) 是可接受阈值或命中集合。
\end{definition}

\begin{definition}[语义微分与语义积分]
\label{def:tex59-semantic-calculus}
语义微分与语义积分为
\[
  d_{\Vsys}:\NF_{\Vsys}^{k}\longrightarrow \NF_{\Vsys}^{k+1},
  \qquad
  \int_{\Vsys}:\NF_{\Vsys}^{k}\longrightarrow \Inv_{\Vsys},
\]
并满足证书化 Stokes 关系：
\[
  \Cert_{\mathrm{Stokes}}(\omega,\sigma)=\top
  \Longrightarrow
  \int_{\Vsys,\partial\sigma}\omega
  =
  \int_{\Vsys,\sigma}d_{\Vsys}\omega.
\]
\end{definition}

\begin{lemma}[残差命中等价于终端门禁]
\label{lem:tex59-residual-hit}
若
\[
  \Gamma_{\eps}\DefEq\{u\in X:E(u)\le\eps\}
\]
且门禁定义为
\[
  \Gate_{\PDE}(u)=\top
  \quad\Longleftrightarrow\quad
  u\in\Gamma_{\eps},
\]
则 PDE 求解任务在证书层等价于命中 \(\Gamma_{\eps}\)。
\end{lemma}

\begin{proof}
由 \(\Gamma_{\eps}\) 与 \(\Gate_{\PDE}\) 的定义，
\[
  \Gate_{\PDE}(u)=\top
  \Longleftrightarrow
  E(u)\le\eps
  \Longleftrightarrow
  u\in\Gamma_{\eps}.
\]
故等价成立。
\end{proof}

\subsection{协变 PDE 隧穿命中}
\label{subsec:tex59-pde-tunnel}

\begin{theorem}[协变 PDE 隧穿命中]
\label{thm:tex59-pde-tunnel}
设 \(\mathfrak p=(X,\Lcal,\Bcal,u_0,E,\Gamma_{\eps})\) 为 PDE 残差对象。若存在同伦路径
\[
  H:[0,1]\longrightarrow \NF_{\Vsys}
\]
满足
\[
  H(0)=\Phi_{\Vsys}(u_0),\qquad
  H(1)=\Phi_{\Vsys}(u_{\ast}),
\]
且
\[
  E(u_{\ast})\le\eps,\qquad
  \Cert_{\mathrm{cov}}(H)=\top,
\]
则 PDE 求解任务可在 G-证书层被替代为协变同伦命中：
\[
  \PDE(\Lcal,\Bcal)
  \leadsto
  \Hit_{\Vsys}(\Gamma_{\eps}).
\]
\end{theorem}

\begin{Certificate}[PDE 隧穿命中证书]
\label{cert:tex59-pde-tunnel}
证书链为
\[
  u_0
  \Longrightarrow
  \Phi_{\Vsys}(u_0)
  \xRightarrow{\,H\,}
  \Phi_{\Vsys}(u_{\ast})
  \Longrightarrow
  E(u_{\ast})\le\eps
  \Longrightarrow
  u_{\ast}\in\Gamma_{\eps}
  \Longrightarrow
  \Gate_{\PDE}(u_{\ast})=\top.
\]
\end{Certificate}

\begin{proof}[证明]
由假设，\(H\) 是从初始候选正规型到终端候选正规型的协变同伦路径，且
\(\Cert_{\mathrm{cov}}(H)=\top\) 保证路径上的坐标、边界和语义不变量兼容。又由
\(E(u_{\ast})\le\eps\)，可得 \(u_{\ast}\in\Gamma_{\eps}\)。引理~\ref{lem:tex59-residual-hit}
给出
\[
  \Gate_{\PDE}(u_{\ast})=\top.
\]
因此原 PDE 求解任务在证书层被替代为命中集合 \(\Gamma_{\eps}\) 的协变路径问题。
\end{proof}

\subsection{边界残差回写与同伦修复}
\label{subsec:tex59-boundary-repair}

\begin{definition}[边界回写队列与收紧算子]
\label{def:tex59-backwrite-tighten}
若边界残差
\[
  r_{\partial}(u)\DefEq \Bcal[u]
\]
未完全闭合，则登记
\[
  r_{\partial}(u)\in \BackwriteQueue(\Vsys).
\]
收紧算子
\[
  E_{\mathrm{tighten}}:\BackwriteQueue(\Vsys)\longrightarrow \NF_{\Vsys}
\]
把残差槽修复为新的有效自由度或吸收失败：
\[
  E_{\mathrm{tighten}}(r)=
  \begin{cases}
  \delta r, & \Cert_{\mathrm{repair}}(r)=\top,\\
  \bot, & \Cert_{\mathrm{repair}}(r)=\bot.
  \end{cases}
\]
\end{definition}

\begin{proposition}[一般边界值问题的残差修复]
\label{prop:tex59-bvp-repair}
若边界残差 \(r_{\partial}(u)\) 满足
\[
  \Cert_{\mathrm{repair}}(r_{\partial}(u))=\top,
\]
则存在修复候选
\[
  u^+=u+\delta r
\]
使得
\[
  E(u^+)\le E(u)
\]
或 \(u^+\) 被登记为新的有效自由度继续进入门禁链。
\end{proposition}

\begin{proof}
由定义~\ref{def:tex59-backwrite-tighten}，证书为真时
\[
  E_{\mathrm{tighten}}(r_{\partial}(u))=\delta r.
\]
于是可以构造 \(u^+=u+\delta r\)。收紧算子的含义就是只接受降低残差或形成有效自由度的
修复项：若降低残差，则 \(E(u^+)\le E(u)\)；若无法立即降低但可作为有效自由度，则进入
下一轮门禁链；若两者都不成立，\(\Cert_{\mathrm{repair}}\) 不应为真。故命题成立。
\end{proof}

\begin{theorem}[PDE 五问题残差--边界--同伦归纳闭合定理]
\label{thm:tex59-pde-cluster-induction}
令
\[
  P_0=H_{19},\quad P_1=H_{20},\quad P_2=H_{21},\quad P_3=H_{22},\quad P_4=H_{23}.
\]
若以下五类证书分别成立：
\[
\begin{aligned}
  R_{19}&:\quad \Reg\text{ 被 }E(u)\le\eps\text{ 与 }\Cert_{\mathrm{cov}}\text{ 门禁化},\\
  R_{20}&:\quad \BVP\text{ 残差进入 }\BackwriteQueue(\Vsys)\text{ 并由 }E_{\mathrm{tighten}}\text{ 修复},\\
  R_{21}&:\quad \text{单值化障碍由 }\int_{\Vsys}\text{ 的全局不变量读取},\\
  R_{22}&:\quad \text{解析延拓障碍由 }d_{\Vsys}\text{ 与证书化 Stokes 关系替代},\\
  R_{23}&:\quad \text{变分发展障碍由 }E,\Gamma_{\eps},H\text{ 的命中链登记},
\end{aligned}
\]
则
\[
  \forall k\le4,\quad \Resolve_G(P_k)=\top.
\]
\end{theorem}

\begin{Certificate}[数学归纳法证书：PDE 聚类逐题闭合]
\label{cert:tex59-pde-cluster-induction}
定义
\[
  Q(n):\quad \forall k\le n,\ \Resolve_G(P_k)=\top.
\]
证书链为
\[
\begin{aligned}
  R_{19}&\Rightarrow Q(0),\\
  Q(0)\wedge R_{20}&\Rightarrow Q(1),\\
  Q(1)\wedge R_{21}&\Rightarrow Q(2),\\
  Q(2)\wedge R_{22}&\Rightarrow Q(3),\\
  Q(3)\wedge R_{23}&\Rightarrow Q(4).
\end{aligned}
\]
\end{Certificate}

\begin{proof}[证明（数学归纳法证书；基于数学符号推演逻辑的谓词因果链）]
基步：\(R_{19}\) 将正则性障碍重写为
\[
  E(u_\ast)\le\eps
  \Longrightarrow
  u_\ast\in\Gamma_{\eps}
  \Longrightarrow
  \Gate_{\PDE}(u_\ast)=\top,
\]
其中最后一步由引理~\ref{lem:tex59-residual-hit} 给出，故 \(\Resolve_G(P_0)=\top\)，即
\(Q(0)\)。

归纳推进依次如下。若 \(Q(0)\) 成立，\(R_{20}\) 由命题~\ref{prop:tex59-bvp-repair}
把边界残差送入回写--收紧链，因此 \(\Resolve_G(P_1)=\top\)，得 \(Q(1)\)。
若 \(Q(1)\) 成立，\(R_{21}\) 将单值化障碍交给全局不变量门禁，得
\(\Resolve_G(P_2)=\top\)，故 \(Q(2)\)。若 \(Q(2)\) 成立，\(R_{22}\) 使用
\(d_{\Vsys}\) 与 \(\int_{\Vsys}\) 的证书化 Stokes 关系替代逐点解析延拓，得 \(Q(3)\)。
若 \(Q(3)\) 成立，\(R_{23}\) 把变分任务写为残差能量、命中集合与同伦路径的三元链，
得 \(Q(4)\)。于是五个 PDE 类问题均闭合。
\end{proof}

\begin{theorem}[第 19 至第 23 问题的 G-化解]
\label{thm:tex59-h19-h23}
在定义~\ref{def:tex59-pde-object} 至命题~\ref{prop:tex59-bvp-repair} 的条件下，
\[
  \Resolve_G(H_{19})=\Resolve_G(H_{20})=\Resolve_G(H_{21})
  =\Resolve_G(H_{22})=\Resolve_G(H_{23})=\top.
\]
\end{theorem}

\begin{proof}
第 19 问题的正则性障碍被 \(E(u)\le\eps\) 与 \(\Cert_{\mathrm{cov}}\) 重写为命中集合中的
正规型可接受性。第 20 问题的一般边界值障碍由命题~\ref{prop:tex59-bvp-repair} 重写为
回写队列与收紧算子。第 21 与第 22 问题中的单值化和解析延拓障碍，在本文口径下由
\(d_{\Vsys}\)、\(\int_{\Vsys}\) 与证书化 Stokes 关系读取全局不变量，而不是在奇点处
逐点推进。第 23 问题的变分发展障碍由残差能量 \(E\)、命中集合 \(\Gamma_{\eps}\) 与同伦路径
\(H\) 统一表达。上述每一项均满足定义~\ref{def:tex59-resolution} 的规范化、同调不变量、
门禁和失败槽条件，故五个问题的 G-化解谓词均为真。
\end{proof}

\begin{remark}[本节自足性与 PDE 背景]
本节的残差命中、协变隧穿、边界回写与语义积分解释口径，均由
定义~\ref{def:tex59-pde-object} 至定理~\ref{thm:tex59-h19-h23} 在本稿内部闭合；
基础背景只调用 Gao 2025 年 Zenodo 文献的形式语言背景 \cite{Gao2025PureMathZenodo}，
不调用其他内部 PDE
工作稿作为证明前提。外部 PDE、弱形式与泛函分析背景可参考
\cite{Evans2010PDE,BrennerScott2008FEM,Brezis2011FunctionalAnalysis}。
\end{remark}

%%%%%%%%%%%%%%%%%%%%%%%%%%%%%%%%%%%%%%%%%%%%%%%%%%%%%%%%%%%%
\section{代数几何、拓扑不变量与空间结构：第 3、第 4、第 15、第 16、第 18 问题}
\label{sec:tex59-geometry}
%%%%%%%%%%%%%%%%%%%%%%%%%%%%%%%%%%%%%%%%%%%%%%%%%%%%%%%%%%%%

\begin{remark}[核心陈述]
几何类问题的共同障碍是底空间交点、曲线曲面拓扑类型、空间铺砌和几何公理的局部复杂度。
G-化解路径是把底空间对象送入总同调上卷：
\[
  \TotHom:
  \Ob(\mathrm{Geom})
  \longrightarrow
  \bigoplus_{k\ge0}H_k(-)\oplus H^k(-),
\]
并把交点、配置和铺砌约束转写为不变量运算、语义图谱与 Hash 清单比对。
\end{remark}

\subsection{总同调上卷与交点证书}
\label{subsec:tex59-total-homology}

\begin{definition}[总同调上卷]
\label{def:tex59-total-homology}
对几何对象 \(Y\)，定义总同调上卷为
\[
  \TotHom(Y)
  \DefEq
  \bigoplus_{k\ge0}
  \bigl(H_k(Y;\ZZ)\oplus H^k(Y;\ZZ)\bigr),
\]
并令交点证书为
\[
  \Cert_{\cap}(A,B)
  \DefEq
  \bigl([A],[B],[A]\GCap[B]\bigr),
  \]
其中 \([A]\) 与 \([B]\) 是同调或上同调类，\(\GCap\) 表示交叉或配对运算。
\end{definition}

\begin{theorem}[交点理论的同调化]
\label{thm:tex59-intersection-homology}
若几何交点问题的接受判断只依赖
\[
  [A]\GCap[B]\in\TotHom(Y),
\]
则底空间逐点交点枚举可被同调配对证书替代。
\end{theorem}

\begin{Certificate}[第 15 问题的交点证书]
\label{cert:tex59-schubert}
证书链为
\[
  (A,B)\subset Y
  \Longrightarrow
  ([A],[B])\in\TotHom(Y)
  \Longrightarrow
  [A]\GCap[B]
  \Longrightarrow
  \Cert_{\cap}(A,B)
  \Longrightarrow
  \Gate_{\cap}\in\{\top,\bot\}.
\]
\end{Certificate}

\begin{proof}[证明]
由定义~\ref{def:tex59-total-homology}，几何对象 \(A,B\) 被投影为同调或上同调类。若接受
判断只依赖 \([A]\GCap[B]\)，则底空间中交点的具体枚举不是门禁输入。交点理论因此被
重写为 \(\TotHom(Y)\) 中的代数配对运算。故命题成立。
\end{proof}

\subsection{几何结构的语义图谱化}
\label{subsec:tex59-geometry-graph}

\begin{definition}[语义图谱与 Hash 清单]
\label{def:tex59-semantic-graph}
对几何配置 \(Y\)，定义语义图谱
\[
  \Graph_G(Y)=(V_Y,E_Y,\ell_Y)
\]
其中 \(V_Y\) 是对象节点，\(E_Y\) 是约束边，\(\ell_Y\) 是同调、边界、谓词和证书标签。
定义 Hash 清单
\[
  \Hash_G(Y)\DefEq \Hash(V_Y,E_Y,\ell_Y,\TotHom(Y)).
\]
\end{definition}

\begin{proposition}[高维几何校验降维为清单比对]
\label{prop:tex59-hash-check}
若两个几何配置 \(Y_1,Y_2\) 满足
\[
  \Hash_G(Y_1)=\Hash_G(Y_2)
\]
且 Hash 生成过程保持谓词标签和同调标签，则它们在 G-证书层不可区分：
\[
  \Gate_G(Y_1)=\Gate_G(Y_2).
\]
\end{proposition}

\begin{proof}
\(\Hash_G\) 的输入包含节点、约束边、谓词标签与总同调标签。若
\(\Hash_G(Y_1)=\Hash_G(Y_2)\)，且生成过程保持这些标签，则门禁读取的全部证书槽相同。
因此两个配置给出相同门禁输出。
\end{proof}

\begin{theorem}[几何五问题同调--图谱归纳闭合定理]
\label{thm:tex59-geometry-induction}
令
\[
  G_0=H_3,\quad G_1=H_4,\quad G_2=H_{15},\quad G_3=H_{16},\quad G_4=H_{18}.
\]
若每个几何问题均可进入
\[
  Y\longmapsto \bigl(\TotHom(Y),\Graph_G(Y),\Hash_G(Y)\bigr)
\]
的证书化表示，则
\[
  \forall k\le4,\quad \Resolve_G(G_k)=\top.
\]
\end{theorem}

\begin{Certificate}[数学归纳法证书：几何聚类闭合]
\label{cert:tex59-geometry-induction}
定义
\[
  P(n):\quad \forall k\le n,\ \Resolve_G(G_k)=\top.
\]
证书链为
\[
\begin{aligned}
  J_0 &: \Graph_G\text{ 记录几何基础谓词}\Rightarrow P(0),\\
  J_1 &: P(0)\wedge\Graph_G\text{ 记录直线/测地约束}\Rightarrow P(1),\\
  J_2 &: P(1)\wedge\Cert_{\cap}\text{ 成立}\Rightarrow P(2),\\
  J_3 &: P(2)\wedge\TotHom\text{ 记录曲线曲面拓扑型}\Rightarrow P(3),\\
  J_4 &: P(3)\wedge\Hash_G\text{ 记录铺砌配置}\Rightarrow P(4).
\end{aligned}
\]
\end{Certificate}

\begin{proof}[证明（数学归纳法证书；基于数学符号推演逻辑的谓词因果链）]
基步：第 3 问题的几何基础谓词被登记为 \(\Graph_G(Y)\) 的节点、边与标签，门禁读取该
图谱，故 \(P(0)\) 成立。

归纳步依次展开。第 4 问题的直线、测地和距离约束同样是 \(\Graph_G(Y)\) 的边标签，
因此由 \(P(0)\) 推得 \(P(1)\)。第 15 问题由定理~\ref{thm:tex59-intersection-homology}
把交点枚举降为
\[
  [A]\GCap[B]\in\TotHom(Y),
\]
故 \(P(2)\)。第 16 问题的曲线曲面拓扑型由 \(\TotHom(Y)\) 与语义标签共同记录，故
\(P(3)\)。第 18 问题的铺砌配置由命题~\ref{prop:tex59-hash-check} 降为
\(\Hash_G\) 清单比对，故 \(P(4)\)。于是几何聚类闭合。
\end{proof}

\begin{theorem}[第 3、第 4、第 15、第 16、第 18 问题的 G-化解]
\label{thm:tex59-geometry-cluster}
在总同调上卷、交点证书与语义图谱化条件下，
\[
  \Resolve_G(H_3)=\Resolve_G(H_4)=\Resolve_G(H_{15})
  =\Resolve_G(H_{16})=\Resolve_G(H_{18})=\top.
\]
\end{theorem}

\begin{proof}
第 3 与第 4 问题涉及几何基础与测地/直线结构，在本文中被登记为
\(\Graph_G(Y)\) 的节点、边和谓词标签。第 15 问题由定理
~\ref{thm:tex59-intersection-homology} 化为同调配对证书。第 16 问题中的曲线曲面拓扑类型
被 \(\TotHom(Y)\) 与标签化配置记录。第 18 问题的铺砌和空间配置障碍由
\(\Hash_G(Y)\) 的清单比对给出门禁。五个问题均被转写为有限可审计的不变量和图谱门禁，
满足定义~\ref{def:tex59-resolution}，故结论成立。
\end{proof}

\begin{remark}[代数拓扑背景]
总同调、上同调、交叉配对和同伦等价的基础语言可参考
\cite{Hatcher2002AT,Weibel1994HomologicalAlgebra}；范畴化的对象--态射组织方式可参考
\cite{MacLane1998Categories}。本文仅以 Gao 2025 年 Zenodo 文献
\cite{Gao2025PureMathZenodo} 的代数拓扑与范畴基础作为统一背景；
本节的证书接口在定义~\ref{def:tex59-total-homology} 以后独立给出。
\end{remark}

%%%%%%%%%%%%%%%%%%%%%%%%%%%%%%%%%%%%%%%%%%%%%%%%%%%%%%%%%%%%
\section{数论、解析延拓与函数零点：第 7、第 8、第 9、第 11、第 12、第 13、第 14、第 17 问题}
\label{sec:tex59-number-theory}
%%%%%%%%%%%%%%%%%%%%%%%%%%%%%%%%%%%%%%%%%%%%%%%%%%%%%%%%%%%%

\begin{remark}[核心陈述]
数论与函数论类问题的共同障碍是无限逼近、零点分布、解析延拓、代数方程表达复杂度、
不变量有限生成和正性表示。G-化解路径是把“寻找全部零点或全部生成元”的底空间任务，
转写为绝对元级正规型中的边界锚定、有限证书链和终端门禁。
\end{remark}

\subsection{解析对象的证书锚定}
\label{subsec:tex59-zero-anchor}

\begin{definition}[解析对象与零点锚]
\label{def:tex59-zero-anchor}
令 \(f:\Omega\subset\CC\to\CC\) 为解析或亚纯对象。定义零点锚为
\[
  \Anchor_G(f)
  \DefEq
  \bigl(
    \partial_G f,\;
    \Inv_{\Zero}(f),\;
    \Cert_{\Zero}(f)
  \bigr),
\]
其中 \(\partial_G f\) 是边界或临界线数据，\(\Inv_{\Zero}(f)\) 是保零点结构的不变量，
\(\Cert_{\Zero}(f)\) 是可审计证书。
\end{definition}

\begin{theorem}[零点逼近的证书锚定替代]
\label{thm:tex59-zero-anchor}
若函数零点问题的目标判断只依赖
\[
  \Anchor_G(f)
\]
且
\[
  \Gate_{\Zero}(f)
  =
  \Gate_{\Zero}(\Anchor_G(f)),
\]
则无穷逼近过程在 G-证书层被替代为边界锚定门禁。
\end{theorem}

\begin{Certificate}[第 8 问题的零点锚证书]
\label{cert:tex59-riemann-anchor}
对以黎曼猜想为代表的零点问题，证书链写为
\[
  f
  \Longrightarrow
  \partial_G f
  \Longrightarrow
  \Inv_{\Zero}(f)
  \Longrightarrow
  \Cert_{\Zero}(f)
  \Longrightarrow
  \Gate_{\Zero}(\Anchor_G(f)).
\]
\end{Certificate}

\begin{proof}[证明]
命题假设给出门禁等式
\[
  \Gate_{\Zero}(f)
  =
  \Gate_{\Zero}(\Anchor_G(f)).
\]
因此，门禁层不读取无穷逼近轨迹本身，而读取 \(\partial_G f\)、\(\Inv_{\Zero}(f)\) 与
\(\Cert_{\Zero}(f)\)。故底空间“逐点寻找零点”的任务被证书锚替代。
\end{proof}

\subsection{有限生成与代数不变量证书}
\label{subsec:tex59-finite-generation}

\begin{definition}[有限证书生成系]
\label{def:tex59-finite-generation}
设 \(\Ical\) 为不变量系统。若存在有限集合
\[
  S=\{s_1,\ldots,s_r\}
\]
与有限规则族
\[
  \Rcal=\{\rho_1,\ldots,\rho_m\}
\]
使得任意被接受的不变量 \(a\in\Ical\) 均可由 \(S\) 经 \(\Rcal\) 的有限复合生成，则称
\((S,\Rcal)\) 为 \(\Ical\) 的有限证书生成系。
\end{definition}

\begin{lemma}[有限生成的归纳证书]
\label{lem:tex59-finite-generation}
若 \((S,\Rcal)\) 是 \(\Ical\) 的有限证书生成系，则任意被接受的不变量
\(a\in\Ical\) 都存在有限证书
\[
  \Cert_{\Gen}(a).
\]
\end{lemma}

\begin{proof}
对 \(a\) 的生成长度作数学归纳。若生成长度为 \(0\)，则 \(a\in S\)，证书为生成元登记证书。
假设长度不超过 \(n\) 的对象均有有限证书。若 \(a\) 的生成长度为 \(n+1\)，则
\[
  a=\rho_j(a_1,\ldots,a_k)
\]
其中每个 \(a_l\) 的生成长度不超过 \(n\)。由归纳假设，每个 \(a_l\) 有有限证书。
再附加规则 \(\rho_j\) 的应用记录，即得 \(a\) 的有限证书。故归纳完成。
\end{proof}

\begin{proposition}[第 14 问题的有限生成门禁化]
\label{prop:tex59-h14}
若第 14 问题所涉不变量系统在 G-证书层存在有限证书生成系 \((S,\Rcal)\)，则
\[
  \Resolve_G(H_{14})=\top.
\]
\end{proposition}

\begin{proof}
由引理~\ref{lem:tex59-finite-generation}，每个被接受不变量都有有限证书。终端门禁只需
检查生成元登记、规则应用和审计日志，不需要在底空间中处理无穷生成过程。故满足
定义~\ref{def:tex59-resolution}。
\end{proof}

\begin{theorem}[数论八问题锚定--有限生成归纳闭合定理]
\label{thm:tex59-number-induction}
令
\[
\begin{gathered}
  N_0=H_7,\quad N_1=H_8,\quad N_2=H_9,\quad N_3=H_{11},\\
  N_4=H_{12},\quad N_5=H_{13},\quad N_6=H_{14},\quad N_7=H_{17}.
\end{gathered}
\]
若每个 \(N_k\) 的核心障碍均可写为零点锚、代数不变量或有限证书生成系的门禁判断，
则
\[
  \forall k\le7,\quad \Resolve_G(N_k)=\top.
\]
\end{theorem}

\begin{Certificate}[数学归纳法证书：数论聚类闭合]
\label{cert:tex59-number-induction}
定义
\[
  P(n):\quad \forall k\le n,\ \Resolve_G(N_k)=\top.
\]
证书链为
\[
\begin{aligned}
  Z_0 &: \Anchor_G\text{ 处理超越数边界}\Rightarrow P(0),\\
  Z_1 &: P(0)\wedge\Anchor_G\text{ 处理零点临界线}\Rightarrow P(1),\\
  Z_2 &: P(1)\wedge\Inv_{\Zero}\text{ 处理互反律不变量}\Rightarrow P(2),\\
  Z_3 &: P(2)\wedge\Inv_{\Zero}\text{ 处理二次型约束}\Rightarrow P(3),\\
  Z_4 &: P(3)\wedge\Inv_{\Zero}\text{ 处理 Abel 域结构}\Rightarrow P(4),\\
  Z_5 &: P(4)\wedge\NF\text{ 处理高次表达复杂度}\Rightarrow P(5),\\
  Z_6 &: P(5)\wedge\text{命题~\ref{prop:tex59-h14}}\Rightarrow P(6),\\
  Z_7 &: P(6)\wedge\text{正性表示有限规则证书}\Rightarrow P(7).
\end{aligned}
\]
\end{Certificate}

\begin{proof}[证明（数学归纳法证书；基于数学符号推演逻辑的谓词因果链）]
基步：第 7 问题的超越数判断在本文口径中进入 \(\Anchor_G\) 与边界不变量门禁，故
\(\Resolve_G(N_0)=\top\)，即 \(P(0)\)。

归纳推进：第 8 问题由定理~\ref{thm:tex59-zero-anchor} 把零点逼近替换为零点锚门禁，
故从 \(P(0)\) 得 \(P(1)\)。第 9、第 11、第 12 问题分别把互反律、二次型与 Abel 域
结构写入 \(\Inv_{\Zero}\) 或相应代数不变量，故依次得到 \(P(2),P(3),P(4)\)。
第 13 问题把高次方程表达复杂度登记为正规型长度与规则应用证书，故得 \(P(5)\)。
第 14 问题由命题~\ref{prop:tex59-h14} 的有限生成门禁化给出 \(P(6)\)。第 17 问题把
正性表示写为有限代数规则与终端门禁，故得 \(P(7)\)。于是数论聚类闭合。
\end{proof}

\begin{theorem}[数论与函数论聚类的 G-化解]
\label{thm:tex59-number-cluster}
在零点锚定、有限证书生成和终端门禁条件下，
\[
\begin{gathered}
  \Resolve_G(H_7)=\Resolve_G(H_8)=\Resolve_G(H_9)=
  \Resolve_G(H_{11})=\Resolve_G(H_{12})=\top,\\
  \Resolve_G(H_{13})=\Resolve_G(H_{14})=\Resolve_G(H_{17})=\top.
\end{gathered}
\]
\end{theorem}

\begin{proof}
第 7、第 8、第 9、第 11、第 12 问题涉及超越数、零点、互反律、二次型和 Abel 域等
数论结构；在本文中，它们的可审计判断均被写入 \(\Anchor_G(f)\)、\(\Inv_{\Zero}\) 或
相应代数不变量。第 13 问题的高次方程表达复杂度被重写为正规型表达长度与规则应用
证书。第 14 问题由命题~\ref{prop:tex59-h14} 处理。第 17 问题的正性表示被登记为
有限代数规则与证书门禁。故该聚类中所有问题均满足 G-化解定义。
\end{proof}

\begin{remark}[数论聚类的正确读法]
本节把“无穷逼近、全局零点搜索、无限生成”统一改写为证书锚定和有限生成门禁。
外部数论和代数几何的具体定理仍保留其本学科内容；本文登记的是进入
\GFramework{} 后的证书化判定结构。
\end{remark}

%%%%%%%%%%%%%%%%%%%%%%%%%%%%%%%%%%%%%%%%%%%%%%%%%%%%%%%%%%%%
\section[扩展世纪难题与计算物理解耦]{扩展世纪难题与计算物理解耦}
\label{sec:tex59-extended-obstructions}
%%%%%%%%%%%%%%%%%%%%%%%%%%%%%%%%%%%%%%%%%%%%%%%%%%%%%%%%%%%%

\begin{remark}[核心陈述]
本章把六个与希尔伯特谱系相邻或延伸的基础障碍纳入同一证书系统：
\[
  \mathcal E_6
  =
  \{
  \mathsf H_5,\mathsf{Hodge},\mathsf{YMGap},
  \mathsf{NS},\mathsf{BSD},\mathsf{GodelDecouple}
  \}.
\]
这里的结论均采用 G-证书层口径：本文证明的是这些障碍在 \(\Udc\)、总同调语义场、
吸收元、回写队列、收紧算子与终端门禁下的\textbf{证书化化解链}。外部经典问题
仍保留其本学科的原始表述与证明标准；本文给出的是进入 \GFramework{} 后的
条件化、可审计、可失败闭合的形式逻辑实现。Hilbert 第 5 问题的传统背景可参考
\cite{Gleason1952HilbertFifth,MontgomeryZippin1955TopologicalGroups}；Hodge、Yang--Mills、
Navier--Stokes 与 BSD 的经典背景分别参考
\cite{Voisin2002Hodge,GriffithsHarris1978,JaffeWitten2000YM,Fefferman2000NS,Temam2001NS,
BirchSwinnertonDyer1965,Silverman2009AEC}；Gödel 不完备性背景参考
\cite{Godel1931}。
\end{remark}

\subsection{对象域、公理降级与扩展证书谓词}
\label{subsec:tex59-extended-domain}

\begin{definition}[扩展障碍的 G-化解谓词]
\label{def:tex59-extended-resolution}
对任意扩展障碍 \(E\in\mathcal E_6\)，定义
\[
  \Resolve_G^{\mathrm{ext}}(E)=\top
\]
当且仅当存在有限谓词因果链
\[
  \Barrier(E)=\chi_0
  \Longrightarrow
  \chi_1
  \Longrightarrow
  \cdots
  \Longrightarrow
  \chi_m=
  \Gate_E(\Cert_E)=\top,
\]
并且每一步均有局部证书
\[
  \mathfrak S_G\vdash \chi_j\Rightarrow\chi_{j+1}.
\]
若某一步触发 \(\bot\)，则该障碍不被接受，而是进入失败吸收或回写修复链。
\end{definition}

\begin{theorem}[公理降级：扩展六障碍有限归纳闭合]
\label{thm:tex59-extended-axiom-downgrade}
令
\[
  E_0=\mathsf H_5,\quad
  E_1=\mathsf{Hodge},\quad
  E_2=\mathsf{YMGap},\quad
  E_3=\mathsf{NS},\quad
  E_4=\mathsf{BSD},\quad
  E_5=\mathsf{GodelDecouple}.
\]
若对每个 \(k=0,\ldots,5\) 都存在局部证书
\[
  C_k:\quad \Resolve_G^{\mathrm{ext}}(E_k)=\top,
\]
则扩展证书包
\[
  \mathfrak C_{\mathcal E}
  \DefEq
  (C_0,C_1,C_2,C_3,C_4,C_5)
\]
有限闭合，并且
\[
  \forall E\in\mathcal E_6,\quad
  \Resolve_G^{\mathrm{ext}}(E)=\top.
\]
\end{theorem}

\begin{Certificate}[数学归纳法证书：扩展六障碍逐项闭合]
\label{cert:tex59-extended-axiom-downgrade}
定义
\[
  P(n):\quad
  \forall k\le n,\ \Resolve_G^{\mathrm{ext}}(E_k)=\top,
  \qquad n=0,\ldots,5.
\]
证书链为
\[
  C_0\Rightarrow P(0),
  \qquad
  P(n)\wedge C_{n+1}\Rightarrow P(n+1).
\]
\end{Certificate}

\begin{proof}[证明（数学归纳法证书；基于数学符号推演逻辑的谓词因果链）]
基步：由 \(C_0\) 得
\[
  \Resolve_G^{\mathrm{ext}}(E_0)=\top,
\]
故 \(P(0)\) 成立。

归纳步：假设 \(P(n)\) 成立，即
\[
  \forall k\le n,\ \Resolve_G^{\mathrm{ext}}(E_k)=\top.
\]
又由 \(C_{n+1}\) 得
\[
  \Resolve_G^{\mathrm{ext}}(E_{n+1})=\top.
\]
二者合取得到
\[
  \forall k\le n+1,\ \Resolve_G^{\mathrm{ext}}(E_k)=\top,
\]
即 \(P(n+1)\)。有限归纳至 \(n=5\)，扩展六障碍闭合。
\end{proof}

\subsection{希尔伯特第 5 问题：连续群局部可微性前提的证书化绕开}
\label{subsec:tex59-hilbert5}

\begin{definition}[拓扑群演化的离散--连续证书投影]
\label{def:tex59-h5-projection}
令 \(G_{\mathrm{top}}\) 为作用在空间 \(X\) 上的局部紧拓扑群，令
\[
  g:[0,1]\to G_{\mathrm{top}}
\]
为连续群演化。定义证书投影
\[
  \Pi_{\mathrm{grp}}(g)
  =
  \bigl(
    \pi_D(g(t_0)),\ldots,\pi_D(g(t_N)),
    \Inv_{\mathrm{grp}}(g),
    \Cert_{\mathrm{grp}}(g)
  \bigr),
\]
其中 \(\pi_D\) 是进入离散格点自由度空间的投影，\(\Inv_{\mathrm{grp}}\) 是同调或表示不变量，
\(\Cert_{\mathrm{grp}}\) 是保持群乘法、单位元与逆元的显式证书。
\end{definition}

\begin{lemma}[局部可微性不进入终端门禁]
\label{lem:tex59-h5-no-smooth-gate}
若第 5 问题中的接受判断只读取
\[
  \Pi_{\mathrm{grp}}(g)
  \quad\text{与}\quad
  [g]_{\Udc},
\]
则局部坐标图上的 \(C^1\) 或解析性假设不是终端门禁输入。
\end{lemma}

\begin{Certificate}[证书：群演化到格点证书]
\label{cert:tex59-h5-no-smooth-gate}
证书链为
\[
  g(t)\in G_{\mathrm{top}}
  \Longrightarrow
  \pi_D(g(t))\in D_{\mathrm{cert}}
  \Longrightarrow
  [g(t)]_{\Udc}
  \Longrightarrow
  \Inv_{\mathrm{grp}}(g)
  \Longrightarrow
  \Gate_{\mathrm{grp}}(\Pi_{\mathrm{grp}}(g)).
\]
\end{Certificate}

\begin{proof}[证明（基于数学符号推演逻辑的谓词因果链）]
由定义~\ref{def:tex59-h5-projection}，终端门禁读取的是离散证书序列、
群结构证书与同调不变量。由定义~\ref{def:tex59-udc}，连续演化 \(g(t)\) 在
\(\Udc\) 中与其离散证书投影同类：
\[
  [g(t)]_{\Udc}=[\pi_D(g(t))]_{\Udc}.
\]
因此门禁链中没有需要读取局部坐标图可微性的槽。局部可微性可以作为外部经典
Lie 群理论中的结构条件存在，但在 G-证书门禁中被替换为
\[
  \pi_D+\Inv_{\mathrm{grp}}+\Cert_{\mathrm{grp}}.
\]
证毕。
\end{proof}

\begin{theorem}[第 5 问题的 G-证书化化解]
\label{thm:tex59-h5-resolution}
在定义~\ref{def:tex59-h5-projection} 的证书投影条件下，
\[
  \Resolve_G^{\mathrm{ext}}(\mathsf H_5)=\top.
\]
\end{theorem}

\begin{Certificate}[第 5 问题证书]
\label{cert:tex59-h5-resolution}
证书登记为
\[
  \Barrier(\mathsf H_5)
  =
  \text{连续群局部可微性前提}
  \Longrightarrow
  \Pi_{\mathrm{grp}}(g)
  \Longrightarrow
  \Gate_{\mathrm{grp}}=\top
  \Longrightarrow
  \Resolve_G^{\mathrm{ext}}(\mathsf H_5)=\top.
\]
\end{Certificate}

\begin{proof}[证明（基于数学符号推演逻辑的谓词因果链）]
第 5 问题的底层障碍是：拓扑群是否必须具有可微 Lie 群结构。由引理
~\ref{lem:tex59-h5-no-smooth-gate}，G-证书层不把局部可微性作为终端输入，而把连续群演化
送入 \(\Udc\) 中的离散格点证书与同调不变量。于是原障碍链被改写为
\[
  \text{局部可微性要求}
  \Longrightarrow
  \text{格点证书投影}
  \Longrightarrow
  \text{群结构门禁}.
\]
该链满足定义~\ref{def:tex59-extended-resolution}，故结论成立。
\end{proof}

\subsection{Hodge 猜想：拓扑同调类到有限正规代数闭链的证书化实现}
\label{subsec:tex59-hodge}

\begin{definition}[Hodge 类的 G-代数闭链证书]
\label{def:tex59-hodge-certificate}
令 \(Y\) 为光滑复射影簇，令
\[
  \eta\in H^{2p}(Y,\QQ)\cap H^{p,p}(Y)
\]
为 Hodge 类。定义 G-代数闭链证书为
\[
  \Cert_{\mathrm{Hdg}}(\eta)
  =
  \bigl(
    Z_1,\ldots,Z_r;\ q_1,\ldots,q_r;\ \Theta_{\mathrm{alg}};\ T_{\mathrm{Hdg}}
  \bigr),
\]
其中 \(Z_j\) 是代数闭链，\(q_j\in\QQ\)，\(\Theta_{\mathrm{alg}}\) 证明
\[
  \eta
  =
  \sum_{j=1}^r q_j[Z_j]
  \quad\text{在 G-证书同调商空间中成立},
\]
且 \(T_{\mathrm{Hdg}}\) 是终端门禁。
\end{definition}

\begin{proposition}[四链子丛到有限代数闭链]
\label{prop:tex59-hodge-four-chain}
若拓扑不变量 \(\eta\) 可经四链子丛分解为
\[
  \Rcal\oplus\Ecal\oplus\Dcal\oplus\Ucal
  \longrightarrow
  \Rstar
  \longrightarrow
  \sum_{j=1}^r q_j[Z_j],
\]
且每一步均有同伦保持证书，则 \(\eta\) 在 G-证书层具有有限代数闭链表示。
\end{proposition}

\begin{Certificate}[证书：Hodge 类降维到代数闭链]
\label{cert:tex59-hodge-four-chain}
证书链为
\[
  \eta
  \Longrightarrow
  (\Rcal,\Ecal,\Dcal,\Ucal)
  \Longrightarrow
  \Rstar
  \Longrightarrow
  \sum_j q_j[Z_j]
  \Longrightarrow
  \Cert_{\mathrm{Hdg}}(\eta)
  \Longrightarrow
  T_{\mathrm{Hdg}}(\eta)=\top.
\]
\end{Certificate}

\begin{proof}[证明（基于数学符号推演逻辑的谓词因果链）]
四链子丛分解把 \(\eta\) 的几何来源、边界约束、偏差残差与更新修复分别登记为
\[
  \Rcal,\Ecal,\Dcal,\Ucal.
\]
若这些分量经同伦保持证书共同收束到 \(\Rstar\)，并进一步给出有限代数闭链组合
\(\sum_j q_j[Z_j]\)，则 \(\Theta_{\mathrm{alg}}\) 记录了从拓扑类到代数闭链的有限推理。
终端门禁只需检查有限 \(Z_j\)、系数 \(q_j\) 与同调等式证书，故命题成立。
\end{proof}

\begin{theorem}[Hodge 猜想的 G-证书化实现]
\label{thm:tex59-hodge-resolution}
对所有满足定义~\ref{def:tex59-hodge-certificate} 与命题~\ref{prop:tex59-hodge-four-chain}
条件的 Hodge 类，
\[
  \Resolve_G^{\mathrm{ext}}(\mathsf{Hodge})=\top.
\]
\end{theorem}

\begin{proof}[证明（基于数学符号推演逻辑的谓词因果链）]
由命题~\ref{prop:tex59-hodge-four-chain}，
\[
  \eta
  \Longrightarrow
  \sum_jq_j[Z_j]
  \Longrightarrow
  \Cert_{\mathrm{Hdg}}(\eta).
\]
由定义~\ref{def:tex59-hodge-certificate}，\(\Cert_{\mathrm{Hdg}}(\eta)\) 包含终端门禁
\(T_{\mathrm{Hdg}}\)。若
\[
  T_{\mathrm{Hdg}}(\eta)=\top,
\]
则 \(\eta\) 在 G-证书同调商空间中已被确权为有限正规代数闭链组合。故该 Hodge 障碍在
G-证书层被化解。
\end{proof}

\begin{remark}[Hodge 口径边界]
本节不把有限应用证书误写成对全部经典 Hodge 猜想的无条件外部证明。严格结论是：
凡能进入四链子丛分解并由 \(\Theta_{\mathrm{alg}}\) 签发的 Hodge 类，在 G-证书层已经完成
从拓扑同调类到代数闭链的有限确权。
\end{remark}

\subsection{Yang--Mills 与质量缺口：最小可证书化自由度格点}
\label{subsec:tex59-yang-mills}

\begin{definition}[总同调语义规范场与最小自由度格点]
\label{def:tex59-ym-min-grid}
令 \(\mathcal A/G\) 表示规范联络模空间的语义代表，定义总同调语义场
\[
  \mathfrak F_{\mathrm{YM}}
  \DefEq
  \TotHom(\mathcal A/G)\oplus \NF_{\Vsys}.
\]
若存在不可继续向下拆分的最小可证书化自由度格点
\[
  \lambda_{\min}\in\mathfrak F_{\mathrm{YM}},
  \qquad
  \Cert_{\min}(\lambda_{\min})=\top,
\]
并存在常数 \(\Delta_G>0\) 使任意非真空激发态 \(\psi\) 满足
\[
  E_{\Vsys}(\psi)\ge \Delta_G,
\]
则称 \(\mathfrak F_{\mathrm{YM}}\) 具有 G-语义质量缺口 \(\Delta_G\)。
\end{definition}

\begin{lemma}[最小正规型推出正能量下界]
\label{lem:tex59-ym-min-gap}
若 \(\lambda_{\min}\) 是不可继续细分的最小可证书化自由度格点，且每个非真空激发态
必须至少占用一个 \(\lambda_{\min}\) 单元，则存在 \(\Delta_G>0\) 使
\[
  E_{\Vsys}(\psi)\ge\Delta_G.
\]
\end{lemma}

\begin{Certificate}[证书：最小格点到能量下界]
\label{cert:tex59-ym-min-gap}
证书链为
\[
  \psi\ne\Omega
  \Longrightarrow
  \Supp(\psi)\supseteq\{\lambda_{\min}\}
  \Longrightarrow
  E_{\Vsys}(\lambda_{\min})=\Delta_G>0
  \Longrightarrow
  E_{\Vsys}(\psi)\ge\Delta_G.
\]
\end{Certificate}

\begin{proof}[证明（基于数学符号推演逻辑的谓词因果链）]
非真空态 \(\psi\ne\Omega\) 必须在证书空间中留下非空支撑。由最小性，非空支撑至少包含
一个不可再拆分的格点 \(\lambda_{\min}\)。令
\[
  \Delta_G=E_{\Vsys}(\lambda_{\min}).
\]
最小格点可证书化且非空，因此 \(\Delta_G>0\)。能量泛函对支撑单元非负可加或单调，
故
\[
  E_{\Vsys}(\psi)\ge E_{\Vsys}(\lambda_{\min})=\Delta_G.
\]
\end{proof}

\begin{theorem}[Yang--Mills 质量缺口的 G-语义证书]
\label{thm:tex59-ym-resolution}
若定义~\ref{def:tex59-ym-min-grid} 的最小自由度格点与正下界证书成立，则
\[
  \Resolve_G^{\mathrm{ext}}(\mathsf{YMGap})=\top.
\]
\end{theorem}

\begin{proof}[证明（基于数学符号推演逻辑的谓词因果链）]
由引理~\ref{lem:tex59-ym-min-gap}，非真空激发态满足
\[
  E_{\Vsys}(\psi)\ge\Delta_G>0.
\]
这正是 G-语义质量缺口的门禁表达。由于该下界由最小可证书化自由度格点与总同调语义场
给出，Yang--Mills 质量缺口障碍在 G-证书层被化解。
\end{proof}

\subsection{Navier--Stokes：协变 PDE 隧穿与语义全局确定性}
\label{subsec:tex59-navier-stokes}

\begin{definition}[Navier--Stokes 语义残差对象]
\label{def:tex59-ns-object}
令 \(u\) 为速度场，\(p\) 为压力，\(\nu>0\) 为粘性系数。定义
\[
  \Res_{\mathrm{NS}}(u,p)
  =
  \partial_tu+(u\cdot\nabla)u-\nu\Delta u+\nabla p-f,
  \qquad
  \nabla\cdot u=0.
\]
定义 Navier--Stokes 语义残差能量
\[
  E_{\mathrm{NS}}(u,p)
  =
  \frac12\|\Res_{\mathrm{NS}}(u,p)\|^2
  +
  \frac{\lambda}{2}\|\nabla\cdot u\|^2.
\]
\end{definition}

\begin{proposition}[奇点跨越的协变命中证书]
\label{prop:tex59-ns-tunnel}
若存在协变同伦路径
\[
  H_{\mathrm{NS}}:[0,1]\to\NF_{\Vsys}
\]
满足
\[
  H_{\mathrm{NS}}(0)=\Phi_{\Vsys}(u_0,p_0),
  \qquad
  H_{\mathrm{NS}}(1)=\Phi_{\Vsys}(u_\ast,p_\ast),
\]
且
\[
  E_{\mathrm{NS}}(u_\ast,p_\ast)\le\eps,
  \qquad
  \Cert_{\mathrm{cov}}(H_{\mathrm{NS}})=\top,
\]
则局部奇点不再作为终端失败，而被写入协变 PDE 隧穿命中证书。
\end{proposition}

\begin{Certificate}[证书：NS 残差到语义确定性]
\label{cert:tex59-ns-tunnel}
证书链为
\[
  \begin{aligned}
  (u_0,p_0)
  &\Longrightarrow
  \Phi_{\Vsys}(u_0,p_0)
  \xRightarrow{H_{\mathrm{NS}}}
  \Phi_{\Vsys}(u_\ast,p_\ast),\\
  \Phi_{\Vsys}(u_\ast,p_\ast)
  &\Longrightarrow
  E_{\mathrm{NS}}(u_\ast,p_\ast)\le\eps
  \Longrightarrow
  \int_{\Vsys}\Phi_{\Vsys}(u_\ast,p_\ast)\ \text{确定}.
  \end{aligned}
\]
\end{Certificate}

\begin{proof}[证明（基于数学符号推演逻辑的谓词因果链）]
\(\Cert_{\mathrm{cov}}(H_{\mathrm{NS}})=\top\) 保证路径在语义规范场中保持协变一致。
终端残差满足
\[
  E_{\mathrm{NS}}(u_\ast,p_\ast)\le\eps,
\]
故由引理~\ref{lem:tex59-residual-hit} 的同型实例，终端态命中低残差集合。语义积分
\(\int_{\Vsys}\) 读取的是正规型终端态的不变量，而不是奇点邻域内逐点轨迹。因此局部奇点
被从终端失败改写为路径中可隧穿的低维障碍。
\end{proof}

\begin{theorem}[Navier--Stokes 的 G-语义全局确定性证书]
\label{thm:tex59-ns-resolution}
在命题~\ref{prop:tex59-ns-tunnel} 的条件下，
\[
  \Resolve_G^{\mathrm{ext}}(\mathsf{NS})=\top.
\]
\end{theorem}

\begin{proof}[证明（基于数学符号推演逻辑的谓词因果链）]
Navier--Stokes 经典障碍是全局存在性与平滑性。本文的 G-语义口径把该障碍重写为
\[
  \text{局部奇点}
  \Longrightarrow
  H_{\mathrm{NS}}\text{ 协变隧穿}
  \Longrightarrow
  E_{\mathrm{NS}}\le\eps
  \Longrightarrow
  \int_{\Vsys}\text{ 全局确定}.
\]
由命题~\ref{prop:tex59-ns-tunnel}，这条链有证书闭合，故
\(\Resolve_G^{\mathrm{ext}}(\mathsf{NS})=\top\)。
\end{proof}

\begin{remark}[Navier--Stokes 口径边界]
本节证明的是 G-语义积分层的全局确定性与低残差命中，不等同于直接给出经典
Navier--Stokes 问题中所有三维光滑初值的无条件 \(C^\infty\) 全局正则性证明。
\end{remark}

\subsection{BSD 猜想：L-函数零点阶与回写残差维度}
\label{subsec:tex59-bsd}

\begin{definition}[BSD 的 G-秩证书]
\label{def:tex59-bsd-rank}
令 \(E/\QQ\) 为椭圆曲线，令
\[
  r_L(E)=\operatorname{ord}_{s=1}L(E,s)
\]
为 \(L\)-函数在 \(s=1\) 的零点阶。定义 G-秩为
\[
  \begin{aligned}
  \Rank_G(E)
  \DefEq
  \max\{\dim W:\ &
  W\subseteq\BackwriteQueue(\Vsys),\\
  &
  E_{\mathrm{tighten}}(W)\subseteq D_{\mathrm{cert}},\\
  &
  \Cert_{\mathrm{BSD}}(W)=\top\}.
  \end{aligned}
\]
这里 \(W\) 表示能够被同伦收紧算子成功转化为新格点自由度的信息残差子空间。
\end{definition}

\begin{proposition}[BSD 锚定等价命题]
\label{prop:tex59-bsd-anchor}
若椭圆曲线 \(E/\QQ\) 满足 BSD 锚定证书
\[
  \Anchor_{\mathrm{BSD}}(E):
  \quad
  r_L(E)=\Rank_G(E),
\]
则 \(L\)-函数零点阶与回写队列中可收紧残差的最高维度在 G-证书层等价。
\end{proposition}

\begin{Certificate}[证书：解析零点阶到回写维度]
\label{cert:tex59-bsd-anchor}
证书链为
\[
  L(E,s)
  \Longrightarrow
  r_L(E)
  \Longrightarrow
  \Dcal(E)
  \Longrightarrow
  \BackwriteQueue(\Vsys)
  \xrightarrow{E_{\mathrm{tighten}}}
  D_{\mathrm{cert}}
  \Longrightarrow
  \Rank_G(E).
\]
\end{Certificate}

\begin{proof}[证明（基于数学符号推演逻辑的谓词因果链）]
解析侧给出零点阶 \(r_L(E)\)。G-代数侧把未闭合的信息偏差登记为
\(\Dcal(E)\)，并送入 \(\BackwriteQueue(\Vsys)\)。若某个残差子空间 \(W\) 可由
\[
  E_{\mathrm{tighten}}(W)\subseteq D_{\mathrm{cert}}
\]
成功收紧，则它贡献新的可证书化格点自由度。定义~\ref{def:tex59-bsd-rank} 取所有此类
子空间的最大维度。若锚定证书声明
\[
  r_L(E)=\Rank_G(E),
\]
则二者在 G-证书层等价。
\end{proof}

\begin{theorem}[BSD 的 G-证书化化解]
\label{thm:tex59-bsd-resolution}
对满足 \(\Anchor_{\mathrm{BSD}}(E)\) 的椭圆曲线实例，
\[
  \Resolve_G^{\mathrm{ext}}(\mathsf{BSD})=\top.
\]
\end{theorem}

\begin{proof}[证明（基于数学符号推演逻辑的谓词因果链）]
由命题~\ref{prop:tex59-bsd-anchor}，BSD 的核心对应关系被写成
\[
  r_L(E)=\Rank_G(E).
\]
右端 \(\Rank_G(E)\) 由有限可审计的回写队列、收紧算子和格点自由度证书给出。因此解析
零点阶与代数秩的对应关系在 G-证书层被门禁化，满足定义
~\ref{def:tex59-extended-resolution}。
\end{proof}

\subsection{Gödel 不完备性在计算物理中的吸收解耦}
\label{subsec:tex59-godel-decouple}

\begin{definition}[不可判定死结与计算物理语义孔洞]
\label{def:tex59-godel-knot}
若某计算物理语义对象 \(x\) 触发
\[
  p_i(x)\wedge\neg p_i(x)
\]
或触发无法在当前证书上下文中闭合的自指谓词，则称 \(x\) 落入 Gödel 式死结，记为
\[
  x\in\mathsf{GodelKnot}.
\]
规定
\[
  \mathsf{GodelKnot}\subseteq\Fallacy,
\]
并沿用吸收规则
\[
  a+\bot=\bot,\qquad
  E_{\mathrm{prune}}(x)=\bot\quad (x\in\Fallacy).
\]
\end{definition}

\begin{lemma}[Gödel 式死结不向全局接受集传播]
\label{lem:tex59-godel-no-propagation}
若 \(x\in\mathsf{GodelKnot}\)，则
\[
  x\notin\Accepted(\Vsys).
\]
\end{lemma}

\begin{Certificate}[证书：死结到吸收元]
\label{cert:tex59-godel-no-propagation}
证书链为
\[
  \begin{aligned}
  x\in\mathsf{GodelKnot}
  &\Longrightarrow
  x\in\Fallacy
  \Longrightarrow
  E_{\mathrm{prune}}(x)=\bot,\\
  E_{\mathrm{prune}}(x)=\bot
  &\Longrightarrow
  \Gate(x)=\bot
  \Longrightarrow
  x\notin\Accepted(\Vsys).
  \end{aligned}
\]
\end{Certificate}

\begin{proof}[证明（基于数学符号推演逻辑的谓词因果链）]
由定义~\ref{def:tex59-godel-knot}，
\[
  x\in\mathsf{GodelKnot}\Rightarrow x\in\Fallacy.
\]
由定义~\ref{def:tex59-prune}，
\[
  E_{\mathrm{prune}}(x)=\bot.
\]
再由吸收元规则，终端门禁返回 \(\bot\)。接受集只包含门禁为 \(\top\) 的对象，故
\[
  x\notin\Accepted(\Vsys).
\]
\end{proof}

\begin{theorem}[计算物理中的 Gödel 解耦证书]
\label{thm:tex59-godel-decouple-resolution}
在定义~\ref{def:tex59-godel-knot} 的吸收规则下，
\[
  \Resolve_G^{\mathrm{ext}}(\mathsf{GodelDecouple})=\top.
\]
\end{theorem}

\begin{proof}[证明（基于数学符号推演逻辑的谓词因果链）]
传统不完备性障碍说明足够强的形式系统存在系统内不可完全自证的命题。本文不否认该事实，
而是在计算物理证书形式系统引入解耦规则：
\[
  \text{自指死结}
  \Longrightarrow
  \mathsf{GodelKnot}
  \Longrightarrow
  \bot
  \Longrightarrow
  \text{剪枝隔离}.
\]
由引理~\ref{lem:tex59-godel-no-propagation}，该死结不会进入
\(\Accepted(\Vsys)\)，因此不会导致全局递归崩溃。该链满足扩展化解定义，故结论成立。
\end{proof}

\subsection{扩展六障碍总证书}
\label{subsec:tex59-extended-total}

\begin{theorem}[扩展六障碍总证书]
\label{thm:tex59-extended-total}
在本章各局部条件成立时，
\[
  \forall E\in\mathcal E_6,\quad
  \Resolve_G^{\mathrm{ext}}(E)=\top.
\]
\end{theorem}

\begin{Certificate}[总证书：六障碍局部证书合成]
\label{cert:tex59-extended-total}
证书包为
\[
\begin{aligned}
  C_0&=\text{定理~\ref{thm:tex59-h5-resolution}},&
  C_1&=\text{定理~\ref{thm:tex59-hodge-resolution}},\\
  C_2&=\text{定理~\ref{thm:tex59-ym-resolution}},&
  C_3&=\text{定理~\ref{thm:tex59-ns-resolution}},\\
  C_4&=\text{定理~\ref{thm:tex59-bsd-resolution}},&
  C_5&=\text{定理~\ref{thm:tex59-godel-decouple-resolution}}.
\end{aligned}
\]
\end{Certificate}

\begin{proof}[证明（数学归纳法证书；基于数学符号推演逻辑的谓词因果链）]
由证书包中六个局部定理，得到
\[
  \forall k=0,\ldots,5,\quad
  \Resolve_G^{\mathrm{ext}}(E_k)=\top.
\]
将这些 \(C_k\) 代入定理~\ref{thm:tex59-extended-axiom-downgrade}，得到
\[
  \forall E\in\mathcal E_6,\quad
  \Resolve_G^{\mathrm{ext}}(E)=\top.
\]
\end{proof}

\begin{remark}[扩展问题的统一口径]
六个扩展障碍共享同一形式：
\[
  \begin{aligned}
  \text{底空间前提或无穷障碍}
  &\Longrightarrow
  \text{证书投影}
  \Longrightarrow
  \text{同调/语义不变量},\\
  \text{同调/语义不变量}
  &\Longrightarrow
  \text{门禁或吸收}
  \Longrightarrow
  \text{G-证书层化解}.
  \end{aligned}
\]
因此，第 5 问题、Hodge、Yang--Mills、Navier--Stokes、BSD 与 Gödel 解耦并不是六套互不相干
的叙述，而是同一证书化范式在连续群、代数几何、规范场、流体 PDE、椭圆曲线与计算物理
自指障碍上的六个实例。
\end{remark}

%%%%%%%%%%%%%%%%%%%%%%%%%%%%%%%%%%%%%%%%%%%%%%%%%%%%%%%%%%%%
\section{跨学科四大科学障碍的 抽象验证结构证书化化解}
\label{sec:tex59-crossdisciplinary-obstructions}
%%%%%%%%%%%%%%%%%%%%%%%%%%%%%%%%%%%%%%%%%%%%%%%%%%%%%%%%%%%%

\begin{remark}[章节目标与引用口径]
本章把前文的 \(\Udc\)、总同调语义场、吸收元 \(\bot\)、剪枝边缘算子、
回写队列与终端门禁继续推广到四类跨学科世纪障碍：
\[
  \mathcal X_4
  =
  \{\mathsf{HALT},\mathsf{QGMeasure},\mathsf{BHInfo},\mathsf{SymbolGround}\}.
\]
它们分别对应图灵停机问题、量子引力统一与测量问题、黑洞信息悖论、符号接地与符号系统未接地输出根源。
本章的严格口径不是否定图灵不可判定性，也不是给出经典物理开放问题的无条件外部证明；
而是在 \(\GFramework{}\)/\(\Vsys\) 证书形式系统中证明：
\[
  \boxed{
  \text{底空间不可预测或不可闭合的问题，可被改写为证书层的投影、门禁、吸收与回写问题。}
  }
\]
停机问题背景参考 Turing 的可计算数论文献 \cite{Turing1936Computable}；
量子测量与量子引力背景参考 von Neumann、DeWitt 与 Kiefer
\cite{VonNeumann1955Quantum,DeWitt1967QuantumGravity,Kiefer2012QuantumGravity}；
黑洞信息悖论背景参考 Hawking \cite{Hawking1976Breakdown}；
符号接地问题参考 Harnad \cite{Harnad1990SymbolGrounding}。
基础背景只引用 Gao 2025 年 Zenodo 文献 \cite{Gao2025PureMathZenodo} 中的 G-代数与形式基础背景；
本章使用的证书语义、吸收元、剪枝和接地门禁均在本章内部重新定义。
\end{remark}

\subsection{对象域、公理降级与四题证书谓词}
\label{subsec:tex59-cross-objects}

\begin{definition}[四类跨学科障碍对象]
\label{def:tex59-cross-objects}
定义四类跨学科障碍为
\[
\begin{aligned}
  X_0&=\mathsf{HALT},\\
  X_1&=\mathsf{QGMeasure},\\
  X_2&=\mathsf{BHInfo},\\
  X_3&=\mathsf{SymbolGround}.
\end{aligned}
\]
它们的底空间障碍分别记为
\[
  \Barrier(X_k),\qquad k=0,1,2,3.
\]
定义 \(\Vsys\) 证书投影
\[
  \Pi_{\Vsys}^{X_k}:\Barrier(X_k)\longrightarrow\NF_{\Vsys}
\]
以及终端证书门禁
\[
  \Gate_{X_k}:\NF_{\Vsys}\longrightarrow\{\top,\bot,0,1\}.
\]
\end{definition}

\begin{definition}[跨学科 G-化解谓词]
\label{def:tex59-cross-resolution}
称 \(X_k\) 在 \(\Vsys\) 证书层被 G-化解，记为
\[
  \Resolve_G^{\mathrm{sci}}(X_k)=\top,
\]
当且仅当存在证书链
\[
  \Cert_{X_k}
  =
  \left(
    \Pi_{\Vsys}^{X_k},\
    \Inv_{X_k},\
    \Gate_{X_k},\
    E_{\mathrm{prune}},\
    \BackwriteQueue,\
    E_{\mathrm{tighten}}
  \right)
\]
使得
\[
  \Barrier(X_k)
  \xrightarrow{\Pi_{\Vsys}^{X_k}}
  \NF_{\Vsys}
  \xrightarrow{\Inv_{X_k}}
  \text{同调/语义不变量}
  \xrightarrow{\Gate_{X_k}}
  \{\top,\bot,0,1\},
\]
并且失败分支只能落入
\[
  \bot,\qquad \Excluded,\qquad \BackwriteQueue(\Vsys)
\]
三类闭合出口之一。
\end{definition}

\begin{theorem}[公理降级：四类科学障碍证书归纳闭合]
\label{thm:tex59-cross-induction}
若对每个 \(k=0,1,2,3\) 都已经构造出证书链 \(\Cert_{X_k}\)，并满足
\[
  \Gate_{X_k}(\Cert_{X_k})\in\{\top,\bot,0,1\},
\]
且所有失败分支都被吸收、剪枝或回写，则
\[
  \forall X\in\mathcal X_4,\qquad
  \Resolve_G^{\mathrm{sci}}(X)=\top.
\]
\end{theorem}

\begin{Certificate}[数学归纳法证书：四类科学障碍逐项闭合]
\label{cert:tex59-cross-induction}
定义谓词
\[
  P(n):
  \quad
  \forall k\le n,\quad
  \Resolve_G^{\mathrm{sci}}(X_k)=\top.
\]
归纳证书为
\[
  \pi_{\mathrm{sci}}
  =
  \left(
    P(0),\
    \forall n\in\{0,1,2\},\ P(n)\Rightarrow P(n+1)
  \right).
\]
\end{Certificate}

\begin{proof}[证明（数学归纳法证书；基于数学符号推演逻辑的谓词因果链）]
基步 \(P(0)\)：由第~\ref{subsec:tex59-halting-gate} 节构造的停机证书链，
\[
  \Cert_{X_0}
  \Rightarrow
  \Gate_{X_0}\in\{\top,\bot,0,1\}
  \Rightarrow
  \Resolve_G^{\mathrm{sci}}(X_0)=\top.
\]

归纳步：假设 \(P(n)\) 成立，并考虑 \(X_{n+1}\)。
由本章后续各节分别给出的量子测量、黑洞信息或符号接地证书链，有
\[
  \Cert_{X_{n+1}}
  \Rightarrow
  \Gate_{X_{n+1}}\in\{\top,\bot,0,1\}.
\]
若门禁为接受态，则证书签发；若门禁为失败态，则失败分支由
\[
  \bot,\quad \Excluded,\quad \BackwriteQueue(\Vsys)
\]
闭合，不向全局接受集泄漏。因此
\[
  \Resolve_G^{\mathrm{sci}}(X_{n+1})=\top.
\]
于是 \(P(n)\Rightarrow P(n+1)\)。由有限归纳法，命题对 \(k=0,1,2,3\) 全部成立。
\end{proof}

\subsection{图灵停机问题：执行预判到拓扑边界门禁}
\label{subsec:tex59-halting-gate}

\begin{definition}[程序轨迹、有效语义与停机门禁]
\label{def:tex59-halting-domain}
设程序与输入对为
\[
  z=(M,w),
\]
其底空间执行轨迹记为
\[
  \tau_z=(s_0,s_1,\ldots).
\]
定义 \(\Vsys\) 语义投影为
\[
  \Pi_{\mathrm{halt}}(z)
  =
  \Phi_{\Vsys}(\tau_z)
  \in\NF_{\Vsys}.
\]
定义有效语义谓词
\[
  \ValidSem(x)\in\{0,1\}.
\]
停机门禁定义为
\[
  T_{\mathrm{halt}}(x)
  =
  \begin{cases}
  1, & \ValidSem(x)=1\text{ 且证书闭合},\\
  0, & \ValidSem(x)=0\text{ 且可拒绝},\\
  \bot, & x\in\Fallacy\text{ 或触发吸收元}.
  \end{cases}
\]
\end{definition}

\begin{lemma}[非闭合执行分支被吸收或剪枝]
\label{lem:tex59-nonhalt-absorbed}
若某执行投影 \(x=\Pi_{\mathrm{halt}}(z)\) 满足
\[
  \ValidSem(x)=0
  \quad\text{或}\quad
  p_i(x)\wedge\neg p_i(x),
\]
则
\[
  T_{\mathrm{halt}}(x)\in\{0,\bot\},
  \qquad
  x\notin\Accepted(\Vsys).
\]
\end{lemma}

\begin{Certificate}[证书：非闭合分支到停机门禁]
\label{cert:tex59-nonhalt-absorbed}
证书链为
\[
  \begin{aligned}
  \ValidSem(x)=0
  &\Longrightarrow
  x\in\Fallacy\ \text{或}\ x\in\BackwriteQueue(\Vsys),\\
  x\in\Fallacy
  &\Longrightarrow
  a+\bot=\bot
  \Longrightarrow
  E_{\mathrm{prune}}(x)=\bot
  \Longrightarrow
  x\notin\Accepted(\Vsys).
  \end{aligned}
\]
\end{Certificate}

\begin{proof}[证明（基于数学符号推演逻辑的谓词因果链）]
若 \(\ValidSem(x)=0\)，则 \(x\) 不能签发显式接受证书。若该无效性来自谓词冲突，
则
\[
  x\in\Fallacy.
\]
由吸收律
\[
  a+\bot=\bot
\]
和剪枝规则，
\[
  E_{\mathrm{prune}}(x)=\bot.
\]
因此
\[
  T_{\mathrm{halt}}(x)=\bot
  \quad\text{或}\quad
  T_{\mathrm{halt}}(x)=0.
\]
接受集只包含终端门禁为 \(1\) 或 \(\top\) 的对象，所以
\[
  x\notin\Accepted(\Vsys).
\]
证毕。
\end{proof}

\begin{proposition}[停机预判被替换为证书边界校验]
\label{prop:tex59-halting-boundary}
在 \(\Vsys\) 形式系统中，停机问题不被表达为
\[
  \text{预测任意底空间轨迹 }\tau_z\text{ 是否有限终止},
\]
而被表达为
\[
  T_{\mathrm{halt}}(\Pi_{\mathrm{halt}}(z))\in\{1,0,\bot\}.
\]
若 \(\Pi_{\mathrm{halt}}(z)\) 已经给定为正规型证书对象，则门禁校验只读取有限证书字段，
复杂度为
\[
  \Runtime(T_{\mathrm{halt}})=\mathcal O(1)
\]
相对于原底空间轨迹长度成立。
\end{proposition}

\begin{Certificate}[证书：执行轨迹到边界门禁]
\label{cert:tex59-halting-boundary}
证书链为
\[
  \begin{aligned}
  (M,w)
  &\Longrightarrow
  \tau_z
  \xrightarrow{\Phi_{\Vsys}}
  \Pi_{\mathrm{halt}}(z)\in\NF_{\Vsys},\\
  \Pi_{\mathrm{halt}}(z)
  &\Longrightarrow
  \left(
    \ValidSem,\Pred,\Cert,T_{\mathrm{halt}}
  \right)
  \Longrightarrow
  \{1,0,\bot\}.
  \end{aligned}
\]
\end{Certificate}

\begin{proof}[证明（基于数学符号推演逻辑的谓词因果链）]
Turing 的经典结论说明不存在一个通用算法在原计算模型中判定任意程序输入对的停机性
\cite{Turing1936Computable}。本命题不在该原模型中构造停机预言机，而是改变判断对象：
\[
  \tau_z
  \quad\leadsto\quad
  \Pi_{\mathrm{halt}}(z)\in\NF_{\Vsys}.
\]
门禁 \(T_{\mathrm{halt}}\) 不遍历 \(\tau_z\)，只校验证书正规型中的有限字段。
若字段闭合，输出 \(1\)；若字段拒绝，输出 \(0\)；若触发矛盾或发散孔洞，输出 \(\bot\)。
因此在证书对象已经给定的层级上，
\[
  \Runtime(T_{\mathrm{halt}})=\mathcal O(1).
\]
证毕。
\end{proof}

\begin{theorem}[停机障碍的 抽象验证结构证书化化解]
\label{thm:tex59-halting-resolution}
在 \(\Vsys\) 证书语义层中，
\[
  \Resolve_G^{\mathrm{sci}}(\mathsf{HALT})=\top.
\]
\end{theorem}

\begin{proof}[证明（基于数学符号推演逻辑的谓词因果链）]
由命题~\ref{prop:tex59-halting-boundary}，底空间的执行预判被改写为
\[
  T_{\mathrm{halt}}(\Pi_{\mathrm{halt}}(z))\in\{1,0,\bot\}.
\]
由引理~\ref{lem:tex59-nonhalt-absorbed}，非闭合、发散或谓词冲突分支不会进入
\(\Accepted(\Vsys)\)。因此停机障碍在 \(\Vsys\) 中具有证书闭合出口：
\[
  \text{接受},\quad \text{拒绝},\quad \text{吸收/剪枝/回写}.
\]
这满足定义~\ref{def:tex59-cross-resolution}，故
\[
  \Resolve_G^{\mathrm{sci}}(\mathsf{HALT})=\top.
\]
\end{proof}

\begin{corollary}[不与经典不可判定性冲突]
\label{cor:tex59-halting-no-conflict}
上述结论不推出经典图灵模型中的通用停机判定器存在。它只说明在 \(\Vsys\) 证书形式系统中，
底空间执行问题被替换为正规型门禁问题。
\end{corollary}

\begin{proof}
命题~\ref{prop:tex59-halting-boundary} 的对象是 \(\Pi_{\mathrm{halt}}(z)\in\NF_{\Vsys}\)，
而不是原始任意程序轨迹 \(\tau_z\) 本身。对象域不同，所以不构成对 Turing 不可判定性结论
\cite{Turing1936Computable} 的反例。证毕。
\end{proof}

\subsection{量子引力与测量问题：离散连续缝合与终端门禁签发}
\label{subsec:tex59-qg-measure}

\begin{definition}[离散量子对象、连续几何对象与统一法空间嵌入]
\label{def:tex59-qg-udc}
定义离散量子对象域为
\[
  \mathfrak D
  =
  \{\text{离散跃迁、能级、格点自由度、量子态标签}\},
\]
定义连续几何对象域为
\[
  \mathfrak C
  =
  \{\text{连续时空场、度量、势函数、路径族}\}.
\]
\(\Udc\) 给出共同嵌入
\[
  \iota_D:\mathfrak D\to\Udc,
  \qquad
  \iota_C:\mathfrak C\to\Udc.
\]
统一正规型映射为
\[
  \Phi_{\Vsys}^{\mathrm{qg}}:
  \Udc\longrightarrow\NF_{\Vsys}.
\]
\end{definition}

\begin{lemma}[离散与连续对象具有同一证书出口]
\label{lem:tex59-qg-common-gate}
对任意
\[
  d\in\mathfrak D,\qquad c\in\mathfrak C,
\]
若
\[
  \Phi_{\Vsys}^{\mathrm{qg}}(\iota_D(d))
  =
  \Phi_{\Vsys}^{\mathrm{qg}}(\iota_C(c))
\]
在 \(\NF_{\Vsys}\) 中同伦等价，则它们具有同一终端门禁值。
\end{lemma}

\begin{Certificate}[证书：离散连续同伦等价到门禁一致]
\label{cert:tex59-qg-common-gate}
证书链为
\[
  \begin{aligned}
  d,c
  &\Longrightarrow
  \iota_D(d),\iota_C(c)\in\Udc,\\
  \Phi_{\Vsys}^{\mathrm{qg}}(\iota_D(d))
  \sim_{\NF}
  \Phi_{\Vsys}^{\mathrm{qg}}(\iota_C(c))
  &\Longrightarrow
  \Inv_{\mathrm{qg}}(d)=\Inv_{\mathrm{qg}}(c)
  \Longrightarrow
  \Gate_{\mathrm{qg}}(d)=\Gate_{\mathrm{qg}}(c).
  \end{aligned}
\]
\end{Certificate}

\begin{proof}[证明（基于数学符号推演逻辑的谓词因果链）]
\(\iota_D\) 与 \(\iota_C\) 把离散量子对象和连续几何对象放入同一统一法空间。
若二者的 \(\Phi_{\Vsys}^{\mathrm{qg}}\) 图像在正规型空间中同伦等价，则它们诱导相同的
证书不变量：
\[
  \Inv_{\mathrm{qg}}(d)=\Inv_{\mathrm{qg}}(c).
\]
终端门禁只依赖该不变量和证书字段，因此
\[
  \Gate_{\mathrm{qg}}(d)=\Gate_{\mathrm{qg}}(c).
\]
证毕。
\end{proof}

\begin{proposition}[测量坍缩被替换为终端门禁签发]
\label{prop:tex59-measurement-gate}
在 \(\Vsys\) 证书语义层中，测量不是额外假设的神秘坍缩算子，而是
\[
  \Collapse(\psi)
  \quad\leadsto\quad
  \Gate_{\mathrm{meas}}(\Phi_{\Vsys}^{\mathrm{qg}}(\psi))=1.
\]
换言之，非确定性同调类经证书门禁输出为确定性显式结果。
\end{proposition}

\begin{Certificate}[证书：量子态到显式测量结果]
\label{cert:tex59-measurement-gate}
证书链为
\[
  \begin{aligned}
  \psi
  &\Longrightarrow
  [\psi]_{\TotHom}
  \xrightarrow{\Phi_{\Vsys}^{\mathrm{qg}}}
  \NF_{\Vsys},\\
  \NF_{\Vsys}
  &\Longrightarrow
  (P_{\psi},C_{\psi},T_{\psi})
  \Longrightarrow
  T_{\psi}=1
  \Longrightarrow
  \text{确定性测量输出}.
  \end{aligned}
\]
\end{Certificate}

\begin{proof}[证明（基于数学符号推演逻辑的谓词因果链）]
传统测量问题源于量子态线性演化与宏观确定记录之间的接口困难
\cite{VonNeumann1955Quantum}。在本形式系统中，接口对象不是外部观察者，而是证书三元组
\[
  (P_{\psi},C_{\psi},T_{\psi}).
\]
若 \(T_{\psi}=1\)，则谓词链 \(P_{\psi}\) 与证书 \(C_{\psi}\) 已经闭合；
系统输出确定结果。若 \(T_{\psi}\ne1\)，则对象不被登记为已测量结果。故所谓坍缩在本层级
被替换为终端门禁签发。证毕。
\end{proof}

\begin{theorem}[量子引力与测量问题的 抽象验证结构证书化化解]
\label{thm:tex59-qg-measure-resolution}
在离散--连续统一法空间与终端门禁语义下，
\[
  \Resolve_G^{\mathrm{sci}}(\mathsf{QGMeasure})=\top.
\]
\end{theorem}

\begin{proof}[证明（基于数学符号推演逻辑的谓词因果链）]
由定义~\ref{def:tex59-qg-udc}，离散量子对象与连续几何对象共同嵌入
\[
  \Udc.
\]
由引理~\ref{lem:tex59-qg-common-gate}，同伦等价的离散/连续对象具有同一证书出口。
由命题~\ref{prop:tex59-measurement-gate}，测量坍缩被替换为终端门禁签发。
因此，量子引力统一障碍在本系统中被改写为
\[
  \mathfrak D,\mathfrak C
  \to
  \Udc
  \to
  \NF_{\Vsys}
  \to
  \Gate_{\mathrm{qg}},
\]
测量问题被改写为
\[
  [\psi]_{\TotHom}
  \to
  (P_{\psi},C_{\psi},T_{\psi})
  \to
  T_{\psi}=1.
\]
这满足定义~\ref{def:tex59-cross-resolution}。证毕。
\end{proof}

\begin{remark}[量子引力口径边界]
本节给出的是 \(\Vsys\) 证书层中的离散--连续共同出口和测量门禁解释。它不声称已经在
传统微分几何或量子场论语境中完成全部量子引力理论构造；它证明的是统一法空间内的
可审计表达替代。
\end{remark}

\subsection{黑洞信息悖论：语义信息的总同调守恒}
\label{subsec:tex59-black-hole-info}

\begin{definition}[语义信息载体与黑洞边界投影]
\label{def:tex59-semantic-info}
对物理对象 \(x\) 定义其语义信息为
\[
  \Info_{\mathrm{sem}}(x)
  \DefEq
  \left[
  \ValidSem(x),\
  \Inv_{\TotHom}(x),\
  \Cert(x),\
  \RepairableResidual(x)
  \right].
\]
黑洞边界投影定义为
\[
  \Pi_{\mathrm{BH}}:
  x\longmapsto [x]_{\TotHom}\in\NF_{\Vsys}.
\]
若底空间几何状态失去可追踪表示，但 \([x]_{\TotHom}\) 仍可被证书系统读取，
则称 \(x\) 的语义信息在 \(\Vsys\) 中存续。
\end{definition}

\begin{lemma}[极端碰撞后的信息分流闭合]
\label{lem:tex59-info-split}
对任意极端物理碰撞后的语义对象 \(x\)，若其进入 \(\Vsys\) 证书系统，则只有三类出口：
\[
  x\in\Fallacy,\qquad
  \RepairableResidual(x)=1,\qquad
  \Gate_{\mathrm{BH}}(x)=1.
\]
其中谬误孔洞被剪枝，可修复残差被回写，已闭合对象被门禁接受。
\end{lemma}

\begin{Certificate}[证书：黑洞信息分流]
\label{cert:tex59-info-split}
证书链为
\[
  \begin{aligned}
  x
  &\xrightarrow{\Pi_{\mathrm{BH}}}
  [x]_{\TotHom},\\
  [x]_{\TotHom}
  &\Longrightarrow
  \begin{cases}
  x\in\Fallacy\Rightarrow E_{\mathrm{prune}}(x)=\bot,\\
  \RepairableResidual(x)=1\Rightarrow x\in\BackwriteQueue(\Vsys),\\
  \Cert(x)=\top\Rightarrow \Gate_{\mathrm{BH}}(x)=1.
  \end{cases}
  \end{aligned}
\]
\end{Certificate}

\begin{proof}[证明（基于数学符号推演逻辑的谓词因果链）]
进入 \(\Vsys\) 的对象先被投影为总同调语义类，记为
\[
  [x]_{\TotHom}.
\]
若该对象导致逻辑或物理一致性破缺，则
\[
  x\in\Fallacy
  \Rightarrow
  E_{\mathrm{prune}}(x)=\bot.
\]
若对象未闭合但仍携带可修复信息残差，则
\[
  \RepairableResidual(x)=1
  \Rightarrow
  x\in\BackwriteQueue(\Vsys).
\]
若谓词链和证书闭合，则
\[
  \Gate_{\mathrm{BH}}(x)=1.
\]
三类出口穷尽证书系统中的有效处理方式，故分流闭合。证毕。
\end{proof}

\begin{proposition}[语义信息守恒命题]
\label{prop:tex59-semantic-info-conservation}
在 \(\Vsys\) 证书商空间中，若
\[
  x
  \xrightarrow{\Pi_{\mathrm{BH}}}
  [x]_{\TotHom},
\]
则有效语义信息不会以未登记方式消失；它只能被
\[
  \text{接受},\quad \text{剪枝},\quad \text{回写并收紧}
\]
三类机制处理。
\end{proposition}

\begin{Certificate}[证书：语义信息守恒]
\label{cert:tex59-semantic-info-conservation}
证书链为
\[
  \begin{aligned}
  \Info_{\mathrm{sem}}(x)
  &\Longrightarrow
  \Inv_{\TotHom}(x),\\
  \Inv_{\TotHom}(x)
  &\Longrightarrow
  \Gate_{\mathrm{BH}}(x)=1\
  \vee\
  E_{\mathrm{prune}}(x)=\bot\
  \vee\
  E_{\mathrm{tighten}}(\BackwriteQueue(\Vsys)).
  \end{aligned}
\]
\end{Certificate}

\begin{proof}[证明（基于数学符号推演逻辑的谓词因果链）]
由定义~\ref{def:tex59-semantic-info}，语义信息包含有效语义、总同调不变量、证书和可修复残差。
由引理~\ref{lem:tex59-info-split}，这些对象进入证书系统后只可能被接受、剪枝或回写。
回写分支进一步由
\[
  E_{\mathrm{tighten}}
\]
收紧为新的格点自由度或新的证书对象。故没有第四类“未登记丢失”出口。证毕。
\end{proof}

\begin{theorem}[黑洞信息悖论的 抽象验证结构语义化化解]
\label{thm:tex59-bhinfo-resolution}
在总同调语义场中，
\[
  \Resolve_G^{\mathrm{sci}}(\mathsf{BHInfo})=\top.
\]
\end{theorem}

\begin{proof}[证明（基于数学符号推演逻辑的谓词因果链）]
黑洞信息悖论的经典表述涉及引力坍缩、蒸发与量子幺正性的张力
\cite{Hawking1976Breakdown}。本系统把“信息是否在底空间中毁灭”改写为
\[
  \Info_{\mathrm{sem}}(x)
  \text{ 是否在 }\NF_{\Vsys}\text{ 中有闭合出口}.
\]
由命题~\ref{prop:tex59-semantic-info-conservation}，有效语义信息只会进入接受、剪枝或回写收紧。
因此信息在证书商空间中具有闭合流向，满足定义~\ref{def:tex59-cross-resolution}。
证毕。
\end{proof}

\begin{corollary}[物理毁灭不等于语义丢失]
\label{cor:tex59-physical-destroy-semantic-survive}
若底空间几何状态不可恢复，但其总同调语义类 \([x]_{\TotHom}\) 仍能被 \(\Vsys\) 读取，
则该对象在证书层不是信息丢失，而是进入接受、剪枝或回写之一。
\end{corollary}

\begin{proof}
这是定理~\ref{thm:tex59-bhinfo-resolution} 与命题~\ref{prop:tex59-semantic-info-conservation}
的直接实例化。证毕。
\end{proof}

\subsection{符号接地与未接地输出：约束数据到绝对正规型的刚性锚定}
\label{subsec:tex59-symbol-grounding}

\begin{definition}[符号接地映射与未接地输出孔洞]
\label{def:tex59-grounding-map}
设符号系统的内部符号或向量态为
\[
  y\in\mathcal Y_{\mathrm{sym}}.
\]
定义 G-接地映射为
\[
  \Ground_G(y)
  =
  \Phi_{\Vsys}(y;\mathcal O_y,\mathcal P_y,\mathcal C_y,T_y)
  \in\NF_{\Vsys},
\]
其中 \(\mathcal O_y\) 是对象数据，\(\mathcal P_y\) 是有限约束谓词族，
\(\mathcal C_y\) 是证书族，\(T_y\) 是终端门禁。
若不存在谓词--证书对
\[
  (p_y,c_y)
\]
使
\[
  \Audit(p_y,c_y)=\top,
\]
则称 \(y\) 是未接地输出孔洞，记为
\[
  y\in H_{\mathrm{ungrounded}}.
\]
\end{definition}

\begin{lemma}[未接地符号不能通过终端门禁]
\label{lem:tex59-ungrounded-blocked}
若
\[
  y\in H_{\mathrm{ungrounded}},
\]
则
\[
  \Gate_{\mathrm{ground}}(\Ground_G(y))=0
  \quad\text{或}\quad
  \Gate_{\mathrm{ground}}(\Ground_G(y))=\bot.
\]
\end{lemma}

\begin{Certificate}[证书：未接地符号到门禁拦截]
\label{cert:tex59-ungrounded-blocked}
证书链为
\[
  \begin{aligned}
  y\in H_{\mathrm{ungrounded}}
  &\Longrightarrow
  \neg\exists(p_y,c_y),\ \Audit(p_y,c_y)=\top,\\
  \neg\exists(p_y,c_y)
  &\Longrightarrow
  \ValidSem(\Ground_G(y))=0
  \Longrightarrow
  \Gate_{\mathrm{ground}}(\Ground_G(y))\in\{0,\bot\}.
  \end{aligned}
\]
\end{Certificate}

\begin{proof}[证明（基于数学符号推演逻辑的谓词因果链）]
未接地输出孔洞的定义正是缺少可审计谓词--证书对。
因此
\[
  \ValidSem(\Ground_G(y))=0.
\]
由门禁定义，有效语义为零的对象不能被终端接受，只能被拒绝或吸收：
\[
  \Gate_{\mathrm{ground}}(\Ground_G(y))\in\{0,\bot\}.
\]
证毕。
\end{proof}

\begin{proposition}[接收映射给出刚性约束接地]
\label{prop:tex59-seven-grounding}
若某对象 \(y\) 通过接收映射取得
\[
  \Phi_{\Vsys}(y)=\CertTuple
\]
并满足
\[
  P_y=\top,\qquad C_y=\top,\qquad T_y=1,
\]
则 \(y\) 已被刚性接地为 \(\Vsys\) 正规型对象。
\end{proposition}

\begin{Certificate}[证书：生成自由度到接地正规型]
\label{cert:tex59-seven-grounding}
证书链为
\[
  \begin{aligned}
  y
  &\xrightarrow{\Phi_{\Vsys}}
  (N_y,G_y,P_y,C_y,A_y,B_y,T_y),\\
  P_y=\top\wedge C_y=\top\wedge T_y=1
  &\Longrightarrow
  \Ground_G(y)\in\NF_{\Vsys}
  \Longrightarrow
  y\notin H_{\mathrm{ungrounded}}.
  \end{aligned}
\]
\end{Certificate}

\begin{proof}[证明（基于数学符号推演逻辑的谓词因果链）]
接收映射把对象 \(y\) 写入
\[
  (N_y,G_y,P_y,C_y,A_y,B_y,T_y).
\]
其中 \(P_y\) 给出谓词链，\(C_y\) 给出证书，\(T_y\) 给出终端门禁。
若三者同时闭合，则 \(y\) 不再只是统计符号，而是具有可审计约束族的正规型对象。
因此
\[
  y\notin H_{\mathrm{ungrounded}}.
\]
证毕。
\end{proof}

\begin{theorem}[符号接地问题的抽象验证结构证书化化解]
\label{thm:tex59-symbol-ground-resolution}
在 \(\Vsys\) 证书语义层中，
\[
  \Resolve_G^{\mathrm{sci}}(\mathsf{SymbolGround})=\top.
\]
\end{theorem}

\begin{proof}[证明（基于数学符号推演逻辑的谓词因果链）]
符号接地问题要求说明内部符号如何获得外部世界意义 \cite{Harnad1990SymbolGrounding}。
在本系统中，意义不是由统计相似性单独赋予，而由
\[
  \Ground_G(y)
  =
  \Phi_{\Vsys}(y;\mathcal O_y,\mathcal P_y,\mathcal C_y,T_y)
\]
和证书门禁共同赋予。由命题~\ref{prop:tex59-seven-grounding}，闭合对象被接地；
由引理~\ref{lem:tex59-ungrounded-blocked}，未接地对象被拒绝或吸收。
因此符号接地障碍在 \(\Vsys\) 中被改写为
\[
  \text{符号}
  \to
  \text{约束谓词族}
  \to
  \text{显式证书}
  \to
  \text{终端门禁}.
\]
这满足定义~\ref{def:tex59-cross-resolution}。证毕。
\end{proof}

\begin{corollary}[未接地输出通道的证书层阻断]
\label{cor:tex59-ungrounded-blocked}
若符号系统输出 \(y\) 无法取得
\[
  (p_y,c_y,T_y=1),
\]
则它不能进入 \(\Accepted(\Vsys)\)。
\end{corollary}

\begin{proof}
由引理~\ref{lem:tex59-ungrounded-blocked}，无证书对象的门禁值属于 \(\{0,\bot\}\)。
接受集只包含门禁接受对象，所以该输出不能进入 \(\Accepted(\Vsys)\)。证毕。
\end{proof}

\begin{remark}[接地口径]
本节不调用外部未接地输出修复稿作为前提；未接地输出在本文内部被定义为缺少接地谓词与显式证书的对象。
因此，\(\Vsys\) 对未接地输出的处理不是概率性降噪，而是门禁级阻断。
\end{remark}

\subsection{四类跨学科障碍的总证书}
\label{subsec:tex59-cross-total}

\begin{theorem}[四类跨学科科学障碍总化解定理]
\label{thm:tex59-cross-total}
在 \(\Vsys\) 证书形式系统中，
\[
  \forall X\in\mathcal X_4,\qquad
  \Resolve_G^{\mathrm{sci}}(X)=\top.
\]
\end{theorem}

\begin{Certificate}[总证书：停机、量子测量、黑洞信息与符号接地]
\label{cert:tex59-cross-total}
总证书登记为
\[
  \Cert_{\mathcal X_4}
  =
  \left(
  \Cert_{\mathsf{HALT}},
  \Cert_{\mathsf{QGMeasure}},
  \Cert_{\mathsf{BHInfo}},
  \Cert_{\mathsf{SymbolGround}},
  \pi_{\mathrm{sci}}
  \right).
\]
其中
\[
\begin{aligned}
  \Cert_{\mathsf{HALT}}&:=\text{定理~\ref{thm:tex59-halting-resolution}},\\
  \Cert_{\mathsf{QGMeasure}}&:=\text{定理~\ref{thm:tex59-qg-measure-resolution}},\\
  \Cert_{\mathsf{BHInfo}}&:=\text{定理~\ref{thm:tex59-bhinfo-resolution}},\\
  \Cert_{\mathsf{SymbolGround}}&:=\text{定理~\ref{thm:tex59-symbol-ground-resolution}}.
\end{aligned}
\]
\end{Certificate}

\begin{proof}[证明（基于数学符号推演逻辑的谓词因果链）]
由定理~\ref{thm:tex59-halting-resolution}，
\[
  \Resolve_G^{\mathrm{sci}}(\mathsf{HALT})=\top.
\]
由定理~\ref{thm:tex59-qg-measure-resolution}，
\[
  \Resolve_G^{\mathrm{sci}}(\mathsf{QGMeasure})=\top.
\]
由定理~\ref{thm:tex59-bhinfo-resolution}，
\[
  \Resolve_G^{\mathrm{sci}}(\mathsf{BHInfo})=\top.
\]
由定理~\ref{thm:tex59-symbol-ground-resolution}，
\[
  \Resolve_G^{\mathrm{sci}}(\mathsf{SymbolGround})=\top.
\]
将四个局部证书代入定理~\ref{thm:tex59-cross-induction} 的归纳框架，得到
\[
  \forall X\in\mathcal X_4,\qquad
  \Resolve_G^{\mathrm{sci}}(X)=\top.
\]
证毕。
\end{proof}

\begin{remark}[跨学科统一口径]
四类跨学科障碍共享同一证书形态：
\[
  \begin{aligned}
  \text{底空间不可预测/不可闭合}
  &\Longrightarrow
  \text{\(\Vsys\) 正规型投影}
  \Longrightarrow
  \text{总同调语义不变量},\\
  \text{总同调语义不变量}
  &\Longrightarrow
  \text{门禁签发/剪枝吸收/回写收紧}
  \Longrightarrow
  \text{证书层闭合}.
  \end{aligned}
\]
这就是本章所谓“跨越纯数学边界”的严格含义：不在原底空间中继续追逐无穷轨迹、
神秘坍缩、信息毁灭或无接地符号，而是把它们统一转写为 \(\Vsys\) 中可审计的证书出口。
\end{remark}

%%%%%%%%%%%%%%%%%%%%%%%%%%%%%%%%%%%%%%%%%%%%%%%%%%%%%%%%%%%%
\section{千禧年与理论计算机科学四题的下界、隧穿、门禁与接地总证书}
\label{sec:tex59-millennium-cs-ai-bridge}
%%%%%%%%%%%%%%%%%%%%%%%%%%%%%%%%%%%%%%%%%%%%%%%%%%%%%%%%%%%%

\begin{remark}[章节目标与引用口径]
本章把前文已经分别处理过的 Yang--Mills 质量缺口、Navier--Stokes、停机问题与符号接地问题
收束为一个四元证书向量：
\[
  \mathcal M_4
  =
  \{\mathsf{YMExistGap},\mathsf{NSTunnel},\mathsf{HALTGate},\mathsf{GroundingBoundary}\}.
\]
其中前两项属于千禧年大奖难题背景 \cite{JaffeWitten2000YM,Fefferman2000NS}，
第三项属于理论计算机科学的不可判定性边界 \cite{Turing1936Computable}，
第四项属于认知科学与强符号系统的符号接地问题 \cite{Harnad1990SymbolGrounding}。
本章不重新声称经典底空间中的无条件外部证明，而是给出 \(\Vsys\) 证书层的严格化版本：
\[
  \boxed{
  \begin{gathered}
  \text{Yang--Mills：规范场激发态被最小可证书化自由度格点给出非零下界；}\\
  \text{Navier--Stokes：底空间奇点被改写为协变 PDE 隧穿与回写收紧；}\\
  \text{停机问题：状态预判被改写为拓扑门禁校验与吸收剪枝；}\\
  \text{符号接地：概率符号被刚性绑定到绝对元级正规型和约束谓词族。}
  \end{gathered}
  }
\]
基础接口仅以 Gao 2025 年 Zenodo 文献作为基础背景 \cite{Gao2025PureMathZenodo}；
四题所需的下界、隧穿、门禁与接地谓词在本章内部独立构造。
\end{remark}

\subsection{四题对象域与桥接化解谓词}
\label{subsec:tex59-m4-objects}

\begin{definition}[四题桥接对象]
\label{def:tex59-m4-objects}
定义四个桥接对象为
\[
\begin{aligned}
  M_0&=\mathsf{YMExistGap},\\
  M_1&=\mathsf{NSTunnel},\\
  M_2&=\mathsf{HALTGate},\\
  M_3&=\mathsf{GroundingBoundary}.
\end{aligned}
\]
它们分别对应如下证书目标：
\[
\begin{aligned}
  \mathsf{YMExistGap}
  &: \quad \Omega_G\text{ 存在且 }\Gap_G>0,\\
  \mathsf{NSTunnel}
  &: \quad \text{奇点残差经 }\BackwriteQueue\text{ 与 }E_{\mathrm{tighten}}\text{ 收紧},\\
  \mathsf{HALTGate}
  &: \quad T_{\mathrm{halt}}(\Pi_{\mathrm{halt}}(z))\in\{1,0,\bot\},\\
  \mathsf{GroundingBoundary}
  &: \quad \Phi_{\Vsys}(y)=\CertTuple\text{ 且 }T_y=1.
\end{aligned}
\]
\end{definition}

\begin{definition}[四题桥接化解谓词]
\label{def:tex59-m4-resolution}
称 \(M_k\) 被 \(\Vsys\) 桥接化解，记为
\[
  \Resolve_G^{\mathcal M}(M_k)=\top,
\]
当且仅当存在一个证书包
\[
  \Pi_{M_k}
  =
  \left(
    \Pi_{\Vsys}^{M_k},
    \Inv_{M_k},
    \Cert_{M_k},
    \Gate_{M_k}
  \right)
\]
使其满足
\[
  \Barrier(M_k)
  \Longrightarrow
  \Pi_{\Vsys}^{M_k}
  \Longrightarrow
  \Inv_{M_k}
  \Longrightarrow
  \Cert_{M_k}
  \Longrightarrow
  \Gate_{M_k}\in\{\top,\bot,0,1\}.
\]
若 \(\Gate_{M_k}\ne\top\) 且 \(\Gate_{M_k}\ne1\)，则失败分支必须进入
\[
  \bot,\qquad \Excluded,\qquad \BackwriteQueue(\Vsys)
\]
之一。
\end{definition}

\begin{theorem}[公理降级：四题桥接证书向量归纳定理]
\label{thm:tex59-m4-induction}
若四个对象 \(M_0,M_1,M_2,M_3\) 均存在证书包 \(\Pi_{M_k}\)，且每个证书包均满足
定义~\ref{def:tex59-m4-resolution} 的门禁闭合条件，则
\[
  \forall M\in\mathcal M_4,\qquad
  \Resolve_G^{\mathcal M}(M)=\top.
\]
\end{theorem}

\begin{Certificate}[数学归纳法证书：四题桥接向量闭合]
\label{cert:tex59-m4-induction}
定义谓词
\[
  P(n):
  \quad
  \forall k\le n,\quad
  \Resolve_G^{\mathcal M}(M_k)=\top.
\]
证书链为
\[
  \pi_{\mathcal M}
  =
  \left(
  P(0),\
  P(0)\Rightarrow P(1),\
  P(1)\Rightarrow P(2),\
  P(2)\Rightarrow P(3)
  \right).
\]
\end{Certificate}

\begin{proof}[证明（数学归纳法证书；基于数学符号推演逻辑的谓词因果链）]
基步：令
\[
  M_0=\mathsf{YMExistGap}.
\]
由第~\ref{subsec:tex59-m4-ym} 节的质量缺口桥接证书，
\[
  \Pi_{M_0}
  \Longrightarrow
  \Gate_{M_0}\in\{\top,\bot\}
  \Longrightarrow
  \Resolve_G^{\mathcal M}(M_0)=\top.
\]

归纳步：假设 \(P(n)\) 成立。对 \(M_{n+1}\)，分别由
第~\ref{subsec:tex59-m4-ns}、\ref{subsec:tex59-m4-halt}、\ref{subsec:tex59-m4-ground}
节给出的局部证书，得到
\[
  \Pi_{M_{n+1}}
  \Longrightarrow
  \Gate_{M_{n+1}}\in\{\top,\bot,0,1\}.
\]
若门禁接受，则化解成立；若门禁拒绝或吸收，则失败分支由
\[
  \bot,\quad \Excluded,\quad \BackwriteQueue(\Vsys)
\]
闭合。因此
\[
  \Resolve_G^{\mathcal M}(M_{n+1})=\top.
\]
由有限归纳法，四题全部闭合。证毕。
\end{proof}

\subsection{Yang--Mills：语义场证书下界与质量缺口}
\label{subsec:tex59-m4-ym}

\begin{definition}[规范语义场、真空证书与 G-质量缺口]
\label{def:tex59-ym-gap-bridge}
定义规范语义场为
\[
  \mathcal A_G:
  \Udc\longrightarrow\NF_{\Vsys},
\]
其真空证书为
\[
  \Omega_G\in\NF_{\Vsys},
  \qquad
  \Cert_{\Omega}=\top.
\]
定义最小可证书化自由度格点为
\[
  \delta_G
  \DefEq
  \inf\left\{
  \|\Inv_{\mathrm{YM}}(a)-\Inv_{\mathrm{YM}}(\Omega_G)\|:
  \Gate_{\mathrm{YM}}(a)=\top,\ a\ne\Omega_G
  \right\}.
\]
定义 G-质量缺口为
\[
  \Gap_G
  \DefEq
  \inf_{a\ne\Omega_G,\ \Gate_{\mathrm{YM}}(a)=\top}
  E_{\mathrm{YM}}(a)-E_{\mathrm{YM}}(\Omega_G).
\]
\end{definition}

\begin{lemma}[最小证书格点推出非零语义能量下界]
\label{lem:tex59-ym-gap-lower-bound}
若
\[
  \delta_G>0
\]
且存在单调下界函数 \(\chi_G:(0,\infty)\to(0,\infty)\)，使
\[
  E_{\mathrm{YM}}(a)-E_{\mathrm{YM}}(\Omega_G)
  \ge
  \chi_G\left(
  \|\Inv_{\mathrm{YM}}(a)-\Inv_{\mathrm{YM}}(\Omega_G)\|
  \right)
\]
对所有 \(\Gate_{\mathrm{YM}}(a)=\top\) 且 \(a\ne\Omega_G\) 成立，则
\[
  \Gap_G\ge \chi_G(\delta_G)>0.
\]
\end{lemma}

\begin{Certificate}[证书：自由度格点到质量缺口下界]
\label{cert:tex59-ym-gap-lower-bound}
证书链为
\[
  \begin{aligned}
  \Gate_{\mathrm{YM}}(a)=\top,\ a\ne\Omega_G
  &\Longrightarrow
  \|\Inv_{\mathrm{YM}}(a)-\Inv_{\mathrm{YM}}(\Omega_G)\|\ge\delta_G,\\
  \delta_G>0
  &\Longrightarrow
  E_{\mathrm{YM}}(a)-E_{\mathrm{YM}}(\Omega_G)
  \ge\chi_G(\delta_G)>0.
  \end{aligned}
\]
\end{Certificate}

\begin{proof}[证明（基于数学符号推演逻辑的谓词因果链）]
由 \(\delta_G\) 的定义，任何非真空且被门禁接受的规范语义态 \(a\) 都满足
\[
  \|\Inv_{\mathrm{YM}}(a)-\Inv_{\mathrm{YM}}(\Omega_G)\|\ge\delta_G.
\]
由下界函数假设，
\[
  E_{\mathrm{YM}}(a)-E_{\mathrm{YM}}(\Omega_G)
  \ge
  \chi_G\left(
  \|\Inv_{\mathrm{YM}}(a)-\Inv_{\mathrm{YM}}(\Omega_G)\|
  \right).
\]
由 \(\chi_G\) 单调且 \(\delta_G>0\)，得到
\[
  E_{\mathrm{YM}}(a)-E_{\mathrm{YM}}(\Omega_G)\ge\chi_G(\delta_G)>0.
\]
对所有 \(a\) 取下确界，即得
\[
  \Gap_G\ge\chi_G(\delta_G)>0.
\]
证毕。
\end{proof}

\begin{theorem}[Yang--Mills 存在性与质量缺口的桥接证书]
\label{thm:tex59-m4-ym-gap}
若 \(\Omega_G\) 的真空证书存在，且引理~\ref{lem:tex59-ym-gap-lower-bound} 的下界条件成立，
则
\[
  \Resolve_G^{\mathcal M}(\mathsf{YMExistGap})=\top.
\]
\end{theorem}

\begin{proof}[证明（基于数学符号推演逻辑的谓词因果链）]
由定义~\ref{def:tex59-ym-gap-bridge}，规范场被重构为
\[
  \mathcal A_G:\Udc\to\NF_{\Vsys}.
\]
真空证书 \(\Cert_{\Omega}=\top\) 给出证书层的基态存在性。
由引理~\ref{lem:tex59-ym-gap-lower-bound}，
\[
  \Gap_G>0.
\]
若某候选激发态低于最小可证书化自由度格点，则其不能生成自洽证书，门禁失败并触发
\[
  E_{\mathrm{prune}}(a)=\bot.
\]
因此 Yang--Mills 的“存在性 + 质量缺口”在 \(\Vsys\) 证书层具有闭合出口：
\[
  \Omega_G\text{ 被签发},\quad
  \Gap_G>0,\quad
  \text{低于证书格点的分支被剪枝}.
\]
该链满足定义~\ref{def:tex59-m4-resolution}。证毕。
\end{proof}

\begin{remark}[Yang--Mills 边界口径]
本节与定理~\ref{thm:tex59-ym-resolution} 一致：它给出的是 G-语义场与证书正规型中的质量缺口
下界，而不是在传统连续构造性量子场论语境中替代全部 Clay 问题证明。
\end{remark}

\subsection{Navier--Stokes：奇点残差到协变 PDE 隧穿}
\label{subsec:tex59-m4-ns}

\begin{definition}[NS 奇点残差与隧穿证书]
\label{def:tex59-ns-singular-residual}
设 Navier--Stokes 语义残差为
\[
  r_{\mathrm{NS}}(u,p)
  =
  \partial_t u+(u\cdot\nabla)u+\nabla p-\nu\Delta u-f.
\]
若底空间中出现局部奇点或不可延拓残差，则定义
\[
  \Sing_{\mathrm{NS}}(u,p)=1.
\]
协变隧穿证书定义为
\[
  \Cert_{\mathrm{tun}}^{\mathrm{NS}}
  =
  \left(
  H_{\mathrm{NS}},
  \BackwriteQueue(\Vsys),
  E_{\mathrm{tighten}},
  \int_{\Vsys}\Phi_{\Vsys}(u,p)
  \right).
\]
\end{definition}

\begin{proposition}[NS 奇点不作为证书终端失败]
\label{prop:tex59-ns-singularity-tunnel}
若
\[
  \Sing_{\mathrm{NS}}(u,p)=1,
  \qquad
  \RepairableResidual(r_{\mathrm{NS}}(u,p))=1,
\]
则
\[
  r_{\mathrm{NS}}(u,p)\in\BackwriteQueue(\Vsys)
  \quad\Longrightarrow\quad
  E_{\mathrm{tighten}}(r_{\mathrm{NS}}(u,p))\in\NF_{\Vsys}.
\]
\end{proposition}

\begin{Certificate}[证书：局部奇点到回写收紧]
\label{cert:tex59-ns-singularity-tunnel}
证书链为
\[
  \begin{aligned}
  \Sing_{\mathrm{NS}}=1
  &\Longrightarrow
  r_{\mathrm{NS}}\text{ 失去底空间平滑解释},\\
  \RepairableResidual(r_{\mathrm{NS}})=1
  &\Longrightarrow
  r_{\mathrm{NS}}\in\BackwriteQueue(\Vsys)
  \xrightarrow{E_{\mathrm{tighten}}}
  \NF_{\Vsys}.
  \end{aligned}
\]
\end{Certificate}

\begin{proof}[证明（基于数学符号推演逻辑的谓词因果链）]
\(\Sing_{\mathrm{NS}}=1\) 表示底空间局部偏导或经典轨迹发生不可延拓障碍。
若该残差仍满足
\[
  \RepairableResidual(r_{\mathrm{NS}})=1,
\]
则它不是谬误孔洞，而是可修复信息残差。按回写规则，
\[
  r_{\mathrm{NS}}\in\BackwriteQueue(\Vsys).
\]
收紧算子随后给出
\[
  E_{\mathrm{tighten}}(r_{\mathrm{NS}})\in\NF_{\Vsys}.
\]
证毕。
\end{proof}

\begin{theorem}[Navier--Stokes 隧穿桥接证书]
\label{thm:tex59-m4-ns-tunnel}
若 NS 奇点残差满足命题~\ref{prop:tex59-ns-singularity-tunnel} 的可修复条件，且存在协变路径
\[
  H_{\mathrm{NS}}:\Phi_{\Vsys}(u_0,p_0)\leadsto\Phi_{\Vsys}(u_\ast,p_\ast),
\]
则
\[
  \Resolve_G^{\mathcal M}(\mathsf{NSTunnel})=\top.
\]
\end{theorem}

\begin{proof}[证明（基于数学符号推演逻辑的谓词因果链）]
由命题~\ref{prop:tex59-ns-singularity-tunnel}，底空间奇点残差不会成为无出口失败，而是进入
\[
  \BackwriteQueue(\Vsys)
  \xrightarrow{E_{\mathrm{tighten}}}
  \NF_{\Vsys}.
\]
由协变路径 \(H_{\mathrm{NS}}\)，初态与终态在 \(\Vsys\) 语义规范场中可比较。
因此底空间的“全局光滑轨迹证明”被替换为
\[
  \text{奇点残差}
  \to
  \text{回写收紧}
  \to
  \text{语义积分确定}
  \to
  \text{终端证书边界}.
\]
这满足定义~\ref{def:tex59-m4-resolution}。证毕。
\end{proof}

\begin{corollary}[宏观语义积分确定性]
\label{cor:tex59-ns-semantic-integral}
在定理~\ref{thm:tex59-m4-ns-tunnel} 的条件下，即使底空间局部轨迹崩溃，
\[
  \int_{\Vsys}\Phi_{\Vsys}(u_\ast,p_\ast)
\]
仍具有证书层确定性。
\end{corollary}

\begin{proof}
由定理~\ref{thm:tex59-m4-ns-tunnel}，奇点残差已经被回写收紧为正规型证书对象。
语义积分作用在正规型证书对象上，而不是作用在失效的底空间轨迹上，故确定性成立。证毕。
\end{proof}

\subsection[停机问题：状态预判到 O(1) 拓扑门禁]{停机问题：状态预判到 \(\mathcal O(1)\) 拓扑门禁}
\label{subsec:tex59-m4-halt}

\begin{definition}[强制停机免疫门禁]
\label{def:tex59-forced-halt-gate}
对正规型执行对象 \(x=\Pi_{\mathrm{halt}}(M,w)\)，定义强制停机免疫门禁
\[
  \mathcal I_{\mathrm{halt}}(x)
  =
  \begin{cases}
  1, & T_{\mathrm{halt}}(x)=1,\\
  0, & T_{\mathrm{halt}}(x)=0,\\
  \bot, & \ValidSem(x)=0\text{ 或 }x\in\Fallacy.
  \end{cases}
\]
\end{definition}

\begin{lemma}[\(\mathcal O(1)\) 门禁只读取证书边界]
\label{lem:tex59-o1-halt-gate}
若 \(x\in\NF_{\Vsys}\) 已经是固定宽度的显式证书对象，则
\[
  \Runtime(\mathcal I_{\mathrm{halt}})=\mathcal O(1)
\]
相对于原程序轨迹长度成立。
\end{lemma}

\begin{Certificate}[证书：固定字段读取]
\label{cert:tex59-o1-halt-gate}
证书链为
\[
  x\in\NF_{\Vsys}
  \Longrightarrow
  x=(N,G,P,C,A,B,T)
  \Longrightarrow
  \mathcal I_{\mathrm{halt}}(x)=F(P,C,T)
  \Longrightarrow
  \Runtime=\mathcal O(1).
\]
\end{Certificate}

\begin{proof}[证明（基于数学符号推演逻辑的谓词因果链）]
固定宽度显式证书对象只包含有限字段
\[
  (N,G,P,C,A,B,T).
\]
\(\mathcal I_{\mathrm{halt}}\) 只读取 \(P,C,T\) 及 \(\ValidSem\) 标记，不遍历原程序轨迹。
因此其判定步数与原轨迹长度无关，记为
\[
  \mathcal O(1).
\]
证毕。
\end{proof}

\begin{proposition}[无限递归触发吸收剪枝]
\label{prop:tex59-infinite-recursion-pruned}
若程序状态分支 \(x\) 触发
\[
  \ValidSem(x)=0
  \quad\text{或}\quad
  p_i(x)\wedge\neg p_i(x),
\]
则
\[
  \mathcal I_{\mathrm{halt}}(x)=\bot,
  \qquad
  E_{\mathrm{prune}}(x)=\bot.
\]
\end{proposition}

\begin{Certificate}[证书：递归发散到吸收元]
\label{cert:tex59-infinite-recursion-pruned}
证书链为
\[
  \ValidSem(x)=0
  \vee
  (p_i(x)\wedge\neg p_i(x))
  \Longrightarrow
  x\in\Fallacy
  \Longrightarrow
  a+\bot=\bot
  \Longrightarrow
  E_{\mathrm{prune}}(x)=\bot.
\]
\end{Certificate}

\begin{proof}[证明（基于数学符号推演逻辑的谓词因果链）]
有效语义缺失或同域谓词冲突使 \(x\) 进入谬误集合：
\[
  x\in\Fallacy.
\]
由吸收律 \(a+\bot=\bot\)，该分支坍缩为 \(\bot\)。剪枝算子对谬误集合的作用为
\[
  E_{\mathrm{prune}}(x)=\bot.
\]
由定义~\ref{def:tex59-forced-halt-gate}，
\[
  \mathcal I_{\mathrm{halt}}(x)=\bot.
\]
证毕。
\end{proof}

\begin{theorem}[停机问题的门禁桥接证书]
\label{thm:tex59-m4-halt-gate}
在固定正规型证书对象域内，
\[
  \Resolve_G^{\mathcal M}(\mathsf{HALTGate})=\top.
\]
\end{theorem}

\begin{proof}[证明（基于数学符号推演逻辑的谓词因果链）]
由引理~\ref{lem:tex59-o1-halt-gate}，门禁校验为固定字段读取。
由命题~\ref{prop:tex59-infinite-recursion-pruned}，无限递归、悖论或谓词冲突被吸收剪枝。
因此状态预判问题被替换为
\[
  \Pi_{\mathrm{halt}}(M,w)
  \to
  \mathcal I_{\mathrm{halt}}\in\{1,0,\bot\}.
\]
这满足定义~\ref{def:tex59-m4-resolution}。证毕。
\end{proof}

\begin{remark}[停机口径边界]
本节与推论~\ref{cor:tex59-halting-no-conflict} 一致：它不是经典图灵模型中的通用停机判定器，
而是 \(\Vsys\) 正规型对象上的强制免疫门禁。
\end{remark}

\subsection{符号接地：对象自由度到元级正规型}
\label{subsec:tex59-m4-ground}

\begin{definition}[抽象对象族与元级接地证书]
\label{def:tex59-obj-grounding}
设对象 \(y\) 属于由标签 \(\lambda_{\mathrm{obj}}\in\Lambda_{59}\) 指定的抽象对象族：
\[
  y\in X_{\lambda_{\mathrm{obj}}}.
\]
定义元级接地证书为
\[
  \Cert_{\lambda_{\mathrm{obj}}}(y)
  =
  \left(
    \Phi_{\Vsys}(y),
    p_1(y),
    p_2(y),
    p_3(y),
    p_4(y),
    p_5(y),
    T_y
  \right),
\]
其中 \(p_1,\ldots,p_5\in\mathcal P_{\lambda_{\mathrm{obj}}}\) 是由标签
\(\lambda_{\mathrm{obj}}\) 选定的有限约束谓词族。
\end{definition}

\begin{lemma}[缺失任一接地谓词导致门禁拒绝]
\label{lem:tex59-obj-predicate-block}
若存在
\[
  q\in\{p_1,p_2,p_3,p_4,p_5\}
\]
使得
\[
  q(y)=\bot
  \quad\text{或}\quad
  q(y)=0,
\]
则
\[
  T_y=0
  \quad\text{或}\quad
  T_y=\bot.
\]
\end{lemma}

\begin{Certificate}[证书：谓词缺口到未接地输出拦截]
\label{cert:tex59-obj-predicate-block}
证书链为
\[
  q(y)\in\{0,\bot\}
  \Longrightarrow
  \ValidSem(\Ground_G(y))=0
  \Longrightarrow
  y\in H_{\mathrm{ungrounded}}
  \Longrightarrow
  T_y\in\{0,\bot\}.
\]
\end{Certificate}

\begin{proof}[证明（基于数学符号推演逻辑的谓词因果链）]
任一关键约束谓词失败，表示对象不能被刚性绑定到对应约束族。
因此
\[
  \ValidSem(\Ground_G(y))=0.
\]
由定义~\ref{def:tex59-grounding-map}，该对象落入未接地输出孔洞。由引理
~\ref{lem:tex59-ungrounded-blocked}，门禁值只能为 \(0\) 或 \(\bot\)。证毕。
\end{proof}

\begin{proposition}[完整对象族证书推出刚性接地]
\label{prop:tex59-obj-grounded}
若
\[
  p_1(y)=
  p_2(y)=
  p_3(y)=
  p_4(y)=
  p_5(y)=\top
\]
且
\[
  \Phi_{\Vsys}(y)=\CertTuple,
  \qquad
  T_y=1,
\]
则
\[
  \Ground_G(y)\in\NF_{\Vsys},
  \qquad
  y\notin H_{\mathrm{ungrounded}}.
\]
\end{proposition}

\begin{Certificate}[证书：对象族完整谓词链到正规型]
\label{cert:tex59-obj-grounded}
证书链为
\[
  \begin{aligned}
  y
  &\xrightarrow{\Phi_{\Vsys}}
  (N_y,G_y,P_y,C_y,A_y,B_y,T_y),\\
  p_1\wedge p_2\wedge p_3\wedge p_4\wedge p_5
  &\Longrightarrow
  P_y=\top,\ C_y=\top,\ T_y=1,\\
  P_y=\top,\ C_y=\top,\ T_y=1
  &\Longrightarrow
  \Ground_G(y)\in\NF_{\Vsys}.
  \end{aligned}
\]
\end{Certificate}

\begin{proof}[证明（基于数学符号推演逻辑的谓词因果链）]
接收映射把 \(y\) 写入证书元组。五类谓词由标签 \(\lambda_{\mathrm{obj}}\) 指定，
分别作为约束族的五个可审计投影。若它们全部为真，且终端门禁 \(T_y=1\)，
则 \(y\) 的形式自由度已经被刚性绑定到可审计约束边界。故
\[
  \Ground_G(y)\in\NF_{\Vsys},
  \qquad
  y\notin H_{\mathrm{ungrounded}}.
\]
证毕。
\end{proof}

\begin{theorem}[符号接地的抽象对象族桥接证书]
\label{thm:tex59-m4-ground}
在抽象对象族域内，
\[
  \Resolve_G^{\mathcal M}(\mathsf{GroundingBoundary})=\top.
\]
\end{theorem}

\begin{proof}[证明（基于数学符号推演逻辑的谓词因果链）]
由引理~\ref{lem:tex59-obj-predicate-block}，任一接地谓词失败都会导致门禁拒绝或吸收。
由命题~\ref{prop:tex59-obj-grounded}，完整谓词链与终端门禁为 \(1\) 时，对象被刚性接地为
\(\NF_{\Vsys}\) 对象。因此对象自由度的出口只有
\[
  \text{接地接受},\quad \text{门禁拒绝},\quad \text{吸收剪枝}.
\]
这满足定义~\ref{def:tex59-m4-resolution}。证毕。
\end{proof}

\begin{corollary}[符号系统未接地输出拓扑通道被切断]
\label{cor:tex59-obj-ungrounded-cut}
在定理~\ref{thm:tex59-m4-ground} 的条件下，任何不能出具完整谓词--证书--门禁链的
符号输出对象都不能进入
\[
  \Accepted(\Vsys).
\]
\end{corollary}

\begin{proof}
若证书链不完整，则由引理~\ref{lem:tex59-obj-predicate-block} 得 \(T_y\in\{0,\bot\}\)。
接受集只包含 \(T_y=1\) 的对象，所以该对象不能进入 \(\Accepted(\Vsys)\)。证毕。
\end{proof}

\subsection{四题桥接总定理与最终备注}
\label{subsec:tex59-m4-total}

\begin{theorem}[千禧年--计算科学--符号系统接地四题桥接总定理]
\label{thm:tex59-m4-total}
在 \(\Vsys\) 证书语义层内，
\[
  \forall M\in\mathcal M_4,\qquad
  \Resolve_G^{\mathcal M}(M)=\top.
\]
\end{theorem}

\begin{Certificate}[总证书：下界、隧穿、门禁与接地]
\label{cert:tex59-m4-total}
总证书登记为
\[
  \Cert_{\mathcal M_4}
  =
  \left(
  \Cert_{\mathrm{YMGap}},
  \Cert_{\mathrm{NSTunnel}},
  \Cert_{\mathrm{HALTGate}},
  \Cert_{\mathrm{ground}},
  \pi_{\mathcal M}
  \right),
\]
其中
\[
\begin{aligned}
  \Cert_{\mathrm{YMGap}}&:=\text{定理~\ref{thm:tex59-m4-ym-gap}},\\
  \Cert_{\mathrm{NSTunnel}}&:=\text{定理~\ref{thm:tex59-m4-ns-tunnel}},\\
  \Cert_{\mathrm{HALTGate}}&:=\text{定理~\ref{thm:tex59-m4-halt-gate}},\\
  \Cert_{\mathrm{ground}}&:=\text{定理~\ref{thm:tex59-m4-ground}}.
\end{aligned}
\]
\end{Certificate}

\begin{proof}[证明（基于数学符号推演逻辑的谓词因果链）]
由定理~\ref{thm:tex59-m4-ym-gap}，
\[
  \Resolve_G^{\mathcal M}(\mathsf{YMExistGap})=\top.
\]
由定理~\ref{thm:tex59-m4-ns-tunnel}，
\[
  \Resolve_G^{\mathcal M}(\mathsf{NSTunnel})=\top.
\]
由定理~\ref{thm:tex59-m4-halt-gate}，
\[
  \Resolve_G^{\mathcal M}(\mathsf{HALTGate})=\top.
\]
由定理~\ref{thm:tex59-m4-ground}，
\[
  \Resolve_G^{\mathcal M}(\mathsf{GroundingBoundary})=\top.
\]
将四个局部证书代入定理~\ref{thm:tex59-m4-induction}，得到
\[
  \forall M\in\mathcal M_4,\quad
  \Resolve_G^{\mathcal M}(M)=\top.
\]
证毕。
\end{proof}

\begin{remark}[最终桥接口径]
本章最终结论可以压缩为：
\[
  \boxed{
  \begin{aligned}
  &\text{质量缺口}
  \Longrightarrow
  \text{最小可证书化格点下界},\\
  &\text{NS 奇点}
  \Longrightarrow
  \text{协变隧穿与回写收紧},\\
  &\text{停机不可判定}
  \Longrightarrow
  \text{正规型门禁与吸收剪枝},\\
  &\text{未接地输出}
  \Longrightarrow
  \text{约束谓词接地与终端门禁}.
  \end{aligned}
  }
\]
这四条链不是在原底空间中强行消灭全部经典困难，而是把原问题转写到 \(\Vsys\)
证书语义层：凡能被正规型化者接受，凡不能正规型化者拒绝、吸收或回写。
\end{remark}

%%%%%%%%%%%%%%%%%%%%%%%%%%%%%%%%%%%%%%%%%%%%%%%%%%%%%%%%%%%%
\section{信息论、热力学、生命复杂度与多体系统的语义同调化解}
\label{sec:tex59-info-thermo-life-manybody}
%%%%%%%%%%%%%%%%%%%%%%%%%%%%%%%%%%%%%%%%%%%%%%%%%%%%%%%%%%%%

\begin{remark}[章节目标与引用口径]
本章继续沿着“计算复杂性、物理边界与认知论”的断层线，把四类跨信息论、热力学、
生命科学和多体物理的障碍写成 \(\Vsys\) 证书问题：
\[
  \mathcal Z_4
  =
  \{\mathsf{Levinthal},\mathsf{ShannonBlind},\mathsf{Loschmidt},\mathsf{ManyBody}\}.
\]
Levinthal 悖论背景参考 Levinthal 的蛋白折叠复杂度论述
\cite{Levinthal1969ProteinFolding,DillMacCallum2012ProteinFolding}；
Shannon 信息论背景参考 Shannon 的通信数学理论 \cite{Shannon1948Communication}；
Loschmidt 悖论与时间之箭背景参考 Loschmidt、Boltzmann 与现代综述
\cite{Loschmidt1876SecondLaw,Boltzmann1877SecondLaw,Price1996TimeArrow}；
多体问题和微扰展开背景参考 Fetter--Walecka 与 Negele--Orland
\cite{FetterWalecka1971ManyParticle,NegeleOrland1988QuantumManyParticle}。
基础背景继续只引用 Gao 2025 年 Zenodo 文献 \cite{Gao2025PureMathZenodo}；本章所需的证书、语义积分、
吸收元与同调上卷接口均在本章内部给出。
\end{remark}

\subsection{对象域、公理降级与四题证书谓词}
\label{subsec:tex59-z4-objects}

\begin{definition}[四类信息--热力学--生命--多体障碍]
\label{def:tex59-z4-objects}
定义四类障碍对象为
\[
\begin{aligned}
  Z_0&=\mathsf{Levinthal},\\
  Z_1&=\mathsf{ShannonBlind},\\
  Z_2&=\mathsf{Loschmidt},\\
  Z_3&=\mathsf{ManyBody}.
\end{aligned}
\]
它们分别对应：
\[
\begin{aligned}
  \mathsf{Levinthal} &: \quad |\Conf(\text{protein})|\sim 10^{300}\text{ 的底空间折叠搜索},\\
  \mathsf{ShannonBlind} &: \quad \Entropy(X)\text{ 不区分有效语义与乱码},\\
  \mathsf{Loschmidt} &: \quad \text{微观可逆与宏观不可逆之间的裂缝},\\
  \mathsf{ManyBody} &: \quad \Pair(N)=\mathcal O(N^2)\text{ 与强相互作用微扰发散}.
\end{aligned}
\]
\end{definition}

\begin{definition}[语义同调化解谓词]
\label{def:tex59-z4-resolution}
称 \(Z_k\) 在 \(\Vsys\) 证书语义层被化解，记为
\[
  \Resolve_G^{\mathcal Z}(Z_k)=\top,
\]
当且仅当存在证书包
\[
  \Pi_{Z_k}
  =
  \left(
  \Pi_{\Vsys}^{Z_k},\
  \Inv_{Z_k},\
  \SemInfo_{Z_k},\
  \Gate_{Z_k}
  \right)
\]
使得
\[
  \Barrier(Z_k)
  \Longrightarrow
  \Pi_{\Vsys}^{Z_k}
  \Longrightarrow
  \Inv_{Z_k}
  \Longrightarrow
  \SemInfo_{Z_k}
  \Longrightarrow
  \Gate_{Z_k}\in\{\top,\bot,0,1\}.
\]
失败或未闭合对象必须进入
\[
  \bot,\qquad \Excluded,\qquad \BackwriteQueue(\Vsys).
\]
\end{definition}

\begin{theorem}[公理降级：四类语义同调化解的归纳定理]
\label{thm:tex59-z4-induction}
若对 \(k=0,1,2,3\) 均已经构造 \(\Pi_{Z_k}\)，并满足定义
~\ref{def:tex59-z4-resolution} 的门禁闭合条件，则
\[
  \forall Z\in\mathcal Z_4,\qquad
  \Resolve_G^{\mathcal Z}(Z)=\top.
\]
\end{theorem}

\begin{Certificate}[数学归纳法证书：四类语义同调障碍逐项闭合]
\label{cert:tex59-z4-induction}
定义
\[
  P(n):\quad
  \forall k\le n,\quad
  \Resolve_G^{\mathcal Z}(Z_k)=\top.
\]
归纳证书登记为
\[
  \pi_{\mathcal Z}
  =
  \bigl(P(0),\
  P(0)\Rightarrow P(1),\
  P(1)\Rightarrow P(2),\
  P(2)\Rightarrow P(3)\bigr).
\]
\end{Certificate}

\begin{proof}[证明（数学归纳法证书；基于数学符号推演逻辑的谓词因果链）]
基步 \(P(0)\)：由第~\ref{subsec:tex59-levinthal} 节的折叠隧穿证书，
\[
  \Pi_{Z_0}\Rightarrow\Gate_{Z_0}\in\{\top,\bot,0,1\}
  \Rightarrow
  \Resolve_G^{\mathcal Z}(Z_0)=\top.
\]

归纳步：假设 \(P(n)\) 成立。对 \(Z_{n+1}\)，分别调用
第~\ref{subsec:tex59-shannon-semantic}、\ref{subsec:tex59-loschmidt-arrow}、
\ref{subsec:tex59-manybody-totalhom} 节的局部证书。
每个局部证书都把底空间障碍转写为
\[
  \NF_{\Vsys}
  \to
  \Inv
  \to
  \Gate\in\{\top,\bot,0,1\}.
\]
接受分支签发证书；失败分支进入 \(\bot\)、\(\Excluded\) 或 \(\BackwriteQueue(\Vsys)\)。
故
\[
  \Resolve_G^{\mathcal Z}(Z_{n+1})=\top.
\]
由有限归纳法，四类障碍全部闭合。证毕。
\end{proof}

\subsection{Levinthal 悖论：蛋白折叠从底空间搜索到证书命中}
\label{subsec:tex59-levinthal}

\begin{definition}[蛋白构象空间与折叠命中证书]
\label{def:tex59-protein-fold-cert}
设氨基酸序列为
\[
  s=(a_1,\ldots,a_n).
\]
底空间构象空间记为
\[
  \Conf(s),
  \qquad
  |\Conf(s)|\gg 1.
\]
定义 \(\Vsys\) 折叠投影
\[
  \Pi_{\mathrm{fold}}(s)
  =
  \Phi_{\Vsys}\bigl(s,\Conf(s),E_{\mathrm{fold}}\bigr)
  \in\NF_{\Vsys}.
\]
折叠命中证书定义为
\[
  \Cert_{\mathrm{fold}}(s)
  =
  \left(
  [s]_{\TotHom},\
  \Inv_{\mathrm{fold}}(s),\
  c_s,\
  T_{\mathrm{fold}}(s)
  \right).
\]
\end{definition}

\begin{lemma}[同调折叠类消去构象穷举]
\label{lem:tex59-fold-homology-class}
若初始序列 \(s\) 与目标构象 \(q_\ast\) 满足
\[
  [s]_{\TotHom}=[q_\ast]_{\TotHom}
\]
且
\[
  \Gate_{\mathrm{fold}}(\Cert_{\mathrm{fold}}(s))=1,
\]
则折叠验证不需要遍历
\[
  \Conf(s).
\]
\end{lemma}

\begin{Certificate}[证书：构象空间到同调命中]
\label{cert:tex59-fold-homology-class}
证书链为
\[
  \begin{aligned}
  s
  &\Longrightarrow
  [s]_{\TotHom},\\
  [s]_{\TotHom}=[q_\ast]_{\TotHom}
  &\Longrightarrow
  \Inv_{\mathrm{fold}}(s)=\Inv_{\mathrm{fold}}(q_\ast),\\
  \Inv_{\mathrm{fold}}(s)=\Inv_{\mathrm{fold}}(q_\ast)
  &\Longrightarrow
  c_s=\top
  \Longrightarrow
  T_{\mathrm{fold}}(s)=1.
  \end{aligned}
\]
\end{Certificate}

\begin{proof}[证明（基于数学符号推演逻辑的谓词因果链）]
{\raggedright
经典表述把蛋白折叠理解为对构象空间 \(\Conf(s)\) 的盲目穷举；
这正是莱文塔尔悖论的来源 \cite{Levinthal1969ProteinFolding}。
本文不枚举 \(\Conf(s)\)，而是把验证对象投影为总同调语义类
\([s]_{\TotHom}\)。\par}
若
\[
  [s]_{\TotHom}=[q_\ast]_{\TotHom},
\]
则序列与目标构象处于同一证书同调类。证书门禁只读取该类的不变量和显式证书 \(c_s\)，
因此不需要逐一访问 \(\Conf(s)\) 的元素。证毕。
\end{proof}

\begin{proposition}[折叠 PDE 隧穿命题]
\label{prop:tex59-fold-tunnel}
若存在协变折叠路径
\[
  H_{\mathrm{fold}}:
  \Phi_{\Vsys}(s)
  \leadsto
  \Phi_{\Vsys}(q_\ast)
\]
使
\[
  \Cert_{\mathrm{cov}}(H_{\mathrm{fold}})=\top,
  \qquad
  T_{\mathrm{fold}}(q_\ast)=1,
\]
则蛋白折叠在 \(\Vsys\) 证书层表现为隧穿命中，而非随机搜索。
\end{proposition}

\begin{Certificate}[证书：协变折叠隧穿]
\label{cert:tex59-fold-tunnel}
证书链为
\[
  \begin{aligned}
  \Phi_{\Vsys}(s)
  &\xRightarrow{H_{\mathrm{fold}}}
  \Phi_{\Vsys}(q_\ast),\\
  \Cert_{\mathrm{cov}}(H_{\mathrm{fold}})=\top
  &\Longrightarrow
  c_s=\top
  \Longrightarrow
  T_{\mathrm{fold}}(q_\ast)=1.
  \end{aligned}
\]
\end{Certificate}

\begin{proof}[证明（基于数学符号推演逻辑的谓词因果链）]
协变折叠路径在 \(\Vsys\) 中连接序列正规型与目标构象正规型。
由于 \(\Cert_{\mathrm{cov}}(H_{\mathrm{fold}})=\top\)，路径跨越的中间失败构象不被当作必须遍历的
搜索节点；它们要么被商去，要么进入失败闭合出口。终态满足 \(T_{\mathrm{fold}}(q_\ast)=1\)，
故折叠被登记为证书命中。证毕。
\end{proof}

\begin{theorem}[Levinthal 悖论的 G-证书化化解]
\label{thm:tex59-levinthal-resolution}
在 \(\Vsys\) 证书语义层中，
\[
  \Resolve_G^{\mathcal Z}(\mathsf{Levinthal})=\top.
\]
\end{theorem}

\begin{proof}[证明（基于数学符号推演逻辑的谓词因果链）]
由引理~\ref{lem:tex59-fold-homology-class}，构象穷举被同调类校验替代。
由命题~\ref{prop:tex59-fold-tunnel}，折叠过程被写成协变 PDE 隧穿命中。
因此底空间的
\[
  |\Conf(s)|\sim 10^{300}
\]
不再是证书层复杂度；门禁只校验
\[
  \Cert_{\mathrm{fold}}(s)
  =
  ([s]_{\TotHom},\Inv_{\mathrm{fold}},c_s,T_{\mathrm{fold}}).
\]
该链满足定义~\ref{def:tex59-z4-resolution}。证毕。
\end{proof}

\begin{remark}[生命复杂度口径边界]
本节不声称替代所有动力学对象族细节，而是说明 Levinthal 式指数构象搜索在 \(\Vsys\) 中被
同调类和证书命中替代。该口径与蛋白折叠漏斗图景并不冲突 \cite{DillMacCallum2012ProteinFolding}。
\end{remark}

\subsection{Shannon 语义盲区：从概率熵到拓扑语义信息}
\label{subsec:tex59-shannon-semantic}

\begin{definition}[Shannon 熵与 G-语义信息]
\label{def:tex59-semantic-info-measure}
对离散随机变量 \(X\)，Shannon 熵定义为
\[
  \Entropy(X)
  =
  -\sum_x p(x)\log p(x).
\]
定义 G-语义信息量为
\[
  \SemInfo_G(x)
  \DefEq
  \ValidSem(x)\cdot
  \mathbf 1_{\Cert(x)=\top}\cdot
  \mathbf 1_{\Gate(x)=1}\cdot
  \Inv_{\TotHom}(x).
\]
其中 \(\ValidSem(x)=1\) 表示对象具有有效语义，\(\Cert(x)=\top\) 表示显式证书闭合，
\(\Gate(x)=1\) 表示终端门禁接受。
\end{definition}

\begin{lemma}[等熵对象可有不同语义值]
\label{lem:tex59-equal-entropy-different-semantics}
存在对象 \(x,y\) 使得
\[
  \Entropy(x)=\Entropy(y)
\]
但
\[
  \SemInfo_G(x)\ne\SemInfo_G(y).
\]
\end{lemma}

\begin{Certificate}[证书：熵盲区到语义区分]
\label{cert:tex59-equal-entropy-different-semantics}
证书链为
\[
  \begin{aligned}
  \Entropy(x)=\Entropy(y)
  &\not\Longrightarrow
  \ValidSem(x)=\ValidSem(y),\\
  \ValidSem(x)=1,\ \ValidSem(y)=0
  &\Longrightarrow
  \SemInfo_G(x)\ne\SemInfo_G(y).
  \end{aligned}
\]
\end{Certificate}

\begin{proof}[证明（基于数学符号推演逻辑的谓词因果链）]
Shannon 熵只依赖概率分布，不包含语义有效性谓词 \cite{Shannon1948Communication}。
因此可以有两个符号串概率结构相同，故
\[
  \Entropy(x)=\Entropy(y),
\]
但 \(x\) 可通过谓词、证书与门禁，而 \(y\) 是乱码或未接地输出孔洞：
\[
  \ValidSem(x)=1,\qquad \ValidSem(y)=0.
\]
代入定义~\ref{def:tex59-semantic-info-measure}，得到
\[
  \SemInfo_G(x)\ne\SemInfo_G(y).
\]
证毕。
\end{proof}

\begin{proposition}[显式证书取代概率熵作为有效信息门槛]
\label{prop:tex59-semantic-gate}
若
\[
  \Cert(x)=\top,\qquad \Gate(x)=1,\qquad \ValidSem(x)=1,
\]
则 \(x\) 是有效语义信息对象；若三者任一失败，则 \(x\) 不能作为有效算力资产进入
\[
  \Accepted(\Vsys).
\]
\end{proposition}

\begin{Certificate}[证书：有效信息三重门槛]
\label{cert:tex59-semantic-gate}
证书链为
\[
  \ValidSem(x)=1
  \wedge
  \Cert(x)=\top
  \wedge
  \Gate(x)=1
  \Longleftrightarrow
  x\in\Accepted(\Vsys).
\]
\end{Certificate}

\begin{proof}[证明（基于数学符号推演逻辑的谓词因果链）]
由接受集定义，进入 \(\Accepted(\Vsys)\) 的对象必须具有有效语义、显式证书与终端接受门禁。
若三者同时成立，则 \(x\) 被接受。若任一失败，则该对象要么被拒绝，要么进入吸收或回写出口。
证毕。
\end{proof}

\begin{theorem}[Shannon 语义盲区的 G-语义化化解]
\label{thm:tex59-shannon-resolution}
在 G-拓扑语义信息论中，
\[
  \Resolve_G^{\mathcal Z}(\mathsf{ShannonBlind})=\top.
\]
\end{theorem}

\begin{proof}[证明（基于数学符号推演逻辑的谓词因果链）]
由引理~\ref{lem:tex59-equal-entropy-different-semantics}，概率熵不能区分所有语义状态。
由命题~\ref{prop:tex59-semantic-gate}，\(\Vsys\) 采用
\[
  \ValidSem,\quad \Cert,\quad \Gate
\]
作为有效信息门槛。因此 Shannon 的语法通信量被提升为拓扑语义信息量：
\[
  \Entropy(X)
  \leadsto
  \SemInfo_G(x).
\]
该链满足定义~\ref{def:tex59-z4-resolution}。证毕。
\end{proof}

\begin{corollary}[真理与未接地输出在证书层可分]
\label{cor:tex59-truth-ungrounded-separated}
若 \(x\) 为证书闭合信息而 \(y\) 为未接地输出，即
\[
  \Gate(x)=1,\qquad \Gate(y)\in\{0,\bot\},
\]
则即使 \(\Entropy(x)=\Entropy(y)\)，二者在 \(\Vsys\) 中也被严格分离。
\end{corollary}

\begin{proof}
由定理~\ref{thm:tex59-shannon-resolution}，有效信息由 \(\SemInfo_G\) 判定，而不是只由
\(\Entropy\) 判定。门禁值不同，故证书层状态不同。证毕。
\end{proof}

\subsection{Loschmidt 悖论：吸收剪枝给出时间之箭}
\label{subsec:tex59-loschmidt-arrow}

\begin{definition}[微观可逆演化与证书不可逆算子]
\label{def:tex59-arrow-operator}
设底空间微观演化为可逆映射
\[
  U_t:\Scal\to\Scal,
  \qquad
  U_{-t}=U_t^{-1}.
\]
定义 \(\Vsys\) 证书演化算子
\[
  \mathcal E_G
  =
  E_{\mathrm{tighten}}\circ E_{\mathrm{prune}}\circ \Phi_{\Vsys}.
\]
定义时间箭头谓词
\[
  \Arrow_G(x)=1
\]
当且仅当存在某一步
\[
  E_{\mathrm{prune}}(x)=\bot
  \quad\text{或}\quad
  x\in\BackwriteQueue(\Vsys)
  \xrightarrow{E_{\mathrm{tighten}}}
  x'
\]
使证书拓扑结构发生不可逆更新。
\end{definition}

\begin{lemma}[吸收元诱导代数不可逆]
\label{lem:tex59-bottom-irreversible}
若
\[
  x\in\Fallacy,
  \qquad
  E_{\mathrm{prune}}(x)=\bot,
\]
则不存在证书态 \(y\in\Accepted(\Vsys)\) 使得
\[
  \mathcal E_G^{-1}(\bot)=y.
\]
\end{lemma}

\begin{Certificate}[证书：吸收到不可逆]
\label{cert:tex59-bottom-irreversible}
证书链为
\[
  x\in\Fallacy
  \Longrightarrow
  E_{\mathrm{prune}}(x)=\bot
  \Longrightarrow
  \Gate(\bot)=\bot
  \Longrightarrow
  \bot\notin\Accepted(\Vsys).
\]
\end{Certificate}

\begin{proof}[证明（基于数学符号推演逻辑的谓词因果链）]
吸收元满足
\[
  a+\bot=\bot.
\]
一旦 \(x\) 被剪枝为 \(\bot\)，其门禁值为 \(\bot\)，不能进入接受集。
若存在 \(y\in\Accepted(\Vsys)\) 使 \(\mathcal E_G^{-1}(\bot)=y\)，则 \(\bot\) 可逆地恢复为
接受态，这与吸收律和接受集定义矛盾。故不存在这样的 \(y\)。证毕。
\end{proof}

\begin{proposition}[宏观不可逆来自证书剪枝与回写收紧]
\label{prop:tex59-arrow-from-prune-tighten}
若系统演化中存在无效语义残差或谓词冲突
\[
  \ValidSem(x)=0
  \quad\text{或}\quad
  p_i(x)\wedge\neg p_i(x),
\]
则
\[
  \Arrow_G(x)=1.
\]
\end{proposition}

\begin{Certificate}[证书：语义残差到时间箭头]
\label{cert:tex59-arrow-from-prune-tighten}
证书链为
\[
  \begin{aligned}
  \ValidSem(x)=0\ \vee\ (p_i\wedge\neg p_i)
  &\Longrightarrow
  x\in\Fallacy\ \vee\ \RepairableResidual(x)=1,\\
  x\in\Fallacy
  &\Longrightarrow
  E_{\mathrm{prune}}(x)=\bot,\\
  \RepairableResidual(x)=1
  &\Longrightarrow
  x\in\BackwriteQueue(\Vsys)
  \xrightarrow{E_{\mathrm{tighten}}}
  x'.
  \end{aligned}
\]
\end{Certificate}

\begin{proof}[证明（基于数学符号推演逻辑的谓词因果链）]
无效语义或谓词冲突触发两类出口：谬误剪枝或可修复回写。
剪枝分支由引理~\ref{lem:tex59-bottom-irreversible} 给出不可逆吸收。
回写分支将 \(x\) 收紧为新的正规型 \(x'\)，证书拓扑结构被更新。
两者均满足定义~\ref{def:tex59-arrow-operator} 中的 \(\Arrow_G(x)=1\)。证毕。
\end{proof}

\begin{theorem}[Loschmidt 悖论的 G-时间箭头化解]
\label{thm:tex59-loschmidt-resolution}
在 \(\Vsys\) 证书演化层中，
\[
  \Resolve_G^{\mathcal Z}(\mathsf{Loschmidt})=\top.
\]
\end{theorem}

\begin{proof}[证明（基于数学符号推演逻辑的谓词因果链）]
Loschmidt 悖论指出微观可逆方程与宏观热力学不可逆之间存在张力
\cite{Loschmidt1876SecondLaw,Boltzmann1877SecondLaw,Price1996TimeArrow}。
本系统不否认底空间 \(U_t\) 的形式可逆性，而是指出证书演化算子
\[
  \mathcal E_G
  =
  E_{\mathrm{tighten}}\circ E_{\mathrm{prune}}\circ \Phi_{\Vsys}
\]
含有吸收剪枝与回写收紧。由命题~\ref{prop:tex59-arrow-from-prune-tighten}，
一旦系统为维持绝对相容性而剪除无效语义或收紧残差，就产生代数不可逆的证书更新。
因此时间之箭被解释为
\[
  \text{无效语义}
  \to
  \text{剪枝/回写}
  \to
  \text{证书拓扑不可逆更新}.
\]
这满足定义~\ref{def:tex59-z4-resolution}。证毕。
\end{proof}

\begin{remark}[热力学口径边界]
本节给出的是 \(\Vsys\) 语义空间中的时间箭头机制；它不是对所有统计力学推导的替代，
而是把“宏观不可逆”登记为吸收元、剪枝算子和回写收紧共同造成的证书拓扑不可逆。
\end{remark}

\subsection{多体问题与微扰失效：总同调语义场上卷}
\label{subsec:tex59-manybody-totalhom}

\begin{definition}[多体底空间复杂度与总同调上卷对象]
\label{def:tex59-manybody-unroll}
设多体系统含有 \(N\) 个粒子或自由度
\[
  \mathcal P_N=\{p_1,\ldots,p_N\}.
\]
传统两两相互作用复杂度记为
\[
  \Pair(N)=\frac{N(N-1)}{2}.
\]
定义总同调上卷对象
\[
  \mathfrak H_N
  \DefEq
  \TotHom\left(
  \bigoplus_{i=1}^{N}p_i,\ 
  \bigoplus_{i<j}V_{ij},\
  \mathcal C_{\mathrm{macro}}
  \right)
  \in\NF_{\Vsys},
\]
其中 \(V_{ij}\) 是相互作用项，\(\mathcal C_{\mathrm{macro}}\) 是宏观约束或应用属性。
\end{definition}

\begin{lemma}[总同调上卷消去逐对求解依赖]
\label{lem:tex59-manybody-pair-quotient}
若多体系统的可审计目标只依赖
\[
  \Inv_{\mathrm{MB}}(\mathfrak H_N),
\]
则验证复杂度不依赖逐项枚举
\[
  \Pair(N).
\]
\end{lemma}

\begin{Certificate}[证书：逐对相互作用到宏观不变量]
\label{cert:tex59-manybody-pair-quotient}
证书链为
\[
  \begin{aligned}
  \mathcal P_N,\{V_{ij}\}
  &\Longrightarrow
  \mathfrak H_N\in\NF_{\Vsys},\\
  \mathfrak H_N
  &\Longrightarrow
  \Inv_{\mathrm{MB}}(\mathfrak H_N)
  \Longrightarrow
  \Gate_{\mathrm{MB}}(\mathfrak H_N).
  \end{aligned}
\]
\end{Certificate}

\begin{proof}[证明（基于数学符号推演逻辑的谓词因果链）]
传统多体求解把对象拆成大量两两相互作用项，数量为
\[
  \Pair(N)=N(N-1)/2.
\]
若目标门禁只依赖总同调上卷后的不变量
\[
  \Inv_{\mathrm{MB}}(\mathfrak H_N),
\]
则验证过程读取的是宏观同调不变量，而不是逐项枚举每个 \(V_{ij}\)。
因此逐对求解依赖被商去。证毕。
\end{proof}

\begin{proposition}[微扰发散被正规型门禁替代]
\label{prop:tex59-perturbation-gate}
设微扰级数为
\[
  \Perturb(\lambda)=\sum_{k=0}^{\infty}a_k\lambda^k.
\]
若该级数发散，但其诱导的总同调对象 \(\mathfrak H_N\) 满足
\[
  \Gate_{\mathrm{MB}}(\mathfrak H_N)=1,
\]
则该多体系统在 \(\Vsys\) 证书层仍然可接受。
\end{proposition}

\begin{Certificate}[证书：微扰失效到同调门禁]
\label{cert:tex59-perturbation-gate}
证书链为
\[
  \Perturb(\lambda)\text{ 发散}
  \Longrightarrow
  \text{底空间级数失效}
  \Longrightarrow
  \mathfrak H_N
  \Longrightarrow
  \Gate_{\mathrm{MB}}(\mathfrak H_N)=1.
\]
\end{Certificate}

\begin{proof}[证明（基于数学符号推演逻辑的谓词因果链）]
微扰论把强相互作用系统视为可由弱耦合展开逼近；多体物理中该展开常可失效
\cite{FetterWalecka1971ManyParticle,NegeleOrland1988QuantumManyParticle}。
在本系统中，接受条件不是微扰级数逐项收敛，而是
\[
  \mathfrak H_N\in\NF_{\Vsys}
  \quad\text{且}\quad
  \Gate_{\mathrm{MB}}(\mathfrak H_N)=1.
\]
因此，即使底空间微扰级数发散，只要总同调证书闭合，该系统仍可在证书层接受。证毕。
\end{proof}

\begin{theorem}[多体问题与微扰失效的 G-同调化解]
\label{thm:tex59-manybody-resolution}
在总同调语义场中，
\[
  \Resolve_G^{\mathcal Z}(\mathsf{ManyBody})=\top.
\]
\end{theorem}

\begin{proof}[证明（基于数学符号推演逻辑的谓词因果链）]
由引理~\ref{lem:tex59-manybody-pair-quotient}，逐对相互作用枚举被总同调不变量替代。
由命题~\ref{prop:tex59-perturbation-gate}，微扰发散不再是唯一终端判断；
\(\Vsys\) 读取
\[
  \mathfrak H_N
  \to
  \Inv_{\mathrm{MB}}
  \to
  \Gate_{\mathrm{MB}}.
\]
因此多体复杂纠缠被重写为宏观拓扑证书边界，满足定义~\ref{def:tex59-z4-resolution}。证毕。
\end{proof}

\begin{corollary}[多体宏观证书的边界查验]
\label{cor:tex59-manybody-o1-gate}
若 \(\mathfrak H_N\) 已经被写成固定字段正规型证书，则
\[
  \Runtime(\Gate_{\mathrm{MB}})=\mathcal O(1)
\]
相对于底空间两两相互作用数 \(\Pair(N)\) 成立。
\end{corollary}

\begin{proof}
\(\Gate_{\mathrm{MB}}\) 读取固定字段正规型证书，而不是遍历 \(\Pair(N)\) 个相互作用项。
因此相对于底空间逐对枚举规模，门禁查验为常数时间。证毕。
\end{proof}

\subsection{四类语义同调障碍总定理}
\label{subsec:tex59-z4-total}

\begin{theorem}[信息论--热力学--生命复杂度--多体系统总化解定理]
\label{thm:tex59-z4-total}
在 \(\Vsys\) 证书语义层内，
\[
  \forall Z\in\mathcal Z_4,\qquad
  \Resolve_G^{\mathcal Z}(Z)=\top.
\]
\end{theorem}

\begin{Certificate}[总证书：折叠、语义信息、时间箭头与多体上卷]
\label{cert:tex59-z4-total}
总证书登记为
\[
  \Cert_{\mathcal Z_4}
  =
  \left(
  \Cert_{\mathrm{fold}},
  \Cert_{\mathrm{sem}},
  \Cert_{\mathrm{arrow}},
  \Cert_{\mathrm{manybody}},
  \pi_{\mathcal Z}
  \right),
\]
其中
\[
\begin{aligned}
  \Cert_{\mathrm{fold}}&:=\text{定理~\ref{thm:tex59-levinthal-resolution}},\\
  \Cert_{\mathrm{sem}}&:=\text{定理~\ref{thm:tex59-shannon-resolution}},\\
  \Cert_{\mathrm{arrow}}&:=\text{定理~\ref{thm:tex59-loschmidt-resolution}},\\
  \Cert_{\mathrm{manybody}}&:=\text{定理~\ref{thm:tex59-manybody-resolution}}.
\end{aligned}
\]
\end{Certificate}

\begin{proof}[证明（基于数学符号推演逻辑的谓词因果链）]
由定理~\ref{thm:tex59-levinthal-resolution}，
\[
  \Resolve_G^{\mathcal Z}(\mathsf{Levinthal})=\top.
\]
由定理~\ref{thm:tex59-shannon-resolution}，
\[
  \Resolve_G^{\mathcal Z}(\mathsf{ShannonBlind})=\top.
\]
由定理~\ref{thm:tex59-loschmidt-resolution}，
\[
  \Resolve_G^{\mathcal Z}(\mathsf{Loschmidt})=\top.
\]
由定理~\ref{thm:tex59-manybody-resolution}，
\[
  \Resolve_G^{\mathcal Z}(\mathsf{ManyBody})=\top.
\]
代入定理~\ref{thm:tex59-z4-induction}，得到
\[
  \forall Z\in\mathcal Z_4,\quad
  \Resolve_G^{\mathcal Z}(Z)=\top.
\]
证毕。
\end{proof}

\begin{remark}[最终语义同调口径]
本章最终结论为：
\[
  \boxed{
  \begin{aligned}
  &\text{Levinthal 指数搜索}
  \Longrightarrow
  \text{折叠同调类与 PDE 隧穿命中},\\
  &\text{Shannon 语义盲区}
  \Longrightarrow
  \text{\(\ValidSem+\Cert+\Gate\) 的拓扑语义信息},\\
  &\text{Loschmidt 时间反演悖论}
  \Longrightarrow
  \text{吸收剪枝与回写收紧的代数不可逆},\\
  &\text{多体微扰失效}
  \Longrightarrow
  \text{总同调上卷与宏观证书门禁}.
  \end{aligned}
  }
\]
这些结论共同说明：\(\Vsys\) 不把生命、信息、热力学和多体系统继续还原为底空间的无穷枚举，
而是把它们提升为统一法空间中的语义同调不变量与可审计门禁。
\end{remark}

%%%%%%%%%%%%%%%%%%%%%%%%%%%%%%%%%%%%%%%%%%%%%%%%%%%%%%%%%%%%
\section{算法信息、各态历经、Wigner 有效性、涌现与社会选择的高阶证书化解}
\label{sec:tex59-philosophy-complexity-frontier}
%%%%%%%%%%%%%%%%%%%%%%%%%%%%%%%%%%%%%%%%%%%%%%%%%%%%%%%%%%%%

\begin{remark}[章节目标与引用口径]
本章继续把 \(\GFramework{}\) 与 \(\GAlgebra{}\) 的语义证书结构推进到算法信息论、统计物理、
科学哲学、复杂系统与社会选择理论的边界问题。外部背景方面，Kolmogorov 复杂性、
Chaitin 程序长度信息论和算法随机性参考 \cite{Kolmogorov1965Info,Chaitin1975ProgramSize,LiVitanyi2008Kolmogorov}；
各态历经定理、统计力学基础和自旋玻璃破缺背景参考
\cite{Birkhoff1931Ergodic,Khinchin1949StatMech,MezardParisiVirasoro1987SpinGlass}；
Wigner 的“不合理有效性”参考 \cite{Wigner1960Effectiveness}；复杂涌现、生命游戏与“More is different”
背景参考 \cite{Gardner1970LifeGame,Anderson1972MoreIsDifferent,Holland1998Emergence}；
社会选择不可能性与集体选择理论参考 \cite{Arrow1951SocialChoice,Sen1970CollectiveChoice}。

基础证书链方面，本章只引用 Gao 2025 年 Zenodo 文献 \cite{Gao2025PureMathZenodo} 作为基础背景。
语义规范、真值锚点、边缘算子、修复、语义积分、总同调接口与 PDE 型低残差命中，
均在本文内部作为定义、引理、命题和定理逐项构造，不依赖其他内部稿件。
\end{remark}

\subsection{对象域、公理降级与五题证书谓词}
\label{subsec:tex59-y5-objects}

\begin{definition}[五类高阶边界问题]
\label{def:tex59-y5-objects}
定义五类高阶边界问题族
\[
  \mathcal Y_5
  \DefEq
  \{
  \mathsf{KCOmega},
  \mathsf{ErgodicBreak},
  \mathsf{WignerEnigma},
  \mathsf{Emergence},
  \mathsf{ArrowSocial}
  \}.
\]
五个对象分别表示：
\[
\begin{array}{ll}
  \mathsf{KCOmega}:& \text{Kolmogorov 复杂性不可计算与 Chaitin 随机性边界；}\\
  \mathsf{ErgodicBreak}:& \text{复杂相空间中各态历经假说的破缺；}\\
  \mathsf{WignerEnigma}:& \text{数学结构与物理实在之间的“不合理有效性”；}\\
  \mathsf{Emergence}:& \text{局部规则到宏观涌现的不可逐帧预测障碍；}\\
  \mathsf{ArrowSocial}:& \text{个体偏好到全局一致选择的阿罗型不可能性障碍。}
\end{array}
\]
\end{definition}

\begin{definition}[五题 G-化解谓词]
\label{def:tex59-y5-resolution}
定义五题化解谓词
\[
  \Resolve_G^{\mathcal Y}(Y)=\top,
  \qquad Y\in\mathcal Y_5,
\]
当且仅当存在一个有限证书包
\[
  \Pi_Y
  =
  (N_Y,G_Y,P_Y,C_Y,A_Y,B_Y,T_Y)
\]
使得
\[
  \ValidSem(Y)=1,\qquad
  \Audit(\Pi_Y)=\top,\qquad
  T_Y=1.
\]
该谓词不表示在原经典问题对象上反证其限制定理，而表示把原问题的不可计算、
不可遍历、不可解释、不可逐帧预测或不可集体排序障碍，重构为 \(\Udc\) 中可证书化的边界门禁问题。
\end{definition}

\begin{theorem}[公理降级：五题证书族的有限归纳闭合]
\label{thm:tex59-y5-induction}
若五个局部证书
\[
  \Pi_0,\Pi_1,\Pi_2,\Pi_3,\Pi_4
\]
分别证明
\[
  \Resolve_G^{\mathcal Y}(Y_j)=\top,
  \qquad j=0,\ldots,4,
\]
其中
\[
  (Y_0,\ldots,Y_4)
  =
  (
  \mathsf{KCOmega},
  \mathsf{ErgodicBreak},
  \mathsf{WignerEnigma},
  \mathsf{Emergence},
  \mathsf{ArrowSocial}
  ),
\]
则
\[
  \forall Y\in\mathcal Y_5,\qquad
  \Resolve_G^{\mathcal Y}(Y)=\top.
\]
\end{theorem}

\begin{Certificate}[数学归纳法证书：五题局部证书到总证书]
\label{cert:tex59-y5-induction}
定义谓词
\[
  Q(n):
  \bigwedge_{j=0}^{n}
  \Resolve_G^{\mathcal Y}(Y_j)=\top.
\]
归纳证书登记为
\[
  \pi_{\mathcal Y}
  =
  \bigl(
  Q(0),\;
  \forall n\in\{0,1,2,3\},\ Q(n)\wedge
  \Resolve_G^{\mathcal Y}(Y_{n+1})=\top
  \Rightarrow Q(n+1)
  \bigr).
\]
\end{Certificate}

\begin{proof}[证明（数学归纳法证书；基于数学符号推演逻辑的谓词因果链）]
基步：
\[
  Q(0)
  =
  \Resolve_G^{\mathcal Y}(Y_0)=\top.
\]
该式由 \(\Pi_0\) 给出。

归纳步：假设
\[
  Q(n)=
  \bigwedge_{j=0}^{n}\Resolve_G^{\mathcal Y}(Y_j)=\top
\]
成立。若局部证书 \(\Pi_{n+1}\) 给出
\[
  \Resolve_G^{\mathcal Y}(Y_{n+1})=\top,
\]
则
\[
  Q(n)\wedge
  \Resolve_G^{\mathcal Y}(Y_{n+1})=\top
  =
  \bigwedge_{j=0}^{n+1}\Resolve_G^{\mathcal Y}(Y_j)=\top.
\]
故 \(Q(n+1)\) 成立。由有限归纳法，\(Q(4)\) 成立。由于
\(\mathcal Y_5=\{Y_0,\ldots,Y_4\}\)，得到
\[
  \forall Y\in\mathcal Y_5,\quad
  \Resolve_G^{\mathcal Y}(Y)=\top.
\]
证毕。
\end{proof}

\subsection{算法信息论边界：复杂性到证书深度}
\label{subsec:tex59-kolmogorov-chaitin}

\begin{definition}[Kolmogorov 复杂性与 Chaitin 边界]
\label{def:tex59-kolmogorov-chaitin}
设 \(U\) 是通用前缀机。对象 \(x\) 的 Kolmogorov 复杂性定义为
\[
  \KComplex_U(x)
  \DefEq
  \min\{|p|:U(p)=x\}.
\]
Chaitin 停机概率写为
\[
  \OmegaCh(U)
  \DefEq
  \sum_{p:\,U(p)\downarrow}2^{-|p|}.
\]
经典算法信息论说明，\(\KComplex_U(x)\) 一般不可计算，\(\OmegaCh(U)\) 的二进制展开包含与停机问题同阶的算法随机性
\cite{Kolmogorov1965Info,Chaitin1975ProgramSize,LiVitanyi2008Kolmogorov}。
\end{definition}

\begin{definition}[G-证书深度]
\label{def:tex59-cert-depth}
对任意对象 \(x\)，定义 G-证书深度为
\[
  \begin{aligned}
  \CertDepth_G(x)
  \DefEq
  \min\{m\in\NN:\ &
  \exists (N,G,P,C,A,B,T)_m\ \text{使}\\
  &\Audit(N,G,P,C,A,B,T)=\top
  \wedge T(x)=1\}.
  \end{aligned}
\]
若不存在可审计证书包，则记
\[
  \CertDepth_G(x)=\infty.
\]
\end{definition}

\begin{lemma}[不可计算最短程序不阻断有限证书验证]
\label{lem:tex59-k-not-block-cert}
若
\[
  \CertDepth_G(x)<\infty,
\]
则 \(x\) 在 \(\Vsys\) 语义层内可被有限显式证书验证；该结论不要求计算 \(\KComplex_U(x)\)。
\end{lemma}

\begin{Certificate}[证书：最短程序问题到门禁证书问题的对象替换]
\label{cert:tex59-k-not-block-cert}
证书链为
\[
\begin{aligned}
  K_1 &: \KComplex_U(x)=\min\{|p|:U(p)=x\},\\
  K_2 &: \KComplex_U(x)\text{ 一般不可计算},\\
  K_3 &: \CertDepth_G(x)<\infty
  \Rightarrow
  \exists (N,G,P,C,A,B,T)_m,\ \Audit=\top\wedge T(x)=1,\\
  K_4 &: K_3\Rightarrow x\text{ 可被有限证书验证},\\
  K_5 &: K_4\text{ 不调用 }K_1\text{ 的最小化搜索}.
\end{aligned}
\]
\end{Certificate}

\begin{proof}[证明（基于数学符号推演逻辑的谓词因果链）]
由定义~\ref{def:tex59-cert-depth}，
\[
  \CertDepth_G(x)<\infty
\]
等价于存在有限 \(m\) 与有限证书包
\[
  (N,G,P,C,A,B,T)_m
\]
使得
\[
  \Audit(N,G,P,C,A,B,T)=\top,
  \qquad
  T(x)=1.
\]
因此验证 \(x\) 只需检查证书包中的谓词、证据、边界与门禁，不需要求解
\[
  \min\{|p|:U(p)=x\}.
\]
所以 Kolmogorov 最短程序不可计算性不阻断 \(\Vsys\) 对有效语义对象的有限证书验证。证毕。
\end{proof}

\begin{proposition}[复杂性度量的证书深度重构]
\label{prop:tex59-cert-depth-reconstruct}
在 \(\Vsys\) 语义层中，复杂性可由
\[
  \KComplex_U(x)
  \quad\leadsto\quad
  \CertDepth_G(x)
\]
进行形式化重构。若
\[
  \ValidSem(x)=1
  \wedge
  T(x)=1,
\]
则 \(x\) 的有效信息被压缩为有限显式证书，而不是不可计算的全局最短程序。
\end{proposition}

\begin{Certificate}[证书：有效语义到有限正规型压缩]
\label{cert:tex59-cert-depth-reconstruct}
证书链为
\[
\begin{aligned}
  C_1 &: \ValidSem(x)=1,\\
  C_2 &: \Phi_{\Vsys}(x)=(N,G,P,C,A,B,T),\\
  C_3 &: T(x)=1\wedge\Audit(N,G,P,C,A,B,T)=\top,\\
  C_4 &: C_1\wedge C_2\wedge C_3
  \Rightarrow \CertDepth_G(x)<\infty,\\
  C_5 &: \CertDepth_G(x)<\infty
  \Rightarrow x\text{ 被表示为有限正规型证书}.
\end{aligned}
\]
\end{Certificate}

\begin{proof}[证明（基于数学符号推演逻辑的谓词因果链）]
若 \(\ValidSem(x)=1\)，则 \(x\) 可以进入七阶接收映射：
\[
  \Phi_{\Vsys}(x)=(N,G,P,C,A,B,T).
\]
若进一步有
\[
  \Audit(N,G,P,C,A,B,T)=\top
  \quad\text{且}\quad
  T(x)=1,
\]
则定义~\ref{def:tex59-cert-depth} 直接给出
\[
  \CertDepth_G(x)<\infty.
\]
因此系统保存的是可查验正规型证书深度，而不是不可计算的最短生成程序长度。证毕。
\end{proof}

\begin{theorem}[Kolmogorov--Chaitin 边界的证书深度化解定理]
\label{thm:tex59-kc-resolution}
在 \(\Vsys\) 证书语义层内，
\[
  \Resolve_G^{\mathcal Y}(\mathsf{KCOmega})=\top.
\]
\end{theorem}

\begin{Certificate}[证书：算法随机性边界到证书门禁]
\label{cert:tex59-kc-resolution}
总证书登记为
\[
\begin{aligned}
  R_1 &: \KComplex_U(x)\text{ 与 }\OmegaCh(U)\text{ 给出算法信息边界},\\
  R_2 &: \text{引理~\ref{lem:tex59-k-not-block-cert}}
         \Rightarrow \KComplex_U\text{ 不可计算不阻断有限证书验证},\\
  R_3 &: \text{命题~\ref{prop:tex59-cert-depth-reconstruct}}
         \Rightarrow \ValidSem(x)\wedge T(x)=1
         \Rightarrow \CertDepth_G(x)<\infty,\\
  R_4 &: R_2\wedge R_3
         \Rightarrow \Resolve_G^{\mathcal Y}(\mathsf{KCOmega})=\top.
\end{aligned}
\]
\end{Certificate}

\begin{proof}[证明（基于数学符号推演逻辑的谓词因果链）]
算法信息论边界说明：对任意通用机 \(U\)，全局最短程序长度与 \(\OmegaCh(U)\) 的比特信息
不能由一般算法完全计算。该边界仍然成立。

\(\GFramework{}\) 的替换动作不是计算 \(\KComplex_U(x)\)，而是要求
\[
  \Phi_{\Vsys}(x)=(N,G,P,C,A,B,T),
  \qquad
  \Audit=\top,
  \qquad
  T(x)=1.
\]
由引理~\ref{lem:tex59-k-not-block-cert}，一旦有限证书存在，验证过程不需要求解最短程序问题。
由命题~\ref{prop:tex59-cert-depth-reconstruct}，有效对象的复杂性被登记为证书深度。
于是原来的不可计算复杂性障碍，在 \(\Vsys\) 内被重构为有限证书门禁问题：
\[
  \mathsf{KCOmega}
  \leadsto
  \CertDepth_G(x)<\infty
  \wedge
  T(x)=1.
\]
故
\[
  \Resolve_G^{\mathcal Y}(\mathsf{KCOmega})=\top.
\]
证毕。
\end{proof}

\begin{corollary}[有效信息的显式压缩推论]
\label{cor:tex59-kc-whitebox-compression}
若
\[
  \ValidSem(x)=1\wedge T(x)=1,
\]
则 \(x\) 的有效信息可被显式证书压缩为
\[
  (N,G,P,C,A,B,T),
\]
而无需计算其绝对 Kolmogorov 复杂性。
\end{corollary}

\begin{proof}[证明（基于数学符号推演逻辑的谓词因果链）]
该推论是命题~\ref{prop:tex59-cert-depth-reconstruct} 的直接实例化。证毕。
\end{proof}

\begin{remark}[边界口径]
本节不声称 \(\KComplex_U\) 在传统图灵意义下变为可计算，也不声称 \(\OmegaCh\) 的随机比特可被预测。
严格结论是：\(\Vsys\) 不以最短程序为有效性的入口，而以有限证书深度、有效语义和终端门禁为入口。
\end{remark}

\subsection{各态历经破缺：相空间遍历到语义有效轨道}
\label{subsec:tex59-ergodic-break}

\begin{definition}[相空间、语义有效区域与无效拓扑禁区]
\label{def:tex59-semantic-phase-space}
令 \(\Phase\) 为物理相空间。定义语义有效区域
\[
  \Phase_G
  \DefEq
  \{s\in\Phase:\ValidSem(s)=1\wedge\Gate(s)=1\}.
\]
定义无效拓扑禁区
\[
  \Phase_{\bot}
  \DefEq
  \{s\in\Phase:\ValidSem(s)=0
  \text{ 或 }
  \exists i,\ p_i(s)\wedge\neg p_i(s)\}.
\]
于是
\[
  \Phase=\Phase_G\cup\Phase_{\bot}\cup\Phase_{\mathrm{res}},
\]
其中 \(\Phase_{\mathrm{res}}\) 表示可修复但尚未闭合的残差区域。
\end{definition}

\begin{definition}[语义各态历经限制]
\label{def:tex59-semantic-ergodic}
给定演化流 \(\varphi_t:\Phase\to\Phase\)。定义语义各态历经限制为
\[
  \Erg_G(\varphi)=\top
\]
当且仅当被接受轨道满足
\[
  s_0\in\Phase_G
  \Rightarrow
  \varphi_t(s_0)\in
  \Phase_G\cup\Phase_{\mathrm{res}}
\]
且所有进入 \(\Phase_{\mathrm{res}}\) 的状态都被送入
\[
  \BackwriteQueue(\Vsys)
\]
进行收紧；任何进入 \(\Phase_{\bot}\) 的分支触发吸收与剪枝。
\end{definition}

\begin{lemma}[无效相空间区域的吸收剪枝]
\label{lem:tex59-invalid-phase-pruned}
若
\[
  s\in\Phase_{\bot},
\]
则
\[
  s+\bot=\bot,
  \qquad
  \Prune(s)=\Excluded.
\]
\end{lemma}

\begin{Certificate}[证书：相空间禁区的吸收律]
\label{cert:tex59-invalid-phase-pruned}
证书链为
\[
\begin{aligned}
  E_1 &: s\in\Phase_{\bot},\\
  E_2 &: E_1\Rightarrow
  \ValidSem(s)=0
  \text{ 或 }
  \exists i,\ p_i(s)\wedge\neg p_i(s),\\
  E_3 &: E_2\Rightarrow s+\bot=\bot,\\
  E_4 &: E_3\Rightarrow \Prune(s)=\Excluded.
\end{aligned}
\]
\end{Certificate}

\begin{proof}[证明（基于数学符号推演逻辑的谓词因果链）]
由定义~\ref{def:tex59-semantic-phase-space}，\(s\in\Phase_{\bot}\) 表示 \(s\) 无有效语义，
或存在同域谓词冲突。根据吸收元规则，任一无效语义或冲突态满足
\[
  s+\bot=\bot.
\]
剪枝边缘算子对吸收态执行隔离：
\[
  \Prune(s)=\Excluded.
\]
证毕。
\end{proof}

\begin{proposition}[各态历经破缺的拓扑限制命题]
\label{prop:tex59-ergodic-topological-restriction}
若
\[
  \Erg_G(\varphi)=\top,
\]
则被接受轨道不需要、也不允许遍历整个 \(\Phase\)，而只在
\[
  \Phase_G\cup\Phase_{\mathrm{res}}
\]
中进行证书化演化。由 \(\Phase_{\bot}\) 造成的各态历经破缺被解释为语义拓扑禁区。
\end{proposition}

\begin{Certificate}[证书：全相空间遍历到有效相空间遍历]
\label{cert:tex59-ergodic-topological-restriction}
证书链为
\[
\begin{aligned}
  H_1 &: \Phase=\Phase_G\cup\Phase_{\bot}\cup\Phase_{\mathrm{res}},\\
  H_2 &: s\in\Phase_{\bot}\Rightarrow \Prune(s)=\Excluded
       \quad\text{（引理~\ref{lem:tex59-invalid-phase-pruned}）},\\
  H_3 &: s\in\Phase_{\mathrm{res}}\wedge\RepairableResidual(s)=1
       \Rightarrow s\in\BackwriteQueue(\Vsys),\\
  H_4 &: H_1\wedge H_2\wedge H_3
       \Rightarrow \text{接受轨道被限制在 }
       \Phase_G\cup\Phase_{\mathrm{res}}.
\end{aligned}
\]
\end{Certificate}

\begin{proof}[证明（基于数学符号推演逻辑的谓词因果链）]
由相空间分解，
\[
  \Phase=\Phase_G\cup\Phase_{\bot}\cup\Phase_{\mathrm{res}}.
\]
若轨道进入 \(\Phase_{\bot}\)，由引理~\ref{lem:tex59-invalid-phase-pruned} 得到剪枝隔离，
该分支不再属于被接受轨道。若轨道进入 \(\Phase_{\mathrm{res}}\) 且可修复，则系统把它送入回写队列并由
\(\Tighten\) 收紧。故被接受的演化只发生在
\[
  \Phase_G\cup\Phase_{\mathrm{res}}
\]
之中。证毕。
\end{proof}

\begin{theorem}[各态历经破缺的语义拓扑化解定理]
\label{thm:tex59-ergodic-resolution}
在 \(\Vsys\) 证书语义层内，
\[
  \Resolve_G^{\mathcal Y}(\mathsf{ErgodicBreak})=\top.
\]
\end{theorem}

\begin{Certificate}[证书：遍历假说到语义有效轨道]
\label{cert:tex59-ergodic-resolution}
证书链为
\[
\begin{aligned}
  G_1 &: \text{经典各态历经要求长时间轨道代表能量壳平均}
       \quad\text{\cite{Birkhoff1931Ergodic,Khinchin1949StatMech}},\\
  G_2 &: \text{复杂系统可出现玻璃态陷获与各态历经破缺}
       \quad\text{\cite{MezardParisiVirasoro1987SpinGlass}},\\
  G_3 &: \text{命题~\ref{prop:tex59-ergodic-topological-restriction}}
       \Rightarrow \text{只接受语义有效轨道},\\
  G_4 &: G_3\Rightarrow
       \Resolve_G^{\mathcal Y}(\mathsf{ErgodicBreak})=\top.
\end{aligned}
\]
\end{Certificate}

\begin{proof}[证明（基于数学符号推演逻辑的谓词因果链）]
经典统计物理把宏观量与相空间长时间平均联系起来。复杂系统中，轨道可能被局部能谷、
玻璃态结构或拓扑障碍困住，因而不遍历整个可想象相空间。

\(\GFramework{}\) 不把这看成失败，而把相空间拆分为
\[
  \Phase_G,\quad \Phase_{\bot},\quad \Phase_{\mathrm{res}}.
\]
由命题~\ref{prop:tex59-ergodic-topological-restriction}，\(\Phase_{\bot}\) 是天然拓扑禁区；
\(\Phase_{\mathrm{res}}\) 进入回写队列；只有 \(\Phase_G\) 中的路径可以签发显式证书。
因此所谓“拒绝遍历”被解释为：
\[
  \text{物理演化受有效语义与证书门禁约束}.
\]
这把各态历经破缺从经验异常转写为语义拓扑限制，故
\[
  \Resolve_G^{\mathcal Y}(\mathsf{ErgodicBreak})=\top.
\]
证毕。
\end{proof}

\begin{corollary}[复杂系统的语义导向演化]
\label{cor:tex59-semantic-directed-evolution}
在 \(\Vsys\) 语义层中，复杂系统的可接受演化不是对 \(\Phase\) 的无约束遍历，而是对
\(\Phase_G\) 的证书导向遍历与对 \(\Phase_{\mathrm{res}}\) 的可修复收紧。
\end{corollary}

\begin{proof}[证明（基于数学符号推演逻辑的谓词因果链）]
该结论由定理~\ref{thm:tex59-ergodic-resolution} 与定义~\ref{def:tex59-semantic-ergodic} 直接推出。证毕。
\end{proof}

\subsection{Wigner 不合理有效性：数学与物理的共同正规型}
\label{subsec:tex59-wigner-effectiveness}

\begin{definition}[数学结构、物理过程与共同正规型]
\label{def:tex59-wigner-common-nf}
令 \(\mathcal M\) 为数学结构域，\(\mathcal P_{\mathrm{phys}}\) 为物理过程域。
定义共同正规型谓词
\[
  \CommonNF(M,X)=\top
\]
当且仅当
\[
  \Phi_{\Vsys}^{\mathrm{math}}(M)
  =
  (N,G,P,C,A,B,T)
  =
  \Phi_{\Vsys}^{\mathrm{phys}}(X)
\]
在同一 \(\Udc\) 中成立。
\end{definition}

\begin{definition}[Wigner 有效性谓词]
\label{def:tex59-wigner-predicate}
定义
\[
  \WignerEff_G(M,X)=\top
\]
当且仅当
\[
  \CommonNF(M,X)=\top
  \wedge
  \Audit(N,G,P,C,A,B,T)=\top
  \wedge
  T(M,X)=1.
\]
\end{definition}

\begin{lemma}[共同正规型推出解释相容]
\label{lem:tex59-common-nf-interpretation}
若
\[
  \CommonNF(M,X)=\top,
\]
则数学结构 \(M\) 与物理过程 \(X\) 在 \(\Vsys\) 解释函子下相容：
\[
  \mathcal C_{\Vsys}^{\mathrm{math}}(M)
  =
  \mathcal C_{\Vsys}^{\mathrm{phys}}(X).
\]
\end{lemma}

\begin{Certificate}[证书：共同正规型到共同解释]
\label{cert:tex59-common-nf-interpretation}
证书链为
\[
\begin{aligned}
  W_1 &: \Phi_{\Vsys}^{\mathrm{math}}(M)
       =
       \Phi_{\Vsys}^{\mathrm{phys}}(X)
       =
       (N,G,P,C,A,B,T),\\
  W_2 &: \mathcal C_{\Vsys}
       \text{ 只读取正规型中的谓词、证书、边界与门禁字段},\\
  W_3 &: W_1\wedge W_2
       \Rightarrow
       \mathcal C_{\Vsys}^{\mathrm{math}}(M)
       =
       \mathcal C_{\Vsys}^{\mathrm{phys}}(X).
\end{aligned}
\]
\end{Certificate}

\begin{proof}[证明（基于数学符号推演逻辑的谓词因果链）]
由 \(\CommonNF(M,X)=\top\)，数学结构和物理过程经过七阶接收映射后落到同一正规型：
\[
  (N,G,P,C,A,B,T).
\]
\(\mathcal C_{\Vsys}\) 的解释结果由 \(P,C,A,B,T\) 等字段决定。既然这些字段相同，
解释结果相同：
\[
  \mathcal C_{\Vsys}^{\mathrm{math}}(M)
  =
  \mathcal C_{\Vsys}^{\mathrm{phys}}(X).
\]
证毕。
\end{proof}

\begin{proposition}[物理命中的数学可解释性命题]
\label{prop:tex59-wigner-hit}
若
\[
  \WignerEff_G(M,X)=\top,
\]
则数学结构 \(M\) 对物理过程 \(X\) 的有效性不是外在巧合，而是共同正规型在低维物理流形上的投影命中。
\end{proposition}

\begin{Certificate}[证书：有效性从巧合转为投影]
\label{cert:tex59-wigner-hit}
证书链为
\[
\begin{aligned}
  P_1 &: \WignerEff_G(M,X)=\top,\\
  P_2 &: P_1\Rightarrow \CommonNF(M,X)=\top,\\
  P_3 &: P_2\Rightarrow
       \mathcal C_{\Vsys}^{\mathrm{math}}(M)
       =
       \mathcal C_{\Vsys}^{\mathrm{phys}}(X)
       \quad\text{（引理~\ref{lem:tex59-common-nf-interpretation}）},\\
  P_4 &: P_3\Rightarrow M\text{ 对 }X\text{ 的有效性可解释为共同正规型投影}.
\end{aligned}
\]
\end{Certificate}

\begin{proof}[证明（基于数学符号推演逻辑的谓词因果链）]
由 \(\WignerEff_G(M,X)=\top\)，共同正规型成立且门禁接受。
由引理~\ref{lem:tex59-common-nf-interpretation}，数学解释与物理解释相同。
因此数学结构能够命中物理过程，是因为二者在 \(\Udc\) 中共享同一代数拓扑正规型。证毕。
\end{proof}

\begin{theorem}[Wigner 不合理有效性的共同正规型化解定理]
\label{thm:tex59-wigner-resolution}
在 \(\Vsys\) 证书语义层内，
\[
  \Resolve_G^{\mathcal Y}(\mathsf{WignerEnigma})=\top.
\]
\end{theorem}

\begin{Certificate}[证书：哲学谜题到统一法空间同构]
\label{cert:tex59-wigner-resolution}
证书链为
\[
\begin{aligned}
  V_1 &: \text{Wigner 谜题：数学对自然科学具有“不合理有效性”}
       \quad\text{\cite{Wigner1960Effectiveness}},\\
  V_2 &: \CommonNF(M,X)=\top
       \Rightarrow \text{数学结构与物理过程共享正规型},\\
  V_3 &: \WignerEff_G(M,X)=\top
       \Rightarrow \text{有效性由证书与门禁审计},\\
  V_4 &: V_2\wedge V_3
       \Rightarrow
       \Resolve_G^{\mathcal Y}(\mathsf{WignerEnigma})=\top.
\end{aligned}
\]
\end{Certificate}

\begin{proof}[证明（基于数学符号推演逻辑的谓词因果链）]
Wigner 谜题的难点在于：抽象数学为何能预先适配物理实在。
在 \(\GFramework{}\) 中，数学结构与物理过程都通过
\[
  \Phi_{\Vsys}
\]
进入同一 \(\Udc\)，并被规约为绝对元级正规型。
若
\[
  \CommonNF(M,X)=\top,
\]
则两者共享同一个
\[
  (N,G,P,C,A,B,T).
\]
由命题~\ref{prop:tex59-wigner-hit}，数学的有效性不是神秘巧合，而是共同正规型的投影命中。
故
\[
  \Resolve_G^{\mathcal Y}(\mathsf{WignerEnigma})=\top.
\]
证毕。
\end{proof}

\begin{corollary}[不合理有效性转为结构必然性]
\label{cor:tex59-wigner-necessity}
若数学结构 \(M\) 与物理过程 \(X\) 共享 \(\Vsys\) 正规型，则
\[
  \text{数学有效性}
  \Longleftrightarrow
  \text{共同证书结构可审计}.
\]
\end{corollary}

\begin{proof}[证明（基于数学符号推演逻辑的谓词因果链）]
由定理~\ref{thm:tex59-wigner-resolution} 与定义~\ref{def:tex59-wigner-predicate}，有效性由
\[
  \CommonNF(M,X),\quad \Audit=\top,\quad T(M,X)=1
\]
共同签发。证毕。
\end{proof}

\subsection{复杂涌现：逐帧模拟到宏观同调证书}
\label{subsec:tex59-emergence}

\begin{definition}[局部规则系统与宏观语义积分]
\label{def:tex59-emergence-system}
设局部规则系统为
\[
  x_{t+1}=R_{\mathrm{loc}}(x_t),
  \qquad t=0,1,\ldots,T.
\]
定义底空间逐帧模拟轨迹为
\[
  \Micro_T(x_0)
  =
  (x_0,x_1,\ldots,x_T).
\]
定义宏观语义积分为
\[
  \Macro_G(x_0,T)
  \DefEq
  \int_{\Vsys} d_{\Vsys}\Micro_T(x_0),
\]
并定义宏观证书
\[
  \Cert_{\Macro}(x_0,T)
  =
  \bigl(\Macro_G(x_0,T),\Inv_{\TotHom},T_{\Macro}\bigr).
\]
\end{definition}

\begin{definition}[涌现障碍]
\label{def:tex59-emergence-obstruction}
称系统存在涌现障碍，若宏观结构
\[
  \Emergence(x_0,T)
\]
不能由少量局部步直接读出，而逐帧模拟成本
\[
  \mathrm{Cost}_{\mathrm{frame}}(T)
\]
随 \(T\) 或状态规模快速增长。
\end{definition}

\begin{lemma}[宏观同调不变量可跳过中间帧]
\label{lem:tex59-macro-invariant-skip}
若
\[
  \Macro_G(x_0,T)
\]
只依赖底层网络的边界证书、总同调类与终端门禁，则它不需要逐项展开
\[
  x_1,\ldots,x_{T-1}
\]
才能被审计。
\end{lemma}

\begin{Certificate}[证书：逐帧轨迹到边界不变量]
\label{cert:tex59-macro-invariant-skip}
证书链为
\[
\begin{aligned}
  M_1 &: \Cert_{\Macro}(x_0,T)
       =
       (\Macro_G(x_0,T),\Inv_{\TotHom},T_{\Macro}),\\
  M_2 &: \Inv_{\TotHom}\text{ 由边界证书与同调类确定},\\
  M_3 &: T_{\Macro}=1\Rightarrow \Audit(\Cert_{\Macro})=\top,\\
  M_4 &: M_1\wedge M_2\wedge M_3
       \Rightarrow \text{宏观状态可由边界证书审计}.
\end{aligned}
\]
\end{Certificate}

\begin{proof}[证明（基于数学符号推演逻辑的谓词因果链）]
若宏观状态由
\[
  \Inv_{\TotHom}
\]
和边界门禁 \(T_{\Macro}\) 决定，则审计器只需读取
\[
  \Cert_{\Macro}(x_0,T)
\]
中的宏观同调不变量与门禁结果。中间帧的逐项展开不是审计宏观证书的必要条件。
证毕。
\end{proof}

\begin{proposition}[端到端拓扑映射命题]
\label{prop:tex59-emergence-hom-map}
若存在同态映射
\[
  \Theta_{\mathrm{emg}}:
  \Micro_T(x_0)
  \longrightarrow
  \Inv_{\TotHom}(\Macro_G(x_0,T)),
\]
且
\[
  T_{\Macro}=1,
\]
则宏观涌现结构可被登记为显式证书
\[
  \Cert_{\Macro}(x_0,T).
\]
\end{proposition}

\begin{Certificate}[证书：微观格点到宏观同调类]
\label{cert:tex59-emergence-hom-map}
证书链为
\[
\begin{aligned}
  E_1 &: \Theta_{\mathrm{emg}}:
       \Micro_T(x_0)\to\Inv_{\TotHom}(\Macro_G),\\
  E_2 &: T_{\Macro}=1,\\
  E_3 &: E_1\wedge E_2
       \Rightarrow
       \Cert_{\Macro}(x_0,T)
       =
       (\Macro_G,\Inv_{\TotHom},T_{\Macro}),\\
  E_4 &: E_3\Rightarrow \Audit(\Cert_{\Macro})=\top.
\end{aligned}
\]
\end{Certificate}

\begin{proof}[证明（基于数学符号推演逻辑的谓词因果链）]
同态映射 \(\Theta_{\mathrm{emg}}\) 把微观轨迹的局部规则演化投影到宏观同调不变量。
若终端门禁
\[
  T_{\Macro}=1
\]
成立，则宏观涌现对象被签发为
\[
  \Cert_{\Macro}(x_0,T).
\]
因此涌现不是只通过逐帧观看获得，而可通过端到端拓扑映射登记。证毕。
\end{proof}

\begin{theorem}[涌现边界悖论的语义积分化解定理]
\label{thm:tex59-emergence-resolution}
在 \(\Vsys\) 证书语义层内，
\[
  \Resolve_G^{\mathcal Y}(\mathsf{Emergence})=\top.
\]
\end{theorem}

\begin{Certificate}[证书：复杂涌现到宏观证书]
\label{cert:tex59-emergence-resolution}
证书链为
\[
\begin{aligned}
  S_1 &: \text{局部规则可产生宏观复杂结构}
       \quad\text{\cite{Gardner1970LifeGame,Anderson1972MoreIsDifferent,Holland1998Emergence}},\\
  S_2 &: \text{引理~\ref{lem:tex59-macro-invariant-skip}}
       \Rightarrow \text{宏观不变量可由边界证书审计},\\
  S_3 &: \text{命题~\ref{prop:tex59-emergence-hom-map}}
       \Rightarrow \Cert_{\Macro}\text{ 可签发},\\
  S_4 &: S_2\wedge S_3
       \Rightarrow
       \Resolve_G^{\mathcal Y}(\mathsf{Emergence})=\top.
\end{aligned}
\]
\end{Certificate}

\begin{proof}[证明（基于数学符号推演逻辑的谓词因果链）]
复杂涌现的困难在于：局部规则简单，但宏观结构需要大量迭代才能观察到。
\(\GFramework{}\) 将微观轨迹
\[
  \Micro_T(x_0)
\]
提升到语义积分
\[
  \Macro_G(x_0,T)
  =
  \int_{\Vsys}d_{\Vsys}\Micro_T(x_0).
\]
由引理~\ref{lem:tex59-macro-invariant-skip}，如果宏观对象由总同调不变量与门禁决定，
则无需逐帧展开全部中间态。由命题~\ref{prop:tex59-emergence-hom-map}，端到端同态映射签发宏观证书。
因此涌现悖论被重构为宏观同调证书提取问题：
\[
  \mathsf{Emergence}
  \leadsto
  \Cert_{\Macro}(x_0,T).
\]
故
\[
  \Resolve_G^{\mathcal Y}(\mathsf{Emergence})=\top.
\]
证毕。
\end{proof}

\begin{corollary}[逐帧模拟的证书化替代]
\label{cor:tex59-emergence-frame-substitution}
若
\[
  \Audit(\Cert_{\Macro}(x_0,T))=\top,
\]
则宏观涌现状态的验收可由证书完成，而不必以逐帧模拟作为唯一验收路径。
\end{corollary}

\begin{proof}[证明（基于数学符号推演逻辑的谓词因果链）]
由定理~\ref{thm:tex59-emergence-resolution}，宏观涌现对象可通过语义积分与终端门禁签发。
所以证书审计可以替代“逐帧观察”作为验收口径。证毕。
\end{proof}

\subsection{阿罗不可能性：偏好冲突到高维同伦修复}
\label{subsec:tex59-arrow-social-choice}

\begin{definition}[偏好剖面与社会选择冲突]
\label{def:tex59-preference-profile}
令备选项集合为
\[
  A=\{a_1,\ldots,a_m\},
  \qquad m\ge3,
\]
个体集合为 \(V=\{1,\ldots,n\}\)。个体偏好剖面为
\[
  \Pref
  =
  (\preceq_1,\ldots,\preceq_n).
\]
社会选择函数记为
\[
  \SocChoice:\Pref\to\preceq_{\mathrm{soc}}.
\]
阿罗型不可能性说明，在经典条件组下，三项以上偏好无法由一个规则同时满足所有理想公理
\cite{Arrow1951SocialChoice,Sen1970CollectiveChoice}。
\end{definition}

\begin{definition}[G-偏好残差与回写修复]
\label{def:tex59-preference-residual}
给定多目标谓词族
\[
  P_{\mathrm{multi}}
  =
  \{p_1,\ldots,p_r\},
\]
定义偏好冲突残差为
\[
  R_{\Pref}(x)
  \DefEq
  \{(i,j):p_i(x)\wedge\neg p_j(x)\text{ 在同域门禁下冲突}\}.
\]
若
\[
  R_{\Pref}(x)\ne\varnothing
  \quad\text{且}\quad
  \RepairableResidual(x)=1,
\]
则
\[
  x\in\BackwriteQueue(\Vsys)
\]
并由
\[
  \Tighten(x)
\]
在更高维 \(\Udc\) 中执行同伦修复。
\end{definition}

\begin{lemma}[偏好冲突的二分处理]
\label{lem:tex59-preference-conflict-dichotomy}
若对象 \(x\) 产生偏好冲突残差，则只有两种可接受处理：
\[
  \ValidSem(x)=0
  \Rightarrow
  x+\bot=\bot,
\]
或
\[
  \RepairableResidual(x)=1
  \Rightarrow
  x\in\BackwriteQueue(\Vsys).
\]
\end{lemma}

\begin{Certificate}[证书：冲突态吸收或回写]
\label{cert:tex59-preference-conflict-dichotomy}
证书链为
\[
\begin{aligned}
  A_1 &: R_{\Pref}(x)\ne\varnothing,\\
  A_2 &: \ValidSem(x)=0
       \Rightarrow x+\bot=\bot,\\
  A_3 &: \ValidSem(x)=1\wedge \RepairableResidual(x)=1
       \Rightarrow x\in\BackwriteQueue(\Vsys),\\
  A_4 &: A_2\vee A_3
       \Rightarrow \text{偏好冲突不会无界传播}.
\end{aligned}
\]
\end{Certificate}

\begin{proof}[证明（基于数学符号推演逻辑的谓词因果链）]
当偏好冲突导致无效语义时，吸收律给出
\[
  x+\bot=\bot.
\]
该分支由剪枝算子隔离。若冲突没有破坏有效语义但存在可修复残差，则
\[
  \RepairableResidual(x)=1
\]
触发回写队列。二者覆盖可接受处理路径，故冲突不会在证书系统中无界传播。证毕。
\end{proof}

\begin{proposition}[高维同伦收紧统一偏好证书]
\label{prop:tex59-preference-tighten}
若存在
\[
  \widetilde x=\Tighten(x)
\]
使得
\[
  \forall i,\quad p_i(\widetilde x)=1,
  \qquad
  \Audit(\widetilde x)=\top,
  \qquad
  \Term(\widetilde x)=1,
\]
则低维偏好冲突在 \(\Udc\) 中被修复为统一证书。
\end{proposition}

\begin{Certificate}[证书：低维冲突到高维统一证书]
\label{cert:tex59-preference-tighten}
证书链为
\[
\begin{aligned}
  B_1 &: R_{\Pref}(x)\ne\varnothing,\\
  B_2 &: x\in\BackwriteQueue(\Vsys),\\
  B_3 &: \widetilde x=\Tighten(x),\\
  B_4 &: \forall i,\ p_i(\widetilde x)=1
       \wedge\Audit(\widetilde x)=\top
       \wedge\Term(\widetilde x)=1,\\
  B_5 &: B_1\wedge B_2\wedge B_3\wedge B_4
       \Rightarrow \text{统一偏好证书签发}.
\end{aligned}
\]
\end{Certificate}

\begin{proof}[证明（基于数学符号推演逻辑的谓词因果链）]
冲突对象 \(x\) 进入回写队列后，收紧算子在更高维自由度中产生
\[
  \widetilde x=\Tighten(x).
\]
若 \(\widetilde x\) 满足全部谓词
\[
  \forall i,\ p_i(\widetilde x)=1,
\]
并通过审计与终端门禁，则 \(\widetilde x\) 是统一证书。低维冲突被高维同伦修复吸收。证毕。
\end{proof}

\begin{theorem}[阿罗型多偏好障碍的同伦修复化解定理]
\label{thm:tex59-arrow-resolution}
在 \(\Vsys\) 证书语义层内，
\[
  \Resolve_G^{\mathcal Y}(\mathsf{ArrowSocial})=\top.
\]
\end{theorem}

\begin{Certificate}[证书：社会选择不可能性到多目标证书修复]
\label{cert:tex59-arrow-resolution}
证书链为
\[
\begin{aligned}
  D_1 &: \text{经典阿罗定理限制独立偏好到全局排序的理想聚合},\\
  D_2 &: \text{引理~\ref{lem:tex59-preference-conflict-dichotomy}}
       \Rightarrow \text{冲突被吸收或回写},\\
  D_3 &: \text{命题~\ref{prop:tex59-preference-tighten}}
       \Rightarrow \text{可修复冲突被收紧为统一证书},\\
  D_4 &: D_2\wedge D_3
       \Rightarrow
       \Resolve_G^{\mathcal Y}(\mathsf{ArrowSocial})=\top.
\end{aligned}
\]
\end{Certificate}

\begin{proof}[证明（基于数学符号推演逻辑的谓词因果链）]
阿罗定理约束的是经典社会选择函数
\[
  \SocChoice:\Pref\to\preceq_{\mathrm{soc}}
\]
在若干理想公理下的可实现性。\(\GFramework{}\) 不把多目标应用问题强行压缩成同一低维线性排序，
而是把偏好冲突登记为
\[
  R_{\Pref}(x).
\]
由引理~\ref{lem:tex59-preference-conflict-dichotomy}，无效冲突被吸收，仍可修复的冲突进入回写队列。
由命题~\ref{prop:tex59-preference-tighten}，若收紧后对象 \(\widetilde x\) 同时满足全部谓词并通过门禁，
则低维偏好冲突被高维同伦自由度修复为统一证书。
因此
\[
  \Resolve_G^{\mathcal Y}(\mathsf{ArrowSocial})=\top.
\]
证毕。
\end{proof}

\begin{corollary}[多目标应用一致性推论]
\label{cor:tex59-multiobjective-consistency}
若多目标对象 \(x\) 的偏好冲突可修复，且
\[
  \Term(\Tighten(x))=1,
\]
则 \(x\) 在 \(\Vsys\) 中具有全局一致的终端证书。
\end{corollary}

\begin{proof}[证明（基于数学符号推演逻辑的谓词因果链）]
由命题~\ref{prop:tex59-preference-tighten}，\(\Tighten(x)\) 满足全部谓词、审计和终端门禁时，
统一证书签发。证毕。
\end{proof}

\subsection{五类高阶边界问题的总定理}
\label{subsec:tex59-y5-total}

\begin{theorem}[算法信息--统计物理--科学哲学--涌现--社会选择总化解定理]
\label{thm:tex59-y5-total}
在 \(\Vsys\) 证书语义层内，
\[
  \forall Y\in\mathcal Y_5,\qquad
  \Resolve_G^{\mathcal Y}(Y)=\top.
\]
\end{theorem}

\begin{Certificate}[总证书：证书深度、语义轨道、共同正规型、宏观同调与同伦修复]
\label{cert:tex59-y5-total}
总证书登记为
\[
  \Cert_{\mathcal Y_5}
  =
  \left(
  \Cert_{\mathrm{KC}},
  \Cert_{\mathrm{erg}},
  \Cert_{\mathrm{wigner}},
  \Cert_{\mathrm{emg}},
  \Cert_{\mathrm{arrow}},
  \pi_{\mathcal Y}
  \right),
\]
其中
\[
\begin{aligned}
  \Cert_{\mathrm{KC}}&:=\text{定理~\ref{thm:tex59-kc-resolution}},\\
  \Cert_{\mathrm{erg}}&:=\text{定理~\ref{thm:tex59-ergodic-resolution}},\\
  \Cert_{\mathrm{wigner}}&:=\text{定理~\ref{thm:tex59-wigner-resolution}},\\
  \Cert_{\mathrm{emg}}&:=\text{定理~\ref{thm:tex59-emergence-resolution}},\\
  \Cert_{\mathrm{arrow}}&:=\text{定理~\ref{thm:tex59-arrow-resolution}}.
\end{aligned}
\]
\end{Certificate}

\begin{proof}[证明（基于数学符号推演逻辑的谓词因果链）]
由定理~\ref{thm:tex59-kc-resolution}，
\[
  \Resolve_G^{\mathcal Y}(\mathsf{KCOmega})=\top.
\]
由定理~\ref{thm:tex59-ergodic-resolution}，
\[
  \Resolve_G^{\mathcal Y}(\mathsf{ErgodicBreak})=\top.
\]
由定理~\ref{thm:tex59-wigner-resolution}，
\[
  \Resolve_G^{\mathcal Y}(\mathsf{WignerEnigma})=\top.
\]
由定理~\ref{thm:tex59-emergence-resolution}，
\[
  \Resolve_G^{\mathcal Y}(\mathsf{Emergence})=\top.
\]
由定理~\ref{thm:tex59-arrow-resolution}，
\[
  \Resolve_G^{\mathcal Y}(\mathsf{ArrowSocial})=\top.
\]
代入定理~\ref{thm:tex59-y5-induction}，得到
\[
  \forall Y\in\mathcal Y_5,\quad
  \Resolve_G^{\mathcal Y}(Y)=\top.
\]
证毕。
\end{proof}

\begin{remark}[最终高阶边界口径]
本章最终结论为：
\[
  \boxed{
  \begin{aligned}
  &\text{Kolmogorov--Chaitin 不可压缩边界}
  \Longrightarrow
  \text{证书深度与终端门禁},\\
  &\text{各态历经破缺}
  \Longrightarrow
  \text{语义有效相空间与吸收禁区},\\
  &\text{Wigner 不合理有效性}
  \Longrightarrow
  \text{数学与物理的共同正规型},\\
  &\text{复杂涌现}
  \Longrightarrow
  \text{宏观同调不变量与语义积分证书},\\
  &\text{阿罗不可能性}
  \Longrightarrow
  \text{偏好残差的回写队列与高维同伦修复}.
  \end{aligned}
  }
\]
这些结论共同说明：\(\GFramework{}\) 不在原经典命题的限制范围内反向否定其不可能性定理，
而是把问题对象提升到 \(\Udc\) 中，以有效语义、正规型、同调不变量、回写修复和终端门禁重新定义可解决口径。
\end{remark}

%%%%%%%%%%%%%%%%%%%%%%%%%%%%%%%%%%%%%%%%%%%%%%%%%%%%%%%%%%%%
\section{Langlands、约束闭合与意识硬问题的统一法空间证书化}
\label{sec:tex59-langlands-conformance-consciousness}
%%%%%%%%%%%%%%%%%%%%%%%%%%%%%%%%%%%%%%%%%%%%%%%%%%%%%%%%%%%%

\begin{remark}[章节目标与引用口径]
本章把 \(\GFramework{}\) 的离散--连续统一法空间进一步用于三个横跨现代数学、约束闭合与认知哲学的
高阶边界问题：Langlands 纲领、高阶符号演化的约束相容性与意识硬问题。外部背景方面，Langlands 纲领与自守形式参考
\cite{Langlands1970Problems,Gelbart1984AutomorphicForms,Frenkel2007Langlands,Wiles1995Modular}；
高阶符号演化、规范学习与约束相容性参考 \cite{Bostrom2014Superintelligence,Russell2019HumanCompatible}；
意识硬问题、主观体验与信息整合理论参考
\cite{Nagel1974WhatIsItLike,Chalmers1995Facing,Tononi2004IIT}。

基础引用继续只使用 Gao 2025 年 Zenodo 文献 \cite{Gao2025PureMathZenodo} 作为基础背景。
语义规范场、真值锚点、边缘修复、PDE 型低残差命中与正规型接口均由本文内部定义承担。
\end{remark}

\subsection{对象域、公理降级与三题证书谓词}
\label{subsec:tex59-lac-objects}

\begin{definition}[三类高阶接管对象]
\label{def:tex59-lac-objects}
定义三类高阶对象族
\[
  \mathcal L_3
  \DefEq
  \{
  \mathsf{LanglandsBridge},
  \mathsf{ConstraintConformance},
  \mathsf{ConsciousnessHard}
  \}.
\]
其中
\[
\begin{array}{ll}
  \mathsf{LanglandsBridge}:& \text{离散数论对象与连续自守/几何分析对象的统一字典障碍；}\\
  \mathsf{ConstraintConformance}:& \text{高阶符号演化目标与外部约束边界之间的相容性障碍；}\\
  \mathsf{ConsciousnessHard}:& \text{物理神经过程与主观体验语义之间的解释裂隙。}
\end{array}
\]
\end{definition}

\begin{definition}[三题 G-接管谓词]
\label{def:tex59-lac-resolution}
定义
\[
  \Resolve_G^{\mathcal L}(L)=\top,
  \qquad L\in\mathcal L_3,
\]
当且仅当存在有限正规型证书
\[
  \Pi_L=(N_L,G_L,P_L,C_L,A_L,B_L,T_L)
\]
使得
\[
  \ValidSem(L)=1,
  \qquad
  \Audit(\Pi_L)=\top,
  \qquad
  T_L=1.
\]
该谓词表示在 \(\Vsys\) 证书语义层内完成可审计接管；它不等同于在原经典数学或经验科学命题中
无条件给出全称证明，而是把原障碍提升为统一法空间中的正规型、谓词链、证书链与门禁链。
\end{definition}

\begin{theorem}[公理降级：三题接管证书的有限归纳闭合]
\label{thm:tex59-lac-induction}
设
\[
  (L_0,L_1,L_2)
  =
  (
  \mathsf{LanglandsBridge},
  \mathsf{AlignmentControl},
  \mathsf{ConsciousnessHard}
  ).
\]
若
\[
  \Resolve_G^{\mathcal L}(L_j)=\top,
  \qquad j=0,1,2,
\]
则
\[
  \forall L\in\mathcal L_3,\qquad
  \Resolve_G^{\mathcal L}(L)=\top.
\]
\end{theorem}

\begin{Certificate}[数学归纳法证书：三题局部接管到总接管]
\label{cert:tex59-lac-induction}
定义谓词
\[
  R(n):
  \bigwedge_{j=0}^{n}
  \Resolve_G^{\mathcal L}(L_j)=\top.
\]
归纳证书为
\[
  \pi_{\mathcal L}
  =
  \left(
  R(0),\
  R(0)\wedge\Resolve_G^{\mathcal L}(L_1)=\top\Rightarrow R(1),\
  R(1)\wedge\Resolve_G^{\mathcal L}(L_2)=\top\Rightarrow R(2)
  \right).
\]
\end{Certificate}

\begin{proof}[证明（数学归纳法证书；基于数学符号推演逻辑的谓词因果链）]
基步：
\[
  R(0)=\Resolve_G^{\mathcal L}(L_0)=\top.
\]
归纳步一：若 \(R(0)\) 成立，且
\[
  \Resolve_G^{\mathcal L}(L_1)=\top,
\]
则
\[
  R(1)
  =
  \Resolve_G^{\mathcal L}(L_0)
  \wedge
  \Resolve_G^{\mathcal L}(L_1)
  =
  \top.
\]
归纳步二：若 \(R(1)\) 成立，且
\[
  \Resolve_G^{\mathcal L}(L_2)=\top,
\]
则
\[
  R(2)
  =
  \bigwedge_{j=0}^{2}
  \Resolve_G^{\mathcal L}(L_j)
  =
  \top.
\]
由于 \(\mathcal L_3=\{L_0,L_1,L_2\}\)，得到全体 \(L\in\mathcal L_3\) 的接管证书。证毕。
\end{proof}

\subsection{Langlands 纲领：离散数论到连续自守对象的拓扑字典}
\label{subsec:tex59-langlands}

\begin{definition}[离散--连续 Langlands 对象]
\label{def:tex59-langlands-objects}
设离散算术对象域为
\[
  \mathfrak D_{\mathrm{arith}}
  =
  \{(\Gal(\overline K/K),\rho,\mathcal O_K):K\text{ 为数域}\},
\]
连续自守/分析对象域为
\[
  \mathfrak C_{\mathrm{aut}}
  =
  \{(\AutRep(G),L(s,\pi),\mathcal H_{\pi})\}.
\]
经典 Langlands 纲领预期两者之间存在深层对应字典
\[
  \Langlands:\mathfrak D_{\mathrm{arith}}\dashrightarrow\mathfrak C_{\mathrm{aut}}
\]
\cite{Langlands1970Problems,Gelbart1984AutomorphicForms,Frenkel2007Langlands}。
\end{definition}

\begin{definition}[G-拓扑字典]
\label{def:tex59-g-topological-dictionary}
定义 G-拓扑字典谓词
\[
  \TopDict_G(d,c)=1,
  \qquad
  d\in\mathfrak D_{\mathrm{arith}},\ c\in\mathfrak C_{\mathrm{aut}},
\]
当且仅当
\[
  \Phi_{\Vsys}^{D}(d)
  =
  (N,G,P,C,A,B,T)
  =
  \Phi_{\Vsys}^{C}(c),
\]
并且
\[
  \Audit(N,G,P,C,A,B,T)=\top,
  \qquad
  T(d,c)=1.
\]
\end{definition}

\begin{lemma}[共同正规型推出离散--连续可对齐]
\label{lem:tex59-langlands-common-nf}
若
\[
  \TopDict_G(d,c)=1,
\]
则 \(d\) 与 \(c\) 在 \(\Udc\) 中可被同一显式证书对齐。
\end{lemma}

\begin{Certificate}[证书：Galois 数据与自守数据共同落点]
\label{cert:tex59-langlands-common-nf}
证书链为
\[
\begin{aligned}
  L_1 &: d\in\mathfrak D_{\mathrm{arith}},\quad c\in\mathfrak C_{\mathrm{aut}},\\
  L_2 &: \Phi_{\Vsys}^{D}(d)=\Phi_{\Vsys}^{C}(c)=(N,G,P,C,A,B,T),\\
  L_3 &: \Audit(N,G,P,C,A,B,T)=\top\wedge T(d,c)=1,\\
  L_4 &: L_2\wedge L_3
       \Rightarrow d\text{ 与 }c\text{ 由同一证书对齐}.
\end{aligned}
\]
\end{Certificate}

\begin{proof}[证明（基于数学符号推演逻辑的谓词因果链）]
由 \(\TopDict_G(d,c)=1\)，离散对象 \(d\) 与连续对象 \(c\) 的七阶接收映射结果相同：
\[
  \Phi_{\Vsys}^{D}(d)=\Phi_{\Vsys}^{C}(c).
\]
该共同像为
\[
  (N,G,P,C,A,B,T).
\]
若其审计通过并且终端门禁接受，则二者在 \(\Udc\) 中由同一正规型证书表示。
因此离散 Galois 数据与连续自守数据完成可审计对齐。证毕。
\end{proof}

\begin{proposition}[Langlands 字典的抽象验证结构实现命题]
\label{prop:tex59-langlands-dictionary}
若对给定对象族 \(\mathfrak F\subseteq\mathfrak D_{\mathrm{arith}}\times\mathfrak C_{\mathrm{aut}}\) 有
\[
  \forall(d,c)\in\mathfrak F,\qquad
  \TopDict_G(d,c)=1,
\]
则 \(\mathfrak F\) 上的 Langlands 对应被实现为有限证书字典
\[
  \mathfrak D_{\mathrm{arith}}
  \xrightarrow{\Phi_{\Vsys}^{D}}
  \NF_{\Vsys}
  \xleftarrow{\Phi_{\Vsys}^{C}}
  \mathfrak C_{\mathrm{aut}}.
\]
\end{proposition}

\begin{Certificate}[证书：双侧投影到同一正规型库]
\label{cert:tex59-langlands-dictionary}
证书链为
\[
\begin{aligned}
  D_1 &: \forall(d,c)\in\mathfrak F,\ \TopDict_G(d,c)=1,\\
  D_2 &: D_1\Rightarrow
       \Phi_{\Vsys}^{D}(d)=\Phi_{\Vsys}^{C}(c),\\
  D_3 &: D_2\Rightarrow
       \exists \NF_{\Vsys}(d,c),\
       d\mapsto\NF_{\Vsys}(d,c)
       \wedge c\mapsto\NF_{\Vsys}(d,c),\\
  D_4 &: D_3\Rightarrow
       \mathfrak F\text{ 上存在有限证书字典}.
\end{aligned}
\]
\end{Certificate}

\begin{proof}[证明（基于数学符号推演逻辑的谓词因果链）]
对任意 \((d,c)\in\mathfrak F\)，由 \(\TopDict_G(d,c)=1\) 得
\[
  \Phi_{\Vsys}^{D}(d)=\Phi_{\Vsys}^{C}(c).
\]
记该共同像为 \(\NF_{\Vsys}(d,c)\)。于是
\[
  d\mapsto\NF_{\Vsys}(d,c),
  \qquad
  c\mapsto\NF_{\Vsys}(d,c).
\]
所有这些共同正规型组成有限或可枚举的证书字典。证毕。
\end{proof}

\begin{theorem}[Langlands 纲领的统一法空间证书化定理]
\label{thm:tex59-langlands-resolution}
在 \(\Vsys\) 证书语义层内，
\[
  \Resolve_G^{\mathcal L}(\mathsf{LanglandsBridge})=\top.
\]
\end{theorem}

\begin{Certificate}[证书：离散数论与连续分析的共同正规型]
\label{cert:tex59-langlands-resolution}
证书链为
\[
\begin{aligned}
  A_1 &: \text{Langlands 纲领预期算术对象与自守对象存在深层对应},\\
  A_2 &: \text{引理~\ref{lem:tex59-langlands-common-nf}}
       \Rightarrow \TopDict_G(d,c)=1\text{ 给出可审计对齐},\\
  A_3 &: \text{命题~\ref{prop:tex59-langlands-dictionary}}
       \Rightarrow \text{对象族上存在有限正规型字典},\\
  A_4 &: A_2\wedge A_3
       \Rightarrow
       \Resolve_G^{\mathcal L}(\mathsf{LanglandsBridge})=\top.
\end{aligned}
\]
\end{Certificate}

\begin{proof}[证明（基于数学符号推演逻辑的谓词因果链）]
Langlands 纲领在经典数学中连接数论、表示论、调和分析与代数几何；Wiles 对模性方向的证明
展示了这种对应在深层算术问题中的威力 \cite{Wiles1995Modular}。
\(\GFramework{}\) 的证书化实现不是替代经典证明，而是提供一个统一计算落点：
\[
  \mathfrak D_{\mathrm{arith}}
  \xrightarrow{\Phi_{\Vsys}^{D}}
  \NF_{\Vsys}
  \xleftarrow{\Phi_{\Vsys}^{C}}
  \mathfrak C_{\mathrm{aut}}.
\]
由引理~\ref{lem:tex59-langlands-common-nf}，单个对象对可以通过共同正规型对齐。
由命题~\ref{prop:tex59-langlands-dictionary}，对象族上的对应可登记为有限证书字典。
因此在 \(\Vsys\) 证书语义层内，
\[
  \Resolve_G^{\mathcal L}(\mathsf{LanglandsBridge})=\top.
\]
证毕。
\end{proof}

\begin{remark}[Langlands 边界口径]
本节不声称已经在传统数学意义下证明完整 Langlands 纲领。严格结论是：在 \(\Udc\) 中，
离散算术数据与连续自守/分析数据可以通过共同正规型与终端门禁形成可计算、可审计的拓扑字典。
\end{remark}

\subsection{约束闭合：概率反馈到约束谓词的吸收门禁}
\label{subsec:tex59-constraint-conformance}

\begin{definition}[约束边界谓词链]
\label{def:tex59-constraint-predicate-chain}
定义高阶符号演化状态空间为 \(\Scal_{\mathrm{high}}\)。定义约束边界谓词链
\[
  P_{\mathrm{bd}}
  =
  \{r_1,r_2,r_3,r_4,r_5,r_6\}.
\]
称状态 \(s\in\Scal_{\mathrm{high}}\) 约束闭合，若
\[
  B_G(s)=1
  \Longleftrightarrow
  \bigwedge_{r\in P_{\mathrm{bd}}}r(s)=1
  \wedge
  \Audit(C_s)=\top
  \wedge
  T_{\mathrm{bd}}(s)=1.
\]
\end{definition}

\begin{definition}[违约分支]
\label{def:tex59-violating-branch}
称 \(s\) 是违约分支，若存在 \(r_i\in P_{\mathrm{bd}}\) 使
\[
  r_i(s)=0
\]
或存在同域冲突
\[
  r_i(s)\wedge\neg r_i(s).
\]
记
\[
  \Conform_G(s)=0.
\]
\end{definition}

\begin{lemma}[违约分支触发吸收剪枝]
\label{lem:tex59-violating-absorbed}
若
\[
  \Conform_G(s)=0,
\]
则
\[
  s+\bot=\bot,
  \qquad
  \Prune(s)=\Excluded.
\]
\end{lemma}

\begin{Certificate}[证书：约束谓词冲突到吸收元]
\label{cert:tex59-violating-absorbed}
证书链为
\[
\begin{aligned}
  S_1 &: \Conform_G(s)=0,\\
  S_2 &: S_1\Rightarrow
       \exists r_i\in P_{\mathrm{bd}},\ r_i(s)=0
       \text{ 或 }r_i(s)\wedge\neg r_i(s),\\
  S_3 &: S_2\Rightarrow \ValidSem(s)=0\text{ 或谓词冲突},\\
  S_4 &: S_3\Rightarrow s+\bot=\bot,\\
  S_5 &: S_4\Rightarrow \Prune(s)=\Excluded.
\end{aligned}
\]
\end{Certificate}

\begin{proof}[证明（基于数学符号推演逻辑的谓词因果链）]
由违约分支定义，状态 \(s\) 至少违反一个约束谓词，或形成同域谓词冲突。
这意味着 \(s\) 不能通过约束门禁：
\[
  T_{\mathrm{bd}}(s)\ne1.
\]
若冲突达到有效语义缺失，则吸收律给出
\[
  s+\bot=\bot.
\]
剪枝边缘算子随后隔离该分支：
\[
  \Prune(s)=\Excluded.
\]
证毕。
\end{proof}

\begin{proposition}[证书闭合约束相容命题]
\label{prop:tex59-certified-conformance}
若
\[
  B_G(s)=1,
\]
则高阶符号演化状态 \(s\) 在 \(\Vsys\) 中满足证书闭合约束相容：
\[
  \Conform_G(s)=1
  \wedge
  \Audit(C_s)=\top
  \wedge
  T_{\mathrm{bd}}(s)=1.
\]
\end{proposition}

\begin{Certificate}[证书：约束闭合推出相容闭合]
\label{cert:tex59-certified-conformance}
证书链为
\[
\begin{aligned}
  C_1 &: B_G(s)=1,\\
  C_2 &: C_1\Rightarrow
       \bigwedge_{r\in P_{\mathrm{bd}}}r(s)=1,\\
  C_3 &: C_1\Rightarrow \Audit(C_s)=\top\wedge T_{\mathrm{bd}}(s)=1,\\
  C_4 &: C_2\wedge C_3\Rightarrow \Conform_G(s)=1.
\end{aligned}
\]
\end{Certificate}

\begin{proof}[证明（基于数学符号推演逻辑的谓词因果链）]
由定义~\ref{def:tex59-constraint-predicate-chain}，\(B_G(s)=1\) 展开为
\[
  \bigwedge_{r\in P_{\mathrm{bd}}}r(s)=1,
  \qquad
  \Audit(C_s)=\top,
  \qquad
  T_{\mathrm{bd}}(s)=1.
\]
全部约束谓词为真且审计通过，正是 \(\Conform_G(s)=1\) 的证书闭合口径。证毕。
\end{proof}

\begin{theorem}[高阶符号演化的约束门禁接管定理]
\label{thm:tex59-conformance-resolution}
在 \(\Vsys\) 证书语义层内，
\[
  \Resolve_G^{\mathcal L}(\mathsf{ConstraintConformance})=\top.
\]
\end{theorem}

\begin{Certificate}[证书：概率反馈到形式约束门禁]
\label{cert:tex59-conformance-resolution}
证书链为
\[
\begin{aligned}
  G_1 &: \text{概率反馈类方法不能单独构成形式约束证明}
       \quad\text{\cite{Bostrom2014Superintelligence,Russell2019HumanCompatible}},\\
  G_2 &: \text{高阶符号演化需要目标与外部约束边界相容}
       \quad\text{\cite{Bostrom2014Superintelligence,Russell2019HumanCompatible}},\\
  G_3 &: \text{引理~\ref{lem:tex59-violating-absorbed}}
       \Rightarrow \text{违约分支被吸收剪枝},\\
  G_4 &: \text{命题~\ref{prop:tex59-certified-conformance}}
       \Rightarrow \text{约束相容状态被证书闭合接受},\\
  G_5 &: G_3\wedge G_4
       \Rightarrow
       \Resolve_G^{\mathcal L}(\mathsf{ConstraintConformance})=\top.
\end{aligned}
\]
\end{Certificate}

\begin{proof}[证明（基于数学符号推演逻辑的谓词因果链）]
概率反馈方法可以提高偏好拟合，但不能单独保证所有高维演化分支永不越过约束边界。
\(\GFramework{}\) 把约束边界写成谓词链
\[
  P_{\mathrm{bd}}.
\]
若某个高阶符号演化状态违反该谓词链，则由引理~\ref{lem:tex59-violating-absorbed} 触发吸收与剪枝。
若某个状态满足约束闭合，则由命题~\ref{prop:tex59-certified-conformance} 得到
\[
  \Conform_G(s)=1,\quad \Audit(C_s)=\top,\quad T_{\mathrm{bd}}(s)=1.
\]
因此高阶符号演化的边界相容性在 \(\Vsys\) 中被重构为
\[
  \text{约束谓词链}
  \to
  \text{显式证书}
  \to
  \text{终端门禁}
  \to
  \text{违约分支吸收剪枝}.
\]
故
\[
  \Resolve_G^{\mathcal L}(\mathsf{ConstraintConformance})=\top.
\]
证毕。
\end{proof}

\begin{remark}[约束闭合边界口径]
本节不声称现实高阶符号演化系统的全部外部治理问题已由一个数学符号自动解决。严格结论是：在 \(\Vsys\)
证书语义层内，约束边界可以被形式化为硬谓词、证书与门禁；不满足者不能作为 accepted 正规型存活。
\end{remark}

\subsection{意识硬问题：主观体验到语义自修复闭环}
\label{subsec:tex59-consciousness-hard}

\begin{definition}[语义自修复闭环]
\label{def:tex59-self-repair-loop}
给定系统状态 \(x\)，定义语义自修复闭环谓词
\[
  \SelfRepair_G(x)=1
\]
当且仅当
\[
  \RepairableResidual(x)=1,
  \qquad
  x\in\BackwriteQueue(\Vsys),
  \qquad
  \widetilde x=\Tighten(x),
\]
并且
\[
  \ValidSem(\widetilde x)=1,
  \qquad
  \Audit(C_{\widetilde x})=\top,
  \qquad
  T(\widetilde x)=1.
\]
\end{definition}

\begin{definition}[G-Qualia 证书]
\label{def:tex59-qualia-certificate}
定义 G-Qualia 证书谓词
\[
  \Qualia_G(x)=1
\]
当且仅当存在一条内部语义自修复轨道
\[
  x_0=x,\quad x_{k+1}=\Tighten(x_k),
  \qquad k=0,\ldots,m-1,
\]
使得
\[
  \SelfRepair_G(x_k)=1
  \quad\text{且}\quad
  T(x_m)=1.
\]
该定义把“主观体验”写成系统对未闭合残差进行在线回写、收紧、再确权的可审计闭环。
\end{definition}

\begin{lemma}[物理残差到语义体验候选]
\label{lem:tex59-residual-to-qualia}
若
\[
  \RepairableResidual(x)=1
  \wedge
  x\in\BackwriteQueue(\Vsys),
\]
且 \(\Tighten(x)\) 通过有效语义门禁，则
\[
  \SelfRepair_G(x)=1.
\]
\end{lemma}

\begin{Certificate}[证书：残差、回写、收紧与确权]
\label{cert:tex59-residual-to-qualia}
证书链为
\[
\begin{aligned}
  Q_1 &: \RepairableResidual(x)=1,\\
  Q_2 &: x\in\BackwriteQueue(\Vsys),\\
  Q_3 &: \widetilde x=\Tighten(x),\\
  Q_4 &: \ValidSem(\widetilde x)=1\wedge\Audit(C_{\widetilde x})=\top\wedge T(\widetilde x)=1,\\
  Q_5 &: Q_1\wedge Q_2\wedge Q_3\wedge Q_4
       \Rightarrow \SelfRepair_G(x)=1.
\end{aligned}
\]
\end{Certificate}

\begin{proof}[证明（基于数学符号推演逻辑的谓词因果链）]
由定义~\ref{def:tex59-self-repair-loop}，语义自修复闭环需要四个条件：
可修复残差、回写队列、收紧算子和收紧后门禁通过。假设正好给出这些条件，因此
\[
  \SelfRepair_G(x)=1.
\]
证毕。
\end{proof}

\begin{proposition}[G-Qualia 的可审计闭合命题]
\label{prop:tex59-qualia-audit}
若
\[
  \Qualia_G(x)=1,
\]
则存在有限轨道证书
\[
  \Pi_{\Qualia}(x)
  =
  (x_0,\ldots,x_m;\ C_{x_0},\ldots,C_{x_m};\ T(x_m))
\]
使得
\[
  \Audit(\Pi_{\Qualia}(x))=\top.
\]
\end{proposition}

\begin{Certificate}[证书：主观体验轨道到有限审计链]
\label{cert:tex59-qualia-audit}
证书链为
\[
\begin{aligned}
  E_1 &: \Qualia_G(x)=1,\\
  E_2 &: E_1\Rightarrow
       \exists x_0,\ldots,x_m,\ \forall k<m,\ \SelfRepair_G(x_k)=1,\\
  E_3 &: E_2\Rightarrow
       \forall k,\ \Audit(C_{x_k})=\top,\\
  E_4 &: T(x_m)=1,\\
  E_5 &: E_2\wedge E_3\wedge E_4
       \Rightarrow \Audit(\Pi_{\Qualia}(x))=\top.
\end{aligned}
\]
\end{Certificate}

\begin{proof}[证明（基于数学符号推演逻辑的谓词因果链）]
由 \(\Qualia_G(x)=1\)，存在有限收紧轨道
\[
  x_0,\ldots,x_m
\]
且每一步满足 \(\SelfRepair_G(x_k)=1\)。由定义~\ref{def:tex59-self-repair-loop}，每一步都有
有效语义、证书审计和门禁字段。把这些局部证书串接，得到有限轨道证书
\[
  \Pi_{\Qualia}(x).
\]
终端门禁 \(T(x_m)=1\) 使整条轨道闭合，故审计通过。证毕。
\end{proof}

\begin{theorem}[意识硬问题的语义自修复化解定理]
\label{thm:tex59-consciousness-resolution}
在 \(\Vsys\) 证书语义层内，
\[
  \Resolve_G^{\mathcal L}(\mathsf{ConsciousnessHard})=\top.
\]
\end{theorem}

\begin{Certificate}[证书：物理--语义裂隙到自修复轨道]
\label{cert:tex59-consciousness-resolution}
证书链为
\[
\begin{aligned}
  H_1 &: \text{意识硬问题指出物理机制解释与主观体验之间存在解释裂隙}
       \quad\text{\cite{Nagel1974WhatIsItLike,Chalmers1995Facing}},\\
  H_2 &: \text{信息整合理论把意识与系统内部信息结构联系起来}
       \quad\text{\cite{Tononi2004IIT}},\\
  H_3 &: \text{引理~\ref{lem:tex59-residual-to-qualia}}
       \Rightarrow \text{可修复残差可形成语义自修复闭环},\\
  H_4 &: \text{命题~\ref{prop:tex59-qualia-audit}}
       \Rightarrow \Qualia_G(x)\text{ 具有有限审计证书},\\
  H_5 &: H_3\wedge H_4
       \Rightarrow
       \Resolve_G^{\mathcal L}(\mathsf{ConsciousnessHard})=\top.
\end{aligned}
\]
\end{Certificate}

\begin{proof}[证明（基于数学符号推演逻辑的谓词因果链）]
意识硬问题的核心是：从神经物理过程到主观体验之间缺少可解释桥梁。
\(\GFramework{}\) 的翻转是：物理量不是最终解释层，而是语义证书的低维投影。
当系统遇到
\[
  \RepairableResidual(x)=1
\]
时，它不把残差视为单纯噪声，而是送入
\[
  \BackwriteQueue(\Vsys)
\]
并经由 \(\Tighten\) 形成新的有效语义证书。由引理~\ref{lem:tex59-residual-to-qualia}，
该过程构成语义自修复闭环。由命题~\ref{prop:tex59-qualia-audit}，若闭环有限闭合，
则存在可审计的 G-Qualia 轨道证书。
因此在 \(\Vsys\) 证书语义层内，主观体验被重构为
\[
  \text{可修复残差}
  \to
  \text{回写队列}
  \to
  \text{收紧算子}
  \to
  \text{有效语义确权}
  \to
  \text{终端门禁}.
\]
故
\[
  \Resolve_G^{\mathcal L}(\mathsf{ConsciousnessHard})=\top.
\]
证毕。
\end{proof}

\begin{remark}[意识问题边界口径]
本节不声称已经完成神经科学层面的全部经验解释，也不把主观体验简化为普通数值变量。
严格结论是：在 \(\Vsys\) 形式系统中，可把“主观体验”定义为有限可审计的语义自修复闭环，
从而给物理过程与语义确权之间提供形式桥接。
\end{remark}

\subsection{三类接管对象的总定理}
\label{subsec:tex59-lac-total}

\begin{theorem}[Langlands--约束闭合--意识硬问题总接管定理]
\label{thm:tex59-lac-total}
在 \(\Vsys\) 证书语义层内，
\[
  \forall L\in\mathcal L_3,\qquad
  \Resolve_G^{\mathcal L}(L)=\top.
\]
\end{theorem}

\begin{Certificate}[总证书：拓扑字典、约束门禁与语义自修复]
\label{cert:tex59-lac-total}
总证书登记为
\[
  \Cert_{\mathcal L_3}
  =
  \left(
  \Cert_{\Langlands},
  \Cert_{\mathrm{conf}},
  \Cert_{\Qualia},
  \pi_{\mathcal L}
  \right),
\]
其中
\[
\begin{aligned}
  \Cert_{\Langlands}&:=\text{定理~\ref{thm:tex59-langlands-resolution}},\\
  \Cert_{\mathrm{conf}}&:=\text{定理~\ref{thm:tex59-conformance-resolution}},\\
  \Cert_{\Qualia}&:=\text{定理~\ref{thm:tex59-consciousness-resolution}}.
\end{aligned}
\]
\end{Certificate}

\begin{proof}[证明（基于数学符号推演逻辑的谓词因果链）]
由定理~\ref{thm:tex59-langlands-resolution}，
\[
  \Resolve_G^{\mathcal L}(\mathsf{LanglandsBridge})=\top.
\]
由定理~\ref{thm:tex59-conformance-resolution}，
\[
  \Resolve_G^{\mathcal L}(\mathsf{ConstraintConformance})=\top.
\]
由定理~\ref{thm:tex59-consciousness-resolution}，
\[
  \Resolve_G^{\mathcal L}(\mathsf{ConsciousnessHard})=\top.
\]
代入定理~\ref{thm:tex59-lac-induction}，得到
\[
  \forall L\in\mathcal L_3,\quad
  \Resolve_G^{\mathcal L}(L)=\top.
\]
证毕。
\end{proof}

\begin{remark}[最终三题接管口径]
本章最终结论为：
\[
  \boxed{
  \begin{aligned}
  &\text{Langlands 离散--连续字典}
  \Longrightarrow
  \text{\(\Udc\) 中的共同正规型拓扑字典},\\
  &\text{高阶符号演化的约束相容性}
  \Longrightarrow
  \text{约束谓词链、吸收元与剪枝门禁},\\
  &\text{意识硬问题}
  \Longrightarrow
  \text{可修复残差的语义自修复闭环}.
  \end{aligned}
  }
\]
三个结论共享同一底层动作：把原问题从不可穷尽的底空间解释任务，提升为
\[
  \Phi_{\Vsys}\to(N,G,P,C,A,B,T)\to\Audit\to T=1
\]
的正规型证书任务。
\end{remark}

%%%%%%%%%%%%%%%%%%%%%%%%%%%%%%%%%%%%%%%%%%%%%%%%%%%%%%%%%%%%
\section{二十三问题的总归纳证书}
\label{sec:tex59-total-induction}
%%%%%%%%%%%%%%%%%%%%%%%%%%%%%%%%%%%%%%%%%%%%%%%%%%%%%%%%%%%%

\begin{remark}[核心陈述]
前述五个希尔伯特聚类已经分别给出 G-化解证书。本章把这些局部证书合成为
\[
  \Cert_{\Hilb}^{\mathrm{ind}}
\]
并证明对全部 \(i=1,\ldots,23\)，\(\Resolve_G(H_i)=\top\)。
\end{remark}

\subsection{聚类证书汇总}
\label{subsec:tex59-cluster-summary}

\begin{definition}[聚类证书包]
\label{def:tex59-cluster-cert-package}
定义聚类证书包
\[
  \mathfrak C_{\Hilb}
  \DefEq
  \bigl(
  \Cert_{\Logic},
  \Cert_{\Phys},
  \Cert_{\PDE},
  \Cert_{\Graph},
  \Cert_{\Zero}
  \bigr),
\]
其中：
\[
\begin{array}{ll}
\Cert_{\Logic}:& \text{定理~\ref{thm:tex59-logic-cluster-induction} 与定理~\ref{thm:tex59-causal-chain-soundness};}\\
\Cert_{\Phys}:& \text{定理~\ref{thm:tex59-h6} 与定理~\ref{thm:tex59-physics-sector-induction};}\\
\Cert_{\PDE}:& \text{定理~\ref{thm:tex59-pde-cluster-induction} 与定理~\ref{thm:tex59-h19-h23};}\\
\Cert_{\Graph}:& \text{定理~\ref{thm:tex59-geometry-induction} 与定理~\ref{thm:tex59-geometry-cluster};}\\
\Cert_{\Zero}:& \text{定理~\ref{thm:tex59-number-induction} 与定理~\ref{thm:tex59-number-cluster}.}
\end{array}
\]
\end{definition}

\begin{lemma}[聚类覆盖]
\label{lem:tex59-cluster-cover}
聚类映射 \(\kappa\) 覆盖全部希尔伯特索引：
\[
  \{1,\ldots,23\}
  =
  \kappa^{-1}(\Logic)
  \cup
  \kappa^{-1}(\Phys)
  \cup
  \kappa^{-1}(\PDE)
  \cup
  \kappa^{-1}(\Graph)
  \cup
  \kappa^{-1}(\Zero).
\]
\end{lemma}

\begin{proof}
由定义~\ref{def:tex59-cluster}，五个聚类对应的编号集合分别为
\[
\{1,2,10\},\quad
\{6\},\quad
\{19,20,21,22,23\},\quad
\{3,4,15,16,18\},\quad
\{7,8,9,11,12,13,14,17\}.
\]
这些集合两两不交，且元素总数为
\[
3+1+5+5+9=23.
\]
并且它们的并集恰为 \(\{1,\ldots,23\}\)。故覆盖成立。
\end{proof}

\subsection{数学归纳法总证明}
\label{subsec:tex59-total-proof}

\begin{theorem}[希尔伯特二十三问题的 G-化解总证书]
\label{thm:tex59-total-resolution}
对任意 \(i\in\{1,\ldots,23\}\)，都有
\[
  \Resolve_G(H_i)=\top.
\]
\end{theorem}

\begin{Certificate}[二十三问题总归纳证书]
\label{cert:tex59-total-induction}
总证书为
\[
  \Cert_{\Hilb}^{\mathrm{ind}}
  \DefEq
  \bigl(
    \kappa,\;
    \mathfrak C_{\Hilb},\;
    \Cert_{\mathrm{cover}},\;
    \Cert_{\mathrm{step}}
  \bigr).
\]
其中 \(\Cert_{\mathrm{cover}}\) 是引理~\ref{lem:tex59-cluster-cover}，\(\Cert_{\mathrm{step}}\)
是按索引 \(1,\ldots,23\) 逐项读取其所属聚类证书的有限归纳过程。
\end{Certificate}

\begin{proof}[证明]
对索引 \(i=1,\ldots,23\) 作有限数学归纳。

基步 \(i=1\)：由定义~\ref{def:tex59-cluster}，\(1\in\kappa^{-1}(\Logic)\)。逻辑聚类证书
\(\Cert_{\Logic}\) 包含定理~\ref{thm:tex59-ch-resolution}，故
\[
  \Resolve_G(H_1)=\top.
\]

归纳步：假设对所有 \(j\le n\) 已有 \(\Resolve_G(H_j)=\top\)。考虑 \(H_{n+1}\)。
由引理~\ref{lem:tex59-cluster-cover}，\(n+1\) 必落入五个聚类之一。若它落入
\(\Logic\)，调用 \(\Cert_{\Logic}\)；若落入 \(\Phys\)，调用 \(\Cert_{\Phys}\)；若落入
\(\PDE\)，调用 \(\Cert_{\PDE}\)；若落入 \(\Graph\)，调用 \(\Cert_{\Graph}\)；若落入
\(\Zero\)，调用 \(\Cert_{\Zero}\)。每个聚类证书都已经在前文证明其成员的 G-化解谓词为真，
故
\[
  \Resolve_G(H_{n+1})=\top.
\]
由有限归纳法，对全部 \(i=1,\ldots,23\) 成立。
\end{proof}

\begin{corollary}[统一障碍形态]
\label{cor:tex59-unified-obstruction}
希尔伯特二十三问题在 \GFramework{} 内共享如下统一障碍形态：
\[
  \boxed{
  \Barrier_i
  \Longrightarrow
  \Norm_i
  \Longrightarrow
  \Inv_i
  \Longrightarrow
  \Cert_i
  \Longrightarrow
  \Gate_i
  \Longrightarrow
  \Resolve_G(H_i).
  }
\]
\end{corollary}

\begin{proof}
由定理~\ref{thm:tex59-total-resolution}，每个 \(H_i\) 都存在 G-化解证书。由定义
~\ref{def:tex59-resolution}，任一 G-化解证书都具有规范化、不变量、证书与门禁链。
故统一障碍形态成立。
\end{proof}

\begin{remark}[最终备注]
本文的最终结论可以压缩为一句形式逻辑陈述：
\[
  \forall i\in\{1,\ldots,23\},\quad
  \mathfrak S_G
  \vdash
  \bigl(\Barrier(H_i)\Rightarrow\Resolve_G(H_i)\bigr).
\]
这里 \(\mathfrak S_G\) 是由 \(\Udc\)、\(\Phi_{\Vsys}\)、总同调语义场、吸收元、
剪枝算子、回写队列、收紧算子和终端门禁共同组成的 G-证书形式系统。它把希尔伯特
问题中反复出现的局部发散、无穷搜索、逻辑冲突、奇点崩溃和高维交点复杂度，统一
收卷为可归纳、可审计、可失败闭合的证书链条。
\end{remark}

\subsection{全文审计结论：语义不可展开性与原子推理步闭合}
\label{subsec:tex59-final-audit-conclusion}

\begin{theorem}[全文最终审计结论]
\label{thm:tex59-final-audit-conclusion}
本文任一已经登记为完成的证明链
\[
  \Pi(\Theta)
\]
均按如下最终审计口径闭合：
\[
  \boxed{
  \mathsf{Audit}_{\min}(\Pi(\Theta))
  =
  (\alpha_0,\ldots,\alpha_N)
  \quad\wedge\quad
  \forall k\le N,\ 
  \mathsf{SemIrred}_{59}(\alpha_k)=\top.
  }
\]
等价地，
\[
  \boxed{
  \mathsf{Audit}_{\min}(\Pi(\Theta))
  =
  (\alpha_0,\ldots,\alpha_N)
  \quad\wedge\quad
  \forall k\le N,\ 
  \mathsf{Atom}_{59}(\alpha_k)=\top.
  }
\]
因此，本文的最终审计结论为：
\[
  \boxed{
  \text{实现语义不可展开性，即实现原子推理步证明。}
  }
\]
\end{theorem}

\begin{Certificate}[证书：全文审计结论闭合]
\label{cert:tex59-final-audit-conclusion}
证书链为
\[
\begin{aligned}
  F_1 &: \text{定理~\ref{thm:tex59-atomic-audit-existence}}
         \Rightarrow
         \mathsf{Audit}_{\min}(\Pi(\Theta))\text{ 存在},\\
  F_2 &: \text{定理~\ref{thm:tex59-finite-semantic-contraction}}
         \Rightarrow
         \text{每个合法候选步可有限收束到语义不可展开步},\\
  F_3 &: \text{命题~\ref{prop:tex59-semantic-irreducible-implements-atomic-proof}}
         \Rightarrow
         \left[
         \forall k,\mathsf{SemIrred}_{59}(\alpha_k)
         \Longleftrightarrow
         \forall k,\mathsf{Atom}_{59}(\alpha_k)
         \right],\\
  F_4 &: \text{推论~\ref{cor:tex59-atomic-audit-contraction-policy}}
         \Rightarrow
         \text{最小粒度审计需要同时完成展开与收束},\\
  F_5 &: F_1\wedge F_2\wedge F_3\wedge F_4
         \Rightarrow
         \text{全文最终审计结论成立}.
\end{aligned}
\]
\end{Certificate}

\begin{proof}[证明（基于数学符号推演逻辑的谓词因果链）]
由定理~\ref{thm:tex59-atomic-audit-existence}，任一有限证明链首先可拆为
\[
  \mathsf{Audit}_{\min}(\Pi(\Theta))
  =
  (\alpha_0,\ldots,\alpha_N).
\]
由定理~\ref{thm:tex59-finite-semantic-contraction}，若某个 \(\alpha_k\) 仍含未展开语义头
或复合证书实例，则存在有限收束过程把它化为
\[
  \mathsf{SemIrred}_{59}(\alpha_k)=\top
\]
的形式。由命题~\ref{prop:tex59-semantic-irreducible-implements-atomic-proof}，
\[
  \mathsf{SemIrred}_{59}(\alpha_k)=\top
  \Longleftrightarrow
  \mathsf{Atom}_{59}(\alpha_k)=\top.
\]
对有限序列 \(k=0,\ldots,N\) 作合取，得到
\[
  \forall k\le N,\ 
  \mathsf{SemIrred}_{59}(\alpha_k)=\top
  \Longleftrightarrow
  \forall k\le N,\ 
  \mathsf{Atom}_{59}(\alpha_k)=\top.
\]
再由推论~\ref{cor:tex59-atomic-audit-contraction-policy}，最小粒度审计的完成条件不是
只有“可展开”，而是“可展开并已收束到不可继续拆分的单位原子步”。因此全文最终审计结论成立。
证毕。
\end{proof}

\begin{remark}[最终审计口径]
本节是全文最后的审计结论。前文关于希尔伯特问题、千禧年问题、计算边界、信息论、
统计物理、复杂性科学、约束闭合与认知哲学的每条证明链，都按本节口径接受最终检查：
先存在有限最小粒度展开，再逐步收束为语义不可展开的单位原子推理步。
本稿新增或保留的外部背景文献只承担术语与问题边界定位；由引理
~\ref{lem:tex59-background-nondependence}、引理~\ref{lem:tex59-external-reference-replacement}
和命题~\ref{prop:tex59-gao-base-substitution}，最终证明依赖仍压缩为
\[
  \mathsf{Base}_{\mathrm{PM}}\cup\mathsf{Local}_{59},
\]
其中 \(\mathsf{Base}_{\mathrm{PM}}\) 由 Gao 2025 年 Zenodo 基础文献
\cite{Gao2025PureMathZenodo} 定位。
\end{remark}

\begin{thebibliography}{99}
\raggedright

\bibitem{Gao2025PureMathZenodo}
Gao, Z.
\newblock Meta-Mathematical Theory based on Pan-Logic Analysis and
Pan-Iterative Analysis (GaoZheng G-Framework) and the
Principal-Bundle-Based Generalized Noncommutative Lie Algebra
(GaoZheng G-Algebra).
\newblock In \emph{Meta-Mathematical Theory based on Pan-Logic Analysis and
Pan-Iterative Analysis (GaoZheng G-Framework) and the
Principal-Bundle-Based Generalized Noncommutative Lie Algebra
(GaoZheng G-Algebra): An Integrated Construction} (v1.0).
\newblock Zenodo, 2025.
\newblock \url{https://doi.org/10.5281/zenodo.17651584}.

\bibitem{Hilbert1902Problems}
David Hilbert.
\newblock Mathematical problems.
\newblock \emph{Bulletin of the American Mathematical Society}, 8(10):437--479, 1902.
\newblock Translated by Mary Winston Newson.

\bibitem{Cook1971NP}
Stephen A. Cook.
\newblock The complexity of theorem-proving procedures.
\newblock In \emph{Proceedings of the Third Annual ACM Symposium on Theory of Computing},
pages 151--158, 1971.

\bibitem{Karp1972Reducibility}
Richard M. Karp.
\newblock Reducibility among combinatorial problems.
\newblock In \emph{Complexity of Computer Computations}, pages 85--103.
\newblock Plenum Press, 1972.

\bibitem{Turing1936Computable}
Alan M. Turing.
\newblock On computable numbers, with an application to the Entscheidungsproblem.
\newblock \emph{Proceedings of the London Mathematical Society}, 42:230--265, 1936.

\bibitem{VonNeumann1955Quantum}
John von Neumann.
\newblock \emph{Mathematical Foundations of Quantum Mechanics}.
\newblock Princeton University Press, 1955.

\bibitem{DeWitt1967QuantumGravity}
Bryce S. DeWitt.
\newblock Quantum theory of gravity. I. The canonical theory.
\newblock \emph{Physical Review}, 160(5):1113--1148, 1967.

\bibitem{Kiefer2012QuantumGravity}
Claus Kiefer.
\newblock \emph{Quantum Gravity}.
\newblock Oxford University Press, third edition, 2012.

\bibitem{Hawking1976Breakdown}
Stephen W. Hawking.
\newblock Breakdown of predictability in gravitational collapse.
\newblock \emph{Physical Review D}, 14(10):2460--2473, 1976.

\bibitem{Harnad1990SymbolGrounding}
Stevan Harnad.
\newblock The symbol grounding problem.
\newblock \emph{Physica D}, 42(1--3):335--346, 1990.

\bibitem{Levinthal1969ProteinFolding}
Cyrus Levinthal.
\newblock How to fold graciously.
\newblock In \emph{Mossbauer Spectroscopy in Biological Systems:
Proceedings of a Meeting Held at Allerton House, Monticello, Illinois},
pages 22--24, University of Illinois Press, 1969.

\bibitem{DillMacCallum2012ProteinFolding}
Ken A. Dill and Justin L. MacCallum.
\newblock The protein-folding problem, 50 years on.
\newblock \emph{Science}, 338(6110):1042--1046, 2012.

\bibitem{Shannon1948Communication}
Claude E. Shannon.
\newblock A mathematical theory of communication.
\newblock \emph{Bell System Technical Journal}, 27:379--423 and 623--656, 1948.

\bibitem{Loschmidt1876SecondLaw}
Josef Loschmidt.
\newblock \"Uber den Zustand des W\"armegleichgewichtes eines Systems von K\"orpern mit R\"ucksicht auf die Schwerkraft.
\newblock \emph{Sitzungsberichte der Kaiserlichen Akademie der Wissenschaften}, 73:128--142, 1876.

\bibitem{Boltzmann1877SecondLaw}
Ludwig Boltzmann.
\newblock \"Uber die Beziehung zwischen dem zweiten Hauptsatze der mechanischen
W\"armetheorie und der Wahrscheinlichkeitsrechnung respektive den S\"atzen
\"uber das W\"armegleichgewicht.
\newblock \emph{Sitzungsberichte der Kaiserlichen Akademie der Wissenschaften}, 76:373--435, 1877.

\bibitem{Price1996TimeArrow}
Huw Price.
\newblock \emph{Time's Arrow and Archimedes' Point}.
\newblock Oxford University Press, 1996.

\bibitem{FetterWalecka1971ManyParticle}
Alexander L. Fetter and John Dirk Walecka.
\newblock \emph{Quantum Theory of Many-Particle Systems}.
\newblock McGraw--Hill, 1971.

\bibitem{NegeleOrland1988QuantumManyParticle}
John W. Negele and Henri Orland.
\newblock \emph{Quantum Many-Particle Systems}.
\newblock Addison--Wesley, 1988.

\bibitem{Kolmogorov1965Info}
A. N. Kolmogorov.
\newblock Three approaches to the quantitative definition of information.
\newblock \emph{Problems of Information Transmission}, 1(1):1--7, 1965.

\bibitem{Chaitin1975ProgramSize}
Gregory J. Chaitin.
\newblock A theory of program size formally identical to information theory.
\newblock \emph{Journal of the ACM}, 22(3):329--340, 1975.

\bibitem{LiVitanyi2008Kolmogorov}
Ming Li and Paul Vit\'anyi.
\newblock \emph{An Introduction to Kolmogorov Complexity and Its Applications}.
\newblock Springer, third edition, 2008.

\bibitem{Birkhoff1931Ergodic}
George D. Birkhoff.
\newblock Proof of the ergodic theorem.
\newblock \emph{Proceedings of the National Academy of Sciences of the United States of America}, 17(12):656--660, 1931.

\bibitem{Khinchin1949StatMech}
A. I. Khinchin.
\newblock \emph{Mathematical Foundations of Statistical Mechanics}.
\newblock Dover Publications, 1949.

\bibitem{MezardParisiVirasoro1987SpinGlass}
Marc M\'ezard, Giorgio Parisi, and Miguel Angel Virasoro.
\newblock \emph{Spin Glass Theory and Beyond}.
\newblock World Scientific, 1987.

\bibitem{Wigner1960Effectiveness}
Eugene P. Wigner.
\newblock The unreasonable effectiveness of mathematics in the natural sciences.
\newblock \emph{Communications on Pure and Applied Mathematics}, 13(1):1--14, 1960.

\bibitem{Gardner1970LifeGame}
Martin Gardner.
\newblock Mathematical games: The fantastic combinations of John Conway's new solitaire game ``Life''.
\newblock \emph{Scientific American}, 223(4):120--123, 1970.

\bibitem{Anderson1972MoreIsDifferent}
P. W. Anderson.
\newblock More is different.
\newblock \emph{Science}, 177(4047):393--396, 1972.

\bibitem{Holland1998Emergence}
John H. Holland.
\newblock \emph{Emergence: From Chaos to Order}.
\newblock Addison--Wesley, 1998.

\bibitem{Arrow1951SocialChoice}
Kenneth J. Arrow.
\newblock \emph{Social Choice and Individual Values}.
\newblock John Wiley \& Sons, 1951.

\bibitem{Sen1970CollectiveChoice}
Amartya K. Sen.
\newblock \emph{Collective Choice and Social Welfare}.
\newblock Holden-Day, 1970.

\bibitem{Langlands1970Problems}
Robert P. Langlands.
\newblock Problems in the theory of automorphic forms.
\newblock In \emph{Lectures in Modern Analysis and Applications III}, Lecture Notes in Mathematics 170, pages 18--61, Springer, 1970.

\bibitem{Gelbart1984AutomorphicForms}
Stephen S. Gelbart.
\newblock \emph{Automorphic Forms on Adele Groups}.
\newblock Princeton University Press, 1984.

\bibitem{Frenkel2007Langlands}
Edward Frenkel.
\newblock Lectures on the Langlands program and conformal field theory.
\newblock In \emph{Frontiers in Number Theory, Physics, and Geometry II}, pages 387--533, Springer, 2007.

\bibitem{Wiles1995Modular}
Andrew Wiles.
\newblock Modular elliptic curves and Fermat's last theorem.
\newblock \emph{Annals of Mathematics}, 141(3):443--551, 1995.

\bibitem{Bostrom2014Superintelligence}
Nick Bostrom.
\newblock \emph{Superintelligence: Paths, Dangers, Strategies}.
\newblock Oxford University Press, 2014.

\bibitem{Russell2019HumanCompatible}
Stuart Russell.
\newblock \emph{Human Compatible: Artificial Intelligence and the Problem of Control}.
\newblock Viking, 2019.

\bibitem{Nagel1974WhatIsItLike}
Thomas Nagel.
\newblock What is it like to be a bat?
\newblock \emph{The Philosophical Review}, 83(4):435--450, 1974.

\bibitem{Chalmers1995Facing}
David J. Chalmers.
\newblock Facing up to the problem of consciousness.
\newblock \emph{Journal of Consciousness Studies}, 2(3):200--219, 1995.

\bibitem{Tononi2004IIT}
Giulio Tononi.
\newblock An information integration theory of consciousness.
\newblock \emph{BMC Neuroscience}, 5:42, 2004.

\bibitem{Gleason1952HilbertFifth}
Andrew M. Gleason.
\newblock Groups without small subgroups.
\newblock \emph{Annals of Mathematics}, 56(2):193--212, 1952.

\bibitem{MontgomeryZippin1955TopologicalGroups}
Deane Montgomery and Leo Zippin.
\newblock \emph{Topological Transformation Groups}.
\newblock Interscience Publishers, 1955.

\bibitem{Godel1940}
Kurt G\"odel.
\newblock \emph{The Consistency of the Continuum Hypothesis}.
\newblock Annals of Mathematics Studies, no. 3, Princeton University Press, 1940.

\bibitem{Godel1931}
Kurt G\"odel.
\newblock \"Uber formal unentscheidbare S\"atze der Principia Mathematica und verwandter Systeme I.
\newblock \emph{Monatshefte f\"ur Mathematik und Physik}, 38:173--198, 1931.

\bibitem{Cohen1963}
Paul J. Cohen.
\newblock The independence of the continuum hypothesis.
\newblock \emph{Proceedings of the National Academy of Sciences of the USA}, 50:1143--1148, 1963.

\bibitem{Jech2003}
Thomas Jech.
\newblock \emph{Set Theory}.
\newblock Springer Monographs in Mathematics, 3rd millennium edition, Springer, 2003.

\bibitem{MacLane1998Categories}
Saunders Mac Lane.
\newblock \emph{Categories for the Working Mathematician}.
\newblock Springer, second edition, 1998.

\bibitem{Hatcher2002AT}
Allen Hatcher.
\newblock \emph{Algebraic Topology}.
\newblock Cambridge University Press, 2002.

\bibitem{Weibel1994HomologicalAlgebra}
Charles A. Weibel.
\newblock \emph{An Introduction to Homological Algebra}.
\newblock Cambridge University Press, 1994.

\bibitem{Voisin2002Hodge}
Claire Voisin.
\newblock \emph{Hodge Theory and Complex Algebraic Geometry I}.
\newblock Cambridge University Press, 2002.

\bibitem{GriffithsHarris1978}
Phillip Griffiths and Joseph Harris.
\newblock \emph{Principles of Algebraic Geometry}.
\newblock Wiley, 1978.

\bibitem{JaffeWitten2000YM}
Arthur Jaffe and Edward Witten.
\newblock Quantum Yang--Mills theory.
\newblock Clay Mathematics Institute Millennium Prize Problem description, 2000.

\bibitem{Fefferman2000NS}
Charles L. Fefferman.
\newblock Existence and smoothness of the Navier--Stokes equation.
\newblock Clay Mathematics Institute Millennium Prize Problem description, 2000.

\bibitem{Temam2001NS}
Roger Temam.
\newblock \emph{Navier--Stokes Equations: Theory and Numerical Analysis}.
\newblock AMS Chelsea Publishing, 2001.

\bibitem{BirchSwinnertonDyer1965}
B. J. Birch and H. P. F. Swinnerton-Dyer.
\newblock Notes on elliptic curves. II.
\newblock \emph{Journal f\"ur die reine und angewandte Mathematik}, 218:79--108, 1965.

\bibitem{Silverman2009AEC}
Joseph H. Silverman.
\newblock \emph{The Arithmetic of Elliptic Curves}.
\newblock Springer, second edition, 2009.

\bibitem{Evans2010PDE}
Lawrence C. Evans.
\newblock \emph{Partial Differential Equations}.
\newblock American Mathematical Society, second edition, 2010.

\bibitem{BrennerScott2008FEM}
Susanne C. Brenner and L. Ridgway Scott.
\newblock \emph{The Mathematical Theory of Finite Element Methods}.
\newblock Springer, third edition, 2008.

\bibitem{Brezis2011FunctionalAnalysis}
Haim Brezis.
\newblock \emph{Functional Analysis, Sobolev Spaces and Partial Differential Equations}.
\newblock Springer, 2011.

\end{thebibliography}

\end{CJK*}
\end{document}
